% Options for packages loaded elsewhere
\PassOptionsToPackage{unicode}{hyperref}
\PassOptionsToPackage{hyphens}{url}
%
\documentclass[
  11pt,
  letterpaper,
  oneside,
  openany]{scrbook}
\usepackage{amsmath,amssymb}
\usepackage{setspace}
\usepackage{iftex}
\ifPDFTeX
  \usepackage[T1]{fontenc}
  \usepackage[utf8]{inputenc}
  \usepackage{textcomp} % provide euro and other symbols
\else % if luatex or xetex
  \usepackage{unicode-math} % this also loads fontspec
  \defaultfontfeatures{Scale=MatchLowercase}
  \defaultfontfeatures[\rmfamily]{Ligatures=TeX,Scale=1}
\fi
\usepackage{lmodern}
\ifPDFTeX\else
  % xetex/luatex font selection
  \setmainfont[]{Latin Modern Roman}
  \setsansfont[]{Latin Modern Sans}
  \setmonofont[]{Latin Modern Mono}
  \setmathfont[]{Latin Modern Math}
\fi
% Use upquote if available, for straight quotes in verbatim environments
\IfFileExists{upquote.sty}{\usepackage{upquote}}{}
\IfFileExists{microtype.sty}{% use microtype if available
  \usepackage[]{microtype}
  \UseMicrotypeSet[protrusion]{basicmath} % disable protrusion for tt fonts
}{}
\makeatletter
\@ifundefined{KOMAClassName}{% if non-KOMA class
  \IfFileExists{parskip.sty}{%
    \usepackage{parskip}
  }{% else
    \setlength{\parindent}{0pt}
    \setlength{\parskip}{6pt plus 2pt minus 1pt}}
}{% if KOMA class
  \KOMAoptions{parskip=half}}
\makeatother
\usepackage{xcolor}
\usepackage[margin=0.75in]{geometry}
\usepackage{longtable,booktabs,array}
\usepackage{calc} % for calculating minipage widths
% Correct order of tables after \paragraph or \subparagraph
\usepackage{etoolbox}
\makeatletter
\patchcmd\longtable{\par}{\if@noskipsec\mbox{}\fi\par}{}{}
\makeatother
% Allow footnotes in longtable head/foot
\IfFileExists{footnotehyper.sty}{\usepackage{footnotehyper}}{\usepackage{footnote}}
\makesavenoteenv{longtable}
\setlength{\emergencystretch}{3em} % prevent overfull lines
\providecommand{\tightlist}{%
  \setlength{\itemsep}{0pt}\setlength{\parskip}{0pt}}
\setcounter{secnumdepth}{-\maxdimen} % remove section numbering
\usepackage{graphicx}
\usepackage{microtype}
\usepackage{fancyhdr}
\usepackage{needspace}
\usepackage{tcolorbox}
\tcbuselibrary{skins,breakable}
\definecolor{ink}{HTML}{152D38}
\definecolor{accent}{HTML}{126B7A}
\definecolor{muted}{HTML}{52636A}
\definecolor{pale}{HTML}{EFF5F6}
\newsavebox{\grmathbox}
\setlength{\parindent}{0pt}
\setlength{\parskip}{5pt plus 1pt minus 1pt}
\setlength{\emergencystretch}{2em}
\setlength{\LTpre}{7pt}
\setlength{\LTpost}{7pt}
\setlength{\headheight}{15pt}
\setlength{\footskip}{24pt}
\pagestyle{fancy}
\fancyhf{}
\fancyhead[L]{\sffamily\fontsize{8}{10}\selectfont\color{muted}GENERAL RELATIVITY}
\fancyhead[R]{\sffamily\fontsize{8}{10}\selectfont\color{muted}FROM THE INSIDE OUT}
\fancyfoot[R]{\sffamily\small\color{accent}\thepage}
\renewcommand{\headrulewidth}{0.25pt}
\renewcommand{\footrulewidth}{0pt}
\fancypagestyle{plain}{\fancyhf{}\fancyfoot[R]{\sffamily\small\color{accent}\thepage}\renewcommand{\headrulewidth}{0pt}}
\setkomafont{disposition}{\normalfont\sffamily\bfseries\color{ink}}
\setkomafont{chapter}{\fontsize{23}{28}\selectfont}
\setkomafont{section}{\fontsize{15}{19}\selectfont}
\setkomafont{subsection}{\fontsize{12}{15}\selectfont}
\RedeclareSectionCommand[beforeskip=0pt,afterskip=16pt]{chapter}
\RedeclareSectionCommand[beforeskip=18pt,afterskip=8pt]{section}
\RedeclareSectionCommand[beforeskip=12pt,afterskip=6pt]{subsection}
\renewenvironment{quote}{\begin{tcolorbox}[enhanced,breakable,sharp corners,colback=pale,colframe=accent,boxrule=0pt,borderline west={2pt}{0pt}{accent},left=10pt,right=10pt,top=7pt,bottom=7pt]}{\end{tcolorbox}}
\clubpenalty=10000
\widowpenalty=10000
\displaywidowpenalty=10000
\raggedbottom
\AtBeginDocument{\hypersetup{colorlinks=true,linkcolor=accent,urlcolor=accent,citecolor=accent,pdftitle={General Relativity, From the Inside Out},pdfsubject={An educational guide to spacetime geometry and the Einstein-Hilbert action},pdfauthor={},bookmarksopen=false,bookmarksdepth=2}}
\ifLuaTeX
  \usepackage{selnolig}  % disable illegal ligatures
\fi
\IfFileExists{bookmark.sty}{\usepackage{bookmark}}{\usepackage{hyperref}}
\IfFileExists{xurl.sty}{\usepackage{xurl}}{} % add URL line breaks if available
\urlstyle{same}
\hypersetup{
  hidelinks,
  pdfcreator={LaTeX via pandoc}}

\author{}
\date{}

\begin{document}
\frontmatter

\begin{titlepage}
\thispagestyle{empty}
\vspace*{0.7in}
{\sffamily\fontsize{11}{14}\selectfont\color{accent}A COMPLETE GUIDED EXPEDITION\par}
\vspace{0.35in}
{\sffamily\bfseries\fontsize{38}{43}\selectfont\color{ink}GENERAL\\RELATIVITY\par}
\vspace{0.22in}
{\sffamily\fontsize{23}{29}\selectfont\color{ink}From the Inside Out\par}
\vspace{0.32in}
{\color{accent}\rule{1.05in}{2pt}\par}
\vspace{0.3in}
{\fontsize{14}{20}\selectfont From clocks and vectors to curvature,\\the Einstein-Hilbert action,\\and the modern theory of spacetime.\par}
\vspace{0.55in}
\begin{center}
{\large\color{ink}$\displaystyle R_{\mu\nu}-\frac12 Rg_{\mu\nu}+\Lambda g_{\mu\nu}=\frac{8\pi G_N}{c^4}T_{\mu\nu}$}
\end{center}
\vfill
{\sffamily\fontsize{10}{15}\selectfont\color{muted}24 CHAPTERS\quad /\quad 30 SOLVED EXERCISES\par}
\vspace{0.12in}
{\small Intuition. Derivations. Worked examples. The traps worth understanding.\par}
\vspace{0.2in}
{\small\color{muted}For readers with college-level mathematics and physics.\par}
\end{titlepage}
\frontmatter
\setcounter{tocdepth}{0}
{\small\setstretch{1.0}\DeclareTOCStyleEntry[beforeskip=4pt]{tocline}{chapter}\tableofcontents}
\clearpage

\setstretch{1.06}
\hypertarget{before-we-begin}{%
\chapter{Before we begin}\label{before-we-begin}}

\textbf{For a reader with college-level mathematics and physics.} You
should be comfortable with derivatives, integrals, elementary vectors
and matrices, Newton\textquotesingle s laws, and basic energy
conservation. Everything specifically geometric - manifolds, covectors,
connections, covariant derivatives, curvature, and metric variations -
is developed here. Some later sections introduce graduate-level ideas,
but the conceptual staircase remains visible.

The destination is this equation:

\begin{center}
\begingroup
\sbox{\grmathbox}{$\displaystyle 
\boxed{R_{\mu\nu}-\frac12 Rg_{\mu\nu}+\Lambda g_{\mu\nu}
=\frac{8\pi G_N}{c^4}T_{\mu\nu}.}
$}
\typeout{GRDISPLAY|1|\the\wd\grmathbox|\the\linewidth}
\ifdim\wd\grmathbox>\linewidth
\resizebox{\linewidth}{!}{\usebox{\grmathbox}}
\else\usebox{\grmathbox}\fi
\endgroup
\end{center}

The goal is much more interesting than memorizing it. By the end, you
should be able to explain why ordinary derivatives need repairing, what
a connection actually compares, how curvature survives a change of
coordinates, what pressure has to do with gravity, why the coefficient
of \(Rg_{\mu\nu}\) is exactly \(-1/2\), and how a single action produces
the entire field equation.

You should also be able to recognize several seductive mistakes before
they recognize you.

\hypertarget{how-to-travel-through-this-book}{%
\chapter{How to travel through this
book}\label{how-to-travel-through-this-book}}

Read with a pencil and occasionally stop before the next displayed
equation. Predict its indices, its dimensions, or its sign.
Understanding is much easier to counterfeit while reading than while
predicting.

There are three routes:

\begingroup\small\setlength{\tabcolsep}{5pt}\renewcommand{\arraystretch}{1.13}

\begin{longtable}[]{@{}
  >{\raggedright\arraybackslash}p{(\columnwidth - 4\tabcolsep) * \real{0.2500}}
  >{\raggedright\arraybackslash}p{(\columnwidth - 4\tabcolsep) * \real{0.3600}}
  >{\raggedright\arraybackslash}p{(\columnwidth - 4\tabcolsep) * \real{0.3900}}@{}}
\toprule\noalign{}
\begin{minipage}[b]{\linewidth}\raggedright
Route
\end{minipage} & \begin{minipage}[b]{\linewidth}\raggedright
Read
\end{minipage} & \begin{minipage}[b]{\linewidth}\raggedright
What it builds
\end{minipage} \\
\midrule\noalign{}
\endhead
\bottomrule\noalign{}
\endlastfoot
The full expedition & Chapters 1-24, then the exercises & A continuous
foundation from special relativity to modern GR \\
The equation expedition & Chapters 2-15, then 24 & The mathematical and
physical meaning of every term and the full action derivation \\
The physical expedition & Chapters 1, 3-5, 8-12, 16-19 & Clocks, tides,
matter, black holes, waves, and cosmology; return for the intermediate
machinery \\
\end{longtable}

\endgroup

\textbf{Analogies are scaffolding.} A good analogy reveals a
relationship. It does not provide a license to import every feature of
the familiar object. When we use a rubber sheet, an accountant, a map, a
neighboring laboratory, or a rotating compass, we will say where the
comparison breaks.

\textbf{``Derive'' has several meanings.} We can derive a consequence
from assumptions, derive an equation from a chosen action, or motivate
why that action is a good low-energy model. These are different
accomplishments. General relativity is not forced on us by pure logic or
by the equivalence principle alone. Its assumptions must meet
experiment.

\textbf{About the sources.} The explanations, analogies, and worked
calculations are written as an independent tutorial. Links identify
historical evidence, research results, and places to pursue particular
ideas; the book is not a paraphrase of a single textbook. The final
reading guide distinguishes foundational notes from original research.
Exact contemporary parameter estimates and speculative claims are
deliberately unnecessary to the main argument.

\hypertarget{conventions-the-treaty-that-prevents-a-thousand-sign-wars}{%
\chapter{Conventions: the treaty that prevents a thousand sign
wars}\label{conventions-the-treaty-that-prevents-a-thousand-sign-wars}}

Different excellent books use different signs. A disagreement in
notation need not be a disagreement about nature. This book keeps the
following treaty throughout.

\begingroup\small\setlength{\tabcolsep}{5pt}\renewcommand{\arraystretch}{1.13}

\begin{longtable}[]{@{}
  >{\raggedright\arraybackslash}p{(\columnwidth - 2\tabcolsep) * \real{0.3000}}
  >{\raggedright\arraybackslash}p{(\columnwidth - 2\tabcolsep) * \real{0.7000}}@{}}
\toprule\noalign{}
\begin{minipage}[b]{\linewidth}\raggedright
Symbol or convention
\end{minipage} & \begin{minipage}[b]{\linewidth}\raggedright
Meaning
\end{minipage} \\
\midrule\noalign{}
\endhead
\bottomrule\noalign{}
\endlastfoot
Metric signature & \((-,+,+,+)\) \\
Coordinates & \(x^\mu=(ct,x,y,z)\) when using standard SI component
formulas; special charts explicitly say when they use \(t\) \\
Indices & Greek \(\mu,\nu,\ldots=0,1,2,3\); spatial Latin
\(i,j,\ldots=1,2,3\) \\
Summation & A repeated upper-lower pair is summed unless stated
otherwise \\
\(g_{\mu\nu}\) and \(g^{\mu\nu}\) & Metric matrix and its inverse, not
componentwise reciprocals \\
\(g\) & \(\det(g_{\mu\nu})\), negative for a nondegenerate Lorentzian
metric in four dimensions \\
\(G_N\) & Newton\textquotesingle s gravitational constant \\
\(G_{\mu\nu}\) & Einstein tensor, \(R_{\mu\nu}-\tfrac12Rg_{\mu\nu}\) \\
\(\epsilon\) & Rest-frame energy density, in joules per cubic metre \\
\(\rho\) & Mass-equivalent density \(\epsilon/c^2\), when that notation
is used \\
\(p\) & Pressure; it has the same dimensions as energy density \\
\(u^\mu\) & Four-velocity \(dx^\mu/d\tau\), with \(u^\mu u_\mu=-c^2\) \\
\(\Lambda\) & Cosmological constant, with dimensions inverse length
squared \\
Natural units & A section using \(c=1\) says so; powers of \(c\) return
for physical numerical results \\
\end{longtable}

\endgroup

Our curvature convention is

\begin{center}
\begingroup
\sbox{\grmathbox}{$\displaystyle 
R^\rho{}_{\sigma\mu\nu}
=\partial_\mu\Gamma^\rho{}_{\nu\sigma}
-\partial_\nu\Gamma^\rho{}_{\mu\sigma}
+\Gamma^\rho{}_{\mu\lambda}\Gamma^\lambda{}_{\nu\sigma}
-\Gamma^\rho{}_{\nu\lambda}\Gamma^\lambda{}_{\mu\sigma},
$}
\typeout{GRDISPLAY|2|\the\wd\grmathbox|\the\linewidth}
\ifdim\wd\grmathbox>\linewidth
\resizebox{\linewidth}{!}{\usebox{\grmathbox}}
\else\usebox{\grmathbox}\fi
\endgroup
\end{center}

with

\begin{center}
\begingroup
\sbox{\grmathbox}{$\displaystyle 
[\nabla_\mu,\nabla_\nu]V^\rho
=R^\rho{}_{\sigma\mu\nu}V^\sigma,
\qquad
R_{\mu\nu}=R^\rho{}_{\mu\rho\nu}.
$}
\typeout{GRDISPLAY|3|\the\wd\grmathbox|\the\linewidth}
\ifdim\wd\grmathbox>\linewidth
\resizebox{\linewidth}{!}{\usebox{\grmathbox}}
\else\usebox{\grmathbox}\fi
\endgroup
\end{center}

Do not attempt to digest this on arrival. It is here so that later
calculations have a definite address. With these conventions, a round
two-sphere has positive scalar curvature.

The physical point-particle action is \(S=-mc^2\int d\tau\). In
coordinates with \(x^0=ct\), our gravitational action is

\begin{center}
\begingroup
\sbox{\grmathbox}{$\displaystyle 
S_{\mathrm{EH}}=
\frac{c^3}{16\pi G_N}
\int d^4x\,\sqrt{-g}\,(R-2\Lambda),
$}
\typeout{GRDISPLAY|4|\the\wd\grmathbox|\the\linewidth}
\ifdim\wd\grmathbox>\linewidth
\resizebox{\linewidth}{!}{\usebox{\grmathbox}}
\else\usebox{\grmathbox}\fi
\endgroup
\end{center}

and the stress tensor is normalized by

\begin{center}
\begingroup
\sbox{\grmathbox}{$\displaystyle 
\delta S_m=-\frac{1}{2c}\int d^4x\,\sqrt{-g}\,
T_{\mu\nu}\,\delta g^{\mu\nu}.
$}
\typeout{GRDISPLAY|5|\the\wd\grmathbox|\the\linewidth}
\ifdim\wd\grmathbox>\linewidth
\resizebox{\linewidth}{!}{\usebox{\grmathbox}}
\else\usebox{\grmathbox}\fi
\endgroup
\end{center}

The action is called \textbf{Einstein-Hilbert}, after Einstein and David
Hilbert. The determinant, inverse-metric variation, boundary terms, and
factors of \(c\) will all get their own explanations.

A final units trap: angular coordinates are dimensionless. In
\(ds^2=dr^2+r^2d\theta^2\), \(g_{\theta\theta}=r^2\) has dimensions of
length squared. It is the whole line element that must have the correct
units, not every coordinate component separately.

\mainmatter

\hypertarget{chapter-1}{%
\chapter{1. The scandal: gravity changes the measuring
equipment}\label{chapter-1}}

\hypertarget{what-a-theory-of-gravity-must-explain}{%
\section{1.1 What a theory of gravity must
explain}\label{what-a-theory-of-gravity-must-explain}}

Imagine two spacecraft drifting side by side toward Earth with their
engines off. Inside either craft, a released pen floats. The
accelerometer reads zero. Nobody feels a downward force.

Yet the separation between the spacecraft changes. If they are side by
side at the same altitude, their trajectories aim toward the same center
and tend to converge. If one is directly above the other, the lower one
accelerates toward Earth more strongly in the Newtonian description, and
their radial separation tends to grow.

Something gravitational remains after both crews have removed the
experience of weight.

That ``something'' is the first clue to spacetime curvature. \textbf{The
most revealing local gravitational experiment compares neighboring
free-fall trajectories.} A single freely falling laboratory can erase
the connection coefficients at an event by choosing suitable
coordinates. A family of laboratories can detect tidal effects that no
coordinate change removes.

Now imagine standing on the ground. Your accelerometer reads
approximately \(9.8\,\mathrm{m/s^2}\). In relativistic language, the
ground is preventing your natural free-fall motion. The electromagnetic
forces holding the floor together push upward on you.

Your everyday intuition says the person standing still is unaccelerated
and the falling person accelerates. An accelerometer says the opposite,
because it measures \textbf{proper acceleration}: departure from
inertial free fall, rather than the second derivative of a chosen
coordinate.

This does not make Newtonian physics foolish. Near
Earth\textquotesingle s surface, Newtonian coordinates are extremely
useful. It means ``acceleration'' has more than one meaning, and a good
theory must specify which one an instrument reads.

\begin{quote}
\textbf{First trap door:} ``Gravity disappears in free fall'' means the
locally removable inertial-gravitational effects disappear to first
order. Tidal curvature generally remains. A falling laboratory is not a
universe-sized eraser.
\end{quote}

\hypertarget{the-weak-spot-in-the-old-division-of-labor}{%
\section{1.2 The weak spot in the old division of
labor}\label{the-weak-spot-in-the-old-division-of-labor}}

In elementary mechanics, space and time provide an arena. A particle has
a trajectory through that arena. Forces change its motion. Rulers and
clocks are assumed to provide the arena\textquotesingle s fixed
geometry.

Special relativity already complicates this arrangement: different
observers split spacetime into ``space'' and ``time'' differently, and
elapsed time depends on a worldline. But the flat spacetime metric is
still prescribed.

General relativity takes the next step. The metric becomes a dynamical
field.

A metric tells us the interval between infinitesimally separated events:

\begin{center}
\begingroup
\sbox{\grmathbox}{$\displaystyle 
ds^2=g_{\mu\nu}(x)\,dx^\mu dx^\nu.
$}
\typeout{GRDISPLAY|6|\the\wd\grmathbox|\the\linewidth}
\ifdim\wd\grmathbox>\linewidth
\resizebox{\linewidth}{!}{\usebox{\grmathbox}}
\else\usebox{\grmathbox}\fi
\endgroup
\end{center}

For a massive object\textquotesingle s timelike trajectory, it
determines the time on a clock carried with that object:

\begin{center}
\begingroup
\sbox{\grmathbox}{$\displaystyle 
d\tau^2=-\frac{ds^2}{c^2}.
$}
\typeout{GRDISPLAY|7|\the\wd\grmathbox|\the\linewidth}
\ifdim\wd\grmathbox>\linewidth
\resizebox{\linewidth}{!}{\usebox{\grmathbox}}
\else\usebox{\grmathbox}\fi
\endgroup
\end{center}

It also determines which directions are null, and therefore the local
light cones. Those cones organize which events can influence which other
events.

So when the metric becomes dynamical, gravity changes more than paths
through an arena. It changes the physical structure used to assign
lengths, times, and causal relations.

Think of a board game whose pieces influence the ruler used to measure
each move and the clocks used to time each turn. The analogy captures
mutual dependence. Its limitation is that real spacetime is not an
elastic tabletop inside a larger room. The metric is an intrinsic field;
an outside embedding is optional mathematical visualization, not
required physics.

\hypertarget{what-the-einstein-equation-actually-connects}{%
\section{1.3 What the Einstein equation actually
connects}\label{what-the-einstein-equation-actually-connects}}

The equation relates two tensor fields at each event.

The right-hand side describes \textbf{energy density, momentum density,
energy flux, and stress}. Matter includes fields: electromagnetic
radiation contributes even though photons have no rest mass. Pressure
contributes because it is momentum transport.

The left-hand side describes a specific contraction of spacetime
curvature, plus the cosmological constant term. It contains the metric
and its derivatives. Once written in a coordinate chart, the equation
becomes coupled nonlinear partial differential equations for the metric
components.

But ``the matter at a point determines the curvature at that point''
needs an important repair. The stress tensor determines the
\textbf{Ricci part} of curvature through the equation. It does not fix
the entire Riemann tensor point by point. Free gravitational degrees of
freedom, encoded in the Weyl part in four dimensions, depend on initial
and boundary information and can propagate through vacuum.

This is why black-hole exteriors and gravitational waves can have
\(T_{\mu\nu}=0\) while remaining gravitationally interesting.

A useful comparison is electromagnetism. Source-free Maxwell equations
allow electromagnetic waves. They do not say ``no charges here,
therefore no electromagnetic field here.'' Vacuum Einstein equations
deserve the same courtesy.

\begin{quote}
\textbf{Second trap door:} For \(\Lambda=0\), vacuum implies
\(R_{\mu\nu}=0\), not \(R^\rho{}_{\sigma\mu\nu}=0\). ``Ricci-flat'' and
``flat'' are different statements.
\end{quote}

\hypertarget{two-questions-that-must-not-be-merged}{%
\section{1.4 Two questions that must not be
merged}\label{two-questions-that-must-not-be-merged}}

There are two separate jobs:

\begin{enumerate}
\def\labelenumi{\arabic{enumi}.}
\tightlist
\item
  Given a spacetime metric, determine how clocks, light, and freely
  falling bodies behave.
\item
  Determine which metric is produced by matter and gravitational initial
  data.
\end{enumerate}

The first is largely geometry plus the coupling of matter to that
geometry. The second is gravitational dynamics.

The geodesic equation addresses the first job for suitable test bodies.
The Einstein equation addresses the second. The Einstein-Hilbert action
packages the second into a variational principle.

There is also a third, easily forgotten job: determine how the matter
itself evolves. A fluid needs an equation of state and fluid equations.
Electromagnetism needs Maxwell\textquotesingle s equations. A scalar
field needs its field equation. The stress tensor is not usually a
freely prescribed movie that can ignore the geometry it lives in.

A self-consistent solution is a joint history of geometry and matter.

This is more like solving an ecosystem than painting scenery behind an
actor. The actor changes the scenery; the scenery changes the
actor\textquotesingle s possible movements; both must obey their
evolution equations.

\hypertarget{why-the-rubber-sheet-is-both-useful-and-dangerous}{%
\section{1.5 Why the rubber sheet is both useful and
dangerous}\label{why-the-rubber-sheet-is-both-useful-and-dangerous}}

The standard image is a heavy ball dimpling a rubber sheet while smaller
balls roll around it.

It can help you imagine a geometry that differs from a plane. Beyond
that, the image acquires several debts:

\begin{itemize}
\tightlist
\item
  The sheet bends into an outside dimension, while intrinsic curvature
  needs no such dimension.
\item
  The balls roll because ordinary gravity pulls them downward, so the
  picture uses gravity to explain gravity.
\item
  The picture shows curved space at an instant. General relativity
  concerns spacetime; changes in clock rates are indispensable.
\item
  A visual dip does not tell you whether the relevant spacetime
  curvature is positive, negative, vacuum, or matter-sourced.
\end{itemize}

Use the sheet for one idea: distances need not obey Euclidean rules.
Then retire it before it starts teaching unauthorized physics.

A better everyday starting point is a network of clocks exchanging light
signals and free-fall laboratories comparing relative acceleration.
Clocks and light are actual measurement procedures. A rubber universe
suspended over a basement is not.

\hypertarget{a-little-history-without-the-lightning-bolt-mythology}{%
\section{1.6 A little history, without the lightning-bolt
mythology}\label{a-little-history-without-the-lightning-bolt-mythology}}

Einstein\textquotesingle s route to general relativity was a prolonged
attempt to reconcile physical constraints. Recovering Newtonian gravity,
making sense of accelerated motion, and preserving an appropriate
energy-momentum balance all mattered. Einstein and Marcel Grossmann
developed a metric approach in their 1913 \emph{Entwurf} theory, but the
final field equations came only after substantial revision. Historical
work on Einstein\textquotesingle s notebooks shows that candidate
equations close to the successful theory had already appeared in his
earlier calculations.
\href{https://www.mpiwg-berlin.mpg.de/Preprints/P264.PDF}{Janssen and
Renn, \emph{Untying the Knot}}.

This is an encouraging story for a learner. Failure to understand an
equation\textquotesingle s physical interpretation can be a deeper
obstacle than failure to write down the equation. A symbol can be
correct while the story you attach to it is wrong.

Einstein presented the final field equations on 25 November 1915, after
several November communications refining his theory. The recognizable
modern action formulation is associated with Einstein and
Hilbert\textquotesingle s work in this period. For the documented
development and the distinction between a compact retrospective
explanation and the actual path of discovery, see the historical
analysis based on Einstein\textquotesingle s calculations and
correspondence.
\href{https://online.ucpress.edu/hsns/article/14/2/253/47626/How-Einstein-Found-His-Field-Equations-1912-1915}{Norton,
\emph{How Einstein Found His Field Equations: 1912-1915}}.

The surrounding mathematics was a collective inheritance: non-Euclidean
and intrinsic geometry, tensor calculus, curvature, and eventually a
clearer language of connections and parallel transport. Learning GR does
not require reenacting the order in which these tools were historically
discovered. We can use the completed toolkit and explain what problem
each tool solves.

That is our modern approach: begin with operational measurements and
geometric objects, distinguish coordinate freedom from physical freedom,
derive dynamics from an action while stating its assumptions, and
understand GR as both a classical theory and the leading part of a
low-energy description.

\hypertarget{what-understanding-the-equation-will-eventually-mean}{%
\section{1.7 What ``understanding the equation'' will eventually
mean}\label{what-understanding-the-equation-will-eventually-mean}}

When you first see \(R_{\mu\nu}-\tfrac12 Rg_{\mu\nu}\), it may look as
though someone took an already difficult object and subtracted a second
difficult object to ensure job security.

By chapter 14, the two pieces will have distinct origins:

\begin{itemize}
\tightlist
\item
  Varying the curvature scalar produces a Ricci-tensor contribution,
  plus a boundary term.
\item
  Varying the spacetime volume measure produces the
  \(-\tfrac12 Rg_{\mu\nu}\) contribution.
\end{itemize}

The geometry\textquotesingle s own differential identity then ensures
that this combination has vanishing covariant divergence.
Matter\textquotesingle s local energy-momentum balance fits that
identity.

The coefficient \(8\pi G_N/c^4\) will not be decorative. Matching the
weak-field limit to Newton\textquotesingle s Poisson equation fixes it.
The factor of two comes from trace reversal; the \(4\pi\) comes from the
familiar three-dimensional inverse-square-field normalization; the
powers of \(c\) reconcile relativistic energy density with curvature.

And \(\Lambda\) will not be merely a late footnote about cosmology. A
constant scalar term is allowed in the action, and its variation
produces exactly \(\Lambda g_{\mu\nu}\).

The equation will become a compressed record of several ideas you have
actually earned.

\hypertarget{a-promise-and-a-discipline}{%
\section{1.8 A promise and a
discipline}\label{a-promise-and-a-discipline}}

We will repeatedly ask four questions:

\textbf{What object is this?} A scalar, vector, covector, tensor,
connection coefficient, density, or coordinate choice?

\textbf{What comparison does it perform?} Between directions at one
event, fields at neighboring events, nearby free-fall trajectories, or
entire spacetime histories?

\textbf{What would an observer measure?} Proper time, acceleration,
frequency, energy, relative displacement, or a coordinate-dependent
intermediate quantity?

\textbf{Under which assumptions is the statement true?} Vacuum or
matter? Local or global? Weak field or exact? Classical or
semiclassical? A point particle or an extended spinning body?

Those questions are the real prerequisites. The calculus will follow
them.

\hypertarget{chapter-2}{%
\chapter{2. The mathematical survival kit: objects, components, and the
art of changing your mind without changing the
universe}\label{chapter-2}}

\hypertarget{a-vector-is-not-its-spreadsheet}{%
\section{2.1 A vector is not its
spreadsheet}\label{a-vector-is-not-its-spreadsheet}}

Imagine a drone receiving the instruction ``move three meters east and
four meters north.'' Rotate the map on its screen. The
command\textquotesingle s two displayed numbers change, but the intended
displacement does not. If the drone changes its destination because you
rotated a map, you have found a software bug, not a new law of
mechanics.

General relativity elevates this ordinary distinction into an organizing
principle: \textbf{physical objects and relationships must survive
changes in the labels used to describe them.} Coordinates are an
interface. They are not the world.

In a vector space, choose a basis \(\{e_1,e_2\}\). A vector \(V\) is
represented by

\begin{center}
\begingroup
\sbox{\grmathbox}{$\displaystyle 
V=V^1e_1+V^2e_2=V^ie_i.
$}
\typeout{GRDISPLAY|8|\the\wd\grmathbox|\the\linewidth}
\ifdim\wd\grmathbox>\linewidth
\resizebox{\linewidth}{!}{\usebox{\grmathbox}}
\else\usebox{\grmathbox}\fi
\endgroup
\end{center}

The \(V^i\) are components. The \(e_i\) are basis vectors. Together they
specify the vector. Either set by itself is insufficient: ``three'' is
not a displacement until we know three of what, in which direction.

The final expression uses \textbf{Einstein summation}: an index
occurring once upstairs and once downstairs is summed. Here the range is
the two directions of our example. In spacetime, Greek indices run over
\(0,1,2,3\); spatial Latin indices normally run over \(1,2,3\).

Suppose a new basis uses twice as long a first basis vector:

\begin{center}
\begingroup
\sbox{\grmathbox}{$\displaystyle 
e'_1=2e_1,\qquad e'_2=e_2.
$}
\typeout{GRDISPLAY|9|\the\wd\grmathbox|\the\linewidth}
\ifdim\wd\grmathbox>\linewidth
\resizebox{\linewidth}{!}{\usebox{\grmathbox}}
\else\usebox{\grmathbox}\fi
\endgroup
\end{center}

Then

\begin{center}
\begingroup
\sbox{\grmathbox}{$\displaystyle 
V=3e_1+4e_2=\frac32e'_1+4e'_2.
$}
\typeout{GRDISPLAY|10|\the\wd\grmathbox|\the\linewidth}
\ifdim\wd\grmathbox>\linewidth
\resizebox{\linewidth}{!}{\usebox{\grmathbox}}
\else\usebox{\grmathbox}\fi
\endgroup
\end{center}

The first component halves because the measuring stick doubled.
Components and basis changes compensate. This is the small mechanical
fact underneath the intimidating word \emph{contravariant}.

There are also two different operations that diagrams sometimes blur. A
\textbf{passive} transformation changes the description of the same
vector. An \textbf{active} transformation changes the vector while
keeping the descriptive machinery fixed. Rotating the map and rotating
the drone\textquotesingle s actual flight direction are different
experiments, even if the corresponding matrix calculations resemble each
other.

\hypertarget{linear-maps-the-machine-behind-a-matrix}{%
\section{2.2 Linear maps: the machine behind a
matrix}\label{linear-maps-the-machine-behind-a-matrix}}

A map \(A\) is linear when

\begin{center}
\begingroup
\sbox{\grmathbox}{$\displaystyle 
A(aV+bW)=aA(V)+bA(W).
$}
\typeout{GRDISPLAY|11|\the\wd\grmathbox|\the\linewidth}
\ifdim\wd\grmathbox>\linewidth
\resizebox{\linewidth}{!}{\usebox{\grmathbox}}
\else\usebox{\grmathbox}\fi
\endgroup
\end{center}

It respects addition and scaling. A spring\textquotesingle s force law
in its linear regime, a small deformation, and a rotation are familiar
examples. In a chosen basis,

\begin{center}
\begingroup
\sbox{\grmathbox}{$\displaystyle 
(AV)^i=A^i{}_jV^j.
$}
\typeout{GRDISPLAY|12|\the\wd\grmathbox|\the\linewidth}
\ifdim\wd\grmathbox>\linewidth
\resizebox{\linewidth}{!}{\usebox{\grmathbox}}
\else\usebox{\grmathbox}\fi
\endgroup
\end{center}

The lower \(j\) consumes the input component; the upper \(i\) identifies
the output component. A matrix is the table of numbers representing this
map in specified input and output bases.

The identity map has components

\begin{center}
\begingroup
\sbox{\grmathbox}{$\displaystyle 
\delta^i{}_j=
\begin{cases}
1&i=j,\\
0&i\ne j.
\end{cases}
$}
\typeout{GRDISPLAY|13|\the\wd\grmathbox|\the\linewidth}
\ifdim\wd\grmathbox>\linewidth
\resizebox{\linewidth}{!}{\usebox{\grmathbox}}
\else\usebox{\grmathbox}\fi
\endgroup
\end{center}

Thus \(\delta^i{}_jV^j=V^i\). This is not merely notation for ``make two
letters the same.'' It is the component representation of a geometric
operation: do nothing.

An important distinction will matter later. A two-index object can
represent a linear map, a bilinear measurement, or something else
depending on the positions and meanings of its indices. ``It is a
\(4\times4\) matrix'' does not specify which kind of geometric object it
represents. A seating chart and a multiplication table can have the same
dimensions without having interchangeable jobs.

\hypertarget{covectors-questions-you-can-ask-a-vector}{%
\section{2.3 Covectors: questions you can ask a
vector}\label{covectors-questions-you-can-ask-a-vector}}

A \textbf{covector} is a linear function that takes a vector and returns
a number. Denote one by \(\omega\):

\begin{center}
\begingroup
\sbox{\grmathbox}{$\displaystyle 
\omega(V)\in\mathbb R.
$}
\typeout{GRDISPLAY|14|\the\wd\grmathbox|\the\linewidth}
\ifdim\wd\grmathbox>\linewidth
\resizebox{\linewidth}{!}{\usebox{\grmathbox}}
\else\usebox{\grmathbox}\fi
\endgroup
\end{center}

If a vector is a possible infinitesimal move, a covector can be a
question such as ``how much does temperature change under that move?''
Its output depends on the move, and linearity means that doubling a
sufficiently small move doubles the first-order response.

For every basis \(e_i\), there is a \textbf{dual basis} \(\theta^i\)
defined by

\begin{center}
\begingroup
\sbox{\grmathbox}{$\displaystyle 
\theta^i(e_j)=\delta^i{}_j.
$}
\typeout{GRDISPLAY|15|\the\wd\grmathbox|\the\linewidth}
\ifdim\wd\grmathbox>\linewidth
\resizebox{\linewidth}{!}{\usebox{\grmathbox}}
\else\usebox{\grmathbox}\fi
\endgroup
\end{center}

The covector \(\theta^1\) extracts the first component; \(\theta^2\)
extracts the second. Write

\begin{center}
\begingroup
\sbox{\grmathbox}{$\displaystyle 
\omega=\omega_i\theta^i.
$}
\typeout{GRDISPLAY|16|\the\wd\grmathbox|\the\linewidth}
\ifdim\wd\grmathbox>\linewidth
\resizebox{\linewidth}{!}{\usebox{\grmathbox}}
\else\usebox{\grmathbox}\fi
\endgroup
\end{center}

Then

\begin{center}
\begingroup
\sbox{\grmathbox}{$\displaystyle 
\omega(V)=\omega_iV^i.
$}
\typeout{GRDISPLAY|17|\the\wd\grmathbox|\the\linewidth}
\ifdim\wd\grmathbox>\linewidth
\resizebox{\linewidth}{!}{\usebox{\grmathbox}}
\else\usebox{\grmathbox}\fi
\endgroup
\end{center}

This pairing needs no metric, no angles, and no idea of distance. It is
built into the relationship between a vector space and its dual.

Return to \(e'_1=2e_1\). The new extractor must be
\(\theta'^1=\theta^1/2\) so that \(\theta'^1(e'_1)=1\). Consequently
\(\omega'_1=2\omega_1\). A vector\textquotesingle s first component
halves; a covector\textquotesingle s first component doubles; their
product stays the same:

\begin{center}
\begingroup
\sbox{\grmathbox}{$\displaystyle 
\omega'_1V'^1=(2\omega_1)(V^1/2)=\omega_1V^1.
$}
\typeout{GRDISPLAY|18|\the\wd\grmathbox|\the\linewidth}
\ifdim\wd\grmathbox>\linewidth
\resizebox{\linewidth}{!}{\usebox{\grmathbox}}
\else\usebox{\grmathbox}\fi
\endgroup
\end{center}

This is why covectors carry lower indices. Their components change
oppositely to vector components. The difference is operational, not
typographical.

\begin{quote}
\textbf{Gotcha: a gradient begins life as a covector.} The differential
\(df\) tells you the directional change of a scalar field. Turning it
into an arrow called \(\operatorname{grad}f\) requires a metric. In
ordinary Cartesian calculus, the Euclidean metric makes this conversion
numerically invisible, which is how it gets away with hiding for years.
\end{quote}

\hypertarget{the-chain-rule-already-knows-tensor-calculus}{%
\section{2.4 The chain rule already knows tensor
calculus}\label{the-chain-rule-already-knows-tensor-calculus}}

Let \(x^\mu\) and \(x'^\alpha\) label the same events. An infinitesimal
displacement transforms by the chain rule:

\begin{center}
\begingroup
\sbox{\grmathbox}{$\displaystyle 
dx'^\alpha=\frac{\partial x'^\alpha}{\partial x^\mu}\,dx^\mu.
$}
\typeout{GRDISPLAY|19|\the\wd\grmathbox|\the\linewidth}
\ifdim\wd\grmathbox>\linewidth
\resizebox{\linewidth}{!}{\usebox{\grmathbox}}
\else\usebox{\grmathbox}\fi
\endgroup
\end{center}

Define the Jacobian and its inverse,

\begin{center}
\begingroup
\sbox{\grmathbox}{$\displaystyle 
J^\alpha{}_{\mu}=\frac{\partial x'^\alpha}{\partial x^\mu},
\qquad
K^\mu{}_{\alpha}=\frac{\partial x^\mu}{\partial x'^\alpha}.
$}
\typeout{GRDISPLAY|20|\the\wd\grmathbox|\the\linewidth}
\ifdim\wd\grmathbox>\linewidth
\resizebox{\linewidth}{!}{\usebox{\grmathbox}}
\else\usebox{\grmathbox}\fi
\endgroup
\end{center}

Their product satisfies
\(K^\mu{}_{\alpha}J^\alpha{}_{\nu}=\delta^\mu{}_{\nu}\). This follows by
differentiating the identity obtained by composing a coordinate change
with its inverse.

Vectors transform like tangent displacements:

\begin{center}
\begingroup
\sbox{\grmathbox}{$\displaystyle 
V'^\alpha=J^\alpha{}_{\mu}V^\mu.
$}
\typeout{GRDISPLAY|21|\the\wd\grmathbox|\the\linewidth}
\ifdim\wd\grmathbox>\linewidth
\resizebox{\linewidth}{!}{\usebox{\grmathbox}}
\else\usebox{\grmathbox}\fi
\endgroup
\end{center}

Now differentiate a scalar field \(f\). A scalar has a single value at
an event; its two coordinate representations obey \(f'(x')=f(x)\). The
chain rule gives

\begin{center}
\begingroup
\sbox{\grmathbox}{$\displaystyle 
\partial'_\alpha f'
=\frac{\partial x^\mu}{\partial x'^\alpha}\partial_\mu f
=K^\mu{}_{\alpha}\partial_\mu f,
\qquad
\partial_\mu\equiv\frac{\partial}{\partial x^\mu}.
$}
\typeout{GRDISPLAY|22|\the\wd\grmathbox|\the\linewidth}
\ifdim\wd\grmathbox>\linewidth
\resizebox{\linewidth}{!}{\usebox{\grmathbox}}
\else\usebox{\grmathbox}\fi
\endgroup
\end{center}

That is the covector transformation law. In particular,

\begin{center}
\begingroup
\sbox{\grmathbox}{$\displaystyle 
df=(\partial_\mu f)dx^\mu,
\qquad
df(V)=V^\mu\partial_\mu f
$}
\typeout{GRDISPLAY|23|\the\wd\grmathbox|\the\linewidth}
\ifdim\wd\grmathbox>\linewidth
\resizebox{\linewidth}{!}{\usebox{\grmathbox}}
\else\usebox{\grmathbox}\fi
\endgroup
\end{center}

is independent of the coordinates. The symbol \(dx^\mu\) has two closely
related uses: as a component of an infinitesimal displacement and as the
coordinate covector that extracts a vector\textquotesingle s \(\mu\)
component. The context tells you which use is intended.

The Jacobian can depend on position. This is harmless for a vector at
one event, but it causes trouble when differentiating a vector field: a
derivative can hit the Jacobian as well as the vector components.
Chapter 6 will turn that apparently minor nuisance into the reason we
need a connection.

\hypertarget{tensors-are-multilinear-relationships}{%
\section{2.5 Tensors are multilinear
relationships}\label{tensors-are-multilinear-relationships}}

A bilinear object \(B\) takes two vectors and returns a scalar:

\begin{center}
\begingroup
\sbox{\grmathbox}{$\displaystyle 
B(V,W)=B_{\mu\nu}V^\mu W^\nu.
$}
\typeout{GRDISPLAY|24|\the\wd\grmathbox|\the\linewidth}
\ifdim\wd\grmathbox>\linewidth
\resizebox{\linewidth}{!}{\usebox{\grmathbox}}
\else\usebox{\grmathbox}\fi
\endgroup
\end{center}

It is linear in each input separately. Requiring the output to be
unchanged by relabeling determines how its components transform:

\begin{center}
\begingroup
\sbox{\grmathbox}{$\displaystyle 
B'_{\alpha\beta}
=K^\mu{}_{\alpha}K^\nu{}_{\beta}B_{\mu\nu}.
$}
\typeout{GRDISPLAY|25|\the\wd\grmathbox|\the\linewidth}
\ifdim\wd\grmathbox>\linewidth
\resizebox{\linewidth}{!}{\usebox{\grmathbox}}
\else\usebox{\grmathbox}\fi
\endgroup
\end{center}

There is one inverse Jacobian for each lower index. An upper index gets
a forward Jacobian. For example,

\begin{center}
\begingroup
\sbox{\grmathbox}{$\displaystyle 
A'^\alpha{}_{\beta}
=J^\alpha{}_{\mu}K^\nu{}_{\beta}A^\mu{}_{\nu}.
$}
\typeout{GRDISPLAY|26|\the\wd\grmathbox|\the\linewidth}
\ifdim\wd\grmathbox>\linewidth
\resizebox{\linewidth}{!}{\usebox{\grmathbox}}
\else\usebox{\grmathbox}\fi
\endgroup
\end{center}

A tensor of type \((r,s)\) has \(r\) upper and \(s\) lower indices.
Abstractly, it is a multilinear object with the corresponding vector and
covector slots. The word ``rank'' is often used for \(r+s\), but beware:
\textbf{tensor rank in that sense is not matrix rank}, which counts
independent image directions of a linear map.

One way to build tensors is the tensor product. If \(\omega\) and
\(\eta\) are covectors,

\begin{center}
\begingroup
\sbox{\grmathbox}{$\displaystyle 
(\omega\otimes\eta)(V,W)=\omega(V)\eta(W).
$}
\typeout{GRDISPLAY|27|\the\wd\grmathbox|\the\linewidth}
\ifdim\wd\grmathbox>\linewidth
\resizebox{\linewidth}{!}{\usebox{\grmathbox}}
\else\usebox{\grmathbox}\fi
\endgroup
\end{center}

The two inputs remain separate. This is a richer object than one scalar
\(\omega(V)\). The symbol \(\otimes\) is an instruction to preserve
independent slots rather than multiply everything into a single number
prematurely.

\textbf{Contraction} pairs an upper and a lower index and sums them. A
linear map has the trace

\begin{center}
\begingroup
\sbox{\grmathbox}{$\displaystyle 
A^\mu{}_{\mu}.
$}
\typeout{GRDISPLAY|28|\the\wd\grmathbox|\the\linewidth}
\ifdim\wd\grmathbox>\linewidth
\resizebox{\linewidth}{!}{\usebox{\grmathbox}}
\else\usebox{\grmathbox}\fi
\endgroup
\end{center}

Its value does not depend on the basis. For the identity in four
dimensions, \(\delta^\mu{}_{\mu}=4\). Contracting changes the
tensor\textquotesingle s type: one upper and one lower slot disappear.

The index language supplies excellent error detection:

\begingroup\small\setlength{\tabcolsep}{5pt}\renewcommand{\arraystretch}{1.13}

\begin{longtable}[]{@{}
  >{\raggedright\arraybackslash}p{(\columnwidth - 2\tabcolsep) * \real{0.3000}}
  >{\raggedright\arraybackslash}p{(\columnwidth - 2\tabcolsep) * \real{0.7000}}@{}}
\toprule\noalign{}
\begin{minipage}[b]{\linewidth}\raggedright
Expression
\end{minipage} & \begin{minipage}[b]{\linewidth}\raggedright
Meaning or problem
\end{minipage} \\
\midrule\noalign{}
\endhead
\bottomrule\noalign{}
\endlastfoot
\(V^\mu\omega_\mu\) & Scalar contraction; legal. \\
\(A^\mu{}_{\nu}V^\nu=W^\mu\) & Vector equation; the free index \(\mu\)
agrees. \\
\(A^\mu{}_{\nu}=B^\mu{}_{\rho}\) & Malformed unless another operation is
supplied; the free indices disagree. \\
\(V^\mu W^\mu\) & Not a valid Einstein contraction under our convention;
two upper indices need a metric. \\
\(A^\mu{}_{\nu}B^\nu{}_{\rho}C^\rho{}_{\mu}\) & Scalar formed from a
closed sequence of contractions. \\
\end{longtable}

\endgroup

Dummy indices are replaceable labels:
\(V^\mu\omega_\mu=V^\alpha\omega_\alpha\). Free indices are the outputs
of the expression and must match across an equation. An index should not
occur three times in a single product under ordinary Einstein notation.

Symmetrization and antisymmetrization will be useful:

\begin{center}
\begingroup
\sbox{\grmathbox}{$\displaystyle 
B_{(\mu\nu)}=\frac12(B_{\mu\nu}+B_{\nu\mu}),
\qquad
B_{[\mu\nu]}=\frac12(B_{\mu\nu}-B_{\nu\mu}).
$}
\typeout{GRDISPLAY|29|\the\wd\grmathbox|\the\linewidth}
\ifdim\wd\grmathbox>\linewidth
\resizebox{\linewidth}{!}{\usebox{\grmathbox}}
\else\usebox{\grmathbox}\fi
\endgroup
\end{center}

Their sum reconstructs \(B_{\mu\nu}\). The factor \(1/2\) ensures that
symmetrizing an already symmetric tensor does not double it.

\hypertarget{derivatives-approximations-and-a-brief-glimpse-of-forms}{%
\section{2.6 Derivatives, approximations, and a brief glimpse of
forms}\label{derivatives-approximations-and-a-brief-glimpse-of-forms}}

For a smooth scalar,

\begin{center}
\begingroup
\sbox{\grmathbox}{$\displaystyle 
f(x+\delta x)
=f(x)+\partial_\mu f\,\delta x^\mu
+\frac12\partial_\mu\partial_\nu f\,\delta x^\mu\delta x^\nu
+\cdots.
$}
\typeout{GRDISPLAY|30|\the\wd\grmathbox|\the\linewidth}
\ifdim\wd\grmathbox>\linewidth
\resizebox{\linewidth}{!}{\usebox{\grmathbox}}
\else\usebox{\grmathbox}\fi
\endgroup
\end{center}

The first derivative measures the local slope. The second measures how
that slope changes. Later, first metric derivatives will describe
coordinate-dependent gravitational acceleration, while a particular
combination of second derivatives and products of first derivatives will
describe curvature. The qualification ``a particular combination''
matters: a second derivative alone is not automatically a tensor.

A perturbative statement such as
\(F=F_0+\varepsilon F_1+O(\varepsilon^2)\) says that discarded terms are
at least quadratic in a specified small parameter, in the regime under
discussion. Always ask what is small. A metric perturbation being small
in convenient coordinates does not mean every derivative of it is small;
a low-amplitude wave can oscillate rapidly.

One short preview of differential forms will prevent future alarm. An
antisymmetric covariant tensor is called a differential form. A one-form
is a covector field; a two-form has two antisymmetric slots. The wedge
product satisfies

\begin{center}
\begingroup
\sbox{\grmathbox}{$\displaystyle 
(dx\wedge dy)(V,W)=V^xW^y-V^yW^x.
$}
\typeout{GRDISPLAY|31|\the\wd\grmathbox|\the\linewidth}
\ifdim\wd\grmathbox>\linewidth
\resizebox{\linewidth}{!}{\usebox{\grmathbox}}
\else\usebox{\grmathbox}\fi
\endgroup
\end{center}

This computes oriented coordinate area. Exchange the two input vectors
and its sign flips. Two parallel inputs give zero. Forms are naturally
suited to integrating along curves, over surfaces, and through
higher-dimensional regions. Chapter 21 will use them to express geometry
with surprisingly few symbols.

You do not need to memorize an entire branch of mathematics before
proceeding. Carry three questions: \textbf{What object is this? What
inputs does it consume? How do its components change when I relabel the
situation?} These questions are more valuable than being able to
pronounce ``contravariant'' without hesitation. For a formal development
of the tangent-space and dual-space machinery, see
\href{https://davidtong.org/pdfs/teaching/general-relativity/gr2.pdf}{David
Tong\textquotesingle s differential-geometry chapter}.

\hypertarget{chapter-3}{%
\chapter{3. Special relativity: learning what a clock is actually
measuring}\label{chapter-3}}

\hypertarget{events-not-photographs}{%
\section{3.1 Events, not photographs}\label{events-not-photographs}}

An event is something localized in space and time: a flash, a detector
click, two particles meeting. In an inertial coordinate system, label it

\begin{center}
\begingroup
\sbox{\grmathbox}{$\displaystyle 
x^\mu=(ct,x,y,z).
$}
\typeout{GRDISPLAY|32|\the\wd\grmathbox|\the\linewidth}
\ifdim\wd\grmathbox>\linewidth
\resizebox{\linewidth}{!}{\usebox{\grmathbox}}
\else\usebox{\grmathbox}\fi
\endgroup
\end{center}

Writing \(x^0=ct\) makes all four coordinates have length units in this
chart. It does not turn time into ordinary space. The difference lives
in the metric\textquotesingle s sign.

An inertial frame is an ideal network of mutually stationary rulers and
synchronized clocks in flat spacetime, with no acceleration or rotation
of the network. ``The time of a distant event'' is the reading assigned
by this network, not the time at which light from that event reaches
your eye. Relativity of simultaneity remains after light-travel delays
have been accounted for.

Einstein\textquotesingle s 1905 construction combined the equivalence of
inertial frames for the laws of physics with the invariant vacuum speed
of light. The original paper makes the operational treatment of clocks
central to the theory. It is useful to meet that emphasis directly in
\href{https://sites.pitt.edu/~jdnorton/teaching/Einstein_graduate/pdfs/Einstein_STR_1905_English.pdf}{Einstein\textquotesingle s
1905 paper, in English translation}.

\hypertarget{deriving-a-lorentz-boost-without-pulling-a-rabbit-from-a-matrix}{%
\section{3.2 Deriving a Lorentz boost without pulling a rabbit from a
matrix}\label{deriving-a-lorentz-boost-without-pulling-a-rabbit-from-a-matrix}}

Let frame \(S'\) move at speed \(v\) in the positive \(x\) direction
relative to \(S\). Their origins coincide at \(t=t'=0\). We restrict
attention to \(t,x\) and assume spatial homogeneity, time homogeneity,
standard clock synchronization, and matching units. These assumptions
make the transformation between inertial coordinates linear.

The moving origin obeys \(x=vt\) and must have \(x'=0\), so

\begin{center}
\begingroup
\sbox{\grmathbox}{$\displaystyle 
x'=A(x-vt)
$}
\typeout{GRDISPLAY|33|\the\wd\grmathbox|\the\linewidth}
\ifdim\wd\grmathbox>\linewidth
\resizebox{\linewidth}{!}{\usebox{\grmathbox}}
\else\usebox{\grmathbox}\fi
\endgroup
\end{center}

for some factor \(A\) depending on \(v\). Write the most general linear
time transformation as \(t'=B t+C x\).

Now send light to the right, so \(x=ct\) and \(x'=ct'\). Substitution
gives

\begin{center}
\begingroup
\sbox{\grmathbox}{$\displaystyle 
A(c-v)=c(B+Cc).
$}
\typeout{GRDISPLAY|34|\the\wd\grmathbox|\the\linewidth}
\ifdim\wd\grmathbox>\linewidth
\resizebox{\linewidth}{!}{\usebox{\grmathbox}}
\else\usebox{\grmathbox}\fi
\endgroup
\end{center}

Send light to the left, so \(x=-ct\) and \(x'=-ct'\). This gives

\begin{center}
\begingroup
\sbox{\grmathbox}{$\displaystyle 
A(c+v)=c(B-Cc).
$}
\typeout{GRDISPLAY|35|\the\wd\grmathbox|\the\linewidth}
\ifdim\wd\grmathbox>\linewidth
\resizebox{\linewidth}{!}{\usebox{\grmathbox}}
\else\usebox{\grmathbox}\fi
\endgroup
\end{center}

Add the equations: \(2Ac=2cB\), hence \(B=A\). Subtract them:
\(-2Av=2c^2C\), hence \(C=-Av/c^2\). Therefore

\begin{center}
\begingroup
\sbox{\grmathbox}{$\displaystyle 
x'=A(x-vt),
\qquad
t'=A\left(t-\frac{vx}{c^2}\right).
$}
\typeout{GRDISPLAY|36|\the\wd\grmathbox|\the\linewidth}
\ifdim\wd\grmathbox>\linewidth
\resizebox{\linewidth}{!}{\usebox{\grmathbox}}
\else\usebox{\grmathbox}\fi
\endgroup
\end{center}

We have not guessed that time must mix with space. The two light
directions forced it.

Reciprocity and isotropy imply that the inverse transformation has the
same factor with \(v\) replaced by \(-v\). Substituting the
transformations into their inverse must return \(x\), which yields

\begin{center}
\begingroup
\sbox{\grmathbox}{$\displaystyle 
A^2\left(1-\frac{v^2}{c^2}\right)=1.
$}
\typeout{GRDISPLAY|37|\the\wd\grmathbox|\the\linewidth}
\ifdim\wd\grmathbox>\linewidth
\resizebox{\linewidth}{!}{\usebox{\grmathbox}}
\else\usebox{\grmathbox}\fi
\endgroup
\end{center}

Choose the positive root continuously connected to \(A=1\) at \(v=0\):

\begin{center}
\begingroup
\sbox{\grmathbox}{$\displaystyle 
\boxed{\gamma=\frac1{\sqrt{1-v^2/c^2}},\qquad
x'=\gamma(x-vt),\qquad
t'=\gamma\left(t-\frac{vx}{c^2}\right).}
$}
\typeout{GRDISPLAY|38|\the\wd\grmathbox|\the\linewidth}
\ifdim\wd\grmathbox>\linewidth
\resizebox{\linewidth}{!}{\usebox{\grmathbox}}
\else\usebox{\grmathbox}\fi
\endgroup
\end{center}

For this standard boost, \(y'=y\) and \(z'=z\). The \(vx/c^2\) term is
the great conceptual disruption. Two events with \(\Delta t=0\) but
different \(x\) generally have

\begin{center}
\begingroup
\sbox{\grmathbox}{$\displaystyle 
\Delta t'=-\gamma\frac{v\Delta x}{c^2}\ne0.
$}
\typeout{GRDISPLAY|39|\the\wd\grmathbox|\the\linewidth}
\ifdim\wd\grmathbox>\linewidth
\resizebox{\linewidth}{!}{\usebox{\grmathbox}}
\else\usebox{\grmathbox}\fi
\endgroup
\end{center}

An observer\textquotesingle s ``now'' is a slice through spacetime
selected by their state of motion. It is not a layer of cosmic frosting
already spread across the universe.

The everyday limit is sensible. If \(v/c\ll1\), then \(\gamma\approx1\)
and \(vx/c^2\) becomes negligible for ordinary distances and timing
precision. We recover the Galilean approximation \(x'\approx x-vt\),
\(t'\approx t\).

\hypertarget{the-interval-the-quantity-that-refuses-to-change}{%
\section{3.3 The interval: the quantity that refuses to
change}\label{the-interval-the-quantity-that-refuses-to-change}}

Insert the Lorentz formulas into \(-c^2dt'^2+dx'^2\). Expanding both
squares produces cross terms \(+2v\,dt\,dx\) and \(-2v\,dt\,dx\), which
cancel. The remaining factor \(\gamma^2(1-v^2/c^2)\) equals one. Thus

\begin{center}
\begingroup
\sbox{\grmathbox}{$\displaystyle 
ds^2=-c^2dt^2+dx^2+dy^2+dz^2
=\eta_{\mu\nu}dx^\mu dx^\nu,
$}
\typeout{GRDISPLAY|40|\the\wd\grmathbox|\the\linewidth}
\ifdim\wd\grmathbox>\linewidth
\resizebox{\linewidth}{!}{\usebox{\grmathbox}}
\else\usebox{\grmathbox}\fi
\endgroup
\end{center}

where

\begin{center}
\begingroup
\sbox{\grmathbox}{$\displaystyle 
\eta_{\mu\nu}=\operatorname{diag}(-1,1,1,1).
$}
\typeout{GRDISPLAY|41|\the\wd\grmathbox|\the\linewidth}
\ifdim\wd\grmathbox>\linewidth
\resizebox{\linewidth}{!}{\usebox{\grmathbox}}
\else\usebox{\grmathbox}\fi
\endgroup
\end{center}

All inertial frames assign the same interval. They disagree about its
division into temporal and spatial components.

A useful analogy is an ordinary rotation: different observers assign
different horizontal and vertical components to a rod while agreeing on
its length. A Lorentz boost similarly changes temporal and spatial
components while preserving a quadratic form. \textbf{The limitation is
crucial:} the spacetime quadratic form has a minus sign. It is not
Euclidean distance in disguise. Nonzero vectors can have zero norm, and
the geometry distinguishes three causal types.

\begingroup\small\setlength{\tabcolsep}{5pt}\renewcommand{\arraystretch}{1.13}

\begin{longtable}[]{@{}
  >{\raggedright\arraybackslash}p{(\columnwidth - 4\tabcolsep) * \real{0.2500}}
  >{\raggedright\arraybackslash}p{(\columnwidth - 4\tabcolsep) * \real{0.3600}}
  >{\raggedright\arraybackslash}p{(\columnwidth - 4\tabcolsep) * \real{0.3900}}@{}}
\toprule\noalign{}
\begin{minipage}[b]{\linewidth}\raggedright
Separation or tangent
\end{minipage} & \begin{minipage}[b]{\linewidth}\raggedright
Interval sign
\end{minipage} & \begin{minipage}[b]{\linewidth}\raggedright
Physical significance
\end{minipage} \\
\midrule\noalign{}
\endhead
\bottomrule\noalign{}
\endlastfoot
Timelike & \(ds^2<0\) & A sufficiently small displacement can lie on a
massive observer\textquotesingle s worldline. \\
Null & \(ds^2=0\) & A nonzero displacement tangent to a vacuum light ray
in geometric optics. \\
Spacelike & \(ds^2>0\) & No causal signal can connect sufficiently
nearby events with that displacement. \\
\end{longtable}

\endgroup

In Minkowski space these classifications also apply directly to finite
differences between two events. In curved spacetime the metric initially
classifies tangent directions; global causal relationships require
examining actual curves.

The null directions form a \textbf{light cone} at each event. The future
cone contains directions in which physical observers and signals can
proceed. Proper Lorentz transformations preserving time orientation do
not turn a future timelike direction into a past direction. The temporal
order of spacelike-separated events can change between inertial frames;
the causal order of connected events cannot.

Rapidity makes the rotation analogy mathematically exact in its
appropriate sense. Define \(\chi\) by

\begin{center}
\begingroup
\sbox{\grmathbox}{$\displaystyle 
\tanh\chi=\frac vc,\qquad
\cosh\chi=\gamma,\qquad
\sinh\chi=\gamma\frac vc.
$}
\typeout{GRDISPLAY|42|\the\wd\grmathbox|\the\linewidth}
\ifdim\wd\grmathbox>\linewidth
\resizebox{\linewidth}{!}{\usebox{\grmathbox}}
\else\usebox{\grmathbox}\fi
\endgroup
\end{center}

Then

\begin{center}
\begingroup
\sbox{\grmathbox}{$\displaystyle 
\begin{pmatrix}ct'\\x'\end{pmatrix}
=
\begin{pmatrix}\cosh\chi&-\sinh\chi\\-\sinh\chi&\cosh\chi\end{pmatrix}
\begin{pmatrix}ct\\x\end{pmatrix}.
$}
\typeout{GRDISPLAY|43|\the\wd\grmathbox|\the\linewidth}
\ifdim\wd\grmathbox>\linewidth
\resizebox{\linewidth}{!}{\usebox{\grmathbox}}
\else\usebox{\grmathbox}\fi
\endgroup
\end{center}

Collinear boost rapidities add. Velocities consequently combine as

\begin{center}
\begingroup
\sbox{\grmathbox}{$\displaystyle 
v_{\mathrm{combined}}=\frac{v_1+v_2}{1+v_1v_2/c^2}.
$}
\typeout{GRDISPLAY|44|\the\wd\grmathbox|\the\linewidth}
\ifdim\wd\grmathbox>\linewidth
\resizebox{\linewidth}{!}{\usebox{\grmathbox}}
\else\usebox{\grmathbox}\fi
\endgroup
\end{center}

Two successive boosts of \(0.8c\) give \(1.6c/1.64\approx0.976c\), not
\(1.6c\). The speed limit is encoded in the geometry of composition.

\hypertarget{proper-time-your-life-is-a-line-integral}{%
\section{3.4 Proper time: your life is a line
integral}\label{proper-time-your-life-is-a-line-integral}}

Along a timelike worldline, define

\begin{center}
\begingroup
\sbox{\grmathbox}{$\displaystyle 
d\tau=\frac{\sqrt{-ds^2}}c
=dt\sqrt{1-\frac{\mathbf v^2}{c^2}}.
$}
\typeout{GRDISPLAY|45|\the\wd\grmathbox|\the\linewidth}
\ifdim\wd\grmathbox>\linewidth
\resizebox{\linewidth}{!}{\usebox{\grmathbox}}
\else\usebox{\grmathbox}\fi
\endgroup
\end{center}

The proper time \(\tau\) is what an ideal clock carried along that
worldline measures. ``Ideal'' means the clock measures this geometric
quantity without appreciable changes from its construction,
acceleration-induced damage, temperature, or other environmental
effects.

For constant speed,

\begin{center}
\begingroup
\sbox{\grmathbox}{$\displaystyle 
\Delta\tau=\frac{\Delta t}{\gamma}.
$}
\typeout{GRDISPLAY|46|\the\wd\grmathbox|\the\linewidth}
\ifdim\wd\grmathbox>\linewidth
\resizebox{\linewidth}{!}{\usebox{\grmathbox}}
\else\usebox{\grmathbox}\fi
\endgroup
\end{center}

At \(v=0.8c\), \(\gamma=5/3\). Five years of inertial-frame coordinate
time correspond to three years on the moving clock. Both numbers are
legitimate measurements of specified quantities. There is no
contradiction to be resolved by identifying which clock is secretly
defective.

For variable speed,

\begin{center}
\begingroup
\sbox{\grmathbox}{$\displaystyle 
\tau=\int_{t_A}^{t_B}\sqrt{1-\frac{\mathbf v(t)^2}{c^2}}\,dt.
$}
\typeout{GRDISPLAY|47|\the\wd\grmathbox|\the\linewidth}
\ifdim\wd\grmathbox>\linewidth
\resizebox{\linewidth}{!}{\usebox{\grmathbox}}
\else\usebox{\grmathbox}\fi
\endgroup
\end{center}

This is a functional of the path. Two clocks can start together, follow
different worldlines, reunite, and display different accumulated times.
Their disagreement at reunion is an invariant local comparison.

In flat spacetime, choose the inertial frame in which two
timelike-separated reunion events occur at the same spatial position.
The clock at rest between them accumulates \(\Delta t\). Every other
future-directed timelike path has an integrand at most one, so it
accumulates no more. The inertial path maximizes proper time between
those events.

That last sentence often causes a double take. A straight Euclidean path
minimizes length. A straight timelike Minkowski path maximizes elapsed
time. The minus sign is not decorative.

\hypertarget{the-twins-acceleration-explains-the-asymmetry-geometry-computes-the-age}{%
\section{3.5 The twins: acceleration explains the asymmetry, geometry
computes the
age}\label{the-twins-acceleration-explains-the-asymmetry-geometry-computes-the-age}}

One twin remains inertial. The other travels outward at \(0.8c\) for
five years of the first twin\textquotesingle s time and returns at the
same speed for another five years, idealizing the turnaround as brief.
At reunion,

\begin{center}
\begingroup
\sbox{\grmathbox}{$\displaystyle 
\tau_{\mathrm{home}}=10\ \text{years},
\qquad
\tau_{\mathrm{traveler}}=2\times5\sqrt{1-0.8^2}=6\ \text{years}.
$}
\typeout{GRDISPLAY|48|\the\wd\grmathbox|\the\linewidth}
\ifdim\wd\grmathbox>\linewidth
\resizebox{\linewidth}{!}{\usebox{\grmathbox}}
\else\usebox{\grmathbox}\fi
\endgroup
\end{center}

The traveler does not remain in one inertial frame for the whole
experiment, so the histories are not symmetric. But ``acceleration makes
clocks run slowly'' is an unreliable explanation. In the ideal clock
formula, elapsed time is obtained from the worldline\textquotesingle s
tangent, not from an extra direct acceleration term.

The turnaround makes it possible for the traveler to switch between
outward and inward inertial segments and reunite. One can make that
turnaround brief without making the age difference disappear.
Conversely, two differently accelerated clocks can be arranged to
accumulate equal proper time. \textbf{Acceleration identifies an
asymmetry in this experiment; the proper-time integral gives the
answer.}

Special relativity can handle accelerated observers perfectly well. What
it lacks is general curved spacetime dynamics. ``Acceleration requires
general relativity'' confuses a choice of observer with a property of
spacetime.

\hypertarget{four-velocity-and-four-momentum}{%
\section{3.6 Four-velocity and
four-momentum}\label{four-velocity-and-four-momentum}}

For a massive particle,

\begin{center}
\begingroup
\sbox{\grmathbox}{$\displaystyle 
u^\mu\equiv\frac{dx^\mu}{d\tau}
=\gamma(c,\mathbf v).
$}
\typeout{GRDISPLAY|49|\the\wd\grmathbox|\the\linewidth}
\ifdim\wd\grmathbox>\linewidth
\resizebox{\linewidth}{!}{\usebox{\grmathbox}}
\else\usebox{\grmathbox}\fi
\endgroup
\end{center}

Its norm follows immediately:

\begin{center}
\begingroup
\sbox{\grmathbox}{$\displaystyle 
\eta_{\mu\nu}u^\mu u^\nu
=\gamma^2(-c^2+\mathbf v^2)=-c^2.
$}
\typeout{GRDISPLAY|50|\the\wd\grmathbox|\the\linewidth}
\ifdim\wd\grmathbox>\linewidth
\resizebox{\linewidth}{!}{\usebox{\grmathbox}}
\else\usebox{\grmathbox}\fi
\endgroup
\end{center}

The spatial velocity can vary; the four-velocity\textquotesingle s norm
cannot. With constant rest mass \(m\), define four-momentum

\begin{center}
\begingroup
\sbox{\grmathbox}{$\displaystyle 
p^\mu=mu^\mu=\left(\frac Ec,\mathbf p\right),
\qquad
E=\gamma mc^2,
\qquad
\mathbf p=\gamma m\mathbf v.
$}
\typeout{GRDISPLAY|51|\the\wd\grmathbox|\the\linewidth}
\ifdim\wd\grmathbox>\linewidth
\resizebox{\linewidth}{!}{\usebox{\grmathbox}}
\else\usebox{\grmathbox}\fi
\endgroup
\end{center}

Its invariant norm is

\begin{center}
\begingroup
\sbox{\grmathbox}{$\displaystyle 
p_\mu p^\mu=-m^2c^2,
$}
\typeout{GRDISPLAY|52|\the\wd\grmathbox|\the\linewidth}
\ifdim\wd\grmathbox>\linewidth
\resizebox{\linewidth}{!}{\usebox{\grmathbox}}
\else\usebox{\grmathbox}\fi
\endgroup
\end{center}

or, after expanding,

\begin{center}
\begingroup
\sbox{\grmathbox}{$\displaystyle 
\boxed{E^2=\mathbf p^2c^2+m^2c^4.}
$}
\typeout{GRDISPLAY|53|\the\wd\grmathbox|\the\linewidth}
\ifdim\wd\grmathbox>\linewidth
\resizebox{\linewidth}{!}{\usebox{\grmathbox}}
\else\usebox{\grmathbox}\fi
\endgroup
\end{center}

The familiar \(E=mc^2\) is the special case of zero spatial momentum in
the particle\textquotesingle s rest frame. At small speed, the binomial
expansion \(\gamma=1+\frac12v^2/c^2+O(v^4/c^4)\) gives

\begin{center}
\begingroup
\sbox{\grmathbox}{$\displaystyle 
E=mc^2+\frac12mv^2+O(mv^4/c^2).
$}
\typeout{GRDISPLAY|54|\the\wd\grmathbox|\the\linewidth}
\ifdim\wd\grmathbox>\linewidth
\resizebox{\linewidth}{!}{\usebox{\grmathbox}}
\else\usebox{\grmathbox}\fi
\endgroup
\end{center}

Newtonian kinetic energy appears as the first correction to rest energy.

For a photon, \(m=0\) and \(E=c|\mathbf p|\). A photon has no rest frame
and no proper-time four-velocity. Its four-momentum still exists and is
null. Dividing a photon\textquotesingle s displacement by its proper
time would divide by zero; that is not a profound alternative definition
of motion.

\begin{quote}
\textbf{Gotcha: system mass is not the sum of constituent rest masses.}
Consider two photons, each of energy \(E_\gamma\), traveling in opposite
directions. Total momentum is zero and total energy is \(2E_\gamma\), so
the system has invariant mass \(M=2E_\gamma/c^2\). Every constituent is
massless. The system is not. Internal motion, radiation, and binding
energy all matter when accounting for mass-energy.
\end{quote}

\hypertarget{energy-is-a-measurement-made-by-an-observer}{%
\section{3.7 Energy is a measurement made by an
observer}\label{energy-is-a-measurement-made-by-an-observer}}

Let an observer have four-velocity \(U^\mu\), with \(U_\mu U^\mu=-c^2\).
The energy that observer measures for a particle with four-momentum
\(p^\mu\) is

\begin{center}
\begingroup
\sbox{\grmathbox}{$\displaystyle 
\boxed{E_{(U)}=-p_\mu U^\mu.}
$}
\typeout{GRDISPLAY|55|\the\wd\grmathbox|\the\linewidth}
\ifdim\wd\grmathbox>\linewidth
\resizebox{\linewidth}{!}{\usebox{\grmathbox}}
\else\usebox{\grmathbox}\fi
\endgroup
\end{center}

Check the units and sign in the observer\textquotesingle s rest frame:
\(U^\mu=(c,0,0,0)\) and \(p_0=-E/c\), so \(-p_\mu U^\mu=E\). The formula
is not an assertion that energy is invariant between observers. It says
the energy measured by a \textbf{specified} observer is a scalar under
coordinate changes. Change the observer \(U\), and the
scalar\textquotesingle s value can change.

For a massive particle, define the relative Lorentz factor

\begin{center}
\begingroup
\sbox{\grmathbox}{$\displaystyle 
\gamma_{\mathrm{rel}}=-\frac{u_\mu U^\mu}{c^2}.
$}
\typeout{GRDISPLAY|56|\the\wd\grmathbox|\the\linewidth}
\ifdim\wd\grmathbox>\linewidth
\resizebox{\linewidth}{!}{\usebox{\grmathbox}}
\else\usebox{\grmathbox}\fi
\endgroup
\end{center}

Then \(E_{(U)}=\gamma_{\mathrm{rel}}mc^2\). Decompose the momentum into
the observer\textquotesingle s temporal and spatial parts:

\begin{center}
\begingroup
\sbox{\grmathbox}{$\displaystyle 
p^\mu=\frac{E_{(U)}}{c^2}U^\mu+p_\perp^\mu,
\qquad
U_\mu p_\perp^\mu=0.
$}
\typeout{GRDISPLAY|57|\the\wd\grmathbox|\the\linewidth}
\ifdim\wd\grmathbox>\linewidth
\resizebox{\linewidth}{!}{\usebox{\grmathbox}}
\else\usebox{\grmathbox}\fi
\endgroup
\end{center}

The orthogonality condition defines the observer\textquotesingle s
instantaneous three-dimensional rest space. It does not require a global
universal simultaneity surface.

For a photon moving in the \(+x\) direction, take
\(p^\mu=(E/c,E/c,0,0)\) and an observer chasing it with
\(U^\mu=\gamma(c,v,0,0)\). The energy measured by that observer is

\begin{center}
\begingroup
\sbox{\grmathbox}{$\displaystyle 
E_{(U)}=\gamma E(1-v/c)
=E\sqrt{\frac{1-v/c}{1+v/c}}.
$}
\typeout{GRDISPLAY|58|\the\wd\grmathbox|\the\linewidth}
\ifdim\wd\grmathbox>\linewidth
\resizebox{\linewidth}{!}{\usebox{\grmathbox}}
\else\usebox{\grmathbox}\fi
\endgroup
\end{center}

The photon is redshifted. Its local speed remains \(c\). Chasing light
changes the measured frequency, not the vacuum speed of light in the
chasing observer\textquotesingle s local inertial frame.

This observer-projection viewpoint will later make the stress-energy
tensor intelligible. ``Energy density'' is not just a number floating in
spacetime; it means energy density measured by some local observer.

\hypertarget{proper-acceleration-what-an-accelerometer-knows}{%
\section{3.8 Proper acceleration: what an accelerometer
knows}\label{proper-acceleration-what-an-accelerometer-knows}}

In inertial Minkowski coordinates, define four-acceleration

\begin{center}
\begingroup
\sbox{\grmathbox}{$\displaystyle 
a^\mu=\frac{du^\mu}{d\tau}.
$}
\typeout{GRDISPLAY|59|\the\wd\grmathbox|\the\linewidth}
\ifdim\wd\grmathbox>\linewidth
\resizebox{\linewidth}{!}{\usebox{\grmathbox}}
\else\usebox{\grmathbox}\fi
\endgroup
\end{center}

Differentiate \(u_\mu u^\mu=-c^2\). Because the Minkowski metric is
constant,

\begin{center}
\begingroup
\sbox{\grmathbox}{$\displaystyle 
2u_\mu a^\mu=0.
$}
\typeout{GRDISPLAY|60|\the\wd\grmathbox|\the\linewidth}
\ifdim\wd\grmathbox>\linewidth
\resizebox{\linewidth}{!}{\usebox{\grmathbox}}
\else\usebox{\grmathbox}\fi
\endgroup
\end{center}

Four-acceleration is orthogonal to four-velocity. In the instantaneous
rest frame, \(u^\mu=(c,0,0,0)\), so \(a^0=0\) and \(a^\mu\) is purely
spatial. Its magnitude

\begin{center}
\begingroup
\sbox{\grmathbox}{$\displaystyle 
\alpha=\sqrt{a_\mu a^\mu}
$}
\typeout{GRDISPLAY|61|\the\wd\grmathbox|\the\linewidth}
\ifdim\wd\grmathbox>\linewidth
\resizebox{\linewidth}{!}{\usebox{\grmathbox}}
\else\usebox{\grmathbox}\fi
\endgroup
\end{center}

is the \textbf{proper acceleration}, measured by an ideal accelerometer.
There is no square-root sign crisis: a vector orthogonal to a timelike
vector is spacelike or zero.

In curved spacetime, or even in noninertial coordinates on flat
spacetime, \(du^\mu/d\tau\) alone is not the correct vector
acceleration. We will replace it by a covariant derivative. That
replacement is exactly what will let us say, precisely rather than
poetically, that a falling apple is unaccelerated while the floor
beneath you is accelerating.

\hypertarget{chapter-4}{%
\chapter{4. Spacetime as a manifold: maps, rulers, and the geometry
beneath them}\label{chapter-4}}

\hypertarget{a-manifold-is-a-place-where-local-coordinates-work}{%
\section{4.1 A manifold is a place where local coordinates
work}\label{a-manifold-is-a-place-where-local-coordinates-work}}

The surface of Earth can be mapped in neighborhoods, even though no
single ordinary flat map represents the entire surface smoothly and
one-to-one without exclusions or identification rules. A \textbf{smooth
manifold} formalizes that pattern: near every point, there are \(n\)
real coordinates, and overlapping coordinate descriptions are related by
smooth invertible transformations.

A coordinate \textbf{chart} is a map

\begin{center}
\begingroup
\sbox{\grmathbox}{$\displaystyle 
x:U\subset M\longrightarrow x(U)\subset\mathbb R^n
$}
\typeout{GRDISPLAY|62|\the\wd\grmathbox|\the\linewidth}
\ifdim\wd\grmathbox>\linewidth
\resizebox{\linewidth}{!}{\usebox{\grmathbox}}
\else\usebox{\grmathbox}\fi
\endgroup
\end{center}

from an open region of the manifold to an open region of ordinary
coordinate space. An \textbf{atlas} is a collection of compatible charts
covering the manifold. Smoothness of their transition maps lets us
differentiate fields without a change of chart creating fictitious
discontinuities. The standard definition also includes technical
separation and countability conditions that exclude pathological spaces;
the local-chart picture is the ingredient we need for calculations.

The manifold specifies which events exist and how their neighborhoods
fit together. It does \textbf{not} yet specify distances, angles, light
cones, clocks, or gravitational dynamics. Those require further
structure.

A rubber glove and a steel glove can have the same underlying
organization of points while having different physical distances between
them. Think of the manifold as the arrangement and the metric as the
measuring instructions. The analogy is limited: a spacetime metric is
Lorentzian and includes causal structure, not merely elastic stretching
of spatial material.

There is no requirement that four-dimensional spacetime be embedded in
some physical five-dimensional room. Intrinsic geometry works without an
external vantage point. If you ask ``what is spacetime curved into?'',
the theory is entitled to answer: that extra place is not needed to
formulate the question we actually measure.

\hypertarget{coordinate-singularities-are-not-automatically-broken-physics}{%
\section{4.2 Coordinate singularities are not automatically broken
physics}\label{coordinate-singularities-are-not-automatically-broken-physics}}

On the Euclidean plane,

\begin{center}
\begingroup
\sbox{\grmathbox}{$\displaystyle 
x=r\cos\theta,\qquad y=r\sin\theta.
$}
\typeout{GRDISPLAY|63|\the\wd\grmathbox|\the\linewidth}
\ifdim\wd\grmathbox>\linewidth
\resizebox{\linewidth}{!}{\usebox{\grmathbox}}
\else\usebox{\grmathbox}\fi
\endgroup
\end{center}

At \(r=0\), every value of \(\theta\) labels the same point. Polar
coordinates fail there. The plane does not acquire a physical puncture
just because our coordinate system loses its manners.

Likewise, longitude fails to distinguish directions at a
sphere\textquotesingle s poles. One can use another chart near a pole.
The original labels were inadequate; the geometry is regular.

This distinction becomes consequential around black holes. A metric
component blowing up or vanishing might indicate a chart boundary, a
poor coordinate choice, or a physical singularity. The component alone
does not decide. We must inspect quantities and structures that survive
coordinate changes, and sometimes the extendibility of the spacetime
itself.

Another subtlety: an event\textquotesingle s coordinate tuple \(x^\mu\)
is generally \textbf{not} a vector. Under a nonlinear coordinate change,
coordinates transform nonlinearly. Tangent vectors transform with the
Jacobian. In flat affine coordinates, differences of position tuples
happen to behave as vectors, which can conceal the distinction.

\hypertarget{tangent-vectors-live-at-events}{%
\section{4.3 Tangent vectors live at
events}\label{tangent-vectors-live-at-events}}

Take a curve through an event \(p\), written \(x^\mu(\lambda)\). Its
tangent components are

\begin{center}
\begingroup
\sbox{\grmathbox}{$\displaystyle 
V^\mu=\left.\frac{dx^\mu}{d\lambda}\right|_p.
$}
\typeout{GRDISPLAY|64|\the\wd\grmathbox|\the\linewidth}
\ifdim\wd\grmathbox>\linewidth
\resizebox{\linewidth}{!}{\usebox{\grmathbox}}
\else\usebox{\grmathbox}\fi
\endgroup
\end{center}

A geometrically clean definition asks what this tangent does to any
smooth scalar field \(f\):

\begin{center}
\begingroup
\sbox{\grmathbox}{$\displaystyle 
V[f]=\left.\frac{d}{d\lambda}f(x(\lambda))\right|_p
=V^\mu\partial_\mu f\big|_p.
$}
\typeout{GRDISPLAY|65|\the\wd\grmathbox|\the\linewidth}
\ifdim\wd\grmathbox>\linewidth
\resizebox{\linewidth}{!}{\usebox{\grmathbox}}
\else\usebox{\grmathbox}\fi
\endgroup
\end{center}

The vector is a directional derivative operator. This is not an
abstraction designed to remove the arrow from physics. It tells you
exactly how to detect the arrow: put scalar fields in its way and
measure their directional changes.

The operator satisfies linearity and the product rule,

\begin{center}
\begingroup
\sbox{\grmathbox}{$\displaystyle 
V[af+bg]=aV[f]+bV[g],
\qquad
V[fg]=V[f]g(p)+f(p)V[g].
$}
\typeout{GRDISPLAY|66|\the\wd\grmathbox|\the\linewidth}
\ifdim\wd\grmathbox>\linewidth
\resizebox{\linewidth}{!}{\usebox{\grmathbox}}
\else\usebox{\grmathbox}\fi
\endgroup
\end{center}

All tangent vectors at \(p\) form the tangent space \(T_pM\). A
coordinate basis is

\begin{center}
\begingroup
\sbox{\grmathbox}{$\displaystyle 
e_\mu=\partial_\mu\big|_p,
\qquad
V=V^\mu e_\mu.
$}
\typeout{GRDISPLAY|67|\the\wd\grmathbox|\the\linewidth}
\ifdim\wd\grmathbox>\linewidth
\resizebox{\linewidth}{!}{\usebox{\grmathbox}}
\else\usebox{\grmathbox}\fi
\endgroup
\end{center}

Its dual basis is \(dx^\mu|_p\), with
\(dx^\mu(e_\nu)=\delta^\mu{}_{\nu}\).

Now the problem that motivates the next part of the book becomes
unavoidable. A vector at \(p\) belongs to \(T_pM\); a vector at a
neighboring event \(q\) belongs to \(T_qM\). They are elements of
different vector spaces. Subtracting them requires a rule for comparing
those spaces.

On a flat Cartesian grid, an obvious translation rule is already
operating in the background. On a general manifold there is no preferred
rule supplied by the manifold alone. A \textbf{connection} will provide
the missing comparison procedure. It is geometry\textquotesingle s
answer to ``before we compare these measurements, are the instruments
aligned?''

\hypertarget{the-metric-is-a-local-measuring-operation}{%
\section{4.4 The metric is a local measuring
operation}\label{the-metric-is-a-local-measuring-operation}}

At every event, the metric takes two tangent vectors and returns a
scalar:

\begin{center}
\begingroup
\sbox{\grmathbox}{$\displaystyle 
g(V,W)=g_{\mu\nu}V^\mu W^\nu.
$}
\typeout{GRDISPLAY|68|\the\wd\grmathbox|\the\linewidth}
\ifdim\wd\grmathbox>\linewidth
\resizebox{\linewidth}{!}{\usebox{\grmathbox}}
\else\usebox{\grmathbox}\fi
\endgroup
\end{center}

It is symmetric,

\begin{center}
\begingroup
\sbox{\grmathbox}{$\displaystyle 
g_{\mu\nu}=g_{\nu\mu},
$}
\typeout{GRDISPLAY|69|\the\wd\grmathbox|\the\linewidth}
\ifdim\wd\grmathbox>\linewidth
\resizebox{\linewidth}{!}{\usebox{\grmathbox}}
\else\usebox{\grmathbox}\fi
\endgroup
\end{center}

and nondegenerate: no nonzero vector is orthogonal to every vector. In
GR it has Lorentzian signature \((-,+,+,+)\). At a point, there is a
basis in which its components are \(\eta_{\mu\nu}\), with one negative
and three positive eigenvalue directions. The number of positive and
negative directions cannot be changed by a nonsingular real basis
transformation.

Nondegenerate does \textbf{not} mean \(g(V,V)\) is nonzero for every
nonzero \(V\). Null vectors have \(g(V,V)=0\), but a nonzero null vector
still has nonzero pairing with some other vectors. A metric with light
cones is not a singular matrix.

The line element

\begin{center}
\begingroup
\sbox{\grmathbox}{$\displaystyle 
ds^2=g_{\mu\nu}(x)dx^\mu dx^\nu
$}
\typeout{GRDISPLAY|70|\the\wd\grmathbox|\the\linewidth}
\ifdim\wd\grmathbox>\linewidth
\resizebox{\linewidth}{!}{\usebox{\grmathbox}}
\else\usebox{\grmathbox}\fi
\endgroup
\end{center}

encodes the local geometry. For a timelike worldline,
\(d\tau=\sqrt{-ds^2}/c\). For a spacelike curve, its length is obtained
from \(\int\sqrt{ds^2}\). Defining a finite spatial distance between
distant observers requires a specified spacelike slice or measurement
protocol; the metric does not hand everyone the same universal
``distance right now.''

The metric also defines the light cone by \(g(V,V)=0\). Thus it is
simultaneously a ruler, a clock specification, and a causal organizer.
``Gravity changes distances'' describes only part of its job.

In four dimensions, a symmetric \(4\times4\) metric has \(4(4+1)/2=10\)
independent components. Those are ten functions in a coordinate
description, \textbf{not ten independent propagating gravitational
modes}. Coordinate freedom and the structure of the field equations will
substantially change that counting.

\hypertarget{a-complete-metric-calculation-in-flat-polar-coordinates}{%
\section{4.5 A complete metric calculation in flat polar
coordinates}\label{a-complete-metric-calculation-in-flat-polar-coordinates}}

Differentiate the polar transformation:

\begin{center}
\begingroup
\sbox{\grmathbox}{$\displaystyle 
dx=\cos\theta\,dr-r\sin\theta\,d\theta,
\qquad
dy=\sin\theta\,dr+r\cos\theta\,d\theta.
$}
\typeout{GRDISPLAY|71|\the\wd\grmathbox|\the\linewidth}
\ifdim\wd\grmathbox>\linewidth
\resizebox{\linewidth}{!}{\usebox{\grmathbox}}
\else\usebox{\grmathbox}\fi
\endgroup
\end{center}

Square and add. The mixed terms cancel, and
\(\sin^2\theta+\cos^2\theta=1\) gives

\begin{center}
\begingroup
\sbox{\grmathbox}{$\displaystyle 
d\ell^2=dx^2+dy^2=dr^2+r^2d\theta^2.
$}
\typeout{GRDISPLAY|72|\the\wd\grmathbox|\the\linewidth}
\ifdim\wd\grmathbox>\linewidth
\resizebox{\linewidth}{!}{\usebox{\grmathbox}}
\else\usebox{\grmathbox}\fi
\endgroup
\end{center}

Therefore the spatial metric is

\begin{center}
\begingroup
\sbox{\grmathbox}{$\displaystyle 
(g_{ij})=
\begin{pmatrix}1&0\\0&r^2\end{pmatrix}.
$}
\typeout{GRDISPLAY|73|\the\wd\grmathbox|\the\linewidth}
\ifdim\wd\grmathbox>\linewidth
\resizebox{\linewidth}{!}{\usebox{\grmathbox}}
\else\usebox{\grmathbox}\fi
\endgroup
\end{center}

The coefficient \(r^2\) is necessary because one radian of angular
change spans a longer arc farther from the origin. It is not evidence of
curved space. We obtained the metric from the flat Euclidean plane by
changing labels.

Units deserve attention. If \(r\) is measured in meters and \(\theta\)
is dimensionless, \(g_{rr}\) is dimensionless but \(g_{\theta\theta}\)
has units of square meters. It is the full line element that must have
consistent units. The statement ``the metric is dimensionless'' is only
true in suitable coordinate conventions.

Off-diagonal terms also have a simple meaning. On the same plane, choose
oblique coordinates with position

\begin{center}
\begingroup
\sbox{\grmathbox}{$\displaystyle 
\mathbf r(q^1,q^2)=q^1\hat{\mathbf x}
+q^2(\hat{\mathbf x}+\hat{\mathbf y}).
$}
\typeout{GRDISPLAY|74|\the\wd\grmathbox|\the\linewidth}
\ifdim\wd\grmathbox>\linewidth
\resizebox{\linewidth}{!}{\usebox{\grmathbox}}
\else\usebox{\grmathbox}\fi
\endgroup
\end{center}

Then

\begin{center}
\begingroup
\sbox{\grmathbox}{$\displaystyle 
d\ell^2=(dq^1)^2+2dq^1dq^2+2(dq^2)^2,
\qquad
(g_{ij})=\begin{pmatrix}1&1\\1&2\end{pmatrix}.
$}
\typeout{GRDISPLAY|75|\the\wd\grmathbox|\the\linewidth}
\ifdim\wd\grmathbox>\linewidth
\resizebox{\linewidth}{!}{\usebox{\grmathbox}}
\else\usebox{\grmathbox}\fi
\endgroup
\end{center}

The cross term tells you the coordinate basis directions are not
orthogonal. Its coefficient is \(2g_{12}\) because the Einstein sum
includes both \(g_{12}dq^1dq^2\) and \(g_{21}dq^2dq^1\).

\hypertarget{inverse-metrics-and-raising-or-lowering-indices}{%
\section{4.6 Inverse metrics and raising or lowering
indices}\label{inverse-metrics-and-raising-or-lowering-indices}}

The inverse metric is defined by

\begin{center}
\begingroup
\sbox{\grmathbox}{$\displaystyle 
g^{\mu\alpha}g_{\alpha\nu}=\delta^\mu{}_{\nu}.
$}
\typeout{GRDISPLAY|76|\the\wd\grmathbox|\the\linewidth}
\ifdim\wd\grmathbox>\linewidth
\resizebox{\linewidth}{!}{\usebox{\grmathbox}}
\else\usebox{\grmathbox}\fi
\endgroup
\end{center}

It is the matrix inverse, not component-by-component reciprocation
except when the metric is diagonal. For our oblique metric,

\begin{center}
\begingroup
\sbox{\grmathbox}{$\displaystyle 
(g^{ij})=\begin{pmatrix}2&-1\\-1&1\end{pmatrix}.
$}
\typeout{GRDISPLAY|77|\the\wd\grmathbox|\the\linewidth}
\ifdim\wd\grmathbox>\linewidth
\resizebox{\linewidth}{!}{\usebox{\grmathbox}}
\else\usebox{\grmathbox}\fi
\endgroup
\end{center}

Multiplying the two matrices verifies the identity directly.

The metric converts a vector into a covector:

\begin{center}
\begingroup
\sbox{\grmathbox}{$\displaystyle 
V_\mu=g_{\mu\nu}V^\nu.
$}
\typeout{GRDISPLAY|78|\the\wd\grmathbox|\the\linewidth}
\ifdim\wd\grmathbox>\linewidth
\resizebox{\linewidth}{!}{\usebox{\grmathbox}}
\else\usebox{\grmathbox}\fi
\endgroup
\end{center}

The new covector asks another vector \(W\) the question
\(V_\mu W^\mu=g(V,W)\). The inverse conversion is

\begin{center}
\begingroup
\sbox{\grmathbox}{$\displaystyle 
\omega^\mu=g^{\mu\nu}\omega_\nu.
$}
\typeout{GRDISPLAY|79|\the\wd\grmathbox|\the\linewidth}
\ifdim\wd\grmathbox>\linewidth
\resizebox{\linewidth}{!}{\usebox{\grmathbox}}
\else\usebox{\grmathbox}\fi
\endgroup
\end{center}

For the oblique example, \(V^i=(1,2)\) gives \(V_i=(3,5)\), and

\begin{center}
\begingroup
\sbox{\grmathbox}{$\displaystyle 
g(V,V)=V_iV^i=3\cdot1+5\cdot2=13.
$}
\typeout{GRDISPLAY|80|\the\wd\grmathbox|\the\linewidth}
\ifdim\wd\grmathbox>\linewidth
\resizebox{\linewidth}{!}{\usebox{\grmathbox}}
\else\usebox{\grmathbox}\fi
\endgroup
\end{center}

The corresponding ordinary Cartesian vector is
\(3\hat{\mathbf x}+2\hat{\mathbf y}\), whose Euclidean norm squared is
indeed \(9+4=13\). Raising and lowering has translated between two
different kinds of object while preserving a precise relationship; it
has not preserved every numerical component.

In inertial spacetime coordinates,

\begin{center}
\begingroup
\sbox{\grmathbox}{$\displaystyle 
V^\mu=(V^0,V^1,V^2,V^3)
\quad\Longrightarrow\quad
V_\mu=(-V^0,V^1,V^2,V^3).
$}
\typeout{GRDISPLAY|81|\the\wd\grmathbox|\the\linewidth}
\ifdim\wd\grmathbox>\linewidth
\resizebox{\linewidth}{!}{\usebox{\grmathbox}}
\else\usebox{\grmathbox}\fi
\endgroup
\end{center}

For example, \(u_0=-\gamma c\) while \(u^0=\gamma c\). This minus sign
is responsible for many apparently mysterious signs in energy formulas.

The gradient distinction from Chapter 2 now resolves:

\begin{center}
\begingroup
\sbox{\grmathbox}{$\displaystyle 
(df)_\mu=\partial_\mu f,
\qquad
(\operatorname{grad}f)^\mu=g^{\mu\nu}\partial_\nu f.
$}
\typeout{GRDISPLAY|82|\the\wd\grmathbox|\the\linewidth}
\ifdim\wd\grmathbox>\linewidth
\resizebox{\linewidth}{!}{\usebox{\grmathbox}}
\else\usebox{\grmathbox}\fi
\endgroup
\end{center}

The differential was available before the metric. The gradient vector
was not. In Lorentzian geometry it is also unsafe to carry over every
Euclidean slogan about a gradient ``pointing uphill most steeply,''
because the unit-vector set and norm structure are different.

\hypertarget{determinants-tell-integration-how-much-room-there-is}{%
\section{4.7 Determinants tell integration how much room there
is}\label{determinants-tell-integration-how-much-room-there-is}}

On a plane, a narrow polar cell has physical area approximately
\(dr\times r\,d\theta\), so

\begin{center}
\begingroup
\sbox{\grmathbox}{$\displaystyle 
dA=r\,dr\,d\theta=\sqrt{\det g_{ij}}\,dr\,d\theta.
$}
\typeout{GRDISPLAY|83|\the\wd\grmathbox|\the\linewidth}
\ifdim\wd\grmathbox>\linewidth
\resizebox{\linewidth}{!}{\usebox{\grmathbox}}
\else\usebox{\grmathbox}\fi
\endgroup
\end{center}

In general, the metric determinant measures the squared volume scaling
of the coordinate basis. Let

\begin{center}
\begingroup
\sbox{\grmathbox}{$\displaystyle 
g\equiv\det(g_{\mu\nu}).
$}
\typeout{GRDISPLAY|84|\the\wd\grmathbox|\the\linewidth}
\ifdim\wd\grmathbox>\linewidth
\resizebox{\linewidth}{!}{\usebox{\grmathbox}}
\else\usebox{\grmathbox}\fi
\endgroup
\end{center}

With Lorentzian signature \((-,+,+,+)\), \(g<0\) in every regular
coordinate chart. The invariant spacetime volume measure is

\begin{center}
\begingroup
\sbox{\grmathbox}{$\displaystyle 
\boxed{dV_4=\sqrt{-g}\,d^4x.}
$}
\typeout{GRDISPLAY|85|\the\wd\grmathbox|\the\linewidth}
\ifdim\wd\grmathbox>\linewidth
\resizebox{\linewidth}{!}{\usebox{\grmathbox}}
\else\usebox{\grmathbox}\fi
\endgroup
\end{center}

Why exactly this factor? Under \(x\mapsto x'\), the covariant metric
transforms as a matrix according to

\begin{center}
\begingroup
\sbox{\grmathbox}{$\displaystyle 
g'=K^{\mathsf T}gK,
\qquad
\det(g')=(\det K)^2\det(g),
$}
\typeout{GRDISPLAY|86|\the\wd\grmathbox|\the\linewidth}
\ifdim\wd\grmathbox>\linewidth
\resizebox{\linewidth}{!}{\usebox{\grmathbox}}
\else\usebox{\grmathbox}\fi
\endgroup
\end{center}

where the first equation uses \(g\) for the metric matrix and the second
its determinant. Thus the square root picks up \(|\det K|\). Meanwhile
\(d^4x'=|\det J|d^4x\), and \(K=J^{-1}\). The factors cancel:

\begin{center}
\begingroup
\sbox{\grmathbox}{$\displaystyle 
\sqrt{-g'}\,d^4x'=\sqrt{-g}\,d^4x.
$}
\typeout{GRDISPLAY|87|\the\wd\grmathbox|\the\linewidth}
\ifdim\wd\grmathbox>\linewidth
\resizebox{\linewidth}{!}{\usebox{\grmathbox}}
\else\usebox{\grmathbox}\fi
\endgroup
\end{center}

An oriented volume \textbf{form} additionally tracks the sign of
orientation using a wedge product. The positive integration measure
above uses absolute Jacobians. These are compatible viewpoints, but
changing orientation is where their notation must be handled carefully.
A formal treatment is available in
\href{https://davidtong.org/pdfs/teaching/general-relativity/gr3.pdf}{Tong\textquotesingle s
discussion of the metric volume form}.

For flat spacetime in spherical spatial coordinates
\((x^0,r,\theta,\phi)=(ct,r,\theta,\phi)\),

\begin{center}
\begingroup
\sbox{\grmathbox}{$\displaystyle 
ds^2=-(dx^0)^2+dr^2+r^2d\theta^2+r^2\sin^2\theta\,d\phi^2,
$}
\typeout{GRDISPLAY|88|\the\wd\grmathbox|\the\linewidth}
\ifdim\wd\grmathbox>\linewidth
\resizebox{\linewidth}{!}{\usebox{\grmathbox}}
\else\usebox{\grmathbox}\fi
\endgroup
\end{center}

so

\begin{center}
\begingroup
\sbox{\grmathbox}{$\displaystyle 
g=-r^4\sin^2\theta,
\qquad
dV_4=c\,dt\,r^2\sin\theta\,dr\,d\theta\,d\phi.
$}
\typeout{GRDISPLAY|89|\the\wd\grmathbox|\the\linewidth}
\ifdim\wd\grmathbox>\linewidth
\resizebox{\linewidth}{!}{\usebox{\grmathbox}}
\else\usebox{\grmathbox}\fi
\endgroup
\end{center}

The familiar spherical Jacobian was the metric determinant in civilian
clothes. The vanishing determinant at \(r=0\) or at a polar axis signals
failure of this chart there; it does not mean the regular Minkowski
metric becomes physically degenerate.

\hypertarget{intrinsic-curvature-extrinsic-bending-and-laboratory-frames}{%
\section{4.8 Intrinsic curvature, extrinsic bending, and laboratory
frames}\label{intrinsic-curvature-extrinsic-bending-and-laboratory-frames}}

A sheet can be rolled into a cylinder without stretching it. Distances
and angles measured within a sufficiently small patch are unchanged, so
its intrinsic curvature is zero. Its embedding in three-dimensional
space is visibly bent. A sphere cannot be produced from a flat sheet
without stretching or cutting, and its intrinsic geometry differs from a
plane\textquotesingle s.

Intrinsic curvature concerns measurements available to inhabitants of
the geometry. Extrinsic curvature concerns how a chosen submanifold sits
inside a larger geometry. Both ideas occur in relativity: spacetime has
intrinsic curvature, and a spatial slice can have extrinsic curvature
within spacetime. They are different objects answering different
questions.

Coordinate basis vectors need not be unit length or orthogonal. A
physical laboratory instead likes an \textbf{orthonormal frame}
\(e_{\hat a}\) satisfying

\begin{center}
\begingroup
\sbox{\grmathbox}{$\displaystyle 
g(e_{\hat a},e_{\hat b})=\eta_{\hat a\hat b}.
$}
\typeout{GRDISPLAY|90|\the\wd\grmathbox|\the\linewidth}
\ifdim\wd\grmathbox>\linewidth
\resizebox{\linewidth}{!}{\usebox{\grmathbox}}
\else\usebox{\grmathbox}\fi
\endgroup
\end{center}

Hats here label laboratory directions. For an observer of four-velocity
\(U\), choose \(e_{\hat0}=U/c\). The other three frame vectors specify
their instantaneous spatial axes.

On the polar plane, an orthonormal frame is

\begin{center}
\begingroup
\sbox{\grmathbox}{$\displaystyle 
e_{\hat r}=\partial_r,
\qquad
e_{\hat\theta}=\frac1r\partial_\theta.
$}
\typeout{GRDISPLAY|91|\the\wd\grmathbox|\the\linewidth}
\ifdim\wd\grmathbox>\linewidth
\resizebox{\linewidth}{!}{\usebox{\grmathbox}}
\else\usebox{\grmathbox}\fi
\endgroup
\end{center}

Thus \(V^{\hat r}=V^r\) and \(V^{\hat\theta}=rV^\theta\). The physical
tangential component includes the conversion from angle to arc length.

There is a subtle distinction between making the
metric\textquotesingle s components constant in a \textbf{frame} and
finding \textbf{coordinates} in which they are constant throughout a
neighborhood. The first is possible locally for a smooth Lorentzian
metric. The second would make the neighborhood flat. A varying
orthonormal frame generally cannot be the coordinate basis of one
coordinate system; its axes can turn relative to one another from point
to point. The connection keeps track of that turning.

At any ordinary event of a smooth spacetime, one can choose locally
inertial coordinates with

\begin{center}
\begingroup
\sbox{\grmathbox}{$\displaystyle 
g_{\mu\nu}(p)=\eta_{\mu\nu},
\qquad
\partial_\rho g_{\mu\nu}(p)=0.
$}
\typeout{GRDISPLAY|92|\the\wd\grmathbox|\the\linewidth}
\ifdim\wd\grmathbox>\linewidth
\resizebox{\linewidth}{!}{\usebox{\grmathbox}}
\else\usebox{\grmathbox}\fi
\endgroup
\end{center}

This is stronger than merely choosing an orthonormal basis at the point,
but it does not generally remove second-order departures from flatness.
It is the mathematical doorway to the equivalence principle.

\hypertarget{chapter-5}{%
\chapter{5. Free fall, the equivalence principle, and the worldline
action}\label{chapter-5}}

\hypertarget{ask-the-accelerometer-not-the-window}{%
\section{5.1 Ask the accelerometer, not the
window}\label{ask-the-accelerometer-not-the-window}}

Standing on the ground, you assign yourself constant spatial
coordinates. A dropped ball\textquotesingle s coordinates accelerate
downward. Everyday language calls you unaccelerated and the ball
accelerated.

Now ask an accelerometer. Yours reads approximately the local
gravitational acceleration because the floor pushes upward on you. An
ideal accelerometer falling with the ball reads zero, ignoring air
resistance and finite-size effects. Its case and internal proof mass
follow free fall together, so the proof mass does not deflect relative
to the case.

GR organizes its local inertial physics around that second distinction.
\textbf{Free fall is zero proper acceleration.} Remaining at a fixed
altitude near Earth requires a nongravitational force.

This does not make Newton\textquotesingle s description useless or make
gravity imaginary. It distinguishes a coordinate acceleration from a
physical acceleration measured along a worldline. Later, tidal effects
will reveal gravitational structure even when every individual
accelerometer reads zero.

\hypertarget{what-the-equivalence-principle-says---and-what-it-does-not-buy-you}{%
\section{5.2 What the equivalence principle says - and what it does not
buy
you}\label{what-the-equivalence-principle-says---and-what-it-does-not-buy-you}}

In Newtonian notation, allow an inertial mass \(m_{\mathrm{I}}\) and a
passive gravitational mass \(m_{\mathrm{G}}\):

\begin{center}
\begingroup
\sbox{\grmathbox}{$\displaystyle 
m_{\mathrm{I}}\frac{d^2\mathbf x}{dt^2}=-m_{\mathrm{G}}\symbfit\nabla\Phi.
$}
\typeout{GRDISPLAY|93|\the\wd\grmathbox|\the\linewidth}
\ifdim\wd\grmathbox>\linewidth
\resizebox{\linewidth}{!}{\usebox{\grmathbox}}
\else\usebox{\grmathbox}\fi
\endgroup
\end{center}

Universality of free fall says the ratio
\(m_{\mathrm{G}}/m_{\mathrm{I}}\) is independent of a sufficiently small
test body\textquotesingle s composition and internal constitution, under
the appropriate idealizations. A universal proportionality can be
absorbed into the definition of the gravitational coupling, leaving
\(m_{\mathrm{G}}=m_{\mathrm{I}}\). Bodies with the same initial position
and velocity then follow the same trajectory.

The \textbf{Einstein equivalence principle} extends the idea to local
nongravitational physics: freely falling laboratories obey special
relativity locally, and local nongravitational experiments do not
acquire different laws merely from the laboratory\textquotesingle s
velocity or location. The \textbf{strong equivalence principle} extends
the scope to gravitational experiments and self-gravitating bodies. The
distinctions and their experimental roles are carefully organized in
\href{https://arxiv.org/abs/1403.7377}{Clifford Will\textquotesingle s
review of tests of gravitation}.

Three qualifications make these statements stronger intellectually
rather than weaker rhetorically.

First, \emph{local} matters. In freely falling coordinates the metric
can be Minkowskian and its first derivatives zero at an event. Curvature
can still produce measurable effects across a finite laboratory or after
a finite time. Making the laboratory smaller suppresses such effects; it
does not declare curvature nonexistent.

Second, ideal test-body motion neglects the body\textquotesingle s own
gravitational backreaction and finite-size structure. Spinning extended
bodies, bodies with multipole moments, and objects subject to self-force
corrections require more elaborate motion laws. The geodesic
approximation has a domain of validity.

Third, the equivalence principle does not uniquely imply
Einstein\textquotesingle s field equation. Multiple theories can use a
metric, respect the local free-fall picture, and supply different
dynamics for that metric. We have learned how the local measuring system
behaves; we have not yet derived what creates the gravitational field.

\hypertarget{an-accelerating-laboratory-can-have-a-varying-metric-in-flat-spacetime}{%
\section{5.3 An accelerating laboratory can have a varying metric in
flat
spacetime}\label{an-accelerating-laboratory-can-have-a-varying-metric-in-flat-spacetime}}

This worked example prevents a persistent misconception:
gravitational-looking clock rates do not by themselves establish
curvature.

Start in flat spacetime with inertial coordinates \((T,X,Y,Z)\).
Introduce an accelerating chart \((t,X,Y,z)\) through

\begin{center}
\begingroup
\sbox{\grmathbox}{$\displaystyle 
cT=\left(\frac{c^2}{a}+z\right)\sinh\left(\frac{at}{c}\right),
$}
\typeout{GRDISPLAY|94|\the\wd\grmathbox|\the\linewidth}
\ifdim\wd\grmathbox>\linewidth
\resizebox{\linewidth}{!}{\usebox{\grmathbox}}
\else\usebox{\grmathbox}\fi
\endgroup
\end{center}

\begin{center}
\begingroup
\sbox{\grmathbox}{$\displaystyle 
Z=\left(\frac{c^2}{a}+z\right)\cosh\left(\frac{at}{c}\right)-\frac{c^2}{a},
$}
\typeout{GRDISPLAY|95|\the\wd\grmathbox|\the\linewidth}
\ifdim\wd\grmathbox>\linewidth
\resizebox{\linewidth}{!}{\usebox{\grmathbox}}
\else\usebox{\grmathbox}\fi
\endgroup
\end{center}

with \(a>0\) and \(z>-c^2/a\). These are Rindler coordinates on a region
of Minkowski spacetime. We are explicitly using \(t\), rather than
\(ct\), as the time coordinate in this chart.

Differentiate these equations and substitute into
\(-c^2dT^2+dZ^2+dX^2+dY^2\). The mixed \(dt\,dz\) terms cancel, and
\(\cosh^2-\sinh^2=1\) leaves

\begin{center}
\begingroup
\sbox{\grmathbox}{$\displaystyle 
ds^2=-\left(1+\frac{az}{c^2}\right)^2c^2dt^2
+dz^2+dX^2+dY^2.
$}
\typeout{GRDISPLAY|96|\the\wd\grmathbox|\the\linewidth}
\ifdim\wd\grmathbox>\linewidth
\resizebox{\linewidth}{!}{\usebox{\grmathbox}}
\else\usebox{\grmathbox}\fi
\endgroup
\end{center}

A clock at fixed \(z,X,Y\) measures

\begin{center}
\begingroup
\sbox{\grmathbox}{$\displaystyle 
d\tau=\left(1+\frac{az}{c^2}\right)dt.
$}
\typeout{GRDISPLAY|97|\the\wd\grmathbox|\the\linewidth}
\ifdim\wd\grmathbox>\linewidth
\resizebox{\linewidth}{!}{\usebox{\grmathbox}}
\else\usebox{\grmathbox}\fi
\endgroup
\end{center}

Stationary clocks at different heights in this accelerating chart
accumulate different proper times per coordinate time. Yet spacetime is
exactly flat: we constructed the metric by changing coordinates in
Minkowski space.

The observer at \(z=0\) has constant proper acceleration \(a\).
Differentiating the other fixed-\(z\) hyperbolic worldlines gives proper
acceleration

\begin{center}
\begingroup
\sbox{\grmathbox}{$\displaystyle 
\alpha(z)=\frac{a}{1+az/c^2}.
$}
\typeout{GRDISPLAY|98|\the\wd\grmathbox|\the\linewidth}
\ifdim\wd\grmathbox>\linewidth
\resizebox{\linewidth}{!}{\usebox{\grmathbox}}
\else\usebox{\grmathbox}\fi
\endgroup
\end{center}

An extended array that maintains these fixed separations does not have
identical proper acceleration at every height. Relativity puts
restrictions on the notion of an accelerating rigid elevator.

This example is also a warning against identifying ``nonconstant
metric,'' ``clock-rate difference,'' and ``spacetime curvature.'' They
are not equivalent statements. One must calculate curvature or measure
tidal geometry. For a complementary derivation of these accelerating
coordinates, see
\href{https://davidtong.org/pdfs/teaching/general-relativity/gr1.pdf}{Tong\textquotesingle s
treatment of the equivalence principle and Rindler motion}.

\hypertarget{why-a-worldline-action-is-the-right-next-move}{%
\section{5.4 Why a worldline action is the right next
move}\label{why-a-worldline-action-is-the-right-next-move}}

In Newtonian mechanics, a free particle moves along a straight line at
constant speed. In spacetime, the invariant version is that its
worldline has stationary proper time between fixed endpoint events.

For a structureless massive test particle, the standard minimally
coupled action is

\begin{center}
\begingroup
\sbox{\grmathbox}{$\displaystyle 
\boxed{S_{\mathrm{particle}}=-mc^2\int d\tau.}
$}
\typeout{GRDISPLAY|99|\the\wd\grmathbox|\the\linewidth}
\ifdim\wd\grmathbox>\linewidth
\resizebox{\linewidth}{!}{\usebox{\grmathbox}}
\else\usebox{\grmathbox}\fi
\endgroup
\end{center}

The action has units of energy times time. The prefactor and sign ensure
that its slow-motion limit gives the usual positive kinetic-energy term.
The fixed overall factor does not affect the free trajectory for
\(m\ne0\).

Why this form? Proper time is a scalar quantity attached to the path,
and the integral is unchanged if we relabel points along that path. With
no additional internal structure or higher-derivative couplings, it is
the simplest local relativistic free-particle action. This is a physical
modeling choice with extraordinary success, not a logical proof that
every conceivable body in every conceivable gravitational theory must
have this exact action.

``Stationary'' means the first-order change in the action vanishes under
sufficiently small endpoint-preserving deformations of the path. It does
not mean the particle tries every route, predicts the future, and
selects the best review score. A stationary-action statement is a
compact way to encode local differential equations.

For timelike geodesics, sufficiently short segments locally maximize
proper time. Over long segments, conjugate points or global geometry can
spoil the maximizing property. The field equation will not turn
``extremum'' into ``global minimum'' by force of enthusiasm.

\hypertarget{deriving-the-geodesic-equation-one-operation-at-a-time}{%
\section{5.5 Deriving the geodesic equation, one operation at a
time}\label{deriving-the-geodesic-equation-one-operation-at-a-time}}

Choose an arbitrary increasing parameter \(\lambda\) along a timelike
path. A dot in this subsection means \(d/d\lambda\). Define

\begin{center}
\begingroup
\sbox{\grmathbox}{$\displaystyle 
\ell=\sqrt{-g_{\mu\nu}(x)\dot x^\mu\dot x^\nu},
\qquad
d\tau=\frac{\ell}{c}\,d\lambda.
$}
\typeout{GRDISPLAY|100|\the\wd\grmathbox|\the\linewidth}
\ifdim\wd\grmathbox>\linewidth
\resizebox{\linewidth}{!}{\usebox{\grmathbox}}
\else\usebox{\grmathbox}\fi
\endgroup
\end{center}

Then \(S_{\mathrm{particle}}=-mc\int\ell\,d\lambda\). Because the
nonzero constant \(-mc\) does not change the stationary paths, vary
\(I=\int\ell\,d\lambda\) instead.

Perturb the curve by

\begin{center}
\begingroup
\sbox{\grmathbox}{$\displaystyle 
x^\mu(\lambda)\longrightarrow x^\mu(\lambda)+s\,\eta^\mu(\lambda),
$}
\typeout{GRDISPLAY|101|\the\wd\grmathbox|\the\linewidth}
\ifdim\wd\grmathbox>\linewidth
\resizebox{\linewidth}{!}{\usebox{\grmathbox}}
\else\usebox{\grmathbox}\fi
\endgroup
\end{center}

where \(s\) is a small number and \(\eta^\mu\) vanishes at both
endpoints. We take the derivative with respect to \(s\) at \(s=0\) and
denote it by \(\delta\).

Two things change: the metric is sampled at a slightly different point,
and the tangent to the path changes. Specifically,

\begin{center}
\begingroup
\sbox{\grmathbox}{$\displaystyle 
\delta g_{\mu\nu}=\partial_\rho g_{\mu\nu}\,\eta^\rho,
\qquad
\delta\dot x^\mu=\dot\eta^\mu.
$}
\typeout{GRDISPLAY|102|\the\wd\grmathbox|\the\linewidth}
\ifdim\wd\grmathbox>\linewidth
\resizebox{\linewidth}{!}{\usebox{\grmathbox}}
\else\usebox{\grmathbox}\fi
\endgroup
\end{center}

Differentiate the square root using
\(\delta\sqrt{Q}=\delta Q/(2\sqrt Q)\). Metric symmetry combines the two
tangent variations into one term:

\begin{center}
\begingroup
\sbox{\grmathbox}{$\displaystyle 
\delta\ell
=-\frac{1}{2\ell}\partial_\rho g_{\mu\nu}\,
\dot x^\mu\dot x^\nu\eta^\rho
-\frac{g_{\rho\nu}\dot x^\nu}{\ell}\dot\eta^\rho.
$}
\typeout{GRDISPLAY|103|\the\wd\grmathbox|\the\linewidth}
\ifdim\wd\grmathbox>\linewidth
\resizebox{\linewidth}{!}{\usebox{\grmathbox}}
\else\usebox{\grmathbox}\fi
\endgroup
\end{center}

The variation contains a derivative of the arbitrary deformation
\(\eta\). Integration by parts transfers that derivative onto its
coefficient:

\begin{center}
\begingroup
\sbox{\grmathbox}{$\displaystyle 
\delta I
=\left[-\frac{g_{\rho\nu}\dot x^\nu}{\ell}\eta^\rho\right]_{A}^{B}
+\int d\lambda\left[
\frac{d}{d\lambda}\left(\frac{g_{\rho\nu}\dot x^\nu}{\ell}\right)
-\frac{1}{2\ell}\partial_\rho g_{\mu\nu}\dot x^\mu\dot x^\nu
\right]\eta^\rho.
$}
\typeout{GRDISPLAY|104|\the\wd\grmathbox|\the\linewidth}
\ifdim\wd\grmathbox>\linewidth
\resizebox{\linewidth}{!}{\usebox{\grmathbox}}
\else\usebox{\grmathbox}\fi
\endgroup
\end{center}

The boundary term vanishes because the endpoint events are fixed. Inside
the interval, each \(\eta^\rho\) can be chosen freely. The only way the
integral can vanish for every such deformation is for its coefficient to
vanish:

\begin{center}
\begingroup
\sbox{\grmathbox}{$\displaystyle 
\frac{d}{d\lambda}\left(\frac{g_{\rho\nu}\dot x^\nu}{\ell}\right)
-\frac{1}{2\ell}\partial_\rho g_{\mu\nu}\dot x^\mu\dot x^\nu=0.
$}
\typeout{GRDISPLAY|105|\the\wd\grmathbox|\the\linewidth}
\ifdim\wd\grmathbox>\linewidth
\resizebox{\linewidth}{!}{\usebox{\grmathbox}}
\else\usebox{\grmathbox}\fi
\endgroup
\end{center}

That is already the Euler-Lagrange equation in expanded form. Multiply
by \(\ell\) and apply the product rule:

\begin{center}
\begingroup
\sbox{\grmathbox}{$\displaystyle 
g_{\rho\nu}\ddot x^\nu
+\partial_\mu g_{\rho\nu}\dot x^\mu\dot x^\nu
-\frac12\partial_\rho g_{\mu\nu}\dot x^\mu\dot x^\nu
-\frac{\dot\ell}{\ell}g_{\rho\nu}\dot x^\nu=0.
$}
\typeout{GRDISPLAY|106|\the\wd\grmathbox|\the\linewidth}
\ifdim\wd\grmathbox>\linewidth
\resizebox{\linewidth}{!}{\usebox{\grmathbox}}
\else\usebox{\grmathbox}\fi
\endgroup
\end{center}

The product \(\dot x^\mu\dot x^\nu\) is symmetric in \(\mu,\nu\).
Therefore the second term is unchanged if we replace its coefficient by
its symmetrized version:

\begin{center}
\begingroup
\sbox{\grmathbox}{$\displaystyle 
\partial_\mu g_{\rho\nu}\dot x^\mu\dot x^\nu
=\frac12\left(\partial_\mu g_{\rho\nu}
+\partial_\nu g_{\rho\mu}\right)\dot x^\mu\dot x^\nu.
$}
\typeout{GRDISPLAY|107|\the\wd\grmathbox|\the\linewidth}
\ifdim\wd\grmathbox>\linewidth
\resizebox{\linewidth}{!}{\usebox{\grmathbox}}
\else\usebox{\grmathbox}\fi
\endgroup
\end{center}

Multiply by \(g^{\alpha\rho}\) to remove the metric in front of
\(\ddot x^\nu\). The combination of metric derivatives that emerges is
denoted

\begin{center}
\begingroup
\sbox{\grmathbox}{$\displaystyle 
\Gamma^\alpha{}_{\mu\nu}
\equiv\frac12g^{\alpha\rho}
\left(\partial_\mu g_{\rho\nu}
+\partial_\nu g_{\rho\mu}
-\partial_\rho g_{\mu\nu}\right).
$}
\typeout{GRDISPLAY|108|\the\wd\grmathbox|\the\linewidth}
\ifdim\wd\grmathbox>\linewidth
\resizebox{\linewidth}{!}{\usebox{\grmathbox}}
\else\usebox{\grmathbox}\fi
\endgroup
\end{center}

These are the Christoffel symbols of the Levi-Civita connection. For now
they are the coefficients our variation has produced. Chapter 7 will
explain their geometric origin and why they transform differently from
tensor components.

Our result is

\begin{center}
\begingroup
\sbox{\grmathbox}{$\displaystyle 
\ddot x^\alpha+\Gamma^\alpha{}_{\mu\nu}\dot x^\mu\dot x^\nu
=\frac{d\ln\ell}{d\lambda}\dot x^\alpha.
$}
\typeout{GRDISPLAY|109|\the\wd\grmathbox|\the\linewidth}
\ifdim\wd\grmathbox>\linewidth
\resizebox{\linewidth}{!}{\usebox{\grmathbox}}
\else\usebox{\grmathbox}\fi
\endgroup
\end{center}

If we choose \(\lambda=\tau\), then \(\ell=c\), so its derivative
vanishes. We obtain

\begin{center}
\begingroup
\sbox{\grmathbox}{$\displaystyle 
\boxed{\frac{d^2x^\alpha}{d\tau^2}
+\Gamma^\alpha{}_{\mu\nu}
\frac{dx^\mu}{d\tau}\frac{dx^\nu}{d\tau}=0.}
$}
\typeout{GRDISPLAY|110|\the\wd\grmathbox|\the\linewidth}
\ifdim\wd\grmathbox>\linewidth
\resizebox{\linewidth}{!}{\usebox{\grmathbox}}
\else\usebox{\grmathbox}\fi
\endgroup
\end{center}

This is the timelike geodesic equation. It contains an ordinary
coordinate acceleration plus a correction describing how the local
coordinate basis and geometry change along the path. The whole
combination is the covariant acceleration, which is zero in free fall.

Notice what disappeared: the test particle\textquotesingle s mass. That
is the universality of free fall appearing in the variational
description.

\hypertarget{affine-parameters-the-route-and-the-speed-of-the-pen}{%
\section{5.6 Affine parameters: the route and the speed of the
pen}\label{affine-parameters-the-route-and-the-speed-of-the-pen}}

A curve is a route through spacetime. A parameter determines how quickly
your mathematical pen moves along that route. These are not the same
information.

The geodesic equation with zero right-hand side chooses an
\textbf{affine parameter}. For a timelike geodesic, proper time is
affine, and so is

\begin{center}
\begingroup
\sbox{\grmathbox}{$\displaystyle 
\lambda=A\tau+B
$}
\typeout{GRDISPLAY|111|\the\wd\grmathbox|\the\linewidth}
\ifdim\wd\grmathbox>\linewidth
\resizebox{\linewidth}{!}{\usebox{\grmathbox}}
\else\usebox{\grmathbox}\fi
\endgroup
\end{center}

for constants \(A\ne0\) and \(B\). An increasing parameter uses \(A>0\).
A general nonlinear relabeling introduces a term parallel to the
tangent, as the preceding derivation showed. Such a term changes the
rate at which the pen traverses the curve; it does not by itself bend
the route away from a geodesic.

One-dimensional flat-space example: the straight line \(x(s)=s\) has
\(d^2x/ds^2=0\). Relabel it using \(s=e^\lambda\). Now
\(x(\lambda)=e^\lambda\) and \(d^2x/d\lambda^2=dx/d\lambda\ne0\). The
line did not become geometrically curved. Its parameter became
nonaffine.

The same issue occurs if coordinate time \(t\) is used to parameterize a
relativistic geodesic. In a general spacetime it need not be affine, so
one must transform the equation correctly rather than replace every
\(\tau\) by \(t\) and hope. These parameter subtleties and the null case
are discussed in \href{https://arxiv.org/pdf/gr-qc/9712019}{Sean
Carroll\textquotesingle s notes, in the geodesics section}.

\hypertarget{null-curves-zero-clock-time-does-not-mean-zero-geometry}{%
\section{5.7 Null curves: zero clock time does not mean zero
geometry}\label{null-curves-zero-clock-time-does-not-mean-zero-geometry}}

For a light ray, \(ds^2=0\) and \(d\tau=0\). Proper time cannot
parameterize the path. We therefore use an affine parameter \(\lambda\)
and write

\begin{center}
\begingroup
\sbox{\grmathbox}{$\displaystyle 
k^\mu=\frac{dx^\mu}{d\lambda},
\qquad
g_{\mu\nu}k^\mu k^\nu=0,
$}
\typeout{GRDISPLAY|112|\the\wd\grmathbox|\the\linewidth}
\ifdim\wd\grmathbox>\linewidth
\resizebox{\linewidth}{!}{\usebox{\grmathbox}}
\else\usebox{\grmathbox}\fi
\endgroup
\end{center}

with

\begin{center}
\begingroup
\sbox{\grmathbox}{$\displaystyle 
\frac{dk^\mu}{d\lambda}
+\Gamma^\mu{}_{\alpha\beta}k^\alpha k^\beta=0.
$}
\typeout{GRDISPLAY|113|\the\wd\grmathbox|\the\linewidth}
\ifdim\wd\grmathbox>\linewidth
\resizebox{\linewidth}{!}{\usebox{\grmathbox}}
\else\usebox{\grmathbox}\fi
\endgroup
\end{center}

This governs vacuum light rays in the geometric-optics approximation to
Maxwell theory. A wave of finite wavelength has more structure than one
infinitely thin ray, and light in material media requires a different
propagation analysis.

For a null geodesic, changing the affine scale changes \(k^\mu\) without
changing the route. The ray\textquotesingle s geometry alone therefore
does not fix a photon\textquotesingle s energy. To identify the tangent
with a physically normalized wave-vector or momentum, supply frequency
or energy data from an observer.

Simply putting \(m=0\) into \(-mc^2\int d\tau\) makes the action vanish
and yields no equation. A correct variational shortcut is the quadratic
functional

\begin{center}
\begingroup
\sbox{\grmathbox}{$\displaystyle 
I_2=\frac12\int g_{\mu\nu}\dot x^\mu\dot x^\nu\,d\lambda.
$}
\typeout{GRDISPLAY|114|\the\wd\grmathbox|\the\linewidth}
\ifdim\wd\grmathbox>\linewidth
\resizebox{\linewidth}{!}{\usebox{\grmathbox}}
\else\usebox{\grmathbox}\fi
\endgroup
\end{center}

Varying over unrestricted nearby paths with fixed endpoints gives the
affine geodesic equation. Nullness is then selected by a null initial
tangent and is preserved along the resulting geodesic. \textbf{Do not
restrict the entire family of varied curves to have identically zero
integrand first} and then expect varying zero to provide dynamics.

A more systematic massless action introduces an auxiliary
one-dimensional field \(e(\lambda)\):

\begin{center}
\begingroup
\sbox{\grmathbox}{$\displaystyle 
I_0=\frac12\int e^{-1}g_{\mu\nu}\dot x^\mu\dot x^\nu\,d\lambda.
$}
\typeout{GRDISPLAY|115|\the\wd\grmathbox|\the\linewidth}
\ifdim\wd\grmathbox>\linewidth
\resizebox{\linewidth}{!}{\usebox{\grmathbox}}
\else\usebox{\grmathbox}\fi
\endgroup
\end{center}

Varying \(e\) imposes the null condition. Varying \(x\) gives

\begin{center}
\begingroup
\sbox{\grmathbox}{$\displaystyle 
\ddot x^\mu+\Gamma^\mu{}_{\alpha\beta}\dot x^\alpha\dot x^\beta
=\frac{\dot e}{e}\dot x^\mu.
$}
\typeout{GRDISPLAY|116|\the\wd\grmathbox|\the\linewidth}
\ifdim\wd\grmathbox>\linewidth
\resizebox{\linewidth}{!}{\usebox{\grmathbox}}
\else\usebox{\grmathbox}\fi
\endgroup
\end{center}

A reparameterization can make \(e\) constant, producing an affine
parameter. This small example previews a broad modern theme: an
auxiliary field can enforce a constraint without introducing an extra
physical propagating degree of freedom.

Finally, there is no ``photon\textquotesingle s perspective'' in the
sense of a valid inertial rest frame. The vanishing proper time of a
null curve is a statement about Lorentzian geometry, not a license to
assign consciousness or an ordinary clock to a frame traveling at \(c\).

\hypertarget{check-that-the-action-remembers-newton}{%
\section{5.8 Check that the action remembers
Newton}\label{check-that-the-action-remembers-newton}}

Chapter 12 will relate the weak gravitational field to the Newtonian
potential. Preview its slowly varying, weak-field result in a suitable
nearly Cartesian chart:

\begin{center}
\begingroup
\sbox{\grmathbox}{$\displaystyle 
ds^2\simeq-\left(1+\frac{2\Phi}{c^2}\right)c^2dt^2+d\mathbf x^2,
$}
\typeout{GRDISPLAY|117|\the\wd\grmathbox|\the\linewidth}
\ifdim\wd\grmathbox>\linewidth
\resizebox{\linewidth}{!}{\usebox{\grmathbox}}
\else\usebox{\grmathbox}\fi
\endgroup
\end{center}

where \(|\Phi|/c^2\ll1\), particle speeds obey \(v^2/c^2\ll1\), and
spatial-metric corrections affect the calculation below only at higher
combined order. The potential \(\Phi\) has units of velocity squared;
for a localized spherical mass, \(\Phi=-G_NM/r\) when its zero is chosen
at infinity.

The clock rate along a slow timelike path is

\begin{center}
\begingroup
\sbox{\grmathbox}{$\displaystyle 
\frac{d\tau}{dt}
\simeq\sqrt{1+\frac{2\Phi}{c^2}-\frac{v^2}{c^2}}
\simeq1+\frac{\Phi}{c^2}-\frac{v^2}{2c^2}.
$}
\typeout{GRDISPLAY|118|\the\wd\grmathbox|\the\linewidth}
\ifdim\wd\grmathbox>\linewidth
\resizebox{\linewidth}{!}{\usebox{\grmathbox}}
\else\usebox{\grmathbox}\fi
\endgroup
\end{center}

Insert this into the particle action:

\begin{center}
\begingroup
\sbox{\grmathbox}{$\displaystyle 
S_{\mathrm{particle}}
\simeq\int dt\left[-mc^2+\frac12mv^2-m\Phi\right].
$}
\typeout{GRDISPLAY|119|\the\wd\grmathbox|\the\linewidth}
\ifdim\wd\grmathbox>\linewidth
\resizebox{\linewidth}{!}{\usebox{\grmathbox}}
\else\usebox{\grmathbox}\fi
\endgroup
\end{center}

For fixed endpoint times, the constant rest-energy term does not affect
the path variation. The remaining Lagrangian is exactly the Newtonian
expression \(L=K-U\). Its Euler-Lagrange equation is

\begin{center}
\begingroup
\sbox{\grmathbox}{$\displaystyle 
m\frac{d^2\mathbf x}{dt^2}=-m\symbfit\nabla\Phi.
$}
\typeout{GRDISPLAY|120|\the\wd\grmathbox|\the\linewidth}
\ifdim\wd\grmathbox>\linewidth
\resizebox{\linewidth}{!}{\usebox{\grmathbox}}
\else\usebox{\grmathbox}\fi
\endgroup
\end{center}

The relativistic action has combined what Newton separated into kinetic
and potential terms. This is not a derivation of the gravitational field
equation: we supplied the appropriate weak-field metric as a preview. It
is a consistency check that, given that metric, the relativistic motion
law has the expected limit.

One should also resist the slogan ``objects fall toward slower time'' as
a universal replacement for GR. It can convey part of slow motion in a
static weak field. It does not contain spatial curvature, frame
dragging, null propagation in full generality, or the behavior of
time-dependent geometries.

\hypertarget{tides-the-gravity-a-falling-laboratory-cannot-remove}{%
\section{5.9 Tides: the gravity a falling laboratory cannot
remove}\label{tides-the-gravity-a-falling-laboratory-cannot-remove}}

Take two nearby freely falling particles in Newtonian gravity, separated
by \(\xi^i\). Their individual accelerations are approximately
\(a^i=-\delta^{ij}\partial_j\Phi\). Subtract their equations and
Taylor-expand the acceleration of the second particle about the first:

\begin{center}
\begingroup
\sbox{\grmathbox}{$\displaystyle 
\frac{d^2\xi^i}{dt^2}
=-\delta^{ik}\partial_k\partial_j\Phi\,\xi^j
+O(|\symbfit\xi|^2).
$}
\typeout{GRDISPLAY|121|\the\wd\grmathbox|\the\linewidth}
\ifdim\wd\grmathbox>\linewidth
\resizebox{\linewidth}{!}{\usebox{\grmathbox}}
\else\usebox{\grmathbox}\fi
\endgroup
\end{center}

The gradient of \(\Phi\) controls the common acceleration. Its Hessian,
the matrix of second derivatives, controls the relative acceleration. A
falling frame can cancel the former at its origin; it cannot generally
cancel the latter throughout its neighborhood.

For \(\Phi=-G_NM/r\),

\begin{center}
\begingroup
\sbox{\grmathbox}{$\displaystyle 
\frac{d\Phi}{dr}=\frac{G_NM}{r^2},
\qquad
\frac{d^2\Phi}{dr^2}=-\frac{2G_NM}{r^3}.
$}
\typeout{GRDISPLAY|122|\the\wd\grmathbox|\the\linewidth}
\ifdim\wd\grmathbox>\linewidth
\resizebox{\linewidth}{!}{\usebox{\grmathbox}}
\else\usebox{\grmathbox}\fi
\endgroup
\end{center}

Two radially separated falling particles therefore have relative radial
acceleration

\begin{center}
\begingroup
\sbox{\grmathbox}{$\displaystyle 
\frac{d^2\xi^{\hat r}}{dt^2}
\simeq\frac{2G_NM}{r^3}\xi^{\hat r}.
$}
\typeout{GRDISPLAY|123|\the\wd\grmathbox|\the\linewidth}
\ifdim\wd\grmathbox>\linewidth
\resizebox{\linewidth}{!}{\usebox{\grmathbox}}
\else\usebox{\grmathbox}\fi
\endgroup
\end{center}

They stretch apart radially because the lower one falls more strongly.
Neighboring side-by-side particles instead converge toward the central
mass. A cloud of freely falling beads changes shape even though each
bead\textquotesingle s ideal accelerometer reads zero.

This supplies an operational distinction:

\begingroup\small\setlength{\tabcolsep}{5pt}\renewcommand{\arraystretch}{1.13}

\begin{longtable}[]{@{}
  >{\raggedright\arraybackslash}p{(\columnwidth - 4\tabcolsep) * \real{0.2500}}
  >{\raggedright\arraybackslash}p{(\columnwidth - 4\tabcolsep) * \real{0.3600}}
  >{\raggedright\arraybackslash}p{(\columnwidth - 4\tabcolsep) * \real{0.3900}}@{}}
\toprule\noalign{}
\begin{minipage}[b]{\linewidth}\raggedright
Question
\end{minipage} & \begin{minipage}[b]{\linewidth}\raggedright
Instrument or comparison
\end{minipage} & \begin{minipage}[b]{\linewidth}\raggedright
Geometric quantity to come
\end{minipage} \\
\midrule\noalign{}
\endhead
\bottomrule\noalign{}
\endlastfoot
Am I being pushed away from free fall? & One ideal accelerometer &
Proper acceleration of my worldline \\
How do these coordinate components change along my path? & A coordinate
description and comparison rule & Connection coefficients \\
Do neighboring free-fall trajectories converge, diverge, or shear? & A
finite arrangement of freely falling bodies & Spacetime curvature \\
\end{longtable}

\endgroup

The next chapters build the language needed to express the last row
without borrowing a preferred Newtonian frame. The connection will tell
us how to compare directions. Curvature will tell us why those
comparisons can fail to fit together around a loop.
Einstein\textquotesingle s equation will then say how that geometry is
related to matter, energy, momentum, and stress.

\hypertarget{chapter-6}{%
\chapter{6. Differentiation when your measuring axes will not sit
still}\label{chapter-6}}

\hypertarget{the-crime-an-innocent-derivative-produces-a-coordinate-dependent-answer}{%
\section{6.1 The crime: an innocent derivative produces a
coordinate-dependent
answer}\label{the-crime-an-innocent-derivative-produces-a-coordinate-dependent-answer}}

Imagine an ant walking across a flat dinner plate. At each location, a
mischievous waiter draws a new pair of coordinate arrows, rotated
relative to the previous pair. A vector that points toward the kitchen
everywhere acquires changing components in these arrows. If the ant
differentiates only those components, it announces that the vector field
is turning. The kitchen, understandably, denies everything.

This is already a problem in flat space. Curvature makes the problem
deeper, but does not create it. Coordinates can change their scale and
orientation from point to point even on a perfectly ordinary plane.

Let a vector field have components \(V^\alpha\) in coordinates
\(x^\mu\). Under new coordinates \(x'^\alpha(x)\), define the Jacobian
and its inverse by

\begin{center}
\begingroup
\sbox{\grmathbox}{$\displaystyle 
J^\alpha{}_{\beta}=\frac{\partial x'^\alpha}{\partial x^\beta},
\qquad
K^\nu{}_{\mu}=\frac{\partial x^\nu}{\partial x'^\mu}.
$}
\typeout{GRDISPLAY|124|\the\wd\grmathbox|\the\linewidth}
\ifdim\wd\grmathbox>\linewidth
\resizebox{\linewidth}{!}{\usebox{\grmathbox}}
\else\usebox{\grmathbox}\fi
\endgroup
\end{center}

Then \(V'^\alpha=J^\alpha{}_{\beta}V^\beta\). Apply the chain rule and
the ordinary product rule:

\begin{center}
\begingroup
\sbox{\grmathbox}{$\displaystyle 
\partial'_\mu V'^\alpha
=
K^\nu{}_{\mu}J^\alpha{}_{\beta}\partial_\nu V^\beta
+
K^\nu{}_{\mu}(\partial_\nu J^\alpha{}_{\beta})V^\beta.
$}
\typeout{GRDISPLAY|125|\the\wd\grmathbox|\the\linewidth}
\ifdim\wd\grmathbox>\linewidth
\resizebox{\linewidth}{!}{\usebox{\grmathbox}}
\else\usebox{\grmathbox}\fi
\endgroup
\end{center}

The first term is exactly how a tensor with one upper and one lower
index should transform. The second term is the trouble: it involves
derivatives of the coordinate transformation itself. A tensor
transformation changes the description of an object using the Jacobian
at the point; it does not need second derivatives of the map.

Notice when the trouble disappears. A Cartesian rotation or a Lorentz
transformation has a constant Jacobian, so the extra term vanishes. This
is why ordinary component differentiation worked so comfortably in
introductory mechanics and special relativity. Those courses usually
handed you unusually well-behaved rulers.

The solution is to differentiate the geometric vector, accounting for
the changing local basis. We call the resulting operation the
\textbf{covariant derivative}.

\hypertarget{what-a-derivative-must-be-able-to-compare}{%
\section{6.2 What a derivative must be able to
compare}\label{what-a-derivative-must-be-able-to-compare}}

There is a subtlety even before the algebra. A vector at event \(p\)
belongs to \(T_pM\), the tangent space at \(p\). A vector at a
neighboring event \(q\) belongs to \(T_qM\). These are different vector
spaces. Subtracting their component lists does not, by itself, define a
geometric subtraction.

Think of two bank balances reported in currencies whose exchange rate
changes between branches. Subtracting the printed numbers is possible;
assigning an invariant meaning to the answer requires a comparison rule.
Here the rule is a \textbf{connection}. The analogy is limited: a
spacetime connection can mix directions, and comparisons generally
depend on the path taken. There is no universal currency conversion
table for separated tangent spaces.

In a coordinate basis \(e_\nu=\partial/\partial x^\nu\), define
connection coefficients by

\begin{center}
\begingroup
\sbox{\grmathbox}{$\displaystyle 
\nabla_\mu e_\nu=\Gamma^\rho{}_{\mu\nu}e_\rho.
$}
\typeout{GRDISPLAY|126|\the\wd\grmathbox|\the\linewidth}
\ifdim\wd\grmathbox>\linewidth
\resizebox{\linewidth}{!}{\usebox{\grmathbox}}
\else\usebox{\grmathbox}\fi
\endgroup
\end{center}

This does not mean that ordinary subtraction of basis vectors at
different points was secretly available. The symbol \(\nabla\) specifies
the comparison rule that makes such a derivative meaningful.

Apply the product rule to \(V=V^\nu e_\nu\):

\begin{center}
\begingroup
\sbox{\grmathbox}{$\displaystyle 
\begin{aligned}
\nabla_\mu V
&=\partial_\mu V^\nu\,e_\nu+V^\nu\nabla_\mu e_\nu\\
&=\left(\partial_\mu V^\rho+\Gamma^\rho{}_{\mu\nu}V^\nu\right)e_\rho.
\end{aligned}
$}
\typeout{GRDISPLAY|127|\the\wd\grmathbox|\the\linewidth}
\ifdim\wd\grmathbox>\linewidth
\resizebox{\linewidth}{!}{\usebox{\grmathbox}}
\else\usebox{\grmathbox}\fi
\endgroup
\end{center}

Therefore

\begin{center}
\begingroup
\sbox{\grmathbox}{$\displaystyle 
\boxed{\nabla_\mu V^\rho
=\partial_\mu V^\rho+\Gamma^\rho{}_{\mu\nu}V^\nu.}
$}
\typeout{GRDISPLAY|128|\the\wd\grmathbox|\the\linewidth}
\ifdim\wd\grmathbox>\linewidth
\resizebox{\linewidth}{!}{\usebox{\grmathbox}}
\else\usebox{\grmathbox}\fi
\endgroup
\end{center}

The first term measures changing components. The second corrects for the
comparison of local bases. Neither term separately has to be a tensor;
their sum does.

The derivative index \(\mu\) is a lower index because differentiation
asks for a direction as its input. If \(X^\mu\) specifies that
direction, then \(\nabla_XV=X^\mu\nabla_\mu V\) is a vector. Before
choosing \(X\), \(\nabla V\) is a tensor with an extra covector slot
waiting to receive it.

\hypertarget{why-covectors-acquire-a-minus-sign}{%
\section{6.3 Why covectors acquire a minus
sign}\label{why-covectors-acquire-a-minus-sign}}

A covector \(\omega_\nu\) eats a vector and produces a scalar:

\begin{center}
\begingroup
\sbox{\grmathbox}{$\displaystyle 
f=\omega_\nu V^\nu.
$}
\typeout{GRDISPLAY|129|\the\wd\grmathbox|\the\linewidth}
\ifdim\wd\grmathbox>\linewidth
\resizebox{\linewidth}{!}{\usebox{\grmathbox}}
\else\usebox{\grmathbox}\fi
\endgroup
\end{center}

A scalar has no moving basis indices to correct, so
\(\nabla_\mu f=\partial_\mu f\). We also require a derivative to obey
the product rule and respect contractions:

\begin{center}
\begingroup
\sbox{\grmathbox}{$\displaystyle 
\partial_\mu(\omega_\nu V^\nu)
=(\nabla_\mu\omega_\nu)V^\nu
+\omega_\nu\nabla_\mu V^\nu.
$}
\typeout{GRDISPLAY|130|\the\wd\grmathbox|\the\linewidth}
\ifdim\wd\grmathbox>\linewidth
\resizebox{\linewidth}{!}{\usebox{\grmathbox}}
\else\usebox{\grmathbox}\fi
\endgroup
\end{center}

Substitute the vector derivative. The final term contributes
\(+\omega_\nu\Gamma^\nu{}_{\mu\lambda}V^\lambda\). There is no
corresponding connection term in the ordinary derivative of the scalar
on the left. The covector derivative must cancel it, for every possible
vector. Thus

\begin{center}
\begingroup
\sbox{\grmathbox}{$\displaystyle 
\boxed{\nabla_\mu\omega_\nu
=\partial_\mu\omega_\nu-\Gamma^\lambda{}_{\mu\nu}\omega_\lambda.}
$}
\typeout{GRDISPLAY|131|\the\wd\grmathbox|\the\linewidth}
\ifdim\wd\grmathbox>\linewidth
\resizebox{\linewidth}{!}{\usebox{\grmathbox}}
\else\usebox{\grmathbox}\fi
\endgroup
\end{center}

The minus sign is not a mnemonic imposed by an unfriendly textbook. It
is the price of preserving the scalar pairing. Vectors and covectors
change in mutually compensating ways.

For a general tensor, every upper index gets a plus correction and every
lower index gets a minus correction. For example,

\begin{center}
\begingroup
\sbox{\grmathbox}{$\displaystyle 
\begin{aligned}
\nabla_\mu T^{\alpha\beta}{}_{\gamma}
={}&\partial_\mu T^{\alpha\beta}{}_{\gamma}
+\Gamma^\alpha{}_{\mu\lambda}T^{\lambda\beta}{}_{\gamma}
+\Gamma^\beta{}_{\mu\lambda}T^{\alpha\lambda}{}_{\gamma}\\
&-\Gamma^\lambda{}_{\mu\gamma}T^{\alpha\beta}{}_{\lambda}.
\end{aligned}
$}
\typeout{GRDISPLAY|132|\the\wd\grmathbox|\the\linewidth}
\ifdim\wd\grmathbox>\linewidth
\resizebox{\linewidth}{!}{\usebox{\grmathbox}}
\else\usebox{\grmathbox}\fi
\endgroup
\end{center}

Each correction replaces exactly one index with a summed index. Every
term retains the same free indices \(\mu,\alpha,\beta,\gamma\). This is
an excellent way to catch mistakes before they breed.

Respecting contractions means, for example,

\begin{center}
\begingroup
\sbox{\grmathbox}{$\displaystyle 
\nabla_\mu(T^\alpha{}_{\alpha})
=(\nabla_\mu T)^\alpha{}_{\alpha}.
$}
\typeout{GRDISPLAY|133|\the\wd\grmathbox|\the\linewidth}
\ifdim\wd\grmathbox>\linewidth
\resizebox{\linewidth}{!}{\usebox{\grmathbox}}
\else\usebox{\grmathbox}\fi
\endgroup
\end{center}

The two connection corrections cancel after a dummy-index relabeling.
However, \textbf{raising an index is an additional operation involving
the metric}. Commuting differentiation with raising and lowering
requires metric compatibility, \(\nabla g=0\), which we will derive into
the story shortly. Compatibility with vector-covector contraction and
compatibility with the metric are related ideas, but are not identical
assumptions.

\hypertarget{the-scalar-trap-one-derivative-is-easy-two-derivatives-are-not}{%
\section{6.4 The scalar trap: one derivative is easy; two derivatives
are
not}\label{the-scalar-trap-one-derivative-is-easy-two-derivatives-are-not}}

For a scalar field \(f\),

\begin{center}
\begingroup
\sbox{\grmathbox}{$\displaystyle 
\nabla_\mu f=\partial_\mu f.
$}
\typeout{GRDISPLAY|134|\the\wd\grmathbox|\the\linewidth}
\ifdim\wd\grmathbox>\linewidth
\resizebox{\linewidth}{!}{\usebox{\grmathbox}}
\else\usebox{\grmathbox}\fi
\endgroup
\end{center}

It is tempting to conclude that every derivative of a scalar is ordinary
differentiation. The trap is that the first derivative is now a
\textbf{covector}. Differentiating that covector gives

\begin{center}
\begingroup
\sbox{\grmathbox}{$\displaystyle 
\boxed{\nabla_\mu\nabla_\nu f
=\partial_\mu\partial_\nu f
-\Gamma^\lambda{}_{\mu\nu}\partial_\lambda f.}
$}
\typeout{GRDISPLAY|135|\the\wd\grmathbox|\the\linewidth}
\ifdim\wd\grmathbox>\linewidth
\resizebox{\linewidth}{!}{\usebox{\grmathbox}}
\else\usebox{\grmathbox}\fi
\endgroup
\end{center}

This is the covariant Hessian. With a torsion-free connection, it is
symmetric in \(\mu,\nu\). Ordinary second partial derivatives are
symmetric too, but generally fail to transform as a tensor. Symmetry
alone is not a certificate of geometric respectability.

There is one useful special case. At a critical point, where
\(\partial_\lambda f=0\), the correction vanishes. The Hessian computed
there with ordinary second derivatives has a coordinate-independent
meaning as a bilinear form. This is why classifying a stationary point
as a maximum, minimum, or saddle can be done intrinsically despite using
coordinate derivatives.

A related distinction: \(df\), with components \(\partial_\mu f\), is a
covector defined without a metric. The \textbf{gradient vector} is

\begin{center}
\begingroup
\sbox{\grmathbox}{$\displaystyle 
(\operatorname{grad}f)^\mu=g^{\mu\nu}\partial_\nu f,
$}
\typeout{GRDISPLAY|136|\the\wd\grmathbox|\the\linewidth}
\ifdim\wd\grmathbox>\linewidth
\resizebox{\linewidth}{!}{\usebox{\grmathbox}}
\else\usebox{\grmathbox}\fi
\endgroup
\end{center}

and does require a metric. In Euclidean space these objects are often
merged into one mental image. In spacetime, raising a time index even
changes a sign in an orthonormal frame. The distinction earns its keep.

\hypertarget{divergence-and-the-hidden-accounting-of-physical-volume}{%
\section{6.5 Divergence and the hidden accounting of physical
volume}\label{divergence-and-the-hidden-accounting-of-physical-volume}}

Using the Levi-Civita connection introduced in the next chapter, the
divergence of a vector is

\begin{center}
\begingroup
\sbox{\grmathbox}{$\displaystyle 
\nabla_\mu V^\mu
=\partial_\mu V^\mu+\Gamma^\mu{}_{\mu\nu}V^\nu.
$}
\typeout{GRDISPLAY|137|\the\wd\grmathbox|\the\linewidth}
\ifdim\wd\grmathbox>\linewidth
\resizebox{\linewidth}{!}{\usebox{\grmathbox}}
\else\usebox{\grmathbox}\fi
\endgroup
\end{center}

The trace of the connection simplifies dramatically:

\begin{center}
\begingroup
\sbox{\grmathbox}{$\displaystyle 
\Gamma^\mu{}_{\mu\nu}
=\frac12g^{\alpha\beta}\partial_\nu g_{\alpha\beta}
=\partial_\nu\ln\sqrt{-g},
\qquad g=\det(g_{\alpha\beta}).
$}
\typeout{GRDISPLAY|138|\the\wd\grmathbox|\the\linewidth}
\ifdim\wd\grmathbox>\linewidth
\resizebox{\linewidth}{!}{\usebox{\grmathbox}}
\else\usebox{\grmathbox}\fi
\endgroup
\end{center}

The last step uses a matrix identity. For an invertible matrix \(A\), an
infinitesimal change satisfies
\(\delta\ln|\det A|=\operatorname{tr}(A^{-1}\delta A)\). Apply this to
the metric, then take half because of the square root. In a Lorentzian
chart with our signature, \(g<0\).

Combining the two terms by the product rule gives

\begin{center}
\begingroup
\sbox{\grmathbox}{$\displaystyle 
\boxed{\nabla_\mu V^\mu
=\frac{1}{\sqrt{-g}}\partial_\mu\left(\sqrt{-g}\,V^\mu\right).}
$}
\typeout{GRDISPLAY|139|\the\wd\grmathbox|\the\linewidth}
\ifdim\wd\grmathbox>\linewidth
\resizebox{\linewidth}{!}{\usebox{\grmathbox}}
\else\usebox{\grmathbox}\fi
\endgroup
\end{center}

Why should the determinant appear? A coordinate box of side lengths
\(dx^\mu\) represents physical four-volume \(\sqrt{-g}\,d^4x\). A flow
can have changing coordinate components merely because the coordinate
boxes expand or shrink. Divergence measures net outflow per physical
volume, so it must include that change in the measuring boxes.

In an \(n\)-dimensional Riemannian space, replace \(\sqrt{-g}\) by
\(\sqrt{g}\); the general expression uses \(\sqrt{|g|}\). The idea is
the same.

Applying divergence to the gradient of a scalar yields the
curved-spacetime wave operator:

\begin{center}
\begingroup
\sbox{\grmathbox}{$\displaystyle 
\boxed{\Box f
=\nabla_\mu\nabla^\mu f
=\frac{1}{\sqrt{-g}}\partial_\mu
\left(\sqrt{-g}\,g^{\mu\nu}\partial_\nu f\right).}
$}
\typeout{GRDISPLAY|140|\the\wd\grmathbox|\the\linewidth}
\ifdim\wd\grmathbox>\linewidth
\resizebox{\linewidth}{!}{\usebox{\grmathbox}}
\else\usebox{\grmathbox}\fi
\endgroup
\end{center}

In Cartesian Minkowski coordinates with \(x^0=ct\), this becomes
\(-c^{-2}\partial_t^2f+\partial_x^2f+\partial_y^2f+\partial_z^2f\). The
ordinary wave equation was already a geometric equation; its familiar
coordinates hid the machinery.

One warning will matter for stress-energy: a rank-two tensor has another
index to correct. In particular,

\begin{center}
\begingroup
\sbox{\grmathbox}{$\displaystyle 
\nabla_\mu T^{\mu\nu}
=\frac{1}{\sqrt{-g}}\partial_\mu
\left(\sqrt{-g}\,T^{\mu\nu}\right)
+\Gamma^\nu{}_{\mu\lambda}T^{\mu\lambda}.
$}
\typeout{GRDISPLAY|141|\the\wd\grmathbox|\the\linewidth}
\ifdim\wd\grmathbox>\linewidth
\resizebox{\linewidth}{!}{\usebox{\grmathbox}}
\else\usebox{\grmathbox}\fi
\endgroup
\end{center}

Leaving off the last term does not create a clever conservation law. It
creates an incorrect equation.

\hypertarget{chapter-7}{%
\chapter{7. The connection: how neighboring laboratories compare
directions}\label{chapter-7}}

\hypertarget{parallel-transport-is-an-instruction-not-a-picture}{%
\section{7.1 Parallel transport is an instruction, not a
picture}\label{parallel-transport-is-an-instruction-not-a-picture}}

Take a curve \(x^\mu(\lambda)\) and a vector \(V^\mu(\lambda)\) attached
to its points. The covariant derivative along the curve is

\begin{center}
\begingroup
\sbox{\grmathbox}{$\displaystyle 
\frac{DV^\mu}{d\lambda}
=\frac{dV^\mu}{d\lambda}
+\Gamma^\mu{}_{\alpha\beta}
\frac{dx^\alpha}{d\lambda}V^\beta.
$}
\typeout{GRDISPLAY|142|\the\wd\grmathbox|\the\linewidth}
\ifdim\wd\grmathbox>\linewidth
\resizebox{\linewidth}{!}{\usebox{\grmathbox}}
\else\usebox{\grmathbox}\fi
\endgroup
\end{center}

We say \(V\) is \textbf{parallel transported} when this derivative
vanishes:

\begin{center}
\begingroup
\sbox{\grmathbox}{$\displaystyle 
\frac{DV^\mu}{d\lambda}=0.
$}
\typeout{GRDISPLAY|143|\the\wd\grmathbox|\the\linewidth}
\ifdim\wd\grmathbox>\linewidth
\resizebox{\linewidth}{!}{\usebox{\grmathbox}}
\else\usebox{\grmathbox}\fi
\endgroup
\end{center}

Given an initial vector and a smooth specified path, this is a linear
first-order differential equation with a locally unique solution.
Parallel transport is therefore an operational prescription: solve this
equation while moving along the path.

If the transported vector is the tangent to the path itself, the
prescription becomes

\begin{center}
\begingroup
\sbox{\grmathbox}{$\displaystyle 
\frac{D}{d\lambda}\frac{dx^\mu}{d\lambda}=0,
$}
\typeout{GRDISPLAY|144|\the\wd\grmathbox|\the\linewidth}
\ifdim\wd\grmathbox>\linewidth
\resizebox{\linewidth}{!}{\usebox{\grmathbox}}
\else\usebox{\grmathbox}\fi
\endgroup
\end{center}

which is the affinely parameterized geodesic equation. A geodesic
transports its own direction. This supplies a precise meaning of ``as
straight as possible'' that does not require drawing the curve inside a
larger space.

A freely falling, torque-free gyroscope provides a physical way to think
about parallel transport of orientation along a timelike geodesic. An
accelerated observer needs a related rule, Fermi-Walker transport, if
the goal is a nonrotating frame. Simply saying ``a gyroscope always
follows the connection'' skips the effect of the
observer\textquotesingle s acceleration.

\hypertarget{two-requirements-that-select-ordinary-grs-connection}{%
\section{7.2 Two requirements that select ordinary GR\textquotesingle s
connection}\label{two-requirements-that-select-ordinary-grs-connection}}

A manifold can carry many connections. Standard metric GR chooses the
\textbf{Levi-Civita connection}, characterized by two conditions.

First, \textbf{metric compatibility}:

\begin{center}
\begingroup
\sbox{\grmathbox}{$\displaystyle 
\nabla_\lambda g_{\mu\nu}=0.
$}
\typeout{GRDISPLAY|145|\the\wd\grmathbox|\the\linewidth}
\ifdim\wd\grmathbox>\linewidth
\resizebox{\linewidth}{!}{\usebox{\grmathbox}}
\else\usebox{\grmathbox}\fi
\endgroup
\end{center}

This says the connection preserves the metric\textquotesingle s inner
products. If \(V\) and \(W\) are parallel transported along the same
curve, the product rule gives

\begin{center}
\begingroup
\sbox{\grmathbox}{$\displaystyle 
\frac{d}{d\lambda}g(V,W)
=(\nabla_{\dot x}g)(V,W)
+g(\nabla_{\dot x}V,W)+g(V,\nabla_{\dot x}W)=0.
$}
\typeout{GRDISPLAY|146|\the\wd\grmathbox|\the\linewidth}
\ifdim\wd\grmathbox>\linewidth
\resizebox{\linewidth}{!}{\usebox{\grmathbox}}
\else\usebox{\grmathbox}\fi
\endgroup
\end{center}

Thus their lengths and mutual inner product remain fixed under
transport. In Lorentzian geometry, causal character is preserved too:
timelike vectors remain timelike, and null vectors remain null.

The condition does \textbf{not} say \(\partial_\lambda g_{\mu\nu}=0\)
everywhere. Metric components may vary because the coordinates vary.
Compatibility says the connection accounts for that variation
consistently.

Second, \textbf{zero torsion}. For vector fields \(X,Y\), torsion is

\begin{center}
\begingroup
\sbox{\grmathbox}{$\displaystyle 
\mathcal T(X,Y)=\nabla_XY-\nabla_YX-[X,Y],
$}
\typeout{GRDISPLAY|147|\the\wd\grmathbox|\the\linewidth}
\ifdim\wd\grmathbox>\linewidth
\resizebox{\linewidth}{!}{\usebox{\grmathbox}}
\else\usebox{\grmathbox}\fi
\endgroup
\end{center}

where \([X,Y]\) is their Lie bracket. The subtraction removes the
failure of the vector fields themselves to form commuting coordinate
directions. In a coordinate basis, \([\partial_\mu,\partial_\nu]=0\), so

\begin{center}
\begingroup
\sbox{\grmathbox}{$\displaystyle 
\mathcal T^\rho{}_{\mu\nu}
=\Gamma^\rho{}_{\mu\nu}-\Gamma^\rho{}_{\nu\mu}.
$}
\typeout{GRDISPLAY|148|\the\wd\grmathbox|\the\linewidth}
\ifdim\wd\grmathbox>\linewidth
\resizebox{\linewidth}{!}{\usebox{\grmathbox}}
\else\usebox{\grmathbox}\fi
\endgroup
\end{center}

Torsion-free therefore means symmetry of the lower two connection
indices \textbf{in a coordinate basis}. In a noncoordinate frame,
noncommuting basis vectors contribute to torsion too. Forgetting that
qualification becomes painful when tetrads appear later.

An intuitive distinction is useful. Curvature concerns a mismatch in
transported orientation around a loop. Torsion concerns a different
infinitesimal closure defect associated with transporting directions to
build a parallelogram. Neither is ``how twisty a coordinate grid
looks.'' Torsion-free is a structural choice of ordinary GR, not a
mathematical necessity for all theories of gravity.

\hypertarget{deriving-the-christoffel-symbols-instead-of-receiving-them-as-a-curse}{%
\section{7.3 Deriving the Christoffel symbols instead of receiving them
as a
curse}\label{deriving-the-christoffel-symbols-instead-of-receiving-them-as-a-curse}}

Expand metric compatibility:

\begin{center}
\begingroup
\sbox{\grmathbox}{$\displaystyle 
\partial_\mu g_{\nu\sigma}
=g_{\lambda\sigma}\Gamma^\lambda{}_{\mu\nu}
+g_{\nu\lambda}\Gamma^\lambda{}_{\mu\sigma}.
$}
\typeout{GRDISPLAY|149|\the\wd\grmathbox|\the\linewidth}
\ifdim\wd\grmathbox>\linewidth
\resizebox{\linewidth}{!}{\usebox{\grmathbox}}
\else\usebox{\grmathbox}\fi
\endgroup
\end{center}

There are two connection terms because the metric has two lower indices.
Now write the same statement with the indices permuted:

\begin{center}
\begingroup
\sbox{\grmathbox}{$\displaystyle 
\begin{aligned}
\partial_\nu g_{\mu\sigma}
&=g_{\lambda\sigma}\Gamma^\lambda{}_{\nu\mu}
+g_{\mu\lambda}\Gamma^\lambda{}_{\nu\sigma},\\
\partial_\sigma g_{\mu\nu}
&=g_{\lambda\nu}\Gamma^\lambda{}_{\sigma\mu}
+g_{\mu\lambda}\Gamma^\lambda{}_{\sigma\nu}.
\end{aligned}
$}
\typeout{GRDISPLAY|150|\the\wd\grmathbox|\the\linewidth}
\ifdim\wd\grmathbox>\linewidth
\resizebox{\linewidth}{!}{\usebox{\grmathbox}}
\else\usebox{\grmathbox}\fi
\endgroup
\end{center}

Add the first two equations and subtract the third. Why this particular
maneuver? We want to isolate the connection with lower indices
\(\mu\nu\). The unwanted terms pair off because torsion-free symmetry
lets us exchange the lower two indices of \(\Gamma\):

\begin{center}
\begingroup
\sbox{\grmathbox}{$\displaystyle 
\partial_\mu g_{\nu\sigma}
+\partial_\nu g_{\mu\sigma}
-\partial_\sigma g_{\mu\nu}
=2g_{\sigma\lambda}\Gamma^\lambda{}_{\mu\nu}.
$}
\typeout{GRDISPLAY|151|\the\wd\grmathbox|\the\linewidth}
\ifdim\wd\grmathbox>\linewidth
\resizebox{\linewidth}{!}{\usebox{\grmathbox}}
\else\usebox{\grmathbox}\fi
\endgroup
\end{center}

The inverse metric removes \(g_{\sigma\lambda}\). Multiply by
\(\tfrac12g^{\rho\sigma}\):

\begin{center}
\begingroup
\sbox{\grmathbox}{$\displaystyle 
\boxed{
\Gamma^\rho{}_{\mu\nu}
=\frac12g^{\rho\sigma}
\left(
\partial_\mu g_{\nu\sigma}
+\partial_\nu g_{\mu\sigma}
-\partial_\sigma g_{\mu\nu}
\right).
}
$}
\typeout{GRDISPLAY|152|\the\wd\grmathbox|\the\linewidth}
\ifdim\wd\grmathbox>\linewidth
\resizebox{\linewidth}{!}{\usebox{\grmathbox}}
\else\usebox{\grmathbox}\fi
\endgroup
\end{center}

This proves uniqueness: any torsion-free, metric-compatible connection
must have these coefficients. Existence follows by checking the formula:
it is symmetric in \(\mu\nu\), substitution gives \(\nabla g=0\), and
the chain rule supplies the connection transformation law below. There
is exactly one such connection for every smooth nondegenerate metric.

The formula is a three-term balancing act. Two terms measure changes in
the metric along the directions we are comparing; the third prevents
counting the same metric variation twice in the wrong slot. It is less
mysterious when remembered as the solution of three linear equations.

Metric compatibility also gives \(\nabla_\lambda g^{\mu\nu}=0\).
Differentiate \(g^{\mu\alpha}g_{\alpha\nu}=\delta^\mu{}_{\nu}\), use the
product rule and \(\nabla g=0\), and multiply by the inverse metric.
Raising and lowering indices now commute with covariant differentiation.

\hypertarget{why-a-connection-is-geometric-but-its-coefficients-are-not-a-tensor}{%
\section{7.4 Why a connection is geometric but its coefficients are not
a
tensor}\label{why-a-connection-is-geometric-but-its-coefficients-are-not-a-tensor}}

Under a coordinate change,

\begin{center}
\begingroup
\sbox{\grmathbox}{$\displaystyle 
\boxed{
\begin{aligned}
\Gamma'^\alpha{}_{\mu\nu}
={}&\frac{\partial x'^\alpha}{\partial x^\rho}
\frac{\partial x^\sigma}{\partial x'^\mu}
\frac{\partial x^\lambda}{\partial x'^\nu}
\Gamma^\rho{}_{\sigma\lambda}\\
&+\frac{\partial x'^\alpha}{\partial x^\rho}
\frac{\partial^2x^\rho}{\partial x'^\mu\partial x'^\nu}.
\end{aligned}}
$}
\typeout{GRDISPLAY|153|\the\wd\grmathbox|\the\linewidth}
\ifdim\wd\grmathbox>\linewidth
\resizebox{\linewidth}{!}{\usebox{\grmathbox}}
\else\usebox{\grmathbox}\fi
\endgroup
\end{center}

The first line resembles the transformation of a \((1,2)\) tensor. The
second line is the essential extra term.

You can derive it directly from
\(e'_\nu=(\partial x^\lambda/\partial x'^\nu)e_\lambda\). Apply
\(\nabla_{e'_\mu}\) and use the product rule. One contribution
differentiates the old basis through the old connection; the other
differentiates the coordinate-dependent coefficient multiplying that
basis. Express the result in the primed basis. These are exactly the two
lines above.

The inhomogeneous term cancels the unwanted second derivatives in
\(\partial'_\mu V'^\alpha\). A connection\textquotesingle s coefficients
are non-tensorial for a very good reason: they must repair a
non-tensorial partial derivative.

A tensor that vanishes in one coordinate system at a point vanishes in
every coordinate system there. Connection coefficients can vanish at a
point in one system and be nonzero in another. They therefore cannot
themselves be a tensor measuring gravitational curvature.

Nevertheless, the \textbf{difference of two connections} is a tensor. If
\(A^\rho{}_{\mu\nu}=\Gamma^\rho{}_{\mu\nu}-\widetilde\Gamma^\rho{}_{\mu\nu}\),
the second-derivative terms cancel under transformation. This fact
underlies comparisons between a background connection and a perturbed
connection. Also, \(\delta\Gamma\) in a metric variation is tensorial
when comparing connections on the same manifold with the same coordinate
identification. That observation will become useful in the action
derivation.

For additional derivations of connections and their relation to
transport, see
\href{https://ned.ipac.caltech.edu/level5/March01/Carroll3/Carroll3.html}{Sean
Carroll\textquotesingle s university lecture notes, ``Curvature''}. The
calculations here use the conventions stated in this book.

\hypertarget{a-complete-flat-space-laboratory-polar-coordinates}{%
\section{7.5 A complete flat-space laboratory: polar
coordinates}\label{a-complete-flat-space-laboratory-polar-coordinates}}

Let us make the distinction between connection and curvature impossible
to forget. Work on an ordinary Euclidean plane away from the origin,
with

\begin{center}
\begingroup
\sbox{\grmathbox}{$\displaystyle 
x=r\cos\theta,\qquad y=r\sin\theta,
$}
\typeout{GRDISPLAY|154|\the\wd\grmathbox|\the\linewidth}
\ifdim\wd\grmathbox>\linewidth
\resizebox{\linewidth}{!}{\usebox{\grmathbox}}
\else\usebox{\grmathbox}\fi
\endgroup
\end{center}

and

\begin{center}
\begingroup
\sbox{\grmathbox}{$\displaystyle 
ds^2=dr^2+r^2d\theta^2.
$}
\typeout{GRDISPLAY|155|\the\wd\grmathbox|\the\linewidth}
\ifdim\wd\grmathbox>\linewidth
\resizebox{\linewidth}{!}{\usebox{\grmathbox}}
\else\usebox{\grmathbox}\fi
\endgroup
\end{center}

Thus \(g_{rr}=1\), \(g_{\theta\theta}=r^2\), and
\(g^{\theta\theta}=r^{-2}\). An angular coordinate is dimensionless, so
\(g_{\theta\theta}\) has units of length squared. Metric components need
not all have the same units when their coordinates do not.

Only one metric derivative is nonzero:
\(\partial_rg_{\theta\theta}=2r\). The Christoffel formula gives

\begin{center}
\begingroup
\sbox{\grmathbox}{$\displaystyle 
\Gamma^r{}_{\theta\theta}=-r,
\qquad
\Gamma^\theta{}_{r\theta}
=\Gamma^\theta{}_{\theta r}=\frac1r,
$}
\typeout{GRDISPLAY|156|\the\wd\grmathbox|\the\linewidth}
\ifdim\wd\grmathbox>\linewidth
\resizebox{\linewidth}{!}{\usebox{\grmathbox}}
\else\usebox{\grmathbox}\fi
\endgroup
\end{center}

with every other coefficient zero.

These values have simple origins. Moving around a circle changes the
radial direction, while the angular coordinate basis vector
\(\partial_\theta\) has length \(r\). Its scale changes when you move
radially. The connection records both effects.

\textbf{Experiment 1: a vector that is actually constant.} Take
\(V=\partial_x\), a unit vector pointing in the same Cartesian direction
everywhere. Its polar components are

\begin{center}
\begingroup
\sbox{\grmathbox}{$\displaystyle 
V^r=\cos\theta,
\qquad
V^\theta=-\frac{\sin\theta}{r}.
$}
\typeout{GRDISPLAY|157|\the\wd\grmathbox|\the\linewidth}
\ifdim\wd\grmathbox>\linewidth
\resizebox{\linewidth}{!}{\usebox{\grmathbox}}
\else\usebox{\grmathbox}\fi
\endgroup
\end{center}

The ordinary derivatives are not all zero. But the covariant derivatives
are:

\begin{center}
\begingroup
\sbox{\grmathbox}{$\displaystyle 
\begin{aligned}
\nabla_rV^r&=0,\\
\nabla_\theta V^r
&=-\sin\theta+(-r)\left(-\frac{\sin\theta}{r}\right)=0,\\
\nabla_rV^\theta
&=\frac{\sin\theta}{r^2}
+\frac1r\left(-\frac{\sin\theta}{r}\right)=0,\\
\nabla_\theta V^\theta
&=-\frac{\cos\theta}{r}+\frac1r\cos\theta=0.
\end{aligned}
$}
\typeout{GRDISPLAY|158|\the\wd\grmathbox|\the\linewidth}
\ifdim\wd\grmathbox>\linewidth
\resizebox{\linewidth}{!}{\usebox{\grmathbox}}
\else\usebox{\grmathbox}\fi
\endgroup
\end{center}

The connection has successfully acquitted the kitchen-pointing vector.

\textbf{Experiment 2: compatibility with a visibly changing metric.}
Although \(\partial_rg_{\theta\theta}=2r\),

\begin{center}
\begingroup
\sbox{\grmathbox}{$\displaystyle 
\nabla_rg_{\theta\theta}
=2r-2\Gamma^\theta{}_{r\theta}g_{\theta\theta}
=2r-2\frac1r r^2=0.
$}
\typeout{GRDISPLAY|159|\the\wd\grmathbox|\the\linewidth}
\ifdim\wd\grmathbox>\linewidth
\resizebox{\linewidth}{!}{\usebox{\grmathbox}}
\else\usebox{\grmathbox}\fi
\endgroup
\end{center}

There is no contradiction between nonconstant metric components and a
covariantly constant metric.

\textbf{Experiment 3: a straight line with coordinate acceleration.} The
geodesic equations are

\begin{center}
\begingroup
\sbox{\grmathbox}{$\displaystyle 
\ddot r-r\dot\theta^2=0,
\qquad
\ddot\theta+\frac{2\dot r\dot\theta}{r}=0,
$}
\typeout{GRDISPLAY|160|\the\wd\grmathbox|\the\linewidth}
\ifdim\wd\grmathbox>\linewidth
\resizebox{\linewidth}{!}{\usebox{\grmathbox}}
\else\usebox{\grmathbox}\fi
\endgroup
\end{center}

where dots mean differentiation with respect to an affine parameter. The
second equation says \(d(r^2\dot\theta)/d\lambda=0\).

Now describe the Cartesian straight line \(x=v\lambda\), \(y=b\), with
\(b>0\). Then

\begin{center}
\begingroup
\sbox{\grmathbox}{$\displaystyle 
r=\sqrt{v^2\lambda^2+b^2},
\qquad
\dot\theta=-\frac{bv}{r^2},
\qquad
\ddot r=\frac{v^2b^2}{r^3}=r\dot\theta^2.
$}
\typeout{GRDISPLAY|161|\the\wd\grmathbox|\the\linewidth}
\ifdim\wd\grmathbox>\linewidth
\resizebox{\linewidth}{!}{\usebox{\grmathbox}}
\else\usebox{\grmathbox}\fi
\endgroup
\end{center}

Its radial coordinate accelerates even though the path is perfectly
straight. The \(r\dot\theta^2\) term is the coordinate accounting needed
to express zero geometric acceleration. A nonzero
\(d^2x^\mu/d\lambda^2\) is not, by itself, a physical acceleration
measurement.

\textbf{Experiment 4: transporting around a full circle.} Along
\(r=r_0\), parallel transport obeys

\begin{center}
\begingroup
\sbox{\grmathbox}{$\displaystyle 
\frac{dV^r}{d\theta}=r_0V^\theta,
\qquad
\frac{d(r_0V^\theta)}{d\theta}=-V^r.
$}
\typeout{GRDISPLAY|162|\the\wd\grmathbox|\the\linewidth}
\ifdim\wd\grmathbox>\linewidth
\resizebox{\linewidth}{!}{\usebox{\grmathbox}}
\else\usebox{\grmathbox}\fi
\endgroup
\end{center}

These are the equations for a rotating pair of components. The vector
remains fixed in Cartesian space while the polar basis rotates
underneath it. After \(2\pi\), the components return to their initial
values. The loop causes no net geometric rotation.

\textbf{Experiment 5: the familiar polar Laplacian.} Since
\(\sqrt{g}=r\),

\begin{center}
\begingroup
\sbox{\grmathbox}{$\displaystyle 
\nabla_aV^a
=\frac1r\partial_r(rV^r)+\partial_\theta V^\theta,
$}
\typeout{GRDISPLAY|163|\the\wd\grmathbox|\the\linewidth}
\ifdim\wd\grmathbox>\linewidth
\resizebox{\linewidth}{!}{\usebox{\grmathbox}}
\else\usebox{\grmathbox}\fi
\endgroup
\end{center}

where \(a\) here ranges over the two coordinates. The physical angular
component in a unit-length basis is \(V^{\hat\theta}=rV^\theta\).
Replacing \(V^\theta\) by \(V^{\hat\theta}/r\) gives the familiar
angular divergence term \(r^{-1}\partial_\theta V^{\hat\theta}\).

Likewise,

\begin{center}
\begingroup
\sbox{\grmathbox}{$\displaystyle 
\Delta f
=\frac1r\partial_r(r\partial_rf)
+\frac1{r^2}\partial_\theta^2f.
$}
\typeout{GRDISPLAY|164|\the\wd\grmathbox|\the\linewidth}
\ifdim\wd\grmathbox>\linewidth
\resizebox{\linewidth}{!}{\usebox{\grmathbox}}
\else\usebox{\grmathbox}\fi
\endgroup
\end{center}

The additional terms in undergraduate vector calculus are connection
effects. They were differential geometry wearing a cheaper jacket.

We have nonzero Christoffel symbols, changing basis components, and
coordinate acceleration, all in flat space. The final experiment -
checking the curvature itself - belongs to the next chapter.

\hypertarget{chapter-8}{%
\chapter{8. Curvature: what remains after the coordinate excuses run
out}\label{chapter-8}}

\hypertarget{two-differentiations-enter-their-order-matters}{%
\section{8.1 Two differentiations enter; their order
matters}\label{two-differentiations-enter-their-order-matters}}

For an ordinary smooth scalar in ordinary coordinates, mixed partial
derivatives commute:

\begin{center}
\begingroup
\sbox{\grmathbox}{$\displaystyle 
\partial_\mu\partial_\nu f=\partial_\nu\partial_\mu f.
$}
\typeout{GRDISPLAY|165|\the\wd\grmathbox|\the\linewidth}
\ifdim\wd\grmathbox>\linewidth
\resizebox{\linewidth}{!}{\usebox{\grmathbox}}
\else\usebox{\grmathbox}\fi
\endgroup
\end{center}

But a vector transported and compared through changing tangent spaces
contains extra structure. Differentiating first in one direction and
then another need not give the same answer as reversing the order.

Using the torsion-free Levi-Civita connection, define curvature by

\begin{center}
\begingroup
\sbox{\grmathbox}{$\displaystyle 
\boxed{[\nabla_\mu,\nabla_\nu]V^\rho
=R^\rho{}_{\sigma\mu\nu}V^\sigma.}
$}
\typeout{GRDISPLAY|166|\the\wd\grmathbox|\the\linewidth}
\ifdim\wd\grmathbox>\linewidth
\resizebox{\linewidth}{!}{\usebox{\grmathbox}}
\else\usebox{\grmathbox}\fi
\endgroup
\end{center}

The brackets denote a commutator: the first operation order minus the
reverse. The placement of indices encodes the roles. The final pair
\(\mu,\nu\) identifies the two differentiation directions; \(\sigma\)
receives the original vector; \(\rho\) labels the output vector.

Let us derive the components. Because \(\nabla_\nu V^\rho\) has both an
upper index and a lower index,

\begin{center}
\begingroup
\sbox{\grmathbox}{$\displaystyle 
\begin{aligned}
\nabla_\mu\nabla_\nu V^\rho
={}&\partial_\mu\left(\partial_\nu V^\rho
+\Gamma^\rho{}_{\nu\sigma}V^\sigma\right)\\
&+\Gamma^\rho{}_{\mu\lambda}
\left(\partial_\nu V^\lambda
+\Gamma^\lambda{}_{\nu\sigma}V^\sigma\right)\\
&-\Gamma^\lambda{}_{\mu\nu}
\left(\partial_\lambda V^\rho
+\Gamma^\rho{}_{\lambda\sigma}V^\sigma\right).
\end{aligned}
$}
\typeout{GRDISPLAY|167|\the\wd\grmathbox|\the\linewidth}
\ifdim\wd\grmathbox>\linewidth
\resizebox{\linewidth}{!}{\usebox{\grmathbox}}
\else\usebox{\grmathbox}\fi
\endgroup
\end{center}

The last line is the commonly forgotten correction for the derivative
index \(\nu\). Swap \(\mu\) and \(\nu\), then subtract.

The ordinary second derivatives cancel because mixed partials commute.
The final line cancels because \(\Gamma^\lambda{}_{\mu\nu}\) is
symmetric. The terms containing first derivatives of \(V\) also cancel
in pairs. What survives is proportional to \(V\) itself, with no
derivatives of \(V\) remaining:

\begin{center}
\begingroup
\sbox{\grmathbox}{$\displaystyle 
\boxed{
R^\rho{}_{\sigma\mu\nu}
=\partial_\mu\Gamma^\rho{}_{\nu\sigma}
-\partial_\nu\Gamma^\rho{}_{\mu\sigma}
+\Gamma^\rho{}_{\mu\lambda}\Gamma^\lambda{}_{\nu\sigma}
-\Gamma^\rho{}_{\nu\lambda}\Gamma^\lambda{}_{\mu\sigma}.
}
$}
\typeout{GRDISPLAY|168|\the\wd\grmathbox|\the\linewidth}
\ifdim\wd\grmathbox>\linewidth
\resizebox{\linewidth}{!}{\usebox{\grmathbox}}
\else\usebox{\grmathbox}\fi
\endgroup
\end{center}

That cancellation is conceptually important. The mismatch depends on the
vector\textquotesingle s value at the event, not on how you happened to
extend it into a vector field around the event. Curvature is a local
multilinear geometric object: a tensor.

The derivative terms measure variation of the connection; the quadratic
terms account for successive changes of basis acting on each other.
Schematically,

\begin{center}
\begingroup
\sbox{\grmathbox}{$\displaystyle 
R\sim\partial\Gamma+\Gamma\Gamma
\sim g^{-1}\partial^2g+g^{-2}(\partial g)^2,
$}
\typeout{GRDISPLAY|169|\the\wd\grmathbox|\the\linewidth}
\ifdim\wd\grmathbox>\linewidth
\resizebox{\linewidth}{!}{\usebox{\grmathbox}}
\else\usebox{\grmathbox}\fi
\endgroup
\end{center}

where the schematic final expression suppresses index contractions. The
nonlinear term is not decorative. In polar coordinates it is precisely
what will cancel the apparent curvature from the derivative term.

\hypertarget{curvature-as-a-machine-with-three-vector-inputs}{%
\section{8.2 Curvature as a machine with three vector
inputs}\label{curvature-as-a-machine-with-three-vector-inputs}}

A coordinate-free version makes the structure cleaner:

\begin{center}
\begingroup
\sbox{\grmathbox}{$\displaystyle 
R(X,Y)Z
=\nabla_X\nabla_YZ-\nabla_Y\nabla_XZ-\nabla_{[X,Y]}Z.
$}
\typeout{GRDISPLAY|170|\the\wd\grmathbox|\the\linewidth}
\ifdim\wd\grmathbox>\linewidth
\resizebox{\linewidth}{!}{\usebox{\grmathbox}}
\else\usebox{\grmathbox}\fi
\endgroup
\end{center}

Why subtract the Lie-bracket term? If \(X\) and \(Y\) themselves do not
commute, their flows reach slightly different points when taken in
opposite orders. We must remove that displacement effect before
interpreting the remaining mismatch as curvature. In a coordinate basis
the bracket vanishes, giving the component formula above.

You can picture \(R\) as a machine. Feed it two directions defining a
tiny parallelogram and a vector to carry around that parallelogram. It
returns the infinitesimal change associated with the loop. The machine
is linear in each of its three inputs at the point.

For a covector the curvature acts with the opposite sign:

\begin{center}
\begingroup
\sbox{\grmathbox}{$\displaystyle 
[\nabla_\mu,\nabla_\nu]\omega_\rho
=-R^\sigma{}_{\rho\mu\nu}\omega_\sigma.
$}
\typeout{GRDISPLAY|171|\the\wd\grmathbox|\the\linewidth}
\ifdim\wd\grmathbox>\linewidth
\resizebox{\linewidth}{!}{\usebox{\grmathbox}}
\else\usebox{\grmathbox}\fi
\endgroup
\end{center}

For a tensor, one curvature term acts on each index, with the same
upper-plus/lower-minus pattern as covariant differentiation. The scalar
pairing is again the consistency check.

\hypertarget{holonomy-geometry-remembers-your-route}{%
\section{8.3 Holonomy: geometry remembers your
route}\label{holonomy-geometry-remembers-your-route}}

Transport a vector around a small coordinate parallelogram. Let
\(a^\mu\) and \(b^\mu\) be the small side displacements, and traverse
the sides in the order \(+a,+b,-a,-b\). With our curvature convention
and the transport equation \(dV=-\Gamma_\mu V\,dx^\mu\), the returned
vector satisfies

\begin{center}
\begingroup
\sbox{\grmathbox}{$\displaystyle 
\Delta V^\rho
=-R^\rho{}_{\sigma\mu\nu}V^\sigma a^\mu b^\nu
+O(\ell^3),
$}
\typeout{GRDISPLAY|172|\the\wd\grmathbox|\the\linewidth}
\ifdim\wd\grmathbox>\linewidth
\resizebox{\linewidth}{!}{\usebox{\grmathbox}}
\else\usebox{\grmathbox}\fi
\endgroup
\end{center}

where all side lengths scale with a small parameter \(\ell\). Reverse
the loop and the leading sign reverses. Defining the difference by
subtracting the two open-path results in the opposite order also
reverses it. These orientation choices explain many apparent sign
disagreements in pictures of holonomy.

The net transformation obtained around a closed loop is called
\textbf{holonomy}. Because the initial and final vectors live in the
same tangent space, their mismatch is a genuine comparison. There is no
need to argue about how to compare vectors at different endpoints.

A sphere provides a delightful demonstration. Follow a geodesic triangle
from the equator to the north pole, down to another point on the
equator, then back along the equator. Parallel transport can return a
vector rotated relative to its starting direction. For a spherical
triangle on a sphere of radius \(a\), the rotation magnitude is its area
divided by \(a^2\), equivalently its angular excess. A triangle covering
an octant has three right angles and produces a \(\pi/2\) rotation.

This is not an argument that vectors are secretly mechanical arrows
sliding through ambient three-dimensional space. Parallel transport on
an embedded sphere can be visualized by requiring no tangential turning;
the normal part of the ambient change is allowed. The intrinsic
definition is the connection equation. No embedding is required.

A local/global caution: zero curvature makes sufficiently small
contractible loops have trivial holonomy. A flat connection on a space
with nontrivial topology can still have nontrivial holonomy around
noncontractible loops. Local geometry and global topology are different
chapters in the universe\textquotesingle s accounting system.

\hypertarget{the-polar-plane-verdict}{%
\section{8.4 The polar-plane verdict}\label{the-polar-plane-verdict}}

Using the polar connection from Chapter 7, calculate

\begin{center}
\begingroup
\sbox{\grmathbox}{$\displaystyle 
\begin{aligned}
R^r{}_{\theta r\theta}
&=\partial_r\Gamma^r{}_{\theta\theta}
-\partial_\theta\Gamma^r{}_{r\theta}
+\Gamma^r{}_{r\lambda}\Gamma^\lambda{}_{\theta\theta}
-\Gamma^r{}_{\theta\lambda}\Gamma^\lambda{}_{r\theta}\\
&=-1-0+0-(-r)(1/r)\\
&=0.
\end{aligned}
$}
\typeout{GRDISPLAY|173|\the\wd\grmathbox|\the\linewidth}
\ifdim\wd\grmathbox>\linewidth
\resizebox{\linewidth}{!}{\usebox{\grmathbox}}
\else\usebox{\grmathbox}\fi
\endgroup
\end{center}

In two dimensions the Riemann symmetries leave only one independent
curvature component, so this establishes that the plane is flat wherever
the polar chart is valid. The apparent singularity at \(r=0\) is a
failure of the chart; Cartesian coordinates cover the origin smoothly.

Omitting the \(\Gamma\Gamma\) terms would have declared the dinner plate
curved. The nonlinear terms have just saved lunch.

\hypertarget{symmetries-that-turn-256-entries-into-20}{%
\section{8.5 Symmetries that turn 256 entries into
20}\label{symmetries-that-turn-256-entries-into-20}}

Lower the first index with the metric:

\begin{center}
\begingroup
\sbox{\grmathbox}{$\displaystyle 
R_{\rho\sigma\mu\nu}=g_{\rho\lambda}R^\lambda{}_{\sigma\mu\nu}.
$}
\typeout{GRDISPLAY|174|\the\wd\grmathbox|\the\linewidth}
\ifdim\wd\grmathbox>\linewidth
\resizebox{\linewidth}{!}{\usebox{\grmathbox}}
\else\usebox{\grmathbox}\fi
\endgroup
\end{center}

For the Levi-Civita connection,

\begin{center}
\begingroup
\sbox{\grmathbox}{$\displaystyle 
\begin{aligned}
R_{\rho\sigma\mu\nu}&=-R_{\sigma\rho\mu\nu},\\
R_{\rho\sigma\mu\nu}&=-R_{\rho\sigma\nu\mu},\\
R_{\rho\sigma\mu\nu}&=R_{\mu\nu\rho\sigma},\\
R_{\rho\sigma\mu\nu}
+R_{\rho\mu\nu\sigma}
+R_{\rho\nu\sigma\mu}&=0.
\end{aligned}
$}
\typeout{GRDISPLAY|175|\the\wd\grmathbox|\the\linewidth}
\ifdim\wd\grmathbox>\linewidth
\resizebox{\linewidth}{!}{\usebox{\grmathbox}}
\else\usebox{\grmathbox}\fi
\endgroup
\end{center}

The second symmetry follows immediately from exchanging the
differentiation order. The first reflects metric compatibility: an
infinitesimal parallel-transport transformation preserves inner
products, so its generator is antisymmetric after lowering the output
index. The cyclic identity, often called the algebraic or first Bianchi
identity, depends on the torsion-free condition. At a point where
\(\Gamma=0\), its terms cancel directly using the symmetry of the lower
connection indices. Together these properties imply the pair-exchange
symmetry.

In four dimensions an antisymmetric pair has six possibilities. Treating
each index pair as one label, the pair symmetries make curvature a
symmetric \(6\times6\) array: 21 components. The algebraic Bianchi
identity removes one independent component, leaving 20.

Twenty is a count of algebraically independent curvature components at
an event. It is \textbf{not} a count of propagating gravitational
degrees of freedom. Field equations, constraints, and coordinate freedom
enter that separate question; vacuum GR in four dimensions has two local
gravitational wave polarizations.

\hypertarget{a-sphere-computed-from-its-own-measurements}{%
\section{8.6 A sphere computed from its own
measurements}\label{a-sphere-computed-from-its-own-measurements}}

Now take the round two-sphere of radius \(a\), with intrinsic line
element

\begin{center}
\begingroup
\sbox{\grmathbox}{$\displaystyle 
ds^2=a^2d\theta^2+a^2\sin^2\theta\,d\phi^2.
$}
\typeout{GRDISPLAY|176|\the\wd\grmathbox|\the\linewidth}
\ifdim\wd\grmathbox>\linewidth
\resizebox{\linewidth}{!}{\usebox{\grmathbox}}
\else\usebox{\grmathbox}\fi
\endgroup
\end{center}

The nonzero connection coefficients are

\begin{center}
\begingroup
\sbox{\grmathbox}{$\displaystyle 
\Gamma^\theta{}_{\phi\phi}=-\sin\theta\cos\theta,
\qquad
\Gamma^\phi{}_{\theta\phi}
=\Gamma^\phi{}_{\phi\theta}=\cot\theta.
$}
\typeout{GRDISPLAY|177|\the\wd\grmathbox|\the\linewidth}
\ifdim\wd\grmathbox>\linewidth
\resizebox{\linewidth}{!}{\usebox{\grmathbox}}
\else\usebox{\grmathbox}\fi
\endgroup
\end{center}

The calculation resembles the plane\textquotesingle s polar calculation,
but the circumference scale is \(a\sin\theta\), not a linear radial
distance. That nonlinear relation is where real curvature enters.

Compute one component:

\begin{center}
\begingroup
\sbox{\grmathbox}{$\displaystyle 
\begin{aligned}
R^\theta{}_{\phi\theta\phi}
&=\partial_\theta(-\sin\theta\cos\theta)
-\Gamma^\theta{}_{\phi\phi}\Gamma^\phi{}_{\theta\phi}\\
&=(\sin^2\theta-\cos^2\theta)+\cos^2\theta\\
&=\sin^2\theta.
\end{aligned}
$}
\typeout{GRDISPLAY|178|\the\wd\grmathbox|\the\linewidth}
\ifdim\wd\grmathbox>\linewidth
\resizebox{\linewidth}{!}{\usebox{\grmathbox}}
\else\usebox{\grmathbox}\fi
\endgroup
\end{center}

Thus

\begin{center}
\begingroup
\sbox{\grmathbox}{$\displaystyle 
R_{\theta\phi\theta\phi}=a^2\sin^2\theta.
$}
\typeout{GRDISPLAY|179|\the\wd\grmathbox|\the\linewidth}
\ifdim\wd\grmathbox>\linewidth
\resizebox{\linewidth}{!}{\usebox{\grmathbox}}
\else\usebox{\grmathbox}\fi
\endgroup
\end{center}

This component has units of length squared because all its slots refer
to angular coordinate vectors. In a unit orthonormal basis the
corresponding component is \(1/a^2\). The invariant Gaussian curvature
is

\begin{center}
\begingroup
\sbox{\grmathbox}{$\displaystyle 
K=\frac{R_{\theta\phi\theta\phi}}
{g_{\theta\theta}g_{\phi\phi}-g_{\theta\phi}^2}
=\frac1{a^2}.
$}
\typeout{GRDISPLAY|180|\the\wd\grmathbox|\the\linewidth}
\ifdim\wd\grmathbox>\linewidth
\resizebox{\linewidth}{!}{\usebox{\grmathbox}}
\else\usebox{\grmathbox}\fi
\endgroup
\end{center}

The Ricci tensor and scalar, whose general definitions we discuss next,
are

\begin{center}
\begingroup
\sbox{\grmathbox}{$\displaystyle 
R_{\theta\theta}=1,
\qquad R_{\phi\phi}=\sin^2\theta,
\qquad R=\frac2{a^2}>0.
$}
\typeout{GRDISPLAY|181|\the\wd\grmathbox|\the\linewidth}
\ifdim\wd\grmathbox>\linewidth
\resizebox{\linewidth}{!}{\usebox{\grmathbox}}
\else\usebox{\grmathbox}\fi
\endgroup
\end{center}

This also fixes our sign convention operationally: a round sphere has
positive scalar curvature. Some books reverse the sign of the Riemann
tensor. Equations borrowed from them must be translated consistently,
especially geodesic deviation and the Einstein equation.

An inhabitant of the sphere need not see its embedding. Draw a circle of
geodesic radius \(s\) around a point. Its circumference is

\begin{center}
\begingroup
\sbox{\grmathbox}{$\displaystyle 
C(s)=2\pi a\sin(s/a)
=2\pi s\left(1-\frac{s^2}{6a^2}+\cdots\right).
$}
\typeout{GRDISPLAY|182|\the\wd\grmathbox|\the\linewidth}
\ifdim\wd\grmathbox>\linewidth
\resizebox{\linewidth}{!}{\usebox{\grmathbox}}
\else\usebox{\grmathbox}\fi
\endgroup
\end{center}

It is shorter than the Euclidean prediction. Curvature can be discovered
with intrinsic distance measurements alone.

A cylinder illustrates the opposite lesson. A sheet can be rolled into a
cylinder without locally stretching it. Its embedding looks bent, but
its intrinsic metric is locally flat. \textbf{Extrinsic bending}
describes the relation to an ambient space; \textbf{intrinsic curvature}
describes the metric geometry experienced by inhabitants. Spacetime
curvature does not demand an invisible higher-dimensional room into
which spacetime bends.

\hypertarget{chapter-9}{%
\chapter{9. Ricci, Weyl, and Einstein: different questions asked of
curvature}\label{chapter-9}}

\hypertarget{a-contraction-is-a-deliberate-loss-of-information}{%
\section{9.1 A contraction is a deliberate loss of
information}\label{a-contraction-is-a-deliberate-loss-of-information}}

The Riemann tensor is richly directional. Einstein\textquotesingle s
equation uses a particular contraction of it, the Ricci tensor:

\begin{center}
\begingroup
\sbox{\grmathbox}{$\displaystyle 
\boxed{R_{\mu\nu}=R^\rho{}_{\mu\rho\nu}.}
$}
\typeout{GRDISPLAY|183|\the\wd\grmathbox|\the\linewidth}
\ifdim\wd\grmathbox>\linewidth
\resizebox{\linewidth}{!}{\usebox{\grmathbox}}
\else\usebox{\grmathbox}\fi
\endgroup
\end{center}

This identifies the output index with one curvature-direction index and
sums. It is a trace of a curvature map. Because of the Riemann
symmetries, \(R_{\mu\nu}=R_{\nu\mu}\).

Contract again to obtain the Ricci scalar:

\begin{center}
\begingroup
\sbox{\grmathbox}{$\displaystyle 
\boxed{R=g^{\mu\nu}R_{\mu\nu}.}
$}
\typeout{GRDISPLAY|184|\the\wd\grmathbox|\the\linewidth}
\ifdim\wd\grmathbox>\linewidth
\resizebox{\linewidth}{!}{\usebox{\grmathbox}}
\else\usebox{\grmathbox}\fi
\endgroup
\end{center}

A trace combines directional information into an aggregate. Consider a
matrix with eigenvalues \(2,-1,-1\). Its trace vanishes, but the matrix
is very much present: it stretches one direction and compresses two
others. A vanishing Ricci scalar is even less informative than a
vanishing Ricci tensor, and a vanishing Ricci tensor is less informative
than a vanishing Riemann tensor.

The hierarchy is

\begin{center}
\begingroup
\sbox{\grmathbox}{$\displaystyle 
R_{\rho\sigma\mu\nu}=0
\ \Longrightarrow\ R_{\mu\nu}=0
\ \Longrightarrow\ R=0,
$}
\typeout{GRDISPLAY|185|\the\wd\grmathbox|\the\linewidth}
\ifdim\wd\grmathbox>\linewidth
\resizebox{\linewidth}{!}{\usebox{\grmathbox}}
\else\usebox{\grmathbox}\fi
\endgroup
\end{center}

with neither reverse implication valid in general four-dimensional
spacetime.

This is crucial for understanding gravity in empty space. For
\(\Lambda=0\), the vacuum Einstein equation gives \(R_{\mu\nu}=0\). A
black hole exterior can still have large tidal curvature, and
gravitational waves can still propagate through vacuum. Ricci-flat does
not mean Riemann-flat.

\hypertarget{ricci-curvature-as-a-directional-average-of-tidal-effects}{%
\section{9.2 Ricci curvature as a directional average of tidal
effects}\label{ricci-curvature-as-a-directional-average-of-tidal-effects}}

Choose a freely falling observer with unit orthonormal frame and
\(e_{\hat0}=u/c\). The spatial tidal matrix is

\begin{center}
\begingroup
\sbox{\grmathbox}{$\displaystyle 
\mathcal E_{ij}=R_{\hat i\hat0\hat j\hat0}.
$}
\typeout{GRDISPLAY|186|\the\wd\grmathbox|\the\linewidth}
\ifdim\wd\grmathbox>\linewidth
\resizebox{\linewidth}{!}{\usebox{\grmathbox}}
\else\usebox{\grmathbox}\fi
\endgroup
\end{center}

It is symmetric. The next chapter derives the measurable relative
acceleration it produces:

\begin{center}
\begingroup
\sbox{\grmathbox}{$\displaystyle 
\frac{d^2\xi^{\hat i}}{d\tau^2}
=-c^2\mathcal E^i{}_{j}\xi^{\hat j}.
$}
\typeout{GRDISPLAY|187|\the\wd\grmathbox|\the\linewidth}
\ifdim\wd\grmathbox>\linewidth
\resizebox{\linewidth}{!}{\usebox{\grmathbox}}
\else\usebox{\grmathbox}\fi
\endgroup
\end{center}

The trace of this matrix is

\begin{center}
\begingroup
\sbox{\grmathbox}{$\displaystyle 
\delta^{ij}\mathcal E_{ij}=R_{\hat0\hat0}
=\frac1{c^2}R_{\mu\nu}u^\mu u^\nu.
$}
\typeout{GRDISPLAY|188|\the\wd\grmathbox|\the\linewidth}
\ifdim\wd\grmathbox>\linewidth
\resizebox{\linewidth}{!}{\usebox{\grmathbox}}
\else\usebox{\grmathbox}\fi
\endgroup
\end{center}

Thus Ricci curvature evaluated on the observer\textquotesingle s time
direction measures the sum of the three principal tidal effects. One
direction can stretch while another compresses; the trace tells you the
aggregate.

To turn that into a volume statement, release an infinitesimal ball of
freely falling particles initially at rest relative to one another in
the observer\textquotesingle s frame. If its initial volume is
\(\mathcal V_0\), then initially

\begin{center}
\begingroup
\sbox{\grmathbox}{$\displaystyle 
\frac{\ddot{\mathcal V}}{\mathcal V}
=-R_{\mu\nu}u^\mu u^\nu,
$}
\typeout{GRDISPLAY|189|\the\wd\grmathbox|\the\linewidth}
\ifdim\wd\grmathbox>\linewidth
\resizebox{\linewidth}{!}{\usebox{\grmathbox}}
\else\usebox{\grmathbox}\fi
\endgroup
\end{center}

and, for a short proper-time interval,

\begin{center}
\begingroup
\sbox{\grmathbox}{$\displaystyle 
\frac{\mathcal V(\Delta\tau)}{\mathcal V_0}
=1-\frac12R_{\mu\nu}u^\mu u^\nu(\Delta\tau)^2
+O((\Delta\tau)^3).
$}
\typeout{GRDISPLAY|190|\the\wd\grmathbox|\the\linewidth}
\ifdim\wd\grmathbox>\linewidth
\resizebox{\linewidth}{!}{\usebox{\grmathbox}}
\else\usebox{\grmathbox}\fi
\endgroup
\end{center}

Why the trace? At first, each principal edge changes by its own small
fractional amount. Multiplying the three edge lengths, the leading
fractional volume change is the sum of those changes. Products of the
small changes enter at higher order.

The initial conditions matter. A cloud already expanding, shearing, or
rotating has additional contributions to its volume evolution. The
Raychaudhuri equation later gives the exact infinitesimal bookkeeping.
``Ricci controls volume'' refers to a precisely defined curvature
contribution, not a claim that every cloud\textquotesingle s volume is
determined by Ricci alone.

\hypertarget{extracting-the-trace-free-remainder-the-weyl-tensor}{%
\section{9.3 Extracting the trace-free remainder: the Weyl
tensor}\label{extracting-the-trace-free-remainder-the-weyl-tensor}}

In four dimensions the Riemann tensor can be decomposed as

\begin{center}
\begingroup
\sbox{\grmathbox}{$\displaystyle 
\boxed{
\begin{aligned}
R_{\rho\sigma\mu\nu}
={}&C_{\rho\sigma\mu\nu}\\
&+\frac12\left(
 g_{\rho\mu}R_{\sigma\nu}
-g_{\rho\nu}R_{\sigma\mu}
-g_{\sigma\mu}R_{\rho\nu}
+g_{\sigma\nu}R_{\rho\mu}
\right)\\
&-\frac{R}{6}\left(
 g_{\rho\mu}g_{\sigma\nu}
-g_{\rho\nu}g_{\sigma\mu}
\right).
\end{aligned}}
$}
\typeout{GRDISPLAY|191|\the\wd\grmathbox|\the\linewidth}
\ifdim\wd\grmathbox>\linewidth
\resizebox{\linewidth}{!}{\usebox{\grmathbox}}
\else\usebox{\grmathbox}\fi
\endgroup
\end{center}

The tensor \(C_{\rho\sigma\mu\nu}\) is the \textbf{Weyl tensor}. It has
the Riemann symmetries and every metric trace vanishes. It is the
curvature information left after removing Ricci\textquotesingle s
traces.

The factors \(1/2\) and \(1/6\) can be derived rather than memorized.
Write the four-term Ricci combination in parentheses as
\(Q_{\rho\sigma\mu\nu}\) and the two-term metric combination as
\(P_{\rho\sigma\mu\nu}\). Suppose

\begin{center}
\begingroup
\sbox{\grmathbox}{$\displaystyle 
R_{\rho\sigma\mu\nu}
=C_{\rho\sigma\mu\nu}+A Q_{\rho\sigma\mu\nu}
+B R P_{\rho\sigma\mu\nu}.
$}
\typeout{GRDISPLAY|192|\the\wd\grmathbox|\the\linewidth}
\ifdim\wd\grmathbox>\linewidth
\resizebox{\linewidth}{!}{\usebox{\grmathbox}}
\else\usebox{\grmathbox}\fi
\endgroup
\end{center}

Contract the first and third indices. In \(n\) dimensions,

\begin{center}
\begingroup
\sbox{\grmathbox}{$\displaystyle 
\begin{aligned}
g^{\rho\mu}Q_{\rho\sigma\mu\nu}
&=(n-2)R_{\sigma\nu}+g_{\sigma\nu}R,\\
g^{\rho\mu}P_{\rho\sigma\mu\nu}
&=(n-1)g_{\sigma\nu}.
\end{aligned}
$}
\typeout{GRDISPLAY|193|\the\wd\grmathbox|\the\linewidth}
\ifdim\wd\grmathbox>\linewidth
\resizebox{\linewidth}{!}{\usebox{\grmathbox}}
\else\usebox{\grmathbox}\fi
\endgroup
\end{center}

The Weyl contribution must vanish under this contraction, and the left
side must equal \(R_{\sigma\nu}\). Therefore \(A(n-2)=1\) and
\(A+B(n-1)=0\). With \(n=4\), \(A=1/2\) and \(B=-1/6\). The coefficients
are forced by trace removal.

At an event in four dimensions, Ricci carries ten independent components
and Weyl carries the remaining ten. Again, this is algebraic
information, not ten freely propagating fields.

Weyl curvature is often described as ``shape-changing curvature.'' Its
contribution to the tidal matrix is trace-free, so it produces no
leading initial volume acceleration for the initially comoving ball just
described. It can stretch a sphere into an ellipsoid while preserving
volume to that leading order.

Two qualifications prevent this useful slogan from becoming misleading:

\begin{itemize}
\tightlist
\item
  Ricci curvature can also contribute to anisotropic distortion; it is
  not generally a purely isotropic squeeze. The Weyl/Ricci split is a
  decomposition of the full spacetime curvature, while ``shape versus
  volume'' refers to a chosen observer\textquotesingle s spatial tidal
  problem.
\item
  Weyl-induced shear can subsequently affect volume evolution. A
  Ricci-flat cloud is not guaranteed to maintain its initial volume
  forever. Trace-free initial acceleration is weaker than an exact
  volume-preservation theorem.
\end{itemize}

When \(R_{\mu\nu}=0\), the whole Riemann tensor equals Weyl. This is the
curvature of vacuum tidal fields and vacuum gravitational waves. If
\(\Lambda\ne0\), vacuum instead has \(R_{\mu\nu}=\Lambda g_{\mu\nu}\),
and the Ricci part need not vanish.

There is also a useful conformal interpretation. Under
\(g_{\mu\nu}\mapsto\Omega^2g_{\mu\nu}\) with smooth positive \(\Omega\),
the mixed-index Weyl tensor \(C^\rho{}_{\sigma\mu\nu}\) is unchanged.
The all-lowered tensor acquires the metric\textquotesingle s factor
\(\Omega^2\). In four dimensions, vanishing Weyl curvature throughout a
suitable neighborhood characterizes local conformal flatness. This is
stronger than a statement at one isolated event. For further treatment
of the Weyl decomposition and conformal geometry, see
\href{https://davidtong.org/pdfs/teaching/general-relativity/gr.pdf}{David
Tong\textquotesingle s general relativity notes}.

\hypertarget{bianchi-identities-why-the-curvature-components-cannot-vary-independently}{%
\section{9.4 Bianchi identities: why the curvature components cannot
vary
independently}\label{bianchi-identities-why-the-curvature-components-cannot-vary-independently}}

A metric can vary from point to point, but the resulting curvature is
not an arbitrary tensor field with 20 freely assignable functions.
Because it comes from a connection, it obeys a differential identity:

\begin{center}
\begingroup
\sbox{\grmathbox}{$\displaystyle 
\boxed{
\nabla_\lambda R^\rho{}_{\sigma\mu\nu}
+\nabla_\mu R^\rho{}_{\sigma\nu\lambda}
+\nabla_\nu R^\rho{}_{\sigma\lambda\mu}=0.
}
$}
\typeout{GRDISPLAY|194|\the\wd\grmathbox|\the\linewidth}
\ifdim\wd\grmathbox>\linewidth
\resizebox{\linewidth}{!}{\usebox{\grmathbox}}
\else\usebox{\grmathbox}\fi
\endgroup
\end{center}

This is the \textbf{differential or second Bianchi identity}. It is an
identity of Levi-Civita geometry, before Einstein\textquotesingle s
equation or any matter model is introduced.

Here is a local proof that shows why it exists. At an arbitrary point
choose normal coordinates so that \(\Gamma=0\) there. At that point
\(\nabla R=\partial R\). Differentiating
\(R=\partial\Gamma-\partial\Gamma+\Gamma\Gamma-\Gamma\Gamma\) gives
second derivatives of \(\Gamma\); the derivatives of its quadratic
products contain an undifferentiated \(\Gamma\) and vanish at the point.
Add the three cyclic terms. Every second derivative appears twice with
opposite sign, so commutation of ordinary mixed partial derivatives
makes the sum zero.

The expression is tensorial, so if it vanishes in normal coordinates it
vanishes in every chart at that point. The point was arbitrary, so the
identity holds throughout the smooth region. Choosing clever coordinates
is a proof technique, not a restriction on the geometry.

Geometrically, the identity says the infinitesimal loop rotations
surrounding a tiny three-dimensional box fit together consistently.
Algebraically, it is related to the Jacobi identity for
covariant-derivative commutators. Curvature is the failure of pairwise
commutation, but those failures themselves must satisfy a consistency
relation.

\hypertarget{contracting-bianchi-until-einsteins-tensor-appears}{%
\section{9.5 Contracting Bianchi until Einstein\textquotesingle s tensor
appears}\label{contracting-bianchi-until-einsteins-tensor-appears}}

Contracting the differential identity once gives

\begin{center}
\begingroup
\sbox{\grmathbox}{$\displaystyle 
\nabla_\rho R^\rho{}_{\sigma\mu\nu}
=\nabla_\mu R_{\sigma\nu}
-\nabla_\nu R_{\sigma\mu}.
$}
\typeout{GRDISPLAY|195|\the\wd\grmathbox|\the\linewidth}
\ifdim\wd\grmathbox>\linewidth
\resizebox{\linewidth}{!}{\usebox{\grmathbox}}
\else\usebox{\grmathbox}\fi
\endgroup
\end{center}

Now contract with \(g^{\sigma\nu}\). Compatibility allows the metric to
move through the derivative. The left side becomes
\(\nabla_\rho R^\rho{}_{\mu}\); the right side becomes
\(\nabla_\mu R-\nabla_\nu R^\nu{}_{\mu}\). Thus

\begin{center}
\begingroup
\sbox{\grmathbox}{$\displaystyle 
\nabla_\rho R^\rho{}_{\mu}
=\nabla_\mu R-\nabla_\nu R^\nu{}_{\mu}.
$}
\typeout{GRDISPLAY|196|\the\wd\grmathbox|\the\linewidth}
\ifdim\wd\grmathbox>\linewidth
\resizebox{\linewidth}{!}{\usebox{\grmathbox}}
\else\usebox{\grmathbox}\fi
\endgroup
\end{center}

The two Ricci-divergence expressions differ only by the name of their
summed index. Bring them together:

\begin{center}
\begingroup
\sbox{\grmathbox}{$\displaystyle 
\boxed{\nabla^\mu R_{\mu\nu}=\frac12\nabla_\nu R.}
$}
\typeout{GRDISPLAY|197|\the\wd\grmathbox|\the\linewidth}
\ifdim\wd\grmathbox>\linewidth
\resizebox{\linewidth}{!}{\usebox{\grmathbox}}
\else\usebox{\grmathbox}\fi
\endgroup
\end{center}

This tells us exactly how to repair the Ricci tensor\textquotesingle s
divergence. Define

\begin{center}
\begingroup
\sbox{\grmathbox}{$\displaystyle 
\boxed{G_{\mu\nu}=R_{\mu\nu}-\frac12Rg_{\mu\nu}.}
$}
\typeout{GRDISPLAY|198|\the\wd\grmathbox|\the\linewidth}
\ifdim\wd\grmathbox>\linewidth
\resizebox{\linewidth}{!}{\usebox{\grmathbox}}
\else\usebox{\grmathbox}\fi
\endgroup
\end{center}

Then

\begin{center}
\begingroup
\sbox{\grmathbox}{$\displaystyle 
\begin{aligned}
\nabla^\mu G_{\mu\nu}
&=\nabla^\mu R_{\mu\nu}
-\frac12(\nabla^\mu R)g_{\mu\nu}
-\frac12R\nabla^\mu g_{\mu\nu}\\
&=\frac12\nabla_\nu R-\frac12\nabla_\nu R-0\\
&=0.
\end{aligned}
$}
\typeout{GRDISPLAY|199|\the\wd\grmathbox|\the\linewidth}
\ifdim\wd\grmathbox>\linewidth
\resizebox{\linewidth}{!}{\usebox{\grmathbox}}
\else\usebox{\grmathbox}\fi
\endgroup
\end{center}

The coefficient \(1/2\) is now motivated by a geometric consistency
requirement. It is not an arbitrary decoration between Ricci and the
metric.

\textbf{Divergence-free does not mean covariantly constant.} The
identity is \(\nabla^\mu G_{\mu\nu}=0\), involving a contraction. It
does not say \(\nabla_\lambda G_{\mu\nu}=0\) for every choice of
indices. A fluid can have zero net outflow from each small region while
still varying across the room; similarly, divergence-free geometry can
vary.

The result prepares the field equation. A consistent geometric left side
can be matched to a covariantly conserved stress-energy tensor on the
right. A constant multiple of \(g_{\mu\nu}\) is also divergence-free,
allowing the cosmological term \(\Lambda g_{\mu\nu}\). Geometry alone
has not yet fixed the physical coupling, matter content, or the
theory\textquotesingle s full action. Those are the next stage.

\hypertarget{a-dimensional-gotcha-with-a-lesson}{%
\section{9.6 A dimensional gotcha with a
lesson}\label{a-dimensional-gotcha-with-a-lesson}}

For the two-sphere we found \(R_{ab}=a^{-2}g_{ab}\) and \(R=2/a^2\).
Consequently,

\begin{center}
\begingroup
\sbox{\grmathbox}{$\displaystyle 
G_{ab}=R_{ab}-\frac12Rg_{ab}=0
$}
\typeout{GRDISPLAY|200|\the\wd\grmathbox|\the\linewidth}
\ifdim\wd\grmathbox>\linewidth
\resizebox{\linewidth}{!}{\usebox{\grmathbox}}
\else\usebox{\grmathbox}\fi
\endgroup
\end{center}

for that intrinsically curved surface. In fact, the Einstein tensor
vanishes identically for every two-dimensional metric.

This is not a counterexample to the usefulness of
Einstein\textquotesingle s equation in four dimensions. It shows that a
tensor\textquotesingle s information content depends on dimension. In
two dimensions all intrinsic curvature is summarized by one scalar, and
the Einstein combination cancels it. In three dimensions Weyl vanishes
identically and Ricci determines the full Riemann tensor. Four
dimensions are the first in which a nonzero Weyl tensor carries local
curvature information independent of Ricci.

\hypertarget{chapter-10}{%
\chapter{10. Tides: how to measure curvature without seeing spacetime
from outside}\label{chapter-10}}

\hypertarget{one-falling-astronaut-is-not-enough}{%
\section{10.1 One falling astronaut is not
enough}\label{one-falling-astronaut-is-not-enough}}

An ideal accelerometer carried by a freely falling point particle reads
zero. This is true in Minkowski space and beside a black hole, provided
the particle follows a geodesic and is treated as an ideal test body.

Release a \textbf{collection} of particles, however, and their
separations can accelerate. Near a gravitating body, one part of the
collection may be pulled into a different geodesic than another. A
sufficiently extended astronaut cannot follow every nearby geodesic
simultaneously while maintaining an unchanged shape. Internal stresses
arise because the body resists that relative motion.

A single accelerometer measures proper acceleration, the departure of
its worldline from geodesic motion. A gravity gradiometer compares
nearby free motions and measures tidal curvature. Confusing those two
instruments is the source of a remarkable amount of popular-relativity
confusion.

The governing equation is geodesic deviation, also called the Jacobi
equation. We will derive it using only concepts already assembled.

\hypertarget{organize-the-experiment-as-a-family-of-worldlines}{%
\section{10.2 Organize the experiment as a family of
worldlines}\label{organize-the-experiment-as-a-family-of-worldlines}}

Consider a smooth local family of timelike geodesics

\begin{center}
\begingroup
\sbox{\grmathbox}{$\displaystyle 
x^\mu=x^\mu(\tau,s),
$}
\typeout{GRDISPLAY|201|\the\wd\grmathbox|\the\linewidth}
\ifdim\wd\grmathbox>\linewidth
\resizebox{\linewidth}{!}{\usebox{\grmathbox}}
\else\usebox{\grmathbox}\fi
\endgroup
\end{center}

where \(\tau\) is proper time along each geodesic and \(s\) labels
neighboring geodesics. Work in a region where the family defines a
smooth local congruence. Define

\begin{center}
\begingroup
\sbox{\grmathbox}{$\displaystyle 
u^\mu=\frac{\partial x^\mu}{\partial\tau},
\qquad
\xi^\mu=\frac{\partial x^\mu}{\partial s}.
$}
\typeout{GRDISPLAY|202|\the\wd\grmathbox|\the\linewidth}
\ifdim\wd\grmathbox>\linewidth
\resizebox{\linewidth}{!}{\usebox{\grmathbox}}
\else\usebox{\grmathbox}\fi
\endgroup
\end{center}

The first vector moves along one worldline. The second moves across the
family at fixed \(\tau\). More precisely, \(\xi^\mu ds\) describes an
infinitesimal connecting displacement. A finite separation between
distant events does not automatically define a unique tangent vector;
the family and the infinitesimal limit supply the meaning here.

Because \(\tau\) and \(s\) are commuting parameters,

\begin{center}
\begingroup
\sbox{\grmathbox}{$\displaystyle 
[u,\xi]=0.
$}
\typeout{GRDISPLAY|203|\the\wd\grmathbox|\the\linewidth}
\ifdim\wd\grmathbox>\linewidth
\resizebox{\linewidth}{!}{\usebox{\grmathbox}}
\else\usebox{\grmathbox}\fi
\endgroup
\end{center}

With zero torsion, the definition of torsion therefore gives

\begin{center}
\begingroup
\sbox{\grmathbox}{$\displaystyle 
\nabla_u\xi=\nabla_\xi u.
$}
\typeout{GRDISPLAY|204|\the\wd\grmathbox|\the\linewidth}
\ifdim\wd\grmathbox>\linewidth
\resizebox{\linewidth}{!}{\usebox{\grmathbox}}
\else\usebox{\grmathbox}\fi
\endgroup
\end{center}

Read this as a statement about the same smoothly labeled grid of paths.
The rate at which the separation changes along a trajectory equals the
change in velocity across neighboring trajectories. It is the
curved-space version of exchanging \(\partial_\tau\partial_sx\) with
\(\partial_s\partial_\tau x\).

Proper-time parameterization gives \(u^\mu u_\mu=-c^2\). Geodesic motion
gives

\begin{center}
\begingroup
\sbox{\grmathbox}{$\displaystyle 
\nabla_u u=0.
$}
\typeout{GRDISPLAY|205|\the\wd\grmathbox|\the\linewidth}
\ifdim\wd\grmathbox>\linewidth
\resizebox{\linewidth}{!}{\usebox{\grmathbox}}
\else\usebox{\grmathbox}\fi
\endgroup
\end{center}

These assumptions have different jobs: normalization makes \(\tau\) a
physical clock reading; affine geodesic motion removes a tangential
reparameterization term from the acceleration equation.

\hypertarget{deriving-geodesic-deviation-with-every-cancellation-visible}{%
\section{10.3 Deriving geodesic deviation, with every cancellation
visible}\label{deriving-geodesic-deviation-with-every-cancellation-visible}}

The covariant relative acceleration is

\begin{center}
\begingroup
\sbox{\grmathbox}{$\displaystyle 
\frac{D^2\xi}{d\tau^2}=\nabla_u\nabla_u\xi.
$}
\typeout{GRDISPLAY|206|\the\wd\grmathbox|\the\linewidth}
\ifdim\wd\grmathbox>\linewidth
\resizebox{\linewidth}{!}{\usebox{\grmathbox}}
\else\usebox{\grmathbox}\fi
\endgroup
\end{center}

Replace the inner derivative using \(\nabla_u\xi=\nabla_\xi u\):

\begin{center}
\begingroup
\sbox{\grmathbox}{$\displaystyle 
\frac{D^2\xi}{d\tau^2}=\nabla_u\nabla_\xi u.
$}
\typeout{GRDISPLAY|207|\the\wd\grmathbox|\the\linewidth}
\ifdim\wd\grmathbox>\linewidth
\resizebox{\linewidth}{!}{\usebox{\grmathbox}}
\else\usebox{\grmathbox}\fi
\endgroup
\end{center}

Now use the definition of curvature to interchange the two derivatives:

\begin{center}
\begingroup
\sbox{\grmathbox}{$\displaystyle 
\nabla_u\nabla_\xi u
=\nabla_\xi\nabla_u u+R(u,\xi)u+\nabla_{[u,\xi]}u.
$}
\typeout{GRDISPLAY|208|\the\wd\grmathbox|\the\linewidth}
\ifdim\wd\grmathbox>\linewidth
\resizebox{\linewidth}{!}{\usebox{\grmathbox}}
\else\usebox{\grmathbox}\fi
\endgroup
\end{center}

The first term vanishes because every worldline in the family is
geodesic. The final term vanishes because the parameters commute. We are
left with

\begin{center}
\begingroup
\sbox{\grmathbox}{$\displaystyle 
\frac{D^2\xi}{d\tau^2}=R(u,\xi)u.
$}
\typeout{GRDISPLAY|209|\the\wd\grmathbox|\the\linewidth}
\ifdim\wd\grmathbox>\linewidth
\resizebox{\linewidth}{!}{\usebox{\grmathbox}}
\else\usebox{\grmathbox}\fi
\endgroup
\end{center}

In components, this is
\(R^\mu{}_{\alpha\beta\nu}u^\alpha u^\beta\xi^\nu\). Use antisymmetry in
the last pair to place the separation index before the final velocity
index:

\begin{center}
\begingroup
\sbox{\grmathbox}{$\displaystyle 
\boxed{
\frac{D^2\xi^\mu}{d\tau^2}
=-R^\mu{}_{\alpha\nu\beta}
 u^\alpha\xi^\nu u^\beta.
}
$}
\typeout{GRDISPLAY|210|\the\wd\grmathbox|\the\linewidth}
\ifdim\wd\grmathbox>\linewidth
\resizebox{\linewidth}{!}{\usebox{\grmathbox}}
\else\usebox{\grmathbox}\fi
\endgroup
\end{center}

The minus sign follows from the stated curvature convention and this
index ordering. It is not an independently adjustable physical sign.

The equation is exact for a Jacobi field generated by an infinitesimal
variation of geodesics. Using it for particles separated by a finite
distance is a linear approximation in their separation. Higher-order
separation effects involve additional geometric information, including
curvature variation. The same geometric derivation works for null
geodesics with an affine parameter, but a null worldline has no
proper-time parameter or timelike rest frame.

To interpret \(\xi\) as spatial separation for the reference observer,
choose it initially orthogonal to \(u\). This orthogonality persists:

\begin{center}
\begingroup
\sbox{\grmathbox}{$\displaystyle 
\begin{aligned}
\frac{d}{d\tau}(u\cdot\xi)
&=u\cdot\nabla_u\xi\\
&=u\cdot\nabla_\xi u\\
&=\frac12\nabla_\xi(u\cdot u)=0.
\end{aligned}
$}
\typeout{GRDISPLAY|211|\the\wd\grmathbox|\the\linewidth}
\ifdim\wd\grmathbox>\linewidth
\resizebox{\linewidth}{!}{\usebox{\grmathbox}}
\else\usebox{\grmathbox}\fi
\endgroup
\end{center}

The first equality uses geodesic motion; the last uses the equal
normalization of the family. This is why the separation can consistently
be discussed in the reference observer\textquotesingle s instantaneous
rest space.

\hypertarget{put-an-instrument-frame-on-the-reference-geodesic}{%
\section{10.4 Put an instrument frame on the reference
geodesic}\label{put-an-instrument-frame-on-the-reference-geodesic}}

Choose a parallel-transported orthonormal frame along the reference
trajectory, with \(e_{\hat0}=u/c\). In this frame \(u^{\hat0}=c\),
\(u^{\hat i}=0\), and the frame itself has zero covariant rate of change
along the trajectory. Covariant differentiation of vector components
along it becomes ordinary differentiation.

The spatial deviation equation reduces to

\begin{center}
\begingroup
\sbox{\grmathbox}{$\displaystyle 
\boxed{
\frac{d^2\xi^{\hat i}}{d\tau^2}
=-c^2R^{\hat i}{}_{\hat0\hat j\hat0}\xi^{\hat j}.
}
$}
\typeout{GRDISPLAY|212|\the\wd\grmathbox|\the\linewidth}
\ifdim\wd\grmathbox>\linewidth
\resizebox{\linewidth}{!}{\usebox{\grmathbox}}
\else\usebox{\grmathbox}\fi
\endgroup
\end{center}

The matrix \(c^2R_{\hat i\hat0\hat j\hat0}\) has units of inverse time
squared. Its eigenvectors identify the principal tidal directions. With
our sign convention, a positive eigenvalue produces relative
acceleration toward the reference trajectory in that direction; a
negative eigenvalue produces relative acceleration away.

This matrix is an observer-dependent projection of the invariant Riemann
tensor. Another observer with a different four-velocity can obtain a
different tidal matrix. That does not make the effect a coordinate
illusion: they are physically different observers performing different
local experiments. A complete reconstruction of curvature requires
enough independent relative-motion measurements, not merely one
three-dimensional tidal matrix.

\hypertarget{the-newtonian-limit-checks-both-the-meaning-and-the-sign}{%
\section{10.5 The Newtonian limit checks both the meaning and the
sign}\label{the-newtonian-limit-checks-both-the-meaning-and-the-sign}}

Let \(\Phi\) be a weak, slowly varying Newtonian potential. For this
calculation use \(x^0=ct\) and retain leading weak-field terms:

\begin{center}
\begingroup
\sbox{\grmathbox}{$\displaystyle 
g_{00}\simeq-\left(1+\frac{2\Phi}{c^2}\right),
\qquad g_{0i}\simeq0.
$}
\typeout{GRDISPLAY|213|\the\wd\grmathbox|\the\linewidth}
\ifdim\wd\grmathbox>\linewidth
\resizebox{\linewidth}{!}{\usebox{\grmathbox}}
\else\usebox{\grmathbox}\fi
\endgroup
\end{center}

The spatial inverse metric can be replaced by \(\delta^{ij}\) at this
order in the particular expression we need. For an approximately static
field,

\begin{center}
\begingroup
\sbox{\grmathbox}{$\displaystyle 
\Gamma^i{}_{00}
\simeq-\frac12\delta^{ij}\partial_jg_{00}
=\frac1{c^2}\partial^i\Phi.
$}
\typeout{GRDISPLAY|214|\the\wd\grmathbox|\the\linewidth}
\ifdim\wd\grmathbox>\linewidth
\resizebox{\linewidth}{!}{\usebox{\grmathbox}}
\else\usebox{\grmathbox}\fi
\endgroup
\end{center}

Therefore

\begin{center}
\begingroup
\sbox{\grmathbox}{$\displaystyle 
R^i{}_{0j0}
\simeq\partial_j\Gamma^i{}_{00}
=\frac1{c^2}\partial_j\partial^i\Phi.
$}
\typeout{GRDISPLAY|215|\the\wd\grmathbox|\the\linewidth}
\ifdim\wd\grmathbox>\linewidth
\resizebox{\linewidth}{!}{\usebox{\grmathbox}}
\else\usebox{\grmathbox}\fi
\endgroup
\end{center}

Time-derivative terms have been neglected by the static approximation,
and products of weak-field connection coefficients are higher order.
With \(u^0\simeq c\) and slow reference motion, geodesic deviation gives

\begin{center}
\begingroup
\sbox{\grmathbox}{$\displaystyle 
\boxed{
\ddot\xi^i\simeq
-\partial_j\partial^i\Phi\,\xi^j.
}
$}
\typeout{GRDISPLAY|216|\the\wd\grmathbox|\the\linewidth}
\ifdim\wd\grmathbox>\linewidth
\resizebox{\linewidth}{!}{\usebox{\grmathbox}}
\else\usebox{\grmathbox}\fi
\endgroup
\end{center}

This is precisely Newtonian relative acceleration. If a nearby particle
is at \(\mathbf x+\symbfit\xi\), subtract its acceleration from the
reference equation by Taylor expanding:

\begin{center}
\begingroup
\sbox{\grmathbox}{$\displaystyle 
a^i(\mathbf x+\symbfit\xi)-a^i(\mathbf x)
=-\partial^i\Phi(\mathbf x+\symbfit\xi)
+\partial^i\Phi(\mathbf x)
\simeq-\partial_j\partial^i\Phi\,\xi^j.
$}
\typeout{GRDISPLAY|217|\the\wd\grmathbox|\the\linewidth}
\ifdim\wd\grmathbox>\linewidth
\resizebox{\linewidth}{!}{\usebox{\grmathbox}}
\else\usebox{\grmathbox}\fi
\endgroup
\end{center}

The common acceleration disappears. Only the spatial gradient of
acceleration remains. This is the mathematical content of the elevator
argument: free fall removes a shared gravitational acceleration locally;
it does not remove differences in gravitational acceleration across a
finite laboratory.

For a point mass outside its source,

\begin{center}
\begingroup
\sbox{\grmathbox}{$\displaystyle 
\Phi=-\frac{G_NM}{r},
\qquad
\partial_i\partial_j\Phi
=\frac{G_NM}{r^3}(\delta_{ij}-3n_in_j),
\qquad n_i=\frac{x_i}{r}.
$}
\typeout{GRDISPLAY|218|\the\wd\grmathbox|\the\linewidth}
\ifdim\wd\grmathbox>\linewidth
\resizebox{\linewidth}{!}{\usebox{\grmathbox}}
\else\usebox{\grmathbox}\fi
\endgroup
\end{center}

The Hessian has radial eigenvalue \(-2G_NM/r^3\) and two tangential
eigenvalues \(+G_NM/r^3\). The actual relative accelerations carry the
minus sign:

\begin{center}
\begingroup
\sbox{\grmathbox}{$\displaystyle 
\ddot\xi_{\mathrm{radial}}
=\frac{2G_NM}{r^3}\xi_{\mathrm{radial}},
\qquad
\ddot\xi_{\mathrm{tangential}}
=-\frac{G_NM}{r^3}\xi_{\mathrm{tangential}}.
$}
\typeout{GRDISPLAY|219|\the\wd\grmathbox|\the\linewidth}
\ifdim\wd\grmathbox>\linewidth
\resizebox{\linewidth}{!}{\usebox{\grmathbox}}
\else\usebox{\grmathbox}\fi
\endgroup
\end{center}

A falling cloud stretches radially and squeezes sideways. The
eigenvalues of the acceleration map add to zero outside the source. This
is the Newtonian shadow of vacuum Ricci-flatness, while the nonzero
trace-free tidal field is the shadow of Weyl curvature.

The eigenvalue ratio \(2:-1:-1\) is worth keeping. It is the geometric
reason that the phrase ``spaghettification'' has more physics in it than
its culinary dignity suggests.

\hypertarget{what-normal-coordinates-can-erase---and-what-they-cannot}{%
\section{10.6 What normal coordinates can erase - and what they
cannot}\label{what-normal-coordinates-can-erase---and-what-they-cannot}}

At any regular event \(p\), choose Riemann normal coordinates with an
orthonormal basis at the origin. Then

\begin{center}
\begingroup
\sbox{\grmathbox}{$\displaystyle 
g_{\mu\nu}(p)=\eta_{\mu\nu},
\qquad
\partial_\alpha g_{\mu\nu}(p)=0,
\qquad
\Gamma^\rho{}_{\mu\nu}(p)=0.
$}
\typeout{GRDISPLAY|220|\the\wd\grmathbox|\the\linewidth}
\ifdim\wd\grmathbox>\linewidth
\resizebox{\linewidth}{!}{\usebox{\grmathbox}}
\else\usebox{\grmathbox}\fi
\endgroup
\end{center}

The connection transformation law explains why this is possible: second
derivatives of the coordinate transformation can cancel the connection
coefficients at the chosen event. The geometric construction is to label
nearby events by the initial tangent vectors of geodesics launched from
\(p\), using the exponential map. It works in a sufficiently small
normal neighborhood, before the geodesic-labeling construction becomes
ambiguous.

But the derivatives of \(\Gamma\) generally survive. The metric
expansion is

\begin{center}
\begingroup
\sbox{\grmathbox}{$\displaystyle 
\boxed{
g_{\mu\nu}(X)
=\eta_{\mu\nu}
-\frac13R_{\mu\alpha\nu\beta}(p)X^\alpha X^\beta
+O(|X|^3).
}
$}
\typeout{GRDISPLAY|221|\the\wd\grmathbox|\the\linewidth}
\ifdim\wd\grmathbox>\linewidth
\resizebox{\linewidth}{!}{\usebox{\grmathbox}}
\else\usebox{\grmathbox}\fi
\endgroup
\end{center}

The metric is Minkowskian through first order in displacement, while
curvature enters at second order. This is the precise limitation of a
local inertial frame. Making the connection vanish at one point does not
make all second derivatives of the metric vanish there, and does not
make the surrounding neighborhood flat.

Riemann normal coordinates make the geodesics launched from the origin
straight coordinate rays. They do not make every geodesic in the
neighborhood a straight coordinate line. If they did, they would have
erased curvature rather than merely chosen convenient labels.

For a moving freely falling laboratory, \textbf{Fermi normal
coordinates} extend the construction along a reference timelike
geodesic. Parallel transport a nonrotating orthonormal frame along it,
then use short spacelike geodesics orthogonal to the reference worldline
to label nearby events. Locally, the metric is Minkowskian and its first
derivatives vanish on the entire reference geodesic, while transverse
second-order terms contain curvature. This construction and its
quadratic expansion are developed in
\href{https://doi.org/10.1063/1.1724316}{Manasse and
Misner\textquotesingle s original Fermi-coordinate paper}.

For example, with \(X^0=c\tau\) and spatial distances \(X^i\),

\begin{center}
\begingroup
\sbox{\grmathbox}{$\displaystyle 
g_{00}
=-1-R_{\hat0\hat i\hat0\hat j}(\tau)X^iX^j
+O(|\mathbf X|^3).
$}
\typeout{GRDISPLAY|222|\the\wd\grmathbox|\the\linewidth}
\ifdim\wd\grmathbox>\linewidth
\resizebox{\linewidth}{!}{\usebox{\grmathbox}}
\else\usebox{\grmathbox}\fi
\endgroup
\end{center}

At the reference worldline, the time coordinate is the
observer\textquotesingle s clock and the first-order gravitational terms
are absent. Away from it, a quadratic tidal potential remains. The
coefficient differs from the Riemann-normal expansion because these are
different coordinate constructions, one centered on an event and the
other on a worldline.

If the laboratory accelerates, its comoving nonrotating coordinates
acquire acceleration terms already at first order in distance. A rocket
can keep itself fixed in its own coordinates, but it cannot make its
accelerometer reading disappear by relabeling events.

The size of an approximately inertial laboratory is controlled by
quantities such as \(|R_{\hat a\hat b\hat c\hat d}|L^2\ll1\), together
with sufficiently small curvature-variation effects across the region. A
nominal curvature radius alone is not enough if the curvature changes
rapidly. ``Local'' is a controlled approximation, not a magic word
granting permission to neglect every inconvenient term.

\hypertarget{scalar-invariants-are-powerful-but-spacetime-has-a-loophole}{%
\section{10.7 Scalar invariants are powerful, but spacetime has a
loophole}\label{scalar-invariants-are-powerful-but-spacetime-has-a-loophole}}

Coordinates can become singular while geometry remains regular. To
distinguish a coordinate problem from a physical curvature problem, form
scalars such as

\begin{center}
\begingroup
\sbox{\grmathbox}{$\displaystyle 
R,
\qquad
R_{\mu\nu}R^{\mu\nu},
\qquad
\mathcal K=R_{\rho\sigma\mu\nu}R^{\rho\sigma\mu\nu}.
$}
\typeout{GRDISPLAY|223|\the\wd\grmathbox|\the\linewidth}
\ifdim\wd\grmathbox>\linewidth
\resizebox{\linewidth}{!}{\usebox{\grmathbox}}
\else\usebox{\grmathbox}\fi
\endgroup
\end{center}

The last is the Kretschmann scalar. In Schwarzschild spacetime,

\begin{center}
\begingroup
\sbox{\grmathbox}{$\displaystyle 
\mathcal K=\frac{48G_N^2M^2}{c^4r^6}.
$}
\typeout{GRDISPLAY|224|\the\wd\grmathbox|\the\linewidth}
\ifdim\wd\grmathbox>\linewidth
\resizebox{\linewidth}{!}{\usebox{\grmathbox}}
\else\usebox{\grmathbox}\fi
\endgroup
\end{center}

It stays finite at \(r=2G_NM/c^2\) and diverges as \(r\to0\). This helps
show why the standard Schwarzschild-coordinate problem at the horizon is
removable, while the central curvature divergence is not. Establishing
smooth extension through a horizon still requires regular coordinates;
finiteness of one scalar by itself is not an extension theorem.

Now the deeper gotcha: \textbf{even all polynomial scalar curvature
invariants can vanish while the Riemann tensor is nonzero}. Lorentzian
contractions are not positive sums of squares. A nonzero null vector
already demonstrates the basic logic: its norm can be zero without the
vector being zero.

Here is an exact four-dimensional example. Let \(U,V\) be null
coordinates and take

\begin{center}
\begingroup
\sbox{\grmathbox}{$\displaystyle 
ds^2=-2\,dU\,dV+dx^2+dy^2+H(U,x,y)\,dU^2,
$}
\typeout{GRDISPLAY|225|\the\wd\grmathbox|\the\linewidth}
\ifdim\wd\grmathbox>\linewidth
\resizebox{\linewidth}{!}{\usebox{\grmathbox}}
\else\usebox{\grmathbox}\fi
\endgroup
\end{center}

where

\begin{center}
\begingroup
\sbox{\grmathbox}{$\displaystyle 
H=A(U)(x^2-y^2)+2B(U)xy.
$}
\typeout{GRDISPLAY|226|\the\wd\grmathbox|\the\linewidth}
\ifdim\wd\grmathbox>\linewidth
\resizebox{\linewidth}{!}{\usebox{\grmathbox}}
\else\usebox{\grmathbox}\fi
\endgroup
\end{center}

Its curvature includes

\begin{center}
\begingroup
\sbox{\grmathbox}{$\displaystyle 
R_{UiUj}=-\frac12\partial_i\partial_jH,
\qquad i,j\in\{x,y\},
$}
\typeout{GRDISPLAY|227|\the\wd\grmathbox|\the\linewidth}
\ifdim\wd\grmathbox>\linewidth
\resizebox{\linewidth}{!}{\usebox{\grmathbox}}
\else\usebox{\grmathbox}\fi
\endgroup
\end{center}

so it is nonzero whenever the profiles \(A\) or \(B\) are nonzero. Yet

\begin{center}
\begingroup
\sbox{\grmathbox}{$\displaystyle 
R_{UU}=-\frac12(\partial_x^2H+\partial_y^2H)=0,
$}
\typeout{GRDISPLAY|228|\the\wd\grmathbox|\the\linewidth}
\ifdim\wd\grmathbox>\linewidth
\resizebox{\linewidth}{!}{\usebox{\grmathbox}}
\else\usebox{\grmathbox}\fi
\endgroup
\end{center}

and the other Ricci components vanish: this is a vacuum plane
gravitational wave.

Why does \(\mathcal K\) vanish too? Every nonzero curvature component
has lower \(U\) indices, but the inverse metric has \(g^{UU}=0\).
Raising a \(U\) slot pairs it with a \(V\) slot, and there are no
corresponding nonzero curvature components containing \(V\). The full
contraction therefore vanishes even though the tidal tensor does not.
These plane waves belong to the class of geometries with vanishing
scalar polynomial invariants; the broader classification is given by
\href{https://arxiv.org/abs/gr-qc/0209024}{Pravda, Pravdova, Coley, and
Milson}.

A vanishing scalar is a statement about that scalar. It is not a
permission slip to replace the tensor by zero. For difficult spacetime
classification or singularity questions, curvature components in
physically or geometrically specified frames, covariant derivatives,
geodesic behavior, and extension properties may all matter.

\hypertarget{the-chain-of-ideas-you-now-own}{%
\section{10.8 The chain of ideas you now
own}\label{the-chain-of-ideas-you-now-own}}

You can now move through the geometric core without treating any symbol
as a ceremonial object:

\begingroup\small\setlength{\tabcolsep}{5pt}\renewcommand{\arraystretch}{1.13}

\begin{longtable}[]{@{}
  >{\raggedright\arraybackslash}p{(\columnwidth - 4\tabcolsep) * \real{0.2500}}
  >{\raggedright\arraybackslash}p{(\columnwidth - 4\tabcolsep) * \real{0.3600}}
  >{\raggedright\arraybackslash}p{(\columnwidth - 4\tabcolsep) * \real{0.3900}}@{}}
\toprule\noalign{}
\begin{minipage}[b]{\linewidth}\raggedright
Object
\end{minipage} & \begin{minipage}[b]{\linewidth}\raggedright
What it lets you ask
\end{minipage} & \begin{minipage}[b]{\linewidth}\raggedright
What it does not imply by itself
\end{minipage} \\
\midrule\noalign{}
\endhead
\bottomrule\noalign{}
\endlastfoot
\(g_{\mu\nu}\) & What are intervals, inner products, and causal
directions? & Changing components do not by themselves prove
curvature. \\
\(\Gamma^\rho{}_{\mu\nu}\) & How does this connection compare
neighboring directions in this chart? & Nonzero coefficients do not
prove a curved geometry. \\
\(R^\rho{}_{\sigma\mu\nu}\) & How does transport depend on route, and
how do free paths deviate? & Twenty algebraic components do not mean
twenty propagating modes. \\
\(R_{\mu\nu}\) & What directional traces of curvature survive
contraction? & Zero Ricci does not rule out vacuum tides. \\
\(R\) & What fully contracted curvature remains? & Zero scalar curvature
does not mean flatness. \\
\(C_{\rho\sigma\mu\nu}\) & What trace-free curvature remains after Ricci
is removed? & Trace-free initial tides do not preserve a
cloud\textquotesingle s volume forever. \\
\(G_{\mu\nu}\) & Which Ricci combination has identically zero covariant
divergence? & Zero divergence is not constancy in every direction. \\
\end{longtable}

\endgroup

The next step is physical rather than merely geometrical: identify the
tensor that describes matter\textquotesingle s energy, momentum, and
stresses, then find the dynamical equation and action that relate it to
the geometry you have just learned to measure.

\hypertarget{chapter-11}{%
\chapter{11. Energy, momentum, and stress: what gravity listens
to}\label{chapter-11}}

Geometry has acquired a vocabulary. It can describe clocks, freely
falling trajectories, and tidal stretching. But it has not yet acquired
a law telling it which geometry to adopt. We now need the other
protagonist: matter, including radiation and nongravitational fields.

A useful first surprise is that gravity does not ask matter for a single
number called ``how much stuff?'' It asks for a tensor. Energy can move;
moving energy carries momentum; momentum itself can flow. A source
description that ignores those flows would disagree with itself as soon
as two observers passed each other.

\hypertarget{a-tensor-is-a-shipping-manifest-for-energy-and-momentum}{%
\section{11.1 A tensor is a shipping manifest for energy and
momentum}\label{a-tensor-is-a-shipping-manifest-for-energy-and-momentum}}

Start in a tiny laboratory with orthonormal axes, using
\(x^{\hat 0}=ct\). Hats distinguish physically calibrated local axes
from arbitrary coordinate axes. Let

\begin{itemize}
\tightlist
\item
  \(\epsilon\) be energy per volume, in \(\mathrm{J/m^3}\);
\item
  \(S_i\) be energy flux, in \(\mathrm{J/(m^2\,s)}\);
\item
  \(\pi_i\) be momentum per volume, in \(\mathrm{kg/(m^2\,s)}\);
\item
  \(\Pi_{ij}\) be the flux of \(j\)-momentum through a surface normal to
  direction \(i\), in \(\mathrm{N/m^2}\).
\end{itemize}

The stress-energy tensor packages these quantities as

\begin{center}
\begingroup
\sbox{\grmathbox}{$\displaystyle 
T^{\hat\mu\hat\nu}=
\begin{pmatrix}
\epsilon & c\pi_x & c\pi_y & c\pi_z\\
S_x/c & \Pi_{xx} & \Pi_{xy} & \Pi_{xz}\\
S_y/c & \Pi_{yx} & \Pi_{yy} & \Pi_{yz}\\
S_z/c & \Pi_{zx} & \Pi_{zy} & \Pi_{zz}
\end{pmatrix}.
$}
\typeout{GRDISPLAY|229|\the\wd\grmathbox|\the\linewidth}
\ifdim\wd\grmathbox>\linewidth
\resizebox{\linewidth}{!}{\usebox{\grmathbox}}
\else\usebox{\grmathbox}\fi
\endgroup
\end{center}

Every entry has energy-density units. In this display, the second index
identifies which energy-momentum component is being tracked; the first
identifies the direction through which it flows. The time row records
what is present in a spatial volume. The spatial rows record transport
through its faces.

For the symmetric stress tensor appearing in ordinary metric general
relativity,

\begin{center}
\begingroup
\sbox{\grmathbox}{$\displaystyle 
T^{\mu\nu}=T^{\nu\mu},
\qquad
\boxed{\mathbf S=c^2\symbfit\pi.}
$}
\typeout{GRDISPLAY|230|\the\wd\grmathbox|\the\linewidth}
\ifdim\wd\grmathbox>\linewidth
\resizebox{\linewidth}{!}{\usebox{\grmathbox}}
\else\usebox{\grmathbox}\fi
\endgroup
\end{center}

This is a local relationship between energy flow and momentum density.
It is one of the most useful ways to recognize that energy and momentum
belong to one relativistic object. A beam of light has momentum because
it transports energy; a moving lump of matter does too.

Check the packaging by taking a divergence in an inertial Cartesian
laboratory:

\begin{center}
\begingroup
\sbox{\grmathbox}{$\displaystyle 
\partial_\mu T^{\mu\nu}=0.
$}
\typeout{GRDISPLAY|231|\the\wd\grmathbox|\the\linewidth}
\ifdim\wd\grmathbox>\linewidth
\resizebox{\linewidth}{!}{\usebox{\grmathbox}}
\else\usebox{\grmathbox}\fi
\endgroup
\end{center}

Set \(\nu=0\). Since \(\partial_0=c^{-1}\partial_t\),

\begin{center}
\begingroup
\sbox{\grmathbox}{$\displaystyle 
\frac1c\frac{\partial\epsilon}{\partial t}
+\partial_i\left(\frac{S_i}{c}\right)=0,
$}
\typeout{GRDISPLAY|232|\the\wd\grmathbox|\the\linewidth}
\ifdim\wd\grmathbox>\linewidth
\resizebox{\linewidth}{!}{\usebox{\grmathbox}}
\else\usebox{\grmathbox}\fi
\endgroup
\end{center}

or

\begin{center}
\begingroup
\sbox{\grmathbox}{$\displaystyle 
\boxed{\frac{\partial\epsilon}{\partial t}+\symbfit\nabla\cdot\mathbf S=0.}
$}
\typeout{GRDISPLAY|233|\the\wd\grmathbox|\the\linewidth}
\ifdim\wd\grmathbox>\linewidth
\resizebox{\linewidth}{!}{\usebox{\grmathbox}}
\else\usebox{\grmathbox}\fi
\endgroup
\end{center}

Energy disappearing from a little box must flow through its faces. Now
set \(\nu=j\):

\begin{center}
\begingroup
\sbox{\grmathbox}{$\displaystyle 
\boxed{\frac{\partial\pi_j}{\partial t}+\partial_i\Pi_{ij}=0.}
$}
\typeout{GRDISPLAY|234|\the\wd\grmathbox|\the\linewidth}
\ifdim\wd\grmathbox>\linewidth
\resizebox{\linewidth}{!}{\usebox{\grmathbox}}
\else\usebox{\grmathbox}\fi
\endgroup
\end{center}

That is the momentum balance law. A pressure gradient changes a
fluid\textquotesingle s momentum because one face receives a different
momentum flux from the opposite face. The tensor is four conservation
equations sharing one filing system.

\textbf{Gotcha: a coordinate component is not automatically a meter
reading.} In spherical coordinates, \(T^{\theta\theta}\) is not a
pressure read directly from a gauge aligned with an angular ruler. The
coordinate basis has its own normalization. Physical readings come from
contraction with an observer\textquotesingle s orthonormal frame.

\textbf{Another gotcha: index position changes signs.} In an orthonormal
frame with signature \((-,+,+,+)\), lowering a single time index changes
its sign. Thus a fluid at rest has

\begin{center}
\begingroup
\sbox{\grmathbox}{$\displaystyle 
T^{\hat\mu\hat\nu}=\operatorname{diag}(\epsilon,p,p,p),
\qquad
T^{\hat\mu}{}_{\hat\nu}=\operatorname{diag}(-\epsilon,p,p,p).
$}
\typeout{GRDISPLAY|235|\the\wd\grmathbox|\the\linewidth}
\ifdim\wd\grmathbox>\linewidth
\resizebox{\linewidth}{!}{\usebox{\grmathbox}}
\else\usebox{\grmathbox}\fi
\endgroup
\end{center}

The energy has not become negative. You changed the kind of tensor
components you were looking at.

\hypertarget{energy-belongs-to-an-observer-the-tensor-belongs-to-everyone}{%
\section{11.2 Energy belongs to an observer; the tensor belongs to
everyone}\label{energy-belongs-to-an-observer-the-tensor-belongs-to-everyone}}

A particle\textquotesingle s energy depends on who measures it. So does
a field\textquotesingle s energy density. The covariant object is
\(T_{\mu\nu}\); ``energy density'' is one observer\textquotesingle s
projection of it.

Let an observer have four-velocity \(w^\mu\), with \(w_\mu w^\mu=-c^2\),
and define the dimensionless unit timelike vector

\begin{center}
\begingroup
\sbox{\grmathbox}{$\displaystyle 
n^\mu=\frac{w^\mu}{c},\qquad n_\mu n^\mu=-1.
$}
\typeout{GRDISPLAY|236|\the\wd\grmathbox|\the\linewidth}
\ifdim\wd\grmathbox>\linewidth
\resizebox{\linewidth}{!}{\usebox{\grmathbox}}
\else\usebox{\grmathbox}\fi
\endgroup
\end{center}

The observer\textquotesingle s spatial projector is

\begin{center}
\begingroup
\sbox{\grmathbox}{$\displaystyle 
H^\mu{}_{\nu}=\delta^\mu{}_{\nu}+n^\mu n_\nu.
$}
\typeout{GRDISPLAY|237|\the\wd\grmathbox|\the\linewidth}
\ifdim\wd\grmathbox>\linewidth
\resizebox{\linewidth}{!}{\usebox{\grmathbox}}
\else\usebox{\grmathbox}\fi
\endgroup
\end{center}

Why the plus sign? Apply it to \(n^\nu\):

\begin{center}
\begingroup
\sbox{\grmathbox}{$\displaystyle 
H^\mu{}_{\nu}n^\nu=n^\mu+n^\mu(-1)=0.
$}
\typeout{GRDISPLAY|238|\the\wd\grmathbox|\the\linewidth}
\ifdim\wd\grmathbox>\linewidth
\resizebox{\linewidth}{!}{\usebox{\grmathbox}}
\else\usebox{\grmathbox}\fi
\endgroup
\end{center}

It removes the time direction, leaving precisely the
observer\textquotesingle s local three-dimensional rest space. The
observer measures

\begin{center}
\begingroup
\sbox{\grmathbox}{$\displaystyle 
\epsilon_{(n)}=T_{\mu\nu}n^\mu n^\nu,
$}
\typeout{GRDISPLAY|239|\the\wd\grmathbox|\the\linewidth}
\ifdim\wd\grmathbox>\linewidth
\resizebox{\linewidth}{!}{\usebox{\grmathbox}}
\else\usebox{\grmathbox}\fi
\endgroup
\end{center}

\begin{center}
\begingroup
\sbox{\grmathbox}{$\displaystyle 
j^\mu=-H^\mu{}_{\alpha}T^{\alpha\beta}n_\beta,
\qquad
P^{\mu\nu}=H^\mu{}_{\alpha}H^\nu{}_{\beta}T^{\alpha\beta}.
$}
\typeout{GRDISPLAY|240|\the\wd\grmathbox|\the\linewidth}
\ifdim\wd\grmathbox>\linewidth
\resizebox{\linewidth}{!}{\usebox{\grmathbox}}
\else\usebox{\grmathbox}\fi
\endgroup
\end{center}

Here \(j^\mu\) is spatial and has energy-density units: the physical
energy-flux vector is \(c j^\mu\), and momentum density is \(j^\mu/c\).
The tensor \(P^{\mu\nu}\) is the spatial stress seen by that observer.
Together they reconstruct the whole tensor:

\begin{center}
\begingroup
\sbox{\grmathbox}{$\displaystyle 
\boxed{
T^{\mu\nu}=\epsilon_{(n)}n^\mu n^\nu
+n^\mu j^\nu+j^\mu n^\nu+P^{\mu\nu}.
}
$}
\typeout{GRDISPLAY|241|\the\wd\grmathbox|\the\linewidth}
\ifdim\wd\grmathbox>\linewidth
\resizebox{\linewidth}{!}{\usebox{\grmathbox}}
\else\usebox{\grmathbox}\fi
\endgroup
\end{center}

This is an unusually productive way to think about a tensor: choose your
local time direction, and it decomposes into familiar laboratory
quantities. Change observer, and the pieces reshuffle while the
underlying tensor remains the same.

The analogy is slicing a loaf at different angles. The slices look
different because your slicing procedure changed. Its limitation is that
Lorentz transformations mix time and space with a hyperbolic geometry,
so an ordinary Euclidean loaf does not reproduce the quantitative
transformation law. It is a reminder about observer-dependent
decomposition, not a model of spacetime.

Some matter admits a local rest frame with zero energy flux. A single
beam of light does not: catching up with the beam would require a
timelike observer to become null. Never assume that every stress tensor
can be put into a perfect-fluid rest-frame form.

\hypertarget{dust-a-crowd-that-refuses-to-interact}{%
\section{11.3 Dust: a crowd that refuses to
interact}\label{dust-a-crowd-that-refuses-to-interact}}

``Dust'' in relativity means an idealized collection of particles with
negligible pressure and random velocity dispersion. It can represent
cold matter on suitable scales; it does not mean that the universe has
neglected its housekeeping.

In the common rest frame of a small dust element, only the
energy-density component is nonzero. If its four-velocity is \(u^\mu\),
the covariant expression reproducing this is

\begin{center}
\begingroup
\sbox{\grmathbox}{$\displaystyle 
\boxed{
T^{\mu\nu}_{\mathrm{dust}}=\frac{\epsilon}{c^2}u^\mu u^\nu
=\rho u^\mu u^\nu,
\qquad \rho=\frac{\epsilon}{c^2}.
}
$}
\typeout{GRDISPLAY|242|\the\wd\grmathbox|\the\linewidth}
\ifdim\wd\grmathbox>\linewidth
\resizebox{\linewidth}{!}{\usebox{\grmathbox}}
\else\usebox{\grmathbox}\fi
\endgroup
\end{center}

In the rest frame \(u^{\hat\mu}=(c,0,0,0)\), so
\(T^{\hat0\hat0}=\rho c^2=\epsilon\). A moving observer sees energy flux
and spatial momentum flux automatically, because \(u^i\) is then
nonzero.

The curved-spacetime conservation equation is

\begin{center}
\begingroup
\sbox{\grmathbox}{$\displaystyle 
0=\nabla_\mu T^{\mu\nu}
=u^\nu\nabla_\mu(\rho u^\mu)
+\rho u^\mu\nabla_\mu u^\nu.
$}
\typeout{GRDISPLAY|243|\the\wd\grmathbox|\the\linewidth}
\ifdim\wd\grmathbox>\linewidth
\resizebox{\linewidth}{!}{\usebox{\grmathbox}}
\else\usebox{\grmathbox}\fi
\endgroup
\end{center}

Contract this with \(u_\nu\). The second term vanishes because

\begin{center}
\begingroup
\sbox{\grmathbox}{$\displaystyle 
u_\nu\nabla_\mu u^\nu=\frac12\nabla_\mu(u_\nu u^\nu)=0.
$}
\typeout{GRDISPLAY|244|\the\wd\grmathbox|\the\linewidth}
\ifdim\wd\grmathbox>\linewidth
\resizebox{\linewidth}{!}{\usebox{\grmathbox}}
\else\usebox{\grmathbox}\fi
\endgroup
\end{center}

The first term therefore gives

\begin{center}
\begingroup
\sbox{\grmathbox}{$\displaystyle 
\nabla_\mu(\rho u^\mu)=0.
$}
\typeout{GRDISPLAY|245|\the\wd\grmathbox|\the\linewidth}
\ifdim\wd\grmathbox>\linewidth
\resizebox{\linewidth}{!}{\usebox{\grmathbox}}
\else\usebox{\grmathbox}\fi
\endgroup
\end{center}

Insert that result back into the uncontracted equation. Wherever
\(\rho\ne0\),

\begin{center}
\begingroup
\sbox{\grmathbox}{$\displaystyle 
u^\mu\nabla_\mu u^\nu=0.
$}
\typeout{GRDISPLAY|246|\the\wd\grmathbox|\the\linewidth}
\ifdim\wd\grmathbox>\linewidth
\resizebox{\linewidth}{!}{\usebox{\grmathbox}}
\else\usebox{\grmathbox}\fi
\endgroup
\end{center}

The same conservation law has produced both continuity of dust mass and
geodesic motion. Geometry and matter dynamics fit together; geodesic
motion is not an unrelated ornament attached to the field equations.

The qualification matters. Charged matter exchanging momentum with an
electromagnetic field need not have a separately conserved matter stress
tensor. A pressured fluid accelerates because neighboring fluid elements
push it. Extended spinning bodies can have curvature-dependent
corrections to simple geodesic motion. Dust is a controlled
idealization, not a universal description of matter.

\hypertarget{perfect-fluids-derived-from-isotropy}{%
\section{11.4 Perfect fluids, derived from
isotropy}\label{perfect-fluids-derived-from-isotropy}}

A perfect fluid has no heat flux, viscosity, or preferred spatial
direction in its local rest frame. Rotational symmetry therefore
requires its spatial stress to be \(p\delta^{ij}\).

Let

\begin{center}
\begingroup
\sbox{\grmathbox}{$\displaystyle 
H^{\mu\nu}=g^{\mu\nu}+\frac{u^\mu u^\nu}{c^2}
$}
\typeout{GRDISPLAY|247|\the\wd\grmathbox|\the\linewidth}
\ifdim\wd\grmathbox>\linewidth
\resizebox{\linewidth}{!}{\usebox{\grmathbox}}
\else\usebox{\grmathbox}\fi
\endgroup
\end{center}

be the projector onto the fluid\textquotesingle s rest space. Its energy
contribution must be \(\epsilon u^\mu u^\nu/c^2\); its isotropic spatial
stress must be \(pH^{\mu\nu}\). Adding them gives

\begin{center}
\begingroup
\sbox{\grmathbox}{$\displaystyle 
\boxed{
T^{\mu\nu}=(\epsilon+p)\frac{u^\mu u^\nu}{c^2}
+p g^{\mu\nu}.
}
$}
\typeout{GRDISPLAY|248|\the\wd\grmathbox|\the\linewidth}
\ifdim\wd\grmathbox>\linewidth
\resizebox{\linewidth}{!}{\usebox{\grmathbox}}
\else\usebox{\grmathbox}\fi
\endgroup
\end{center}

This derivation explains the apparently mysterious \(\epsilon+p\).
Pressure initially entered as a purely spatial stress. Expressing
``spatial'' covariantly requires a projector containing \(u^\mu u^\nu\),
so pressure also joins the coefficient of the velocity term.

For an observer moving at relative speed \(v\) with respect to the
fluid,

\begin{center}
\begingroup
\sbox{\grmathbox}{$\displaystyle 
\epsilon_{\mathrm{measured}}=(\epsilon+p)\gamma^2-p,
\qquad
\gamma=\frac{1}{\sqrt{1-v^2/c^2}}.
$}
\typeout{GRDISPLAY|249|\the\wd\grmathbox|\the\linewidth}
\ifdim\wd\grmathbox>\linewidth
\resizebox{\linewidth}{!}{\usebox{\grmathbox}}
\else\usebox{\grmathbox}\fi
\endgroup
\end{center}

This follows by contracting \(T_{\mu\nu}\) with the
observer\textquotesingle s unit time vector and using
\(u\cdot w=-c^2\gamma\). Even the energy density measured by a moving
observer knows about pressure.

Projecting \(\nabla_\mu T^{\mu\nu}=0\) parallel to \(u_\nu\) gives

\begin{center}
\begingroup
\sbox{\grmathbox}{$\displaystyle 
\boxed{
u^\mu\nabla_\mu\epsilon+(\epsilon+p)\theta=0,
\qquad \theta=\nabla_\mu u^\mu.}
$}
\typeout{GRDISPLAY|250|\the\wd\grmathbox|\the\linewidth}
\ifdim\wd\grmathbox>\linewidth
\resizebox{\linewidth}{!}{\usebox{\grmathbox}}
\else\usebox{\grmathbox}\fi
\endgroup
\end{center}

Here \(\theta\) measures fractional expansion of a small comoving
volume. Write its proper volume as \(V\) so that \(dV/d\tau=\theta V\).
Multiplying the equation by \(V\) yields

\begin{center}
\begingroup
\sbox{\grmathbox}{$\displaystyle 
\frac{d(\epsilon V)}{d\tau}=-p\frac{dV}{d\tau}.
$}
\typeout{GRDISPLAY|251|\the\wd\grmathbox|\the\linewidth}
\ifdim\wd\grmathbox>\linewidth
\resizebox{\linewidth}{!}{\usebox{\grmathbox}}
\else\usebox{\grmathbox}\fi
\endgroup
\end{center}

It is the familiar pressure-volume work law, now embedded in
relativistic geometry.

Projecting perpendicular to \(u^\mu\) instead gives

\begin{center}
\begingroup
\sbox{\grmathbox}{$\displaystyle 
\boxed{
\frac{\epsilon+p}{c^2}a^\alpha
=-H^{\alpha\mu}\nabla_\mu p,
\qquad a^\alpha=u^\mu\nabla_\mu u^\alpha.
}
$}
\typeout{GRDISPLAY|252|\the\wd\grmathbox|\the\linewidth}
\ifdim\wd\grmathbox>\linewidth
\resizebox{\linewidth}{!}{\usebox{\grmathbox}}
\else\usebox{\grmathbox}\fi
\endgroup
\end{center}

The pressure gradient supplies the force density; \((\epsilon+p)/c^2\)
supplies the relativistic inertial coefficient. In a cold, slow fluid,
\(p\ll\epsilon\simeq\rho c^2\), recovering
\(\rho\mathbf a=-\symbfit\nabla p\) locally.

\textbf{Stress-sign trap.} Relativity\textquotesingle s \(T^{ij}\) is
momentum flux. Some engineering conventions define Cauchy stress as
positive in tension, giving a stationary fluid the mechanical stress
\(-p\delta^{ij}\). The sign difference is bookkeeping, not a
disagreement about which way a pressurized piston moves.

\hypertarget{pressure-gravitates---and-the-photon-box-is-not-a-mass-doubling-machine}{%
\section{11.5 Pressure gravitates - and the photon box is not a
mass-doubling
machine}\label{pressure-gravitates---and-the-photon-box-is-not-a-mass-doubling-machine}}

The field equation developed next implies, in an orthonormal frame
comoving with an isotropic fluid,

\begin{center}
\begingroup
\sbox{\grmathbox}{$\displaystyle 
R_{\hat0\hat0}
=\frac{4\pi G_N}{c^4}(\epsilon+3p)-\Lambda.
$}
\typeout{GRDISPLAY|253|\the\wd\grmathbox|\the\linewidth}
\ifdim\wd\grmathbox>\linewidth
\resizebox{\linewidth}{!}{\usebox{\grmathbox}}
\else\usebox{\grmathbox}\fi
\endgroup
\end{center}

The factor \(3p\) is the sum of the stresses in three spatial
directions. It describes a particular Ricci-curvature projection
relevant to focusing timelike geodesics. It is not a universal
instruction to replace all mass densities everywhere by
\((\epsilon+3p)/c^2\).

For isotropic radiation, \(p=\epsilon/3\). One can see why: a photon
carries momentum \(E/c\), and its momentum crossing a surface brings
another directional factor. Isotropy gives the average
\(\langle\cos^2\theta\rangle=1/3\). Hence each diagonal momentum flux is
one third of the energy density.

Now comes the trap: \(\epsilon+3p=2\epsilon\). Does putting photons into
a box make their gravitational mass twice \(E/c^2\)?

The missing member of that calculation is the box. Radiation pushes
outward; walls develop stresses to hold it in. An isolated static system
must include its supports in its total stress tensor. Wall tensions
compensate the extra integrated radiation-pressure contribution in the
regime of negligible internal self-gravity. This is the resolution
studied explicitly by
\href{https://link.aps.org/doi/10.1103/PhysRev.116.1045}{Misner and
Putnam in ``Active Gravitational Mass''}.

We can expose the mechanism with a short independent calculation. For a
localized stationary system in approximately flat spacetime, total
momentum conservation says

\begin{center}
\begingroup
\sbox{\grmathbox}{$\displaystyle 
\partial_kT^{kj}=0.
$}
\typeout{GRDISPLAY|254|\the\wd\grmathbox|\the\linewidth}
\ifdim\wd\grmathbox>\linewidth
\resizebox{\linewidth}{!}{\usebox{\grmathbox}}
\else\usebox{\grmathbox}\fi
\endgroup
\end{center}

Multiply by \(x^i\) and use the product rule:

\begin{center}
\begingroup
\sbox{\grmathbox}{$\displaystyle 
\partial_k(x^iT^{kj})=T^{ij}.
$}
\typeout{GRDISPLAY|255|\the\wd\grmathbox|\the\linewidth}
\ifdim\wd\grmathbox>\linewidth
\resizebox{\linewidth}{!}{\usebox{\grmathbox}}
\else\usebox{\grmathbox}\fi
\endgroup
\end{center}

Integrate over all space. If the complete system\textquotesingle s
stresses decay sufficiently fast, the surface term vanishes, leaving

\begin{center}
\begingroup
\sbox{\grmathbox}{$\displaystyle 
\int T^{ij}\,d^3x=0.
$}
\typeout{GRDISPLAY|256|\the\wd\grmathbox|\the\linewidth}
\ifdim\wd\grmathbox>\linewidth
\resizebox{\linewidth}{!}{\usebox{\grmathbox}}
\else\usebox{\grmathbox}\fi
\endgroup
\end{center}

The positive pressure stresses of the contents cannot be the whole
story: other stresses must balance their integral. Consequently the
leading static weak-field mass contribution

\begin{center}
\begingroup
\sbox{\grmathbox}{$\displaystyle 
\frac1{c^2}\int\left(T^{00}+\sum_iT^{ii}\right)d^3x
$}
\typeout{GRDISPLAY|257|\the\wd\grmathbox|\the\linewidth}
\ifdim\wd\grmathbox>\linewidth
\resizebox{\linewidth}{!}{\usebox{\grmathbox}}
\else\usebox{\grmathbox}\fi
\endgroup
\end{center}

reduces to total energy divided by \(c^2\) for the complete system under
these assumptions. If you supply additional energy \(E\) without
otherwise changing the total energy accounting, the leading mass
increase is \(E/c^2\). Gravitational binding corrections require the
appropriate relativistic total-energy definition.

The lesson is sharper than ``pressure does not really gravitate.'' It
does. The lesson is: \textbf{do not compute the gravity of a
mechanically incomplete system and mistake the answer for the gravity of
the complete one.}

\hypertarget{two-field-examples-scalar-waves-and-electromagnetism}{%
\section{11.6 Two field examples: scalar waves and
electromagnetism}\label{two-field-examples-scalar-waves-and-electromagnetism}}

For this scalar-field subsection only, use natural units \(c=\hbar=1\).
A canonically normalized real scalar field has Lagrangian

\begin{center}
\begingroup
\sbox{\grmathbox}{$\displaystyle 
\mathcal L_\phi=-\frac12g^{\mu\nu}\partial_\mu\phi\partial_\nu\phi-V(\phi).
$}
\typeout{GRDISPLAY|258|\the\wd\grmathbox|\the\linewidth}
\ifdim\wd\grmathbox>\linewidth
\resizebox{\linewidth}{!}{\usebox{\grmathbox}}
\else\usebox{\grmathbox}\fi
\endgroup
\end{center}

In four dimensions, \(\phi\) has mass dimension one; \(V\) and
\(T_{\mu\nu}\) have mass dimension four. The sign of the kinetic term is
chosen for signature \((-,+,+,+)\), so time-dependent excitations carry
positive kinetic energy.

The metric-variation method of Chapter 13 gives

\begin{center}
\begingroup
\sbox{\grmathbox}{$\displaystyle 
T_{\mu\nu}^{(\phi)}
=\partial_\mu\phi\partial_\nu\phi
-g_{\mu\nu}\left[\frac12(\partial\phi)^2+V(\phi)\right],
$}
\typeout{GRDISPLAY|259|\the\wd\grmathbox|\the\linewidth}
\ifdim\wd\grmathbox>\linewidth
\resizebox{\linewidth}{!}{\usebox{\grmathbox}}
\else\usebox{\grmathbox}\fi
\endgroup
\end{center}

where
\((\partial\phi)^2=g^{\alpha\beta}\partial_\alpha\phi\partial_\beta\phi\).
In a local inertial frame,

\begin{center}
\begingroup
\sbox{\grmathbox}{$\displaystyle 
\epsilon_\phi=\frac12\dot\phi^2+\frac12|\symbfit\nabla\phi|^2+V.
$}
\typeout{GRDISPLAY|260|\the\wd\grmathbox|\the\linewidth}
\ifdim\wd\grmathbox>\linewidth
\resizebox{\linewidth}{!}{\usebox{\grmathbox}}
\else\usebox{\grmathbox}\fi
\endgroup
\end{center}

Time variation, spatial gradients, and potential energy all contribute.
For a homogeneous field in its comoving frame,

\begin{center}
\begingroup
\sbox{\grmathbox}{$\displaystyle 
p_\phi=\frac12\dot\phi^2-V.
$}
\typeout{GRDISPLAY|261|\the\wd\grmathbox|\the\linewidth}
\ifdim\wd\grmathbox>\linewidth
\resizebox{\linewidth}{!}{\usebox{\grmathbox}}
\else\usebox{\grmathbox}\fi
\endgroup
\end{center}

Kinetic energy gives positive pressure; potential energy gives negative
pressure. This prepares the intuition for scalar-field cosmology without
requiring a field to be imagined as a literal fluid of tiny balls.

\textbf{Return to SI units and \(x^0=ct\).} For electromagnetism choose

\begin{center}
\begingroup
\sbox{\grmathbox}{$\displaystyle 
F_{0i}=-\frac{E_i}{c},\qquad F_{ij}=\varepsilon_{ijk}B_k,
$}
\typeout{GRDISPLAY|262|\the\wd\grmathbox|\the\linewidth}
\ifdim\wd\grmathbox>\linewidth
\resizebox{\linewidth}{!}{\usebox{\grmathbox}}
\else\usebox{\grmathbox}\fi
\endgroup
\end{center}

in a local orthonormal frame. Here \(\varepsilon_{ijk}\) is the
three-dimensional antisymmetric symbol, unrelated to energy density
\(\epsilon\). The electromagnetic Lagrangian energy density and stress
tensor are

\begin{center}
\begingroup
\sbox{\grmathbox}{$\displaystyle 
\mathcal L_{\mathrm{EM}}=-\frac1{4\mu_0}F_{\alpha\beta}F^{\alpha\beta},
$}
\typeout{GRDISPLAY|263|\the\wd\grmathbox|\the\linewidth}
\ifdim\wd\grmathbox>\linewidth
\resizebox{\linewidth}{!}{\usebox{\grmathbox}}
\else\usebox{\grmathbox}\fi
\endgroup
\end{center}

\begin{center}
\begingroup
\sbox{\grmathbox}{$\displaystyle 
\boxed{
T_{\mu\nu}^{\mathrm{EM}}
=\frac1{\mu_0}\left(
F_{\mu\alpha}F_\nu{}^\alpha
-\frac14g_{\mu\nu}F_{\alpha\beta}F^{\alpha\beta}
\right).
}
$}
\typeout{GRDISPLAY|264|\the\wd\grmathbox|\the\linewidth}
\ifdim\wd\grmathbox>\linewidth
\resizebox{\linewidth}{!}{\usebox{\grmathbox}}
\else\usebox{\grmathbox}\fi
\endgroup
\end{center}

These conventions produce

\begin{center}
\begingroup
\sbox{\grmathbox}{$\displaystyle 
\epsilon_{\mathrm{EM}}=
\frac{\epsilon_0E^2}{2}+\frac{B^2}{2\mu_0},
\qquad
\mathbf S=\frac1{\mu_0}\mathbf E\times\mathbf B,
\qquad
\epsilon_0\mu_0c^2=1.
$}
\typeout{GRDISPLAY|265|\the\wd\grmathbox|\the\linewidth}
\ifdim\wd\grmathbox>\linewidth
\resizebox{\linewidth}{!}{\usebox{\grmathbox}}
\else\usebox{\grmathbox}\fi
\endgroup
\end{center}

You recognize the energy density and Poynting vector from
electromagnetism. The same tensor includes electric and magnetic
stresses, so a magnetic field can affect geometry even in a region
containing no material particles.

Its classical trace in four spacetime dimensions vanishes:

\begin{center}
\begingroup
\sbox{\grmathbox}{$\displaystyle 
T^\mu{}_{\mu}=0.
$}
\typeout{GRDISPLAY|266|\the\wd\grmathbox|\the\linewidth}
\ifdim\wd\grmathbox>\linewidth
\resizebox{\linewidth}{!}{\usebox{\grmathbox}}
\else\usebox{\grmathbox}\fi
\endgroup
\end{center}

Contract the displayed tensor: the first term gives
\(F_{\mu\alpha}F^{\mu\alpha}\), and the second gives \(4\times\frac14\)
times that same quantity. They cancel. Zero trace does not mean zero
energy, zero stress, or zero gravitational influence.

\hypertarget{chapter-12}{%
\chapter{12. Einstein\textquotesingle s equation: every symbol earns its
place}\label{chapter-12}}

We can now put the two protagonists into one equation:

\begin{center}
\begingroup
\sbox{\grmathbox}{$\displaystyle 
\boxed{
\underbrace{R_{\mu\nu}-\frac12R g_{\mu\nu}}_{G_{\mu\nu}}
+\Lambda g_{\mu\nu}
=\frac{8\pi G_N}{c^4}T_{\mu\nu}.
}
$}
\typeout{GRDISPLAY|267|\the\wd\grmathbox|\the\linewidth}
\ifdim\wd\grmathbox>\linewidth
\resizebox{\linewidth}{!}{\usebox{\grmathbox}}
\else\usebox{\grmathbox}\fi
\endgroup
\end{center}

It is a local differential equation for the spacetime metric, coupled to
the equations governing matter. It does not say that geometry is a
decorative picture pasted over Newtonian gravity. The metric determines
clocks, distances, causal cones, and free-fall trajectories, and this
equation governs that metric.

\hypertarget{reading-the-equation-without-mysticism}{%
\section{12.1 Reading the equation without
mysticism}\label{reading-the-equation-without-mysticism}}

\begingroup\small\setlength{\tabcolsep}{5pt}\renewcommand{\arraystretch}{1.13}

\begin{longtable}[]{@{}
  >{\raggedright\arraybackslash}p{(\columnwidth - 4\tabcolsep) * \real{0.2500}}
  >{\raggedright\arraybackslash}p{(\columnwidth - 4\tabcolsep) * \real{0.3600}}
  >{\raggedright\arraybackslash}p{(\columnwidth - 4\tabcolsep) * \real{0.3900}}@{}}
\toprule\noalign{}
\begin{minipage}[b]{\linewidth}\raggedright
Symbol
\end{minipage} & \begin{minipage}[b]{\linewidth}\raggedright
What it is
\end{minipage} & \begin{minipage}[b]{\linewidth}\raggedright
What job it does
\end{minipage} \\
\midrule\noalign{}
\endhead
\bottomrule\noalign{}
\endlastfoot
\(g_{\mu\nu}\) & Lorentzian metric & Defines intervals, causal
structure, contractions, and the Levi-Civita connection \\
\(R_{\mu\nu}\) & Ricci curvature & Records a contraction of tidal
curvature \\
\(R=g^{\mu\nu}R_{\mu\nu}\) & Scalar curvature & Supplies the curvature
trace \\
\(G_{\mu\nu}\) & Einstein tensor & Combines Ricci curvature and its
trace into an identically divergence-free tensor \\
\(\Lambda\) & Cosmological constant & Adds a permitted curvature scale
even without ordinary matter \\
\(T_{\mu\nu}\) & Nongravitational stress-energy & Supplies energy,
momentum, flux, and stress \\
\(G_N\) & Newton\textquotesingle s gravitational constant & Calibrates
the strength of gravity using the Newtonian limit \\
\(c\) & Speed of light & Relates temporal and spatial units and energy
to mass \\
\(\mu,\nu\) & Free tensor indices & Specify which pair of directions is
being compared \\
\end{longtable}

\endgroup

In local axes with all coordinates measured in length, curvature
components have units \(\mathrm{m^{-2}}\) and \(T_{\mu\nu}\) has units
\(\mathrm{J/m^3}\). Since

\begin{center}
\begingroup
\sbox{\grmathbox}{$\displaystyle 
\left[\frac{G_N}{c^4}\right]=\mathrm{m/J},
$}
\typeout{GRDISPLAY|268|\the\wd\grmathbox|\the\linewidth}
\ifdim\wd\grmathbox>\linewidth
\resizebox{\linewidth}{!}{\usebox{\grmathbox}}
\else\usebox{\grmathbox}\fi
\endgroup
\end{center}

the right-hand side has curvature units. In angular or differently
normalized coordinates, individual component units follow their
coordinate bases; the tensor equation remains dimensionally consistent.

The equation is nonlinear. The inverse metric appears in contractions;
the connection contains \(g^{-1}\partial g\); curvature contains
\(\partial\Gamma+\Gamma\Gamma\). The object being solved for helps
define the differential operator acting on itself.

An orchestra analogy is useful: the musicians affect the acoustics, and
the acoustics affect what the musicians hear. Its limitation is that
spacetime is not an external concert hall with a separate mechanical
material; the metric is the gravitational field itself.

\hypertarget{why-subtract-half-the-trace}{%
\section{12.2 Why subtract half the
trace?}\label{why-subtract-half-the-trace}}

The contracted Bianchi identity derived in Chapter 9 is

\begin{center}
\begingroup
\sbox{\grmathbox}{$\displaystyle 
\nabla^\mu R_{\mu\nu}=\frac12\nabla_\nu R.
$}
\typeout{GRDISPLAY|269|\the\wd\grmathbox|\the\linewidth}
\ifdim\wd\grmathbox>\linewidth
\resizebox{\linewidth}{!}{\usebox{\grmathbox}}
\else\usebox{\grmathbox}\fi
\endgroup
\end{center}

Suppose we seek a particularly simple symmetric curvature tensor of the
form

\begin{center}
\begingroup
\sbox{\grmathbox}{$\displaystyle 
A R_{\mu\nu}+B Rg_{\mu\nu}+C g_{\mu\nu},
$}
\typeout{GRDISPLAY|270|\the\wd\grmathbox|\the\linewidth}
\ifdim\wd\grmathbox>\linewidth
\resizebox{\linewidth}{!}{\usebox{\grmathbox}}
\else\usebox{\grmathbox}\fi
\endgroup
\end{center}

with constant coefficients. Metric compatibility gives \(\nabla g=0\),
so its divergence is

\begin{center}
\begingroup
\sbox{\grmathbox}{$\displaystyle 
\left(\frac A2+B\right)\nabla_\nu R.
$}
\typeout{GRDISPLAY|271|\the\wd\grmathbox|\the\linewidth}
\ifdim\wd\grmathbox>\linewidth
\resizebox{\linewidth}{!}{\usebox{\grmathbox}}
\else\usebox{\grmathbox}\fi
\endgroup
\end{center}

To make this vanish for arbitrary metrics, choose \(B=-A/2\). Rescale
the overall equation to set \(A=1\). The surviving combination is

\begin{center}
\begingroup
\sbox{\grmathbox}{$\displaystyle 
G_{\mu\nu}+\Lambda g_{\mu\nu}.
$}
\typeout{GRDISPLAY|272|\the\wd\grmathbox|\the\linewidth}
\ifdim\wd\grmathbox>\linewidth
\resizebox{\linewidth}{!}{\usebox{\grmathbox}}
\else\usebox{\grmathbox}\fi
\endgroup
\end{center}

The factor \(1/2\) is therefore a geometric compatibility requirement
within this ansatz. It is not a lucky fit to Mercury\textquotesingle s
orbit.

This does \textbf{not} establish that ``the equivalence principle
uniquely proves Einstein\textquotesingle s equation.'' We selected a
metric theory with a particular low-derivative curvature structure. More
general curvature actions, extra fields, independent connections, or
other assumptions can change the dynamics while retaining coordinate
covariance. The equivalence principle guides the local relation between
matter and geometry; it does not provide every dynamical postulate by
itself.

The action principle in Chapter 14 will supply a second route to
precisely the same trace subtraction. Seeing the coefficient arise from
both the Bianchi identity and the variation of volume is one of the
satisfying internal checks of the theory.

\hypertarget{trace-reversal-the-most-useful-algebraic-rearrangement}{%
\section{12.3 Trace reversal: the most useful algebraic
rearrangement}\label{trace-reversal-the-most-useful-algebraic-rearrangement}}

Define

\begin{center}
\begingroup
\sbox{\grmathbox}{$\displaystyle 
\kappa=\frac{8\pi G_N}{c^4},\qquad
T=g^{\mu\nu}T_{\mu\nu}.
$}
\typeout{GRDISPLAY|273|\the\wd\grmathbox|\the\linewidth}
\ifdim\wd\grmathbox>\linewidth
\resizebox{\linewidth}{!}{\usebox{\grmathbox}}
\else\usebox{\grmathbox}\fi
\endgroup
\end{center}

Contract Einstein\textquotesingle s equation with \(g^{\mu\nu}\). In
four dimensions \(g^{\mu\nu}g_{\mu\nu}=4\), so

\begin{center}
\begingroup
\sbox{\grmathbox}{$\displaystyle 
R-\frac12R(4)+4\Lambda=\kappa T.
$}
\typeout{GRDISPLAY|274|\the\wd\grmathbox|\the\linewidth}
\ifdim\wd\grmathbox>\linewidth
\resizebox{\linewidth}{!}{\usebox{\grmathbox}}
\else\usebox{\grmathbox}\fi
\endgroup
\end{center}

Thus

\begin{center}
\begingroup
\sbox{\grmathbox}{$\displaystyle 
\boxed{R=4\Lambda-\kappa T.}
$}
\typeout{GRDISPLAY|275|\the\wd\grmathbox|\the\linewidth}
\ifdim\wd\grmathbox>\linewidth
\resizebox{\linewidth}{!}{\usebox{\grmathbox}}
\else\usebox{\grmathbox}\fi
\endgroup
\end{center}

Substitute this back into the original equation:

\begin{center}
\begingroup
\sbox{\grmathbox}{$\displaystyle 
\boxed{
R_{\mu\nu}=\kappa\left(T_{\mu\nu}-\frac12Tg_{\mu\nu}\right)
+\Lambda g_{\mu\nu}.
}
$}
\typeout{GRDISPLAY|276|\the\wd\grmathbox|\the\linewidth}
\ifdim\wd\grmathbox>\linewidth
\resizebox{\linewidth}{!}{\usebox{\grmathbox}}
\else\usebox{\grmathbox}\fi
\endgroup
\end{center}

This is the \textbf{trace-reversed form}. It is mathematically
equivalent in four dimensions, but often much easier to use. For a
perfect fluid, \(T=-\epsilon+3p\). Therefore

\begin{center}
\begingroup
\sbox{\grmathbox}{$\displaystyle 
R_{\hat0\hat0}
=\kappa\left[\epsilon+\frac12(-\epsilon+3p)\right]-\Lambda
=\frac\kappa2(\epsilon+3p)-\Lambda,
$}
\typeout{GRDISPLAY|277|\the\wd\grmathbox|\the\linewidth}
\ifdim\wd\grmathbox>\linewidth
\resizebox{\linewidth}{!}{\usebox{\grmathbox}}
\else\usebox{\grmathbox}\fi
\endgroup
\end{center}

which explains the pressure result from Chapter 11.

If \(T_{\mu\nu}=0\) and \(\Lambda=0\), then \(R_{\mu\nu}=0\). This does
not force the entire Riemann tensor to vanish: Weyl curvature can
remain. Black-hole exteriors and gravitational waves make excellent use
of that permission.

If \(T=0\) but \(T_{\mu\nu}\ne0\), then \(R=4\Lambda\) while Ricci
curvature still responds to matter. An electromagnetic field is the
standard counterexample to the false statement ``zero scalar curvature
means empty, flat spacetime.''

The coefficient changes with dimension. In \(d>2\) dimensions,

\begin{center}
\begingroup
\sbox{\grmathbox}{$\displaystyle 
R_{\mu\nu}=\kappa\left(T_{\mu\nu}-\frac{T}{d-2}g_{\mu\nu}\right)
+\frac{2\Lambda}{d-2}g_{\mu\nu}.
$}
\typeout{GRDISPLAY|278|\the\wd\grmathbox|\the\linewidth}
\ifdim\wd\grmathbox>\linewidth
\resizebox{\linewidth}{!}{\usebox{\grmathbox}}
\else\usebox{\grmathbox}\fi
\endgroup
\end{center}

The four-dimensional \(1/2\) in trace reversal is dimension dependent.
The \(1/2\) in the definition of the Einstein tensor is not.

\hypertarget{deriving-the-newtonian-limit-including-the-factor-of-two}{%
\section{12.4 Deriving the Newtonian limit, including the factor of
two}\label{deriving-the-newtonian-limit-including-the-factor-of-two}}

We still owe an explanation of \(8\pi G_N/c^4\). Temporarily write an
unknown coupling \(\kappa\) on the right-hand side.

Assume a weak, nearly static field, slow matter and test particles,
negligible pressure compared with rest energy, and a region where
coordinates are approximately Minkowskian. Let \(\Phi\) denote the
ordinary Newtonian gravitational potential, with dimensions of velocity
squared.

First establish what \(\Phi\) has to do with the metric. For a slow
particle, the spatial geodesic equation is dominated by its two temporal
velocities:

\begin{center}
\begingroup
\sbox{\grmathbox}{$\displaystyle 
\frac{d^2x^i}{dt^2}\simeq-c^2\Gamma^i{}_{00}.
$}
\typeout{GRDISPLAY|279|\the\wd\grmathbox|\the\linewidth}
\ifdim\wd\grmathbox>\linewidth
\resizebox{\linewidth}{!}{\usebox{\grmathbox}}
\else\usebox{\grmathbox}\fi
\endgroup
\end{center}

This follows from the proper-time equation using \(dt/d\tau\simeq1\) at
leading order. Terms involving spatial velocities and the associated
nonaffine corrections in \(t\) are higher order in this limit.

With \(g_{0i}\) negligible and time derivatives negligible,

\begin{center}
\begingroup
\sbox{\grmathbox}{$\displaystyle 
\Gamma^i{}_{00}\simeq-\frac12\delta^{ij}\partial_jg_{00}.
$}
\typeout{GRDISPLAY|280|\the\wd\grmathbox|\the\linewidth}
\ifdim\wd\grmathbox>\linewidth
\resizebox{\linewidth}{!}{\usebox{\grmathbox}}
\else\usebox{\grmathbox}\fi
\endgroup
\end{center}

To reproduce Newton\textquotesingle s \(d^2x^i/dt^2=-\partial_i\Phi\),
we require

\begin{center}
\begingroup
\sbox{\grmathbox}{$\displaystyle 
\boxed{g_{00}\simeq-\left(1+\frac{2\Phi}{c^2}\right).}
$}
\typeout{GRDISPLAY|281|\the\wd\grmathbox|\the\linewidth}
\ifdim\wd\grmathbox>\linewidth
\resizebox{\linewidth}{!}{\usebox{\grmathbox}}
\else\usebox{\grmathbox}\fi
\endgroup
\end{center}

This step is kinematical: it identifies the Newtonian potential through
the behavior of slow free fall. It has not yet used
Einstein\textquotesingle s field equation.

Next calculate \(R_{00}\). With our curvature convention,

\begin{center}
\begingroup
\sbox{\grmathbox}{$\displaystyle 
R_{00}
=\partial_\rho\Gamma^\rho{}_{00}
-\partial_0\Gamma^\rho{}_{\rho0}
+\Gamma^\rho{}_{\rho\lambda}\Gamma^\lambda{}_{00}
-\Gamma^\rho{}_{0\lambda}\Gamma^\lambda{}_{\rho0}.
$}
\typeout{GRDISPLAY|282|\the\wd\grmathbox|\the\linewidth}
\ifdim\wd\grmathbox>\linewidth
\resizebox{\linewidth}{!}{\usebox{\grmathbox}}
\else\usebox{\grmathbox}\fi
\endgroup
\end{center}

Staticity removes the time-derivative term; weakness lets us discard
products of first-order connections. What remains is

\begin{center}
\begingroup
\sbox{\grmathbox}{$\displaystyle 
R_{00}\simeq\partial_i\Gamma^i{}_{00}
=\frac{\symbfit\nabla^2\Phi}{c^2}.
$}
\typeout{GRDISPLAY|283|\the\wd\grmathbox|\the\linewidth}
\ifdim\wd\grmathbox>\linewidth
\resizebox{\linewidth}{!}{\usebox{\grmathbox}}
\else\usebox{\grmathbox}\fi
\endgroup
\end{center}

For slow, pressureless matter,

\begin{center}
\begingroup
\sbox{\grmathbox}{$\displaystyle 
T_{00}\simeq\rho c^2,
\qquad
T\simeq-\rho c^2,
\qquad
 g_{00}\simeq-1.
$}
\typeout{GRDISPLAY|284|\the\wd\grmathbox|\the\linewidth}
\ifdim\wd\grmathbox>\linewidth
\resizebox{\linewidth}{!}{\usebox{\grmathbox}}
\else\usebox{\grmathbox}\fi
\endgroup
\end{center}

The trace-reversed source is consequently

\begin{center}
\begingroup
\sbox{\grmathbox}{$\displaystyle 
T_{00}-\frac12Tg_{00}
\simeq\rho c^2-\frac12\rho c^2
=\frac12\rho c^2.
$}
\typeout{GRDISPLAY|285|\the\wd\grmathbox|\the\linewidth}
\ifdim\wd\grmathbox>\linewidth
\resizebox{\linewidth}{!}{\usebox{\grmathbox}}
\else\usebox{\grmathbox}\fi
\endgroup
\end{center}

Setting \(\Lambda=0\) for the local matching calculation,

\begin{center}
\begingroup
\sbox{\grmathbox}{$\displaystyle 
\frac{\symbfit\nabla^2\Phi}{c^2}
=\frac\kappa2\rho c^2.
$}
\typeout{GRDISPLAY|286|\the\wd\grmathbox|\the\linewidth}
\ifdim\wd\grmathbox>\linewidth
\resizebox{\linewidth}{!}{\usebox{\grmathbox}}
\else\usebox{\grmathbox}\fi
\endgroup
\end{center}

Newtonian gravity requires Poisson\textquotesingle s equation

\begin{center}
\begingroup
\sbox{\grmathbox}{$\displaystyle 
\symbfit\nabla^2\Phi=4\pi G_N\rho.
$}
\typeout{GRDISPLAY|287|\the\wd\grmathbox|\the\linewidth}
\ifdim\wd\grmathbox>\linewidth
\resizebox{\linewidth}{!}{\usebox{\grmathbox}}
\else\usebox{\grmathbox}\fi
\endgroup
\end{center}

Comparing coefficients yields

\begin{center}
\begingroup
\sbox{\grmathbox}{$\displaystyle 
\boxed{\kappa=\frac{8\pi G_N}{c^4}.}
$}
\typeout{GRDISPLAY|288|\the\wd\grmathbox|\the\linewidth}
\ifdim\wd\grmathbox>\linewidth
\resizebox{\linewidth}{!}{\usebox{\grmathbox}}
\else\usebox{\grmathbox}\fi
\endgroup
\end{center}

The \(4\pi\) is the familiar three-dimensional Gauss-law normalization:
a spherical surface has area \(4\pi r^2\). The additional factor of two
came from trace reversal. The powers of \(c\) came from relating
temporal curvature to acceleration and energy density to mass density.
No number in the coefficient arrived by ceremonial decree.

Keeping the cosmological constant gives, in this same static weak-field
approximation,

\begin{center}
\begingroup
\sbox{\grmathbox}{$\displaystyle 
\symbfit\nabla^2\Phi=4\pi G_N\rho-\Lambda c^2.
$}
\typeout{GRDISPLAY|289|\the\wd\grmathbox|\the\linewidth}
\ifdim\wd\grmathbox>\linewidth
\resizebox{\linewidth}{!}{\usebox{\grmathbox}}
\else\usebox{\grmathbox}\fi
\endgroup
\end{center}

For example, a local vacuum solution includes
\(\Phi_\Lambda=-\Lambda c^2r^2/6\). Its acceleration is
\(-\symbfit\nabla\Phi_\Lambda=+\Lambda c^2\mathbf r/3\): positive
\(\Lambda\) produces an outward contribution in this approximation. This
is not a Newtonian description valid across an arbitrary cosmological
spacetime.

\hypertarget{a-calculation-trap-that-catches-experienced-students}{%
\section{12.5 A calculation trap that catches experienced
students}\label{a-calculation-trap-that-catches-experienced-students}}

For slow-particle acceleration, knowing \(g_{00}\) is enough at leading
order. For computing \(G_{00}\), it is not.

In the Newtonian regime of general relativity with negligible
anisotropic stress and suitable boundary conditions, a useful weak-field
chart has

\begin{center}
\begingroup
\sbox{\grmathbox}{$\displaystyle 
ds^2\simeq
-\left(1+\frac{2\Phi}{c^2}\right)c^2dt^2
+\left(1-\frac{2\Phi}{c^2}\right)\delta_{ij}dx^i dx^j.
$}
\typeout{GRDISPLAY|290|\the\wd\grmathbox|\the\linewidth}
\ifdim\wd\grmathbox>\linewidth
\resizebox{\linewidth}{!}{\usebox{\grmathbox}}
\else\usebox{\grmathbox}\fi
\endgroup
\end{center}

The spatial metric is perturbed too. Writing \(\varphi=\Phi/c^2\), a
first-order calculation gives

\begin{center}
\begingroup
\sbox{\grmathbox}{$\displaystyle 
R_{00}\simeq\symbfit\nabla^2\varphi,
\quad
R_{ij}\simeq\delta_{ij}\symbfit\nabla^2\varphi,
\quad
R\simeq2\symbfit\nabla^2\varphi,
$}
\typeout{GRDISPLAY|291|\the\wd\grmathbox|\the\linewidth}
\ifdim\wd\grmathbox>\linewidth
\resizebox{\linewidth}{!}{\usebox{\grmathbox}}
\else\usebox{\grmathbox}\fi
\endgroup
\end{center}

so

\begin{center}
\begingroup
\sbox{\grmathbox}{$\displaystyle 
G_{00}\simeq2\symbfit\nabla^2\varphi.
$}
\typeout{GRDISPLAY|292|\the\wd\grmathbox|\the\linewidth}
\ifdim\wd\grmathbox>\linewidth
\resizebox{\linewidth}{!}{\usebox{\grmathbox}}
\else\usebox{\grmathbox}\fi
\endgroup
\end{center}

If you perturb only \(g_{00}\) while artificially keeping the spatial
metric exactly Euclidean, your first-order \(G_{00}\) actually vanishes
for a static perturbation. That ansatz can identify Newtonian
accelerations but is not the complete weak-field solution sourced by
ordinary matter.

The larger lesson is methodological: \textbf{an approximation adequate
for one observable can be inadequate for another equation.} Trace
reversal made the coupling derivation especially clean because
\(R_{00}\) could be evaluated directly from the temporal perturbation at
this order.

\hypertarget{ten-components-with-structure}{%
\section{12.6 Ten components, with
structure}\label{ten-components-with-structure}}

A symmetric four-by-four tensor has ten independent components.
Einstein\textquotesingle s equation supplies ten component equations,
but the Bianchi identity imposes four differential relations among the
geometric expressions. Four coordinate functions are also gauge choices.
The resulting initial-value system contains constraint equations as well
as evolution equations; Chapter 20 will unpack it.

This is why ``solve ten independent wave equations for ten metric
components'' is the wrong computational picture. It is also why
inserting an arbitrary, nonconserved \(T_{\mu\nu}\) generally fails:
geometry\textquotesingle s identities require the matter equations and
sources to fit together.

Finally, \(T_{\mu\nu}\) here does not contain a universal local
gravitational stress tensor added by hand. Gravitational
self-interaction is already present in the nonlinear left-hand side. We
will return to the important distinction between that fact and the
existence of physically meaningful gravitational-wave energy or total
mass.

\hypertarget{chapter-13}{%
\chapter{13. Variational calculus: learning to ask a whole history a
question}\label{chapter-13}}

An equation of motion tells you how a system evolves. An action assigns
a number to an entire candidate history. The physical history makes that
number stationary under appropriate small changes.

For gravity, the ``history'' includes the geometry of spacetime itself.
We are about to vary the rulers and clocks, not merely move a particle
between them.

\hypertarget{from-an-ordinary-derivative-to-a-variation}{%
\section{13.1 From an ordinary derivative to a
variation}\label{from-an-ordinary-derivative-to-a-variation}}

For an ordinary function \(f(x)\), a stationary point satisfies
\(df/dx=0\). Change \(x\) by a small amount, and the first-order change
in \(f\) vanishes.

For mechanics, consider a candidate trajectory \(q(t)\) and its action

\begin{center}
\begingroup
\sbox{\grmathbox}{$\displaystyle 
S[q]=\int_{t_1}^{t_2}L(q,\dot q,t)\,dt.
$}
\typeout{GRDISPLAY|293|\the\wd\grmathbox|\the\linewidth}
\ifdim\wd\grmathbox>\linewidth
\resizebox{\linewidth}{!}{\usebox{\grmathbox}}
\else\usebox{\grmathbox}\fi
\endgroup
\end{center}

A functional such as \(S[q]\) eats a function and returns a number. To
differentiate it, introduce a one-parameter family of nearby histories,

\begin{center}
\begingroup
\sbox{\grmathbox}{$\displaystyle 
q_\lambda(t)=q(t)+\lambda\eta(t),
\qquad
\eta(t_1)=\eta(t_2)=0.
$}
\typeout{GRDISPLAY|294|\the\wd\grmathbox|\the\linewidth}
\ifdim\wd\grmathbox>\linewidth
\resizebox{\linewidth}{!}{\usebox{\grmathbox}}
\else\usebox{\grmathbox}\fi
\endgroup
\end{center}

The function \(\eta(t)\) describes a proposed wiggle, and \(\lambda\)
controls its size. Define

\begin{center}
\begingroup
\sbox{\grmathbox}{$\displaystyle 
\delta S=\left.\frac{d}{d\lambda}S[q_\lambda]\right|_{\lambda=0}.
$}
\typeout{GRDISPLAY|295|\the\wd\grmathbox|\the\linewidth}
\ifdim\wd\grmathbox>\linewidth
\resizebox{\linewidth}{!}{\usebox{\grmathbox}}
\else\usebox{\grmathbox}\fi
\endgroup
\end{center}

The endpoints are held fixed because we are comparing histories
connecting the same endpoint data. Other physical questions can require
other boundary conditions; they must be specified rather than guessed.

Differentiate inside the integral using the chain rule:

\begin{center}
\begingroup
\sbox{\grmathbox}{$\displaystyle 
\delta S=\int_{t_1}^{t_2}
\left(\frac{\partial L}{\partial q}\eta
+\frac{\partial L}{\partial\dot q}\dot\eta\right)dt.
$}
\typeout{GRDISPLAY|296|\the\wd\grmathbox|\the\linewidth}
\ifdim\wd\grmathbox>\linewidth
\resizebox{\linewidth}{!}{\usebox{\grmathbox}}
\else\usebox{\grmathbox}\fi
\endgroup
\end{center}

We want every term proportional to the arbitrary wiggle \(\eta\), not
its derivative. Integration by parts moves a time derivative away from
the wiggle:

\begin{center}
\begingroup
\sbox{\grmathbox}{$\displaystyle 
\delta S=
\left[\frac{\partial L}{\partial\dot q}\eta\right]_{t_1}^{t_2}
+\int_{t_1}^{t_2}
\left[
\frac{\partial L}{\partial q}
-\frac{d}{dt}\frac{\partial L}{\partial\dot q}
\right]\eta\,dt.
$}
\typeout{GRDISPLAY|297|\the\wd\grmathbox|\the\linewidth}
\ifdim\wd\grmathbox>\linewidth
\resizebox{\linewidth}{!}{\usebox{\grmathbox}}
\else\usebox{\grmathbox}\fi
\endgroup
\end{center}

The boundary term vanishes by the endpoint condition. Since \(\eta\) can
be chosen to have support in any small interior interval, the
coefficient must vanish pointwise:

\begin{center}
\begingroup
\sbox{\grmathbox}{$\displaystyle 
\boxed{
\frac{d}{dt}\frac{\partial L}{\partial\dot q}
-\frac{\partial L}{\partial q}=0.
}
$}
\typeout{GRDISPLAY|298|\the\wd\grmathbox|\the\linewidth}
\ifdim\wd\grmathbox>\linewidth
\resizebox{\linewidth}{!}{\usebox{\grmathbox}}
\else\usebox{\grmathbox}\fi
\endgroup
\end{center}

That last inference is the fundamental lemma of the calculus of
variations. If a continuous coefficient were positive somewhere, we
could choose a positive wiggle supported there and make the integral
nonzero. The same argument excludes a negative coefficient. ``For every
wiggle'' is an immensely powerful demand.

For \(L=\frac12m\dot q^2-V(q)\), the result is

\begin{center}
\begingroup
\sbox{\grmathbox}{$\displaystyle 
m\ddot q=-V'(q).
$}
\typeout{GRDISPLAY|299|\the\wd\grmathbox|\the\linewidth}
\ifdim\wd\grmathbox>\linewidth
\resizebox{\linewidth}{!}{\usebox{\grmathbox}}
\else\usebox{\grmathbox}\fi
\endgroup
\end{center}

Newton\textquotesingle s equation has emerged from a statement about a
complete history.

\hypertarget{stationary-does-not-mean-minimal}{%
\section{13.2 Stationary does not mean
minimal}\label{stationary-does-not-mean-minimal}}

The phrase ``least action'' is dangerously memorable. A stationary value
can be a minimum, a maximum, or a saddle. The variation above only tests
the first-order change.

For a harmonic oscillator, \(L=\frac12m\dot q^2-\frac12m\omega^2q^2\).
About a solution, the quadratic change in action is

\begin{center}
\begingroup
\sbox{\grmathbox}{$\displaystyle 
\Delta S=\frac{m\lambda^2}{2}
\int_{t_1}^{t_2}(\dot\eta^2-\omega^2\eta^2)\,dt.
$}
\typeout{GRDISPLAY|300|\the\wd\grmathbox|\the\linewidth}
\ifdim\wd\grmathbox>\linewidth
\resizebox{\linewidth}{!}{\usebox{\grmathbox}}
\else\usebox{\grmathbox}\fi
\endgroup
\end{center}

Let \(T=t_2-t_1\) and choose \(\eta=\sin[\pi(t-t_1)/T]\). Then

\begin{center}
\begingroup
\sbox{\grmathbox}{$\displaystyle 
\Delta S=\frac{m\lambda^2T}{4}
\left(\frac{\pi^2}{T^2}-\omega^2\right).
$}
\typeout{GRDISPLAY|301|\the\wd\grmathbox|\the\linewidth}
\ifdim\wd\grmathbox>\linewidth
\resizebox{\linewidth}{!}{\usebox{\grmathbox}}
\else\usebox{\grmathbox}\fi
\endgroup
\end{center}

For a sufficiently long interval, this variation decreases the action.
Higher-frequency wiggles can increase it. The physical path is then a
saddle, not a minimum.

Nor does the action principle require the universe to preview every
future and run an optimization contest. It is a compact mathematical
formulation equivalent to local differential equations under the stated
assumptions. It becomes especially useful when symmetries and coupled
fields make guessing the equations difficult.

\hypertarget{fields-a-degree-of-freedom-at-every-point}{%
\section{13.3 Fields: a degree of freedom at every
point}\label{fields-a-degree-of-freedom-at-every-point}}

For a scalar field in curved spacetime, use an SI Lagrangian energy
density \(\mathcal L\) and write

\begin{center}
\begingroup
\sbox{\grmathbox}{$\displaystyle 
S_\phi=\frac1c\int d^4x\,\sqrt{-g}\,
\mathcal L(\phi,\nabla_\mu\phi,g^{\mu\nu}).
$}
\typeout{GRDISPLAY|302|\the\wd\grmathbox|\the\linewidth}
\ifdim\wd\grmathbox>\linewidth
\resizebox{\linewidth}{!}{\usebox{\grmathbox}}
\else\usebox{\grmathbox}\fi
\endgroup
\end{center}

The \(1/c\) compensates for \(dx^0=c\,dt\), so the action has units of
energy times time. While varying \(\phi\), hold the metric fixed. Define

\begin{center}
\begingroup
\sbox{\grmathbox}{$\displaystyle 
P^\mu=\frac{\partial\mathcal L}{\partial(\nabla_\mu\phi)}.
$}
\typeout{GRDISPLAY|303|\the\wd\grmathbox|\the\linewidth}
\ifdim\wd\grmathbox>\linewidth
\resizebox{\linewidth}{!}{\usebox{\grmathbox}}
\else\usebox{\grmathbox}\fi
\endgroup
\end{center}

Since \(\phi\) is a scalar, \(\nabla_\mu\phi=\partial_\mu\phi\), and

\begin{center}
\begingroup
\sbox{\grmathbox}{$\displaystyle 
\delta S_\phi=\frac1c\int\sqrt{-g}
\left[\frac{\partial\mathcal L}{\partial\phi}\delta\phi
+P^\mu\nabla_\mu\delta\phi\right]d^4x.
$}
\typeout{GRDISPLAY|304|\the\wd\grmathbox|\the\linewidth}
\ifdim\wd\grmathbox>\linewidth
\resizebox{\linewidth}{!}{\usebox{\grmathbox}}
\else\usebox{\grmathbox}\fi
\endgroup
\end{center}

The covariant product rule gives

\begin{center}
\begingroup
\sbox{\grmathbox}{$\displaystyle 
P^\mu\nabla_\mu\delta\phi
=\nabla_\mu(P^\mu\delta\phi)-(\nabla_\mu P^\mu)\delta\phi.
$}
\typeout{GRDISPLAY|305|\the\wd\grmathbox|\the\linewidth}
\ifdim\wd\grmathbox>\linewidth
\resizebox{\linewidth}{!}{\usebox{\grmathbox}}
\else\usebox{\grmathbox}\fi
\endgroup
\end{center}

The divergence becomes a boundary integral. For a compactly supported
variation, or appropriate fixed boundary data, it vanishes. The field
equation is

\begin{center}
\begingroup
\sbox{\grmathbox}{$\displaystyle 
\boxed{
\frac{\partial\mathcal L}{\partial\phi}-\nabla_\mu P^\mu=0.
}
$}
\typeout{GRDISPLAY|306|\the\wd\grmathbox|\the\linewidth}
\ifdim\wd\grmathbox>\linewidth
\resizebox{\linewidth}{!}{\usebox{\grmathbox}}
\else\usebox{\grmathbox}\fi
\endgroup
\end{center}

In the natural-unit scalar example of Chapter 11,
\(P^\mu=-\nabla^\mu\phi\) and
\(\partial\mathcal L/\partial\phi=-V'(\phi)\), giving

\begin{center}
\begingroup
\sbox{\grmathbox}{$\displaystyle 
\Box\phi-V'(\phi)=0,
\qquad
\Box=\nabla_\mu\nabla^\mu.
$}
\typeout{GRDISPLAY|307|\the\wd\grmathbox|\the\linewidth}
\ifdim\wd\grmathbox>\linewidth
\resizebox{\linewidth}{!}{\usebox{\grmathbox}}
\else\usebox{\grmathbox}\fi
\endgroup
\end{center}

For \(V=\frac12m^2\phi^2\) in units \(c=\hbar=1\), this is the
Klein-Gordon equation. In a local inertial frame
\(\Box=-\partial_t^2+\symbfit\nabla^2\) at the point, with these units
and signature.

\textbf{Derivative trap.} The variation \(\delta\) commutes with
coordinate partial derivatives when comparing fields at the same
coordinates. It does not blindly commute with a covariant derivative
when the metric is varied. For a vector,

\begin{center}
\begingroup
\sbox{\grmathbox}{$\displaystyle 
\delta(\nabla_\mu V^\nu)
=\nabla_\mu\delta V^\nu
+\delta\Gamma^\nu{}_{\mu\rho}V^\rho.
$}
\typeout{GRDISPLAY|308|\the\wd\grmathbox|\the\linewidth}
\ifdim\wd\grmathbox>\linewidth
\resizebox{\linewidth}{!}{\usebox{\grmathbox}}
\else\usebox{\grmathbox}\fi
\endgroup
\end{center}

The connection is part of what changes. Forgetting this term would
remove the central difficulty of gravity by accidentally deleting it.

\hypertarget{varying-the-inverse-metric-prove-the-minus-sign-once}{%
\section{13.4 Varying the inverse metric: prove the minus sign
once}\label{varying-the-inverse-metric-prove-the-minus-sign-once}}

The covariant and inverse metrics are not independent fields. They obey

\begin{center}
\begingroup
\sbox{\grmathbox}{$\displaystyle 
g^{\mu\alpha}g_{\alpha\nu}=\delta^\mu{}_{\nu}.
$}
\typeout{GRDISPLAY|309|\the\wd\grmathbox|\the\linewidth}
\ifdim\wd\grmathbox>\linewidth
\resizebox{\linewidth}{!}{\usebox{\grmathbox}}
\else\usebox{\grmathbox}\fi
\endgroup
\end{center}

Vary this product:

\begin{center}
\begingroup
\sbox{\grmathbox}{$\displaystyle 
(\delta g^{\mu\alpha})g_{\alpha\nu}
+g^{\mu\alpha}\delta g_{\alpha\nu}=0.
$}
\typeout{GRDISPLAY|310|\the\wd\grmathbox|\the\linewidth}
\ifdim\wd\grmathbox>\linewidth
\resizebox{\linewidth}{!}{\usebox{\grmathbox}}
\else\usebox{\grmathbox}\fi
\endgroup
\end{center}

Multiply by \(g^{\nu\beta}\) to isolate the inverse-metric variation:

\begin{center}
\begingroup
\sbox{\grmathbox}{$\displaystyle 
\boxed{
\delta g^{\mu\beta}
=-g^{\mu\alpha}g^{\beta\nu}\delta g_{\alpha\nu}.
}
$}
\typeout{GRDISPLAY|311|\the\wd\grmathbox|\the\linewidth}
\ifdim\wd\grmathbox>\linewidth
\resizebox{\linewidth}{!}{\usebox{\grmathbox}}
\else\usebox{\grmathbox}\fi
\endgroup
\end{center}

Equivalently,

\begin{center}
\begingroup
\sbox{\grmathbox}{$\displaystyle 
\boxed{
\delta g_{\mu\nu}
=-g_{\mu\alpha}g_{\nu\beta}\delta g^{\alpha\beta}.
}
$}
\typeout{GRDISPLAY|312|\the\wd\grmathbox|\the\linewidth}
\ifdim\wd\grmathbox>\linewidth
\resizebox{\linewidth}{!}{\usebox{\grmathbox}}
\else\usebox{\grmathbox}\fi
\endgroup
\end{center}

The matrix identity underneath is
\(\delta(M^{-1})=-M^{-1}(\delta M)M^{-1}\). Enlarging a positive
eigenvalue of a matrix shrinks the corresponding inverse eigenvalue,
which makes the minus sign intuitive. The identity itself holds for
every invertible matrix, including Lorentzian metrics.

A common notation hazard is to define \(k_{\mu\nu}=\delta g_{\mu\nu}\)
and then raise its indices. The raised tensor
\(k^{\mu\nu}=g^{\mu\alpha}g^{\nu\beta}k_{\alpha\beta}\) satisfies

\begin{center}
\begingroup
\sbox{\grmathbox}{$\displaystyle 
k^{\mu\nu}=-\delta g^{\mu\nu}.
$}
\typeout{GRDISPLAY|313|\the\wd\grmathbox|\the\linewidth}
\ifdim\wd\grmathbox>\linewidth
\resizebox{\linewidth}{!}{\usebox{\grmathbox}}
\else\usebox{\grmathbox}\fi
\endgroup
\end{center}

``Raise the metric perturbation'' and ``vary the inverse metric'' differ
by a minus sign. This tiny distinction has ruined many otherwise
pleasant afternoons.

\hypertarget{why-the-volume-element-varies}{%
\section{13.5 Why the volume element
varies}\label{why-the-volume-element-varies}}

A spacetime region\textquotesingle s invariant volume is not simply its
coordinate volume. It is

\begin{center}
\begingroup
\sbox{\grmathbox}{$\displaystyle 
dV_4=\sqrt{-g}\,d^4x,
\qquad g=\det(g_{\mu\nu}).
$}
\typeout{GRDISPLAY|314|\the\wd\grmathbox|\the\linewidth}
\ifdim\wd\grmathbox>\linewidth
\resizebox{\linewidth}{!}{\usebox{\grmathbox}}
\else\usebox{\grmathbox}\fi
\endgroup
\end{center}

For signature \((-,+,+,+)\) the determinant is negative in any real
nonsingular coordinate chart, hence the real positive square root of
\(-g\).

To vary it, first recall Jacobi\textquotesingle s determinant identity:

\begin{center}
\begingroup
\sbox{\grmathbox}{$\displaystyle 
\delta\det M=(\det M)\operatorname{tr}(M^{-1}\delta M).
$}
\typeout{GRDISPLAY|315|\the\wd\grmathbox|\the\linewidth}
\ifdim\wd\grmathbox>\linewidth
\resizebox{\linewidth}{!}{\usebox{\grmathbox}}
\else\usebox{\grmathbox}\fi
\endgroup
\end{center}

Here is a short derivation. Factor

\begin{center}
\begingroup
\sbox{\grmathbox}{$\displaystyle 
\det(M+\lambda\delta M)=\det M\,
\det(I+\lambda M^{-1}\delta M).
$}
\typeout{GRDISPLAY|316|\the\wd\grmathbox|\the\linewidth}
\ifdim\wd\grmathbox>\linewidth
\resizebox{\linewidth}{!}{\usebox{\grmathbox}}
\else\usebox{\grmathbox}\fi
\endgroup
\end{center}

To first order, the determinant of \(I+\lambda A\) is
\(1+\lambda\operatorname{tr}A\). In the determinant expansion, a term
with exactly one perturbation contributes only when it occupies a
diagonal position; off-diagonal contributions need at least two
perturbations. Differentiating at \(\lambda=0\) gives the identity.

Apply it to the metric:

\begin{center}
\begingroup
\sbox{\grmathbox}{$\displaystyle 
\delta g=g\,g^{\mu\nu}\delta g_{\mu\nu}.
$}
\typeout{GRDISPLAY|317|\the\wd\grmathbox|\the\linewidth}
\ifdim\wd\grmathbox>\linewidth
\resizebox{\linewidth}{!}{\usebox{\grmathbox}}
\else\usebox{\grmathbox}\fi
\endgroup
\end{center}

The square-root chain rule then gives

\begin{center}
\begingroup
\sbox{\grmathbox}{$\displaystyle 
\boxed{
\delta\sqrt{-g}
=\frac12\sqrt{-g}\,g^{\mu\nu}\delta g_{\mu\nu}
=-\frac12\sqrt{-g}\,g_{\mu\nu}\delta g^{\mu\nu}.
}
$}
\typeout{GRDISPLAY|318|\the\wd\grmathbox|\the\linewidth}
\ifdim\wd\grmathbox>\linewidth
\resizebox{\linewidth}{!}{\usebox{\grmathbox}}
\else\usebox{\grmathbox}\fi
\endgroup
\end{center}

The final minus sign comes from inverse-metric variation. It does not
mean the Lorentzian volume element is imaginary or negative.

Geometrically, the trace measures the infinitesimal fractional change of
a volume. A trace-free metric deformation can change a little
region\textquotesingle s shape without changing its volume to first
order. In gravity, both curvature and the volume used to integrate that
curvature change under variation. Ignoring the second effect would lose
the \(-\frac12Rg_{\mu\nu}\) term in Einstein\textquotesingle s equation.

\hypertarget{the-stress-tensor-as-the-response-to-changing-geometry}{%
\section{13.6 The stress tensor as the response to changing
geometry}\label{the-stress-tensor-as-the-response-to-changing-geometry}}

A precise definition of the matter stress tensor, with our SI action
convention, is

\begin{center}
\begingroup
\sbox{\grmathbox}{$\displaystyle 
\boxed{
\delta S_m=-\frac1{2c}\int d^4x\,\sqrt{-g}\,
T_{\mu\nu}\,\delta g^{\mu\nu},
}
$}
\typeout{GRDISPLAY|319|\the\wd\grmathbox|\the\linewidth}
\ifdim\wd\grmathbox>\linewidth
\resizebox{\linewidth}{!}{\usebox{\grmathbox}}
\else\usebox{\grmathbox}\fi
\endgroup
\end{center}

when the matter fields themselves are held fixed and appropriate
integrations by parts have been performed. Equivalently,

\begin{center}
\begingroup
\sbox{\grmathbox}{$\displaystyle 
T_{\mu\nu}=-\frac{2c}{\sqrt{-g}}
\frac{\delta S_m}{\delta g^{\mu\nu}}.
$}
\typeout{GRDISPLAY|320|\the\wd\grmathbox|\the\linewidth}
\ifdim\wd\grmathbox>\linewidth
\resizebox{\linewidth}{!}{\usebox{\grmathbox}}
\else\usebox{\grmathbox}\fi
\endgroup
\end{center}

A functional derivative is defined by its role in the integral giving
the first variation. It is the continuum analogue of the coefficients
\(\partial f/\partial x_i\) in
\(\delta f=\sum_i(\partial f/\partial x_i)\delta x_i\).

Metric variations are symmetric. In particular,

\begin{center}
\begingroup
\sbox{\grmathbox}{$\displaystyle 
\frac{\partial g^{\alpha\beta}}{\partial g^{\mu\nu}}
=\delta^\alpha{}_{(\mu}\delta^\beta{}_{\nu)}
=\frac12\left(
\delta^\alpha{}_{\mu}\delta^\beta{}_{\nu}
+\delta^\alpha{}_{\nu}\delta^\beta{}_{\mu}
\right).
$}
\typeout{GRDISPLAY|321|\the\wd\grmathbox|\the\linewidth}
\ifdim\wd\grmathbox>\linewidth
\resizebox{\linewidth}{!}{\usebox{\grmathbox}}
\else\usebox{\grmathbox}\fi
\endgroup
\end{center}

There are ten independent variations, and the resulting metric stress
tensor is symmetric. If one instead varies \(g_{\mu\nu}\), the
definition becomes

\begin{center}
\begingroup
\sbox{\grmathbox}{$\displaystyle 
\delta S_m=+\frac1{2c}\int\sqrt{-g}\,
T^{\mu\nu}\delta g_{\mu\nu}\,d^4x.
$}
\typeout{GRDISPLAY|322|\the\wd\grmathbox|\the\linewidth}
\ifdim\wd\grmathbox>\linewidth
\resizebox{\linewidth}{!}{\usebox{\grmathbox}}
\else\usebox{\grmathbox}\fi
\endgroup
\end{center}

Changing metric variables changes the sign convention in the variation
formula, not the physical stress tensor.

For a matter Lagrangian with no derivatives of the metric, the
determinant identity immediately yields

\begin{center}
\begingroup
\sbox{\grmathbox}{$\displaystyle 
\boxed{
T_{\mu\nu}=-2\frac{\partial\mathcal L_m}{\partial g^{\mu\nu}}
+g_{\mu\nu}\mathcal L_m.
}
$}
\typeout{GRDISPLAY|323|\the\wd\grmathbox|\the\linewidth}
\ifdim\wd\grmathbox>\linewidth
\resizebox{\linewidth}{!}{\usebox{\grmathbox}}
\else\usebox{\grmathbox}\fi
\endgroup
\end{center}

Apply this in natural units to the scalar Lagrangian. Its explicit
metric derivative is \(-\frac12\partial_\mu\phi\partial_\nu\phi\). The
first term in \(T_{\mu\nu}\) becomes
\(\partial_\mu\phi\partial_\nu\phi\), and the volume variation supplies
\(g_{\mu\nu}\mathcal L_\phi\), reproducing Chapter 11\textquotesingle s
tensor.

For the electromagnetic example, return to SI units and hold the
covariant potential \(A_\mu\) fixed. Then \(F_{\mu\nu}\) is fixed under
the metric variation, but its raised components are not. Write its
contraction as

\begin{center}
\begingroup
\sbox{\grmathbox}{$\displaystyle 
F_{\alpha\beta}F^{\alpha\beta}
=g^{\alpha\rho}g^{\beta\sigma}F_{\alpha\beta}F_{\rho\sigma}.
$}
\typeout{GRDISPLAY|324|\the\wd\grmathbox|\the\linewidth}
\ifdim\wd\grmathbox>\linewidth
\resizebox{\linewidth}{!}{\usebox{\grmathbox}}
\else\usebox{\grmathbox}\fi
\endgroup
\end{center}

The two inverse metrics each contribute a variation. Antisymmetry of
\(F\) and renaming dummy indices show that the two contributions are
equal:

\begin{center}
\begingroup
\sbox{\grmathbox}{$\displaystyle 
\delta(F_{\alpha\beta}F^{\alpha\beta})
=2F_{\mu\alpha}F_\nu{}^\alpha\delta g^{\mu\nu}.
$}
\typeout{GRDISPLAY|325|\the\wd\grmathbox|\the\linewidth}
\ifdim\wd\grmathbox>\linewidth
\resizebox{\linewidth}{!}{\usebox{\grmathbox}}
\else\usebox{\grmathbox}\fi
\endgroup
\end{center}

Thus the explicit metric derivative of \(\mathcal L_{\mathrm{EM}}\)
gives the \(F_{\mu\alpha}F_\nu{}^\alpha/\mu_0\) term in its stress
tensor, and the volume variation gives the \(-g_{\mu\nu}F^2/(4\mu_0)\)
term. Even when the field components held fixed do not change, the
metric changes how they are contracted into a physical energy density.

The physical interpretation is beautiful: \textbf{stress-energy measures
how the matter action responds when its spacetime measuring apparatus is
changed.} Spatial deformations reveal stress; temporal deformations
reveal energy; mixed deformations reveal momentum and energy flow. The
familiar idea of stress as a response to strain has become a spacetime
statement.

For matter actions containing curvature or metric derivatives, use the
full functional definition rather than the short partial-derivative
formula. Likewise, a fluid\textquotesingle s energy density depends on
proper volume and other constrained variables. Treating it as a
metric-independent number during a naive variation will generally
produce the wrong fluid stress tensor.

\hypertarget{chapter-14}{%
\chapter{14. The Einstein-Hilbert action, taken apart
completely}\label{chapter-14}}

The name is \textbf{Einstein-Hilbert}, after David Hilbert. We now have
enough tools to derive the field equation without hiding its central
steps behind ``after some algebra.''

Our dynamical variable will be the inverse metric \(g^{\mu\nu}\). The
connection is its Levi-Civita connection, not an independent field in
this derivation. Matter fields are collectively denoted \(\psi\).

\hypertarget{choosing-an-action-and-checking-what-is-assumed}{%
\section{14.1 Choosing an action and checking what is
assumed}\label{choosing-an-action-and-checking-what-is-assumed}}

Take

\begin{center}
\begingroup
\sbox{\grmathbox}{$\displaystyle 
\boxed{
S[g,\psi]=
\frac{c^3}{16\pi G_N}
\int_{\mathcal M}d^4x\,\sqrt{-g}\,(R-2\Lambda)
+S_m[g,\psi]+S_{\mathrm{boundary}}.
}
$}
\typeout{GRDISPLAY|326|\the\wd\grmathbox|\the\linewidth}
\ifdim\wd\grmathbox>\linewidth
\resizebox{\linewidth}{!}{\usebox{\grmathbox}}
\else\usebox{\grmathbox}\fi
\endgroup
\end{center}

Why this structure? An action should not depend on arbitrary coordinate
labels, so integrate a scalar using invariant volume. A constant scalar
gives a cosmological term. The simplest nontrivial scalar built from a
metric and its curvature at two-derivative order is \(R\).

This is a choice of theory with substantial motivation, not a proof that
other terms are forbidden. Scalars such as \(R^2\) and
\(R_{\mu\nu}R^{\mu\nu}\) are also coordinate invariant; Chapter 23
explains why a modern effective theory expects higher-order terms. Here
we derive the dynamics of the Einstein-Hilbert choice.

The coefficient is already calibrated by the Newtonian limit. With
\(x^0=ct\),

\begin{center}
\begingroup
\sbox{\grmathbox}{$\displaystyle 
\left[\frac{c^3}{G_N}\right]=\mathrm{kg/s},
\qquad
\left[\int d^4x\sqrt{-g}\,R\right]=\mathrm{m^2},
$}
\typeout{GRDISPLAY|327|\the\wd\grmathbox|\the\linewidth}
\ifdim\wd\grmathbox>\linewidth
\resizebox{\linewidth}{!}{\usebox{\grmathbox}}
\else\usebox{\grmathbox}\fi
\endgroup
\end{center}

so \(S\) has units \(\mathrm{kg\,m^2/s}=\mathrm{J\,s}\). Under an actual
coordinate change from \(x^0=ct\) to \(t\), the invariant action keeps
its \(c^3\) prefactor: the transformed metric has \(g_{tt}=c^2g_{00}\)
and its determinant contributes \(\sqrt{-g_{(t)}}=c\sqrt{-g_{(x^0)}}\).
Some presentations factor that \(c\) out of the determinant and write a
\(c^4\) prefactor with a reduced determinant convention. These are
consistent bookkeeping choices. Do not transplant their prefactor into
our \(x^0=ct\) formula.

The symbol \(S_{\mathrm{boundary}}\) is not decorative. For deriving the
local bulk equations, we may initially choose variations supported
strictly inside the region. For a finite-region variational problem that
fixes the boundary geometry, an additional boundary action will be
required.

\hypertarget{first-split-curvature-changes-volume-changes}{%
\section{14.2 First split: curvature changes, volume
changes}\label{first-split-curvature-changes-volume-changes}}

Write \(A=c^3/(16\pi G_N)\) and vary the gravitational bulk term:

\begin{center}
\begingroup
\sbox{\grmathbox}{$\displaystyle 
\delta S_g=A\int d^4x\left[
\sqrt{-g}\,\delta R+(R-2\Lambda)\delta\sqrt{-g}
\right].
$}
\typeout{GRDISPLAY|328|\the\wd\grmathbox|\the\linewidth}
\ifdim\wd\grmathbox>\linewidth
\resizebox{\linewidth}{!}{\usebox{\grmathbox}}
\else\usebox{\grmathbox}\fi
\endgroup
\end{center}

Here \(\Lambda\) is a fixed parameter, so \(\delta\Lambda=0\). Since
\(R=g^{\mu\nu}R_{\mu\nu}\), the product rule gives

\begin{center}
\begingroup
\sbox{\grmathbox}{$\displaystyle 
\delta R=R_{\mu\nu}\delta g^{\mu\nu}
+g^{\mu\nu}\delta R_{\mu\nu}.
$}
\typeout{GRDISPLAY|329|\the\wd\grmathbox|\the\linewidth}
\ifdim\wd\grmathbox>\linewidth
\resizebox{\linewidth}{!}{\usebox{\grmathbox}}
\else\usebox{\grmathbox}\fi
\endgroup
\end{center}

The first term comes from changing the inverse metric used to take the
trace. The second comes from changing curvature itself. Substituting the
volume identity,

\begin{center}
\begingroup
\sbox{\grmathbox}{$\displaystyle 
\delta S_g=A\int\sqrt{-g}\left[
\left(R_{\mu\nu}-\frac12Rg_{\mu\nu}+\Lambda g_{\mu\nu}\right)
\delta g^{\mu\nu}
+g^{\mu\nu}\delta R_{\mu\nu}
\right]d^4x.
$}
\typeout{GRDISPLAY|330|\the\wd\grmathbox|\the\linewidth}
\ifdim\wd\grmathbox>\linewidth
\resizebox{\linewidth}{!}{\usebox{\grmathbox}}
\else\usebox{\grmathbox}\fi
\endgroup
\end{center}

Most of Einstein\textquotesingle s equation is already visible. The
trace subtraction came from volume variation, and
\(+\Lambda g_{\mu\nu}\) came from multiplying \(-2\Lambda\) by the
\(-1/2\) in that same variation. All that remains is to understand the
last term honestly.

\hypertarget{how-a-changing-metric-changes-the-connection}{%
\section{14.3 How a changing metric changes the
connection}\label{how-a-changing-metric-changes-the-connection}}

Set

\begin{center}
\begingroup
\sbox{\grmathbox}{$\displaystyle 
k_{\mu\nu}=\delta g_{\mu\nu},
\qquad k=g^{\mu\nu}k_{\mu\nu}.
$}
\typeout{GRDISPLAY|331|\the\wd\grmathbox|\the\linewidth}
\ifdim\wd\grmathbox>\linewidth
\resizebox{\linewidth}{!}{\usebox{\grmathbox}}
\else\usebox{\grmathbox}\fi
\endgroup
\end{center}

Metric compatibility says \(\nabla_\lambda g_{\mu\nu}=0\). Varying it
gives

\begin{center}
\begingroup
\sbox{\grmathbox}{$\displaystyle 
\nabla_\lambda k_{\mu\nu}
=\delta\Gamma^\rho{}_{\lambda\mu}g_{\rho\nu}
+\delta\Gamma^\rho{}_{\lambda\nu}g_{\mu\rho}.
$}
\typeout{GRDISPLAY|332|\the\wd\grmathbox|\the\linewidth}
\ifdim\wd\grmathbox>\linewidth
\resizebox{\linewidth}{!}{\usebox{\grmathbox}}
\else\usebox{\grmathbox}\fi
\endgroup
\end{center}

We have used the fact that varying the covariant derivative also varies
its two connection terms. Now write the analogous equations with cyclic
permutations of the indices. Add the versions with derivative indices
\(\mu\) and \(\nu\), and subtract the version with derivative index
\(\sigma\). Symmetry of the lower two connection indices makes the
unwanted terms cancel. The result is

\begin{center}
\begingroup
\sbox{\grmathbox}{$\displaystyle 
\boxed{
\delta\Gamma^\rho{}_{\mu\nu}
=\frac12g^{\rho\sigma}
\left(
\nabla_\mu k_{\nu\sigma}
+\nabla_\nu k_{\mu\sigma}
-\nabla_\sigma k_{\mu\nu}
\right).
}
$}
\typeout{GRDISPLAY|333|\the\wd\grmathbox|\the\linewidth}
\ifdim\wd\grmathbox>\linewidth
\resizebox{\linewidth}{!}{\usebox{\grmathbox}}
\else\usebox{\grmathbox}\fi
\endgroup
\end{center}

It resembles the Christoffel formula, but ordinary derivatives have
become covariant derivatives of the metric variation.

Why is this a tensor, when a connection is not? Two connections
transform with the same inhomogeneous coordinate term, so their
difference transforms tensorially. \(\delta\Gamma\) is an infinitesimal
difference of connections. You are allowed to take its covariant
derivative as a genuine \((1,2)\) tensor.

\textbf{Useful warning:} this statement compares two connections using
the same coordinate identification of the underlying manifold. A
separate coordinate transformation is a different operation, even though
both can be written with small parameters.

\hypertarget{the-palatini-identity-curvature-variation-becomes-a-derivative}{%
\section{14.4 The Palatini identity: curvature variation becomes a
derivative}\label{the-palatini-identity-curvature-variation-becomes-a-derivative}}

Start with our Ricci convention:

\begin{center}
\begingroup
\sbox{\grmathbox}{$\displaystyle 
R_{\mu\nu}
=\partial_\rho\Gamma^\rho{}_{\nu\mu}
-\partial_\nu\Gamma^\rho{}_{\rho\mu}
+\Gamma^\rho{}_{\rho\lambda}\Gamma^\lambda{}_{\nu\mu}
-\Gamma^\rho{}_{\nu\lambda}\Gamma^\lambda{}_{\rho\mu}.
$}
\typeout{GRDISPLAY|334|\the\wd\grmathbox|\the\linewidth}
\ifdim\wd\grmathbox>\linewidth
\resizebox{\linewidth}{!}{\usebox{\grmathbox}}
\else\usebox{\grmathbox}\fi
\endgroup
\end{center}

Vary every connection. The first two terms give derivatives of
\(\delta\Gamma\); the product terms each give two terms containing one
\(\Gamma\) and one \(\delta\Gamma\). Those product terms are exactly the
connection corrections required to turn the derivatives into covariant
derivatives:

\begin{center}
\begingroup
\sbox{\grmathbox}{$\displaystyle 
\boxed{
\delta R_{\mu\nu}
=\nabla_\rho\delta\Gamma^\rho{}_{\nu\mu}
-\nabla_\nu\delta\Gamma^\rho{}_{\rho\mu}.
}
$}
\typeout{GRDISPLAY|335|\the\wd\grmathbox|\the\linewidth}
\ifdim\wd\grmathbox>\linewidth
\resizebox{\linewidth}{!}{\usebox{\grmathbox}}
\else\usebox{\grmathbox}\fi
\endgroup
\end{center}

This is the \textbf{Palatini identity}. A clean way to verify the
algebra is to choose normal coordinates for the unvaried metric at one
point. There \(\Gamma=0\), all the product-variation terms vanish at
that point, and the identity reduces to the obvious variation of the two
derivative terms. Both sides are tensors, so equality established that
way holds in every chart. Normal coordinates simplify a tensor
calculation; they do not make curvature vanish.

Contract with \(g^{\mu\nu}\). Because \(\nabla g=0\), the inverse metric
moves through the derivatives:

\begin{center}
\begingroup
\sbox{\grmathbox}{$\displaystyle 
g^{\mu\nu}\delta R_{\mu\nu}=\nabla_\rho V^\rho,
$}
\typeout{GRDISPLAY|336|\the\wd\grmathbox|\the\linewidth}
\ifdim\wd\grmathbox>\linewidth
\resizebox{\linewidth}{!}{\usebox{\grmathbox}}
\else\usebox{\grmathbox}\fi
\endgroup
\end{center}

where

\begin{center}
\begingroup
\sbox{\grmathbox}{$\displaystyle 
\boxed{
V^\rho=g^{\mu\nu}\delta\Gamma^\rho{}_{\mu\nu}
-g^{\rho\mu}\delta\Gamma^\nu{}_{\nu\mu}.
}
$}
\typeout{GRDISPLAY|337|\the\wd\grmathbox|\the\linewidth}
\ifdim\wd\grmathbox>\linewidth
\resizebox{\linewidth}{!}{\usebox{\grmathbox}}
\else\usebox{\grmathbox}\fi
\endgroup
\end{center}

Substitute the connection variation to make its content more explicit:

\begin{center}
\begingroup
\sbox{\grmathbox}{$\displaystyle 
V^\rho=\nabla_\mu k^{\mu\rho}-\nabla^\rho k.
$}
\typeout{GRDISPLAY|338|\the\wd\grmathbox|\the\linewidth}
\ifdim\wd\grmathbox>\linewidth
\resizebox{\linewidth}{!}{\usebox{\grmathbox}}
\else\usebox{\grmathbox}\fi
\endgroup
\end{center}

If \(q^{\mu\nu}=\delta g^{\mu\nu}\) and \(q=g_{\mu\nu}q^{\mu\nu}\), then
\(k^{\mu\nu}=-q^{\mu\nu}\) and \(k=-q\), so equivalently

\begin{center}
\begingroup
\sbox{\grmathbox}{$\displaystyle 
\boxed{V^\rho=\nabla^\rho q-\nabla_\mu q^{\mu\rho}.}
$}
\typeout{GRDISPLAY|339|\the\wd\grmathbox|\the\linewidth}
\ifdim\wd\grmathbox>\linewidth
\resizebox{\linewidth}{!}{\usebox{\grmathbox}}
\else\usebox{\grmathbox}\fi
\endgroup
\end{center}

This is the key structural result. The remaining variation of curvature
contributes a total divergence, rather than an additional bulk field
equation.

The Einstein-Hilbert integrand contains second derivatives of the
metric. Nevertheless its bulk Euler-Lagrange equations contain only
second derivatives, not generic fourth derivatives: this special
linear-curvature structure sends the dangerous variation terms to the
boundary. Curvature-squared actions do not generally share that
simplification.

\hypertarget{put-the-pieces-together}{%
\section{14.5 Put the pieces together}\label{put-the-pieces-together}}

We have established

\begin{center}
\begingroup
\sbox{\grmathbox}{$\displaystyle 
\delta S_g=A\int_{\mathcal M}\sqrt{-g}
(G_{\mu\nu}+\Lambda g_{\mu\nu})\delta g^{\mu\nu}\,d^4x
+A\int_{\mathcal M}\sqrt{-g}\,\nabla_\rho V^\rho\,d^4x.
$}
\typeout{GRDISPLAY|340|\the\wd\grmathbox|\the\linewidth}
\ifdim\wd\grmathbox>\linewidth
\resizebox{\linewidth}{!}{\usebox{\grmathbox}}
\else\usebox{\grmathbox}\fi
\endgroup
\end{center}

For a variation supported inside the region, the second integral
vanishes by the divergence theorem. Add the matter variation:

\begin{center}
\begingroup
\sbox{\grmathbox}{$\displaystyle 
\delta S_m=-\frac1{2c}\int_{\mathcal M}\sqrt{-g}\,
T_{\mu\nu}\delta g^{\mu\nu}\,d^4x.
$}
\typeout{GRDISPLAY|341|\the\wd\grmathbox|\the\linewidth}
\ifdim\wd\grmathbox>\linewidth
\resizebox{\linewidth}{!}{\usebox{\grmathbox}}
\else\usebox{\grmathbox}\fi
\endgroup
\end{center}

Therefore

\begin{center}
\begingroup
\sbox{\grmathbox}{$\displaystyle 
\delta S=
\int_{\mathcal M}\sqrt{-g}
\left[
\frac{c^3}{16\pi G_N}(G_{\mu\nu}+\Lambda g_{\mu\nu})
-\frac1{2c}T_{\mu\nu}
\right]\delta g^{\mu\nu}\,d^4x.
$}
\typeout{GRDISPLAY|342|\the\wd\grmathbox|\the\linewidth}
\ifdim\wd\grmathbox>\linewidth
\resizebox{\linewidth}{!}{\usebox{\grmathbox}}
\else\usebox{\grmathbox}\fi
\endgroup
\end{center}

The ten symmetric metric variations can be chosen arbitrarily in the
interior, so their coefficient vanishes:

\begin{center}
\begingroup
\sbox{\grmathbox}{$\displaystyle 
\frac{c^3}{16\pi G_N}(G_{\mu\nu}+\Lambda g_{\mu\nu})
=\frac1{2c}T_{\mu\nu}.
$}
\typeout{GRDISPLAY|343|\the\wd\grmathbox|\the\linewidth}
\ifdim\wd\grmathbox>\linewidth
\resizebox{\linewidth}{!}{\usebox{\grmathbox}}
\else\usebox{\grmathbox}\fi
\endgroup
\end{center}

Multiply by \(16\pi G_N/c^3\):

\begin{center}
\begingroup
\sbox{\grmathbox}{$\displaystyle 
\boxed{
G_{\mu\nu}+\Lambda g_{\mu\nu}
=\frac{8\pi G_N}{c^4}T_{\mu\nu}.
}
$}
\typeout{GRDISPLAY|344|\the\wd\grmathbox|\the\linewidth}
\ifdim\wd\grmathbox>\linewidth
\resizebox{\linewidth}{!}{\usebox{\grmathbox}}
\else\usebox{\grmathbox}\fi
\endgroup
\end{center}

The geometry and matter equations arise by independent variations of the
metric and matter fields. When varying the metric to derive this
equation, matter fields are held fixed as field variables; when varying
matter, the metric is held fixed. A physical solution must satisfy both
resulting sets of equations.

We have now derived the field equation in words as well as symbols:
change the metric, account for the changed trace of curvature, account
for the changed spacetime volume, separate a divergence from the
curvature variation, measure the matter response by its stress tensor,
and require the total first-order response to vanish.

\hypertarget{why-the-boundary-term-vanishes-is-not-always-a-legal-sentence}{%
\section{14.6 Why ``the boundary term vanishes'' is not always a legal
sentence}\label{why-the-boundary-term-vanishes-is-not-always-a-legal-sentence}}

In mechanics with a first-derivative Lagrangian, fixing \(\delta q=0\)
at the endpoints kills the integration-by-parts boundary term. Here
\(V^\rho\) contains derivatives of \(\delta g\). Fixing the metric on a
boundary does not fix its normal derivative there.

The one-dimensional analogy is a function satisfying \(f(0)=0\) while
\(f'(0)\) is completely arbitrary. The function can meet the wall at any
slope. A fixed boundary metric likewise does not prevent its variation
from changing immediately away from the boundary.

There are two different legitimate problems:

\begin{enumerate}
\def\labelenumi{\arabic{enumi}.}
\tightlist
\item
  To derive \textbf{local bulk equations}, use compactly supported
  variations. No boundary term survives.
\item
  To define a \textbf{finite-region Dirichlet variational principle},
  fix the induced boundary geometry and add a term cancelling normal
  derivatives of its variation.
\end{enumerate}

The second problem is solved, for smooth non-null boundaries, by the
Gibbons-Hawking-York term.

\hypertarget{induced-geometry-and-the-gibbons-hawking-york-term}{%
\section{14.7 Induced geometry and the Gibbons-Hawking-York
term}\label{induced-geometry-and-the-gibbons-hawking-york-term}}

Let \(n^\mu\) be the outward-directed unit normal to a smooth boundary
segment. Set

\begin{center}
\begingroup
\sbox{\grmathbox}{$\displaystyle 
s=n_\mu n^\mu=\begin{cases}
+1,&\text{spacelike normal, timelike boundary},\\
-1,&\text{timelike normal, spacelike boundary}.
\end{cases}
$}
\typeout{GRDISPLAY|345|\the\wd\grmathbox|\the\linewidth}
\ifdim\wd\grmathbox>\linewidth
\resizebox{\linewidth}{!}{\usebox{\grmathbox}}
\else\usebox{\grmathbox}\fi
\endgroup
\end{center}

We use \(s\) for this sign so it cannot be confused with energy density
\(\epsilon\). The tensor projecting tangent to the boundary is

\begin{center}
\begingroup
\sbox{\grmathbox}{$\displaystyle 
h_{\mu\nu}=g_{\mu\nu}-s n_\mu n_\nu.
$}
\typeout{GRDISPLAY|346|\the\wd\grmathbox|\the\linewidth}
\ifdim\wd\grmathbox>\linewidth
\resizebox{\linewidth}{!}{\usebox{\grmathbox}}
\else\usebox{\grmathbox}\fi
\endgroup
\end{center}

In boundary coordinates \(y^a\) with tangent vectors
\(e_a^\mu=\partial x^\mu/\partial y^a\), the induced metric is

\begin{center}
\begingroup
\sbox{\grmathbox}{$\displaystyle 
h_{ab}=g_{\mu\nu}e_a^\mu e_b^\nu.
$}
\typeout{GRDISPLAY|347|\the\wd\grmathbox|\the\linewidth}
\ifdim\wd\grmathbox>\linewidth
\resizebox{\linewidth}{!}{\usebox{\grmathbox}}
\else\usebox{\grmathbox}\fi
\endgroup
\end{center}

It measures distances and, for a timelike boundary, times along the
boundary. Its determinant gives the surface measure
\(\sqrt{|h|}\,d^3y\).

The extrinsic curvature is

\begin{center}
\begingroup
\sbox{\grmathbox}{$\displaystyle 
K_{ab}=e_a^\mu e_b^\nu\nabla_\mu n_\nu,
\qquad K=h^{ab}K_{ab}.
$}
\typeout{GRDISPLAY|348|\the\wd\grmathbox|\the\linewidth}
\ifdim\wd\grmathbox>\linewidth
\resizebox{\linewidth}{!}{\usebox{\grmathbox}}
\else\usebox{\grmathbox}\fi
\endgroup
\end{center}

It describes how the normal changes as one moves along the boundary: how
the boundary bends within spacetime. This differs from its intrinsic
curvature. A cylindrical surface, for example, can have nonzero
extrinsic curvature even while its two-dimensional intrinsic geometry is
locally flat.

With this definition of \(K\) and outward normals, the appropriate term
is

\begin{center}
\begingroup
\sbox{\grmathbox}{$\displaystyle 
\boxed{
S_{\mathrm{GHY}}=\frac{c^3}{8\pi G_N}
\int_{\partial\mathcal M}s\sqrt{|h|}\,K\,d^3y.
}
$}
\typeout{GRDISPLAY|349|\the\wd\grmathbox|\the\linewidth}
\ifdim\wd\grmathbox>\linewidth
\resizebox{\linewidth}{!}{\usebox{\grmathbox}}
\else\usebox{\grmathbox}\fi
\endgroup
\end{center}

Sum over boundary segments when needed. The Lorentzian divergence
theorem in these conventions uses

\begin{center}
\begingroup
\sbox{\grmathbox}{$\displaystyle 
d\Sigma_\mu=s n_\mu\sqrt{|h|}\,d^3y,
$}
\typeout{GRDISPLAY|350|\the\wd\grmathbox|\the\linewidth}
\ifdim\wd\grmathbox>\linewidth
\resizebox{\linewidth}{!}{\usebox{\grmathbox}}
\else\usebox{\grmathbox}\fi
\endgroup
\end{center}

so the Einstein-Hilbert boundary variation is
\(A\int s\sqrt{|h|}n_\rho V^\rho d^3y\). Defining \(K\) with an overall
minus sign, or choosing different normal orientations, changes the
displayed boundary-action signs. Always compare definitions before
comparing formulas.

We can see the cancellation explicitly. Near a non-null boundary choose
Gaussian normal coordinates with outward normal coordinate \(r\):

\begin{center}
\begingroup
\sbox{\grmathbox}{$\displaystyle 
ds^2=s\,dr^2+h_{ab}(r,y)dy^a dy^b,
\qquad n=\partial_r,
\qquad K_{ab}=\frac12\partial_rh_{ab}.
$}
\typeout{GRDISPLAY|351|\the\wd\grmathbox|\the\linewidth}
\ifdim\wd\grmathbox>\linewidth
\resizebox{\linewidth}{!}{\usebox{\grmathbox}}
\else\usebox{\grmathbox}\fi
\endgroup
\end{center}

For variations preserving this coordinate form near the boundary, with
\(\delta h_{ab}=0\) on the boundary itself,

\begin{center}
\begingroup
\sbox{\grmathbox}{$\displaystyle 
\delta K=\frac12h^{ab}\partial_r\delta h_{ab},
\qquad
n_\rho V^\rho=-h^{ab}\partial_r\delta h_{ab}=-2\delta K.
$}
\typeout{GRDISPLAY|352|\the\wd\grmathbox|\the\linewidth}
\ifdim\wd\grmathbox>\linewidth
\resizebox{\linewidth}{!}{\usebox{\grmathbox}}
\else\usebox{\grmathbox}\fi
\endgroup
\end{center}

Thus

\begin{center}
\begingroup
\sbox{\grmathbox}{$\displaystyle 
\delta S_{g,\mathrm{boundary}}
=-2A\int s\sqrt{|h|}\,\delta K\,d^3y,
$}
\typeout{GRDISPLAY|353|\the\wd\grmathbox|\the\linewidth}
\ifdim\wd\grmathbox>\linewidth
\resizebox{\linewidth}{!}{\usebox{\grmathbox}}
\else\usebox{\grmathbox}\fi
\endgroup
\end{center}

whereas, because fixed \(h_{ab}\) also means \(\delta\sqrt{|h|}=0\)
there,

\begin{center}
\begingroup
\sbox{\grmathbox}{$\displaystyle 
\delta S_{\mathrm{GHY}}
=+2A\int s\sqrt{|h|}\,\delta K\,d^3y.
$}
\typeout{GRDISPLAY|354|\the\wd\grmathbox|\the\linewidth}
\ifdim\wd\grmathbox>\linewidth
\resizebox{\linewidth}{!}{\usebox{\grmathbox}}
\else\usebox{\grmathbox}\fi
\endgroup
\end{center}

They cancel. This local coordinate check isolates precisely the
normal-derivative terms that made the original variation problematic.

More generally, for a smooth non-null boundary, the remaining metric
boundary variation takes the form

\begin{center}
\begingroup
\sbox{\grmathbox}{$\displaystyle 
A\int_{\partial\mathcal M}s\sqrt{|h|}
(K_{ab}-Kh_{ab})\delta h^{ab}\,d^3y,
$}
\typeout{GRDISPLAY|355|\the\wd\grmathbox|\the\linewidth}
\ifdim\wd\grmathbox>\linewidth
\resizebox{\linewidth}{!}{\usebox{\grmathbox}}
\else\usebox{\grmathbox}\fi
\endgroup
\end{center}

up to boundary-of-boundary contributions and the specified boundary
conventions. It vanishes when the induced metric is fixed. The boundary
action makes the chosen variational problem well posed; it does not
change Einstein\textquotesingle s bulk field equation.

There are deliberate limits to this formula. Corners and joints
generally require extra terms. Null boundaries have no unit normal and a
degenerate induced metric, so the displayed GHY formula cannot simply be
applied to them. Their boundary and joint terms require additional
choices, including the normalization or parametrization of null
generators. A detailed primary treatment is
\href{https://arxiv.org/abs/1609.00207}{Lehner, Myers, Poisson, and
Sorkin, ``Gravitational action with null boundaries''}.

\hypertarget{palatini-identity-versus-palatini-variation}{%
\section{14.8 Palatini identity versus Palatini
variation}\label{palatini-identity-versus-palatini-variation}}

The \textbf{Palatini identity} used above is an identity for the
variation of curvature. We used it while the connection was determined
by the metric.

The \textbf{Palatini formulation} changes the variational problem: it
treats the metric and a connection \(\widetilde\Gamma\) as independent
variables. For the Einstein-Hilbert action, in dimension greater than
two, assuming a torsion-free independent connection and matter that does
not couple to it, the connection equation implies

\begin{center}
\begingroup
\sbox{\grmathbox}{$\displaystyle 
\widetilde\nabla_\lambda\bigl(\sqrt{-g}\,g^{\mu\nu}\bigr)=0,
$}
\typeout{GRDISPLAY|356|\the\wd\grmathbox|\the\linewidth}
\ifdim\wd\grmathbox>\linewidth
\resizebox{\linewidth}{!}{\usebox{\grmathbox}}
\else\usebox{\grmathbox}\fi
\endgroup
\end{center}

after the relevant traced equations are combined. Here
\(\widetilde\nabla\) acts on a tensor density of weight one, including
the connection term appropriate to that density weight. Under these
assumptions its solution is the Levi-Civita connection, and the metric
equation reproduces general relativity.

The equivalence depends on the assumptions. Changing the curvature
action or allowing matter to couple directly to the independent
connection can change the result. Allowing torsion also requires a
separate analysis. ``We used the Palatini identity'' does not mean ``we
made the connection an independent dynamical variable.''

\hypertarget{chapter-15}{%
\chapter{15. Symmetry, conservation, vacuum energy, and the limits of
slogans}\label{chapter-15}}

We have reached the equation, but understanding it requires knowing what
follows from its structure. Why is stress-energy conserved? What kind of
conservation is that? Does the universe have one total energy? And why
does an apparently harmless constant in a matter Lagrangian suddenly
matter to gravity?

These are places where the most memorable slogans are often one
assumption shorter than the truth.

\hypertarget{diffeomorphisms-moving-the-mathematical-description-coherently}{%
\section{15.1 Diffeomorphisms: moving the mathematical description
coherently}\label{diffeomorphisms-moving-the-mathematical-description-coherently}}

A diffeomorphism is a smooth invertible map of the manifold with a
smooth inverse. Infinitesimally it is generated by a vector field
\(\xi^\mu\). Its action on a tensor is described by a Lie derivative.

For the metric,

\begin{center}
\begingroup
\sbox{\grmathbox}{$\displaystyle 
\delta_\xi g_{\mu\nu}=\mathcal L_\xi g_{\mu\nu}
=\nabla_\mu\xi_\nu+\nabla_\nu\xi_\mu.
$}
\typeout{GRDISPLAY|357|\the\wd\grmathbox|\the\linewidth}
\ifdim\wd\grmathbox>\linewidth
\resizebox{\linewidth}{!}{\usebox{\grmathbox}}
\else\usebox{\grmathbox}\fi
\endgroup
\end{center}

For the inverse metric,

\begin{center}
\begingroup
\sbox{\grmathbox}{$\displaystyle 
\delta_\xi g^{\mu\nu}
=-\nabla^\mu\xi^\nu-\nabla^\nu\xi^\mu.
$}
\typeout{GRDISPLAY|358|\the\wd\grmathbox|\the\linewidth}
\ifdim\wd\grmathbox>\linewidth
\resizebox{\linewidth}{!}{\usebox{\grmathbox}}
\else\usebox{\grmathbox}\fi
\endgroup
\end{center}

For a scalar field,

\begin{center}
\begingroup
\sbox{\grmathbox}{$\displaystyle 
\delta_\xi\phi=\mathcal L_\xi\phi=\xi^\mu\partial_\mu\phi.
$}
\typeout{GRDISPLAY|359|\the\wd\grmathbox|\the\linewidth}
\ifdim\wd\grmathbox>\linewidth
\resizebox{\linewidth}{!}{\usebox{\grmathbox}}
\else\usebox{\grmathbox}\fi
\endgroup
\end{center}

The sign convention here uses the pullback generated by \(\xi\). One can
formulate passive coordinate changes with the opposite infinitesimal
sign; all terms must be changed consistently.

The essential idea is that the geometry and matter fields are
transformed together. Changing only the matter field while leaving
geometry fixed is generally a physical alteration, not the gauge
redundancy being discussed.

A relabeling analogy is helpful. If every street sign, map, address, and
navigation instruction is translated coherently, the journey is
unchanged. Translating only the street signs while leaving the
navigation instructions untouched creates a different experience. The
analogy\textquotesingle s limitation is that diffeomorphisms act on
smooth fields and their locations, not merely written names; in gravity
the distinction between gauge transformations and physical boundary
symmetries also depends on boundary conditions.

\hypertarget{noethers-theorem-has-a-local-gauge-version}{%
\section{15.2 Noether\textquotesingle s theorem has a local-gauge
version}\label{noethers-theorem-has-a-local-gauge-version}}

In mechanics, a continuous global symmetry produces a conserved
quantity. Diffeomorphism invariance has an arbitrary spacetime-dependent
vector field as its generator. This local freedom leads to differential
identities among equations of motion - an instance of
Noether\textquotesingle s second theorem.

First consider the matter action. Assume its matter fields satisfy their
equations of motion, so their first-order action variation vanishes
apart from boundary terms. Choose \(\xi^\mu\) with compact support to
remove the boundary terms. Diffeomorphism invariance and the
stress-tensor definition give

\begin{center}
\begingroup
\sbox{\grmathbox}{$\displaystyle 
0=\delta_\xi S_m
=-\frac1{2c}\int\sqrt{-g}\,T_{\mu\nu}\delta_\xi g^{\mu\nu}\,d^4x.
$}
\typeout{GRDISPLAY|360|\the\wd\grmathbox|\the\linewidth}
\ifdim\wd\grmathbox>\linewidth
\resizebox{\linewidth}{!}{\usebox{\grmathbox}}
\else\usebox{\grmathbox}\fi
\endgroup
\end{center}

Insert the inverse-metric Lie derivative. Symmetry of \(T_{\mu\nu}\)
makes its two terms equal:

\begin{center}
\begingroup
\sbox{\grmathbox}{$\displaystyle 
0=\frac1c\int\sqrt{-g}\,T^{\mu\nu}\nabla_\mu\xi_\nu\,d^4x.
$}
\typeout{GRDISPLAY|361|\the\wd\grmathbox|\the\linewidth}
\ifdim\wd\grmathbox>\linewidth
\resizebox{\linewidth}{!}{\usebox{\grmathbox}}
\else\usebox{\grmathbox}\fi
\endgroup
\end{center}

Integrate by parts:

\begin{center}
\begingroup
\sbox{\grmathbox}{$\displaystyle 
0=-\frac1c\int\sqrt{-g}\,(\nabla_\mu T^{\mu\nu})\xi_\nu\,d^4x.
$}
\typeout{GRDISPLAY|362|\the\wd\grmathbox|\the\linewidth}
\ifdim\wd\grmathbox>\linewidth
\resizebox{\linewidth}{!}{\usebox{\grmathbox}}
\else\usebox{\grmathbox}\fi
\endgroup
\end{center}

Since the interior vector field \(\xi_\nu\) was arbitrary,

\begin{center}
\begingroup
\sbox{\grmathbox}{$\displaystyle 
\boxed{\nabla_\mu T^{\mu\nu}=0.}
$}
\typeout{GRDISPLAY|363|\the\wd\grmathbox|\the\linewidth}
\ifdim\wd\grmathbox>\linewidth
\resizebox{\linewidth}{!}{\usebox{\grmathbox}}
\else\usebox{\grmathbox}\fi
\endgroup
\end{center}

This derivation used the matter equations, but not
Einstein\textquotesingle s equation. Matter on a prescribed curved
background can satisfy this conservation law if its action and equations
have the stated covariance and no unaccounted external exchanges.

The phrase \textbf{on shell} means that the relevant equations of motion
are satisfied. \textbf{Off shell} means we consider arbitrary field
configurations. Conservation of the matter stress tensor is generally an
on-shell statement. For the canonical scalar field in natural units, the
more informative off-shell identity is

\begin{center}
\begingroup
\sbox{\grmathbox}{$\displaystyle 
\nabla_\mu T^{\mu\nu}
=(\Box\phi-V'(\phi))\nabla^\nu\phi.
$}
\typeout{GRDISPLAY|364|\the\wd\grmathbox|\the\linewidth}
\ifdim\wd\grmathbox>\linewidth
\resizebox{\linewidth}{!}{\usebox{\grmathbox}}
\else\usebox{\grmathbox}\fi
\endgroup
\end{center}

You can check it directly by differentiating the scalar stress tensor,
using the product rule and commuting derivatives on a scalar. Its
divergence vanishes once the scalar equation holds. Conversely, stress
conservation alone need not imply every matter equation: if
\(\nabla^\nu\phi\) vanishes, that product can be zero without enforcing
the scalar equation.

Apply the same diffeomorphism argument to the gravitational action. Its
metric variation is proportional to \(G_{\mu\nu}+\Lambda g_{\mu\nu}\).
Because this is a pure geometric symmetry, the resulting identity

\begin{center}
\begingroup
\sbox{\grmathbox}{$\displaystyle 
\nabla^\mu(G_{\mu\nu}+\Lambda g_{\mu\nu})=0
$}
\typeout{GRDISPLAY|365|\the\wd\grmathbox|\the\linewidth}
\ifdim\wd\grmathbox>\linewidth
\resizebox{\linewidth}{!}{\usebox{\grmathbox}}
\else\usebox{\grmathbox}\fi
\endgroup
\end{center}

holds off shell for constant \(\Lambda\). It is the contracted Bianchi
identity in variational clothing. Einstein\textquotesingle s equation
joins an identically compatible geometric tensor to a matter tensor
conserved when matter evolves consistently.

This is not circular reasoning. One route derives a geometric identity
from curvature; another derives matter conservation from matter
dynamics; the action explains why the two structures can be coupled.

\hypertarget{local-conservation-is-not-automatically-a-global-energy-account}{%
\section{15.3 Local conservation is not automatically a global energy
account}\label{local-conservation-is-not-automatically-a-global-energy-account}}

Expand the covariant divergence:

\begin{center}
\begingroup
\sbox{\grmathbox}{$\displaystyle 
0=\nabla_\mu T^{\mu\nu}
=\frac1{\sqrt{-g}}\partial_\mu(\sqrt{-g}\,T^{\mu\nu})
+\Gamma^\nu{}_{\mu\lambda}T^{\mu\lambda}.
$}
\typeout{GRDISPLAY|366|\the\wd\grmathbox|\the\linewidth}
\ifdim\wd\grmathbox>\linewidth
\resizebox{\linewidth}{!}{\usebox{\grmathbox}}
\else\usebox{\grmathbox}\fi
\endgroup
\end{center}

For a vector current \(J^\mu\), there is no last free-index connection
term:

\begin{center}
\begingroup
\sbox{\grmathbox}{$\displaystyle 
\nabla_\mu J^\mu
=\frac1{\sqrt{-g}}\partial_\mu(\sqrt{-g}J^\mu).
$}
\typeout{GRDISPLAY|367|\the\wd\grmathbox|\the\linewidth}
\ifdim\wd\grmathbox>\linewidth
\resizebox{\linewidth}{!}{\usebox{\grmathbox}}
\else\usebox{\grmathbox}\fi
\endgroup
\end{center}

That difference matters. A divergence-free vector current can be
integrated using the divergence theorem to obtain a scalar charge
balance. A stress tensor still carries a free vector index, and vectors
at different points cannot generally be added without specifying how
their frames are related. Curvature makes that comparison path
dependent.

A useful mixed-index version is

\begin{center}
\begingroup
\sbox{\grmathbox}{$\displaystyle 
\boxed{
\frac1{\sqrt{-g}}\partial_\mu(\sqrt{-g}\,T^\mu{}_{\nu})
=\frac12T^{\alpha\beta}\partial_\nu g_{\alpha\beta}.
}
$}
\typeout{GRDISPLAY|368|\the\wd\grmathbox|\the\linewidth}
\ifdim\wd\grmathbox>\linewidth
\resizebox{\linewidth}{!}{\usebox{\grmathbox}}
\else\usebox{\grmathbox}\fi
\endgroup
\end{center}

It follows by expanding the covariant derivative and using symmetry of
\(T^{\alpha\beta}\). If the metric depends on the time coordinate, the
corresponding coordinate energy equation has a geometric term on its
right-hand side. It is not generally the ordinary flat-spacetime
continuity equation for one global energy.

In a sufficiently small freely falling laboratory, the connection
vanishes at a chosen event and the law reduces there to the familiar
local conservation equations. Across a finite curved region, comparing
different local laboratories is part of the problem.

Think of keeping accounts in currencies whose conversion depends on both
place and route. Each office can keep perfectly consistent local
accounts, yet adding their bare numerical balances is meaningless
without a conversion prescription. The limitation is that the geometry
is exact and deterministic; no financial uncertainty or exchange-market
mechanism is being asserted.

\hypertarget{a-killing-vector-supplies-the-missing-comparison-rule}{%
\section{15.4 A Killing vector supplies the missing comparison
rule}\label{a-killing-vector-supplies-the-missing-comparison-rule}}

A Killing vector satisfies

\begin{center}
\begingroup
\sbox{\grmathbox}{$\displaystyle 
\nabla_{(\mu}\xi_{\nu)}=0,
\qquad\text{equivalently}\qquad
\mathcal L_\xi g_{\mu\nu}=0.
$}
\typeout{GRDISPLAY|369|\the\wd\grmathbox|\the\linewidth}
\ifdim\wd\grmathbox>\linewidth
\resizebox{\linewidth}{!}{\usebox{\grmathbox}}
\else\usebox{\grmathbox}\fi
\endgroup
\end{center}

It generates a symmetry of the metric itself. For a symmetric conserved
stress tensor define

\begin{center}
\begingroup
\sbox{\grmathbox}{$\displaystyle 
J^\mu=-T^{\mu\nu}\xi_\nu.
$}
\typeout{GRDISPLAY|370|\the\wd\grmathbox|\the\linewidth}
\ifdim\wd\grmathbox>\linewidth
\resizebox{\linewidth}{!}{\usebox{\grmathbox}}
\else\usebox{\grmathbox}\fi
\endgroup
\end{center}

Its divergence is

\begin{center}
\begingroup
\sbox{\grmathbox}{$\displaystyle 
\nabla_\mu J^\mu
=-(\nabla_\mu T^{\mu\nu})\xi_\nu
-T^{\mu\nu}\nabla_\mu\xi_\nu.
$}
\typeout{GRDISPLAY|371|\the\wd\grmathbox|\the\linewidth}
\ifdim\wd\grmathbox>\linewidth
\resizebox{\linewidth}{!}{\usebox{\grmathbox}}
\else\usebox{\grmathbox}\fi
\endgroup
\end{center}

The first term vanishes by matter conservation. Since \(T^{\mu\nu}\) is
symmetric, only the symmetric part of \(\nabla_\mu\xi_\nu\) contributes
to the second term. The Killing equation kills it. Therefore

\begin{center}
\begingroup
\sbox{\grmathbox}{$\displaystyle 
\boxed{\nabla_\mu J^\mu=0.}
$}
\typeout{GRDISPLAY|372|\the\wd\grmathbox|\the\linewidth}
\ifdim\wd\grmathbox>\linewidth
\resizebox{\linewidth}{!}{\usebox{\grmathbox}}
\else\usebox{\grmathbox}\fi
\endgroup
\end{center}

A timelike Killing vector supplies an energy symmetry; a rotational
Killing vector supplies angular-momentum symmetry. The same construction
works for the corresponding charges, with their conventional
normalizations.

To make the energy case concrete, normalize a timelike \(\xi^\mu\) so
its components are dimensionless and it approaches a unit time direction
in an appropriate asymptotically flat region. On a spacelike
hypersurface \(\Sigma\) with future-directed unit normal \(n^\mu\),
define the matter Killing energy

\begin{center}
\begingroup
\sbox{\grmathbox}{$\displaystyle 
\boxed{
E_\xi=\int_\Sigma T_{\mu\nu}n^\mu\xi^\nu\,dV_\Sigma.
}
$}
\typeout{GRDISPLAY|373|\the\wd\grmathbox|\the\linewidth}
\ifdim\wd\grmathbox>\linewidth
\resizebox{\linewidth}{!}{\usebox{\grmathbox}}
\else\usebox{\grmathbox}\fi
\endgroup
\end{center}

In Minkowski spacetime, with \(n^\mu=\xi^\mu=(1,0,0,0)\), this is
\(\int\epsilon\,d^3x\). Apply the divergence theorem to a region between
two such hypersurfaces: the change of \(E_\xi\) equals minus the outward
flux through the intervening side boundary. If there is no side flux,
the charge is conserved.

The result requires the Killing symmetry, the field equations, and
suitable boundaries or falloff. It is a \textbf{matter} Killing charge;
it is not by itself a universal formula for the full gravitating
system\textquotesingle s total energy.

In a static spacetime, let \(N=\sqrt{-\xi^\mu\xi_\mu}\). Static
observers have unit time direction \(\xi^\mu/N\). A photon has conserved
Killing energy

\begin{center}
\begingroup
\sbox{\grmathbox}{$\displaystyle 
E_\xi=N E_{\mathrm{local}}.
$}
\typeout{GRDISPLAY|374|\the\wd\grmathbox|\the\linewidth}
\ifdim\wd\grmathbox>\linewidth
\resizebox{\linewidth}{!}{\usebox{\grmathbox}}
\else\usebox{\grmathbox}\fi
\endgroup
\end{center}

Two observers at different gravitational potentials can therefore assign
different local photon energies while agreeing on the conserved symmetry
charge. Gravitational redshift is consistent with energy conservation
when the correct conserved quantity and observer normalization are used.

\hypertarget{where-gravitational-energy-lives}{%
\section{15.5 Where gravitational energy
lives}\label{where-gravitational-energy-lives}}

Gravitational waves transfer energy to detectors, and black holes have
measurable mass. Yet general relativity provides no universal, unique,
exact local gravitational stress tensor playing the same role as
matter\textquotesingle s \(T_{\mu\nu}\) in every spacetime.

The equivalence principle offers a useful warning: expressions
constructed as a supposed gravitational energy density from connection
coefficients can be changed dramatically by changing coordinates,
including setting the connection to zero at one event. This warning is
not by itself a complete mathematical proof of every nonexistence
statement. Curvature tensors certainly exist. The full issue is finding
an exact local object with all the desired covariance, conservation,
normalization, and physical-energy properties without extra structure.

Several well-defined constructions answer more specific questions:

\begingroup\small\setlength{\tabcolsep}{5pt}\renewcommand{\arraystretch}{1.13}

\begin{longtable}[]{@{}
  >{\raggedright\arraybackslash}p{(\columnwidth - 4\tabcolsep) * \real{0.2500}}
  >{\raggedright\arraybackslash}p{(\columnwidth - 4\tabcolsep) * \real{0.3600}}
  >{\raggedright\arraybackslash}p{(\columnwidth - 4\tabcolsep) * \real{0.3900}}@{}}
\toprule\noalign{}
\begin{minipage}[b]{\linewidth}\raggedright
Construction
\end{minipage} & \begin{minipage}[b]{\linewidth}\raggedright
Extra structure or regime
\end{minipage} & \begin{minipage}[b]{\linewidth}\raggedright
What it describes
\end{minipage} \\
\midrule\noalign{}
\endhead
\bottomrule\noalign{}
\endlastfoot
ADM energy & Suitable asymptotic flatness at spatial infinity & Total
energy of an isolated gravitating system \\
Bondi energy & Suitable asymptotic structure at null infinity & Energy
remaining as radiation escapes to infinity \\
Quasilocal energy & A finite boundary and a chosen prescription & Energy
associated with a bounded region and boundary observers \\
Averaged gravitational-wave stress tensor & A controlled
short-wavelength approximation and averaging & Effective wave-energy
transport relative to a background \\
\end{longtable}

\endgroup

These are not mutually contradictory definitions of one quantity that
was obvious all along. They address different physical settings. Brown
and York\textquotesingle s quasilocal construction, for example, derives
a boundary stress tensor from variation of a gravitational action with
respect to the boundary metric; in the appropriate asymptotically flat
limit it recovers ADM quantities. See their original
\href{https://arxiv.org/abs/gr-qc/9209012}{``Quasilocal Energy and
Conserved Charges Derived from the Gravitational Action''}.

A generic expanding cosmology has no global timelike Killing vector.
Photons redshift, and matter still obeys local covariant conservation.
Asking where every lost photon joule ``went'' presupposes a globally
conserved energy account that the spacetime may not possess.

The comoving relation \(d(\epsilon V)=-p\,dV\) remains useful and exact
for a homogeneous perfect-fluid element under its assumptions. It does
not require inventing an external reservoir into which ``space stores
the missing energy.'' Nor is ``the total energy of every universe is
zero'' a general theorem of Einstein\textquotesingle s equation.

\hypertarget{the-cosmological-constant-as-vacuum-stress-energy}{%
\section{15.6 The cosmological constant as vacuum
stress-energy}\label{the-cosmological-constant-as-vacuum-stress-energy}}

Move the cosmological term to the right-hand side:

\begin{center}
\begingroup
\sbox{\grmathbox}{$\displaystyle 
G_{\mu\nu}=\kappa\left(T_{\mu\nu}+T^{(\Lambda)}_{\mu\nu}\right),
$}
\typeout{GRDISPLAY|375|\the\wd\grmathbox|\the\linewidth}
\ifdim\wd\grmathbox>\linewidth
\resizebox{\linewidth}{!}{\usebox{\grmathbox}}
\else\usebox{\grmathbox}\fi
\endgroup
\end{center}

with

\begin{center}
\begingroup
\sbox{\grmathbox}{$\displaystyle 
\boxed{
T^{(\Lambda)}_{\mu\nu}
=-\frac{\Lambda}{\kappa}g_{\mu\nu}
=-\epsilon_\Lambda g_{\mu\nu},
\qquad
\epsilon_\Lambda=\frac{\Lambda c^4}{8\pi G_N}.
}
$}
\typeout{GRDISPLAY|376|\the\wd\grmathbox|\the\linewidth}
\ifdim\wd\grmathbox>\linewidth
\resizebox{\linewidth}{!}{\usebox{\grmathbox}}
\else\usebox{\grmathbox}\fi
\endgroup
\end{center}

For any unit timelike observer \(n^\mu\),

\begin{center}
\begingroup
\sbox{\grmathbox}{$\displaystyle 
T^{(\Lambda)}_{\mu\nu}n^\mu n^\nu
=-\epsilon_\Lambda(-1)=\epsilon_\Lambda.
$}
\typeout{GRDISPLAY|377|\the\wd\grmathbox|\the\linewidth}
\ifdim\wd\grmathbox>\linewidth
\resizebox{\linewidth}{!}{\usebox{\grmathbox}}
\else\usebox{\grmathbox}\fi
\endgroup
\end{center}

Every such observer measures the same energy density and zero energy
flux. Spatial projection gives isotropic pressure

\begin{center}
\begingroup
\sbox{\grmathbox}{$\displaystyle 
\boxed{p_\Lambda=-\epsilon_\Lambda.}
$}
\typeout{GRDISPLAY|378|\the\wd\grmathbox|\the\linewidth}
\ifdim\wd\grmathbox>\linewidth
\resizebox{\linewidth}{!}{\usebox{\grmathbox}}
\else\usebox{\grmathbox}\fi
\endgroup
\end{center}

This is precisely the perfect-fluid form with \(\epsilon+p=0\): its
velocity-dependent part disappears. Vacuum stress of this form does not
select a preferred local rest frame.

Negative pressure is not merely a verbal synonym for repulsion. Its role
follows from the equations. In the timelike Ricci projection,

\begin{center}
\begingroup
\sbox{\grmathbox}{$\displaystyle 
\epsilon_\Lambda+3p_\Lambda=-2\epsilon_\Lambda.
$}
\typeout{GRDISPLAY|379|\the\wd\grmathbox|\the\linewidth}
\ifdim\wd\grmathbox>\linewidth
\resizebox{\linewidth}{!}{\usebox{\grmathbox}}
\else\usebox{\grmathbox}\fi
\endgroup
\end{center}

For positive \(\Lambda\), that contribution has the opposite sign from
positive-density, nonnegative-pressure matter. The resulting expansion
or focusing behavior also depends on the spacetime and congruence being
studied; Chapter 19 develops the cosmological case.

There is also a thermodynamic check. If \(\epsilon_\Lambda\) is constant
while a comoving volume changes, \(E=\epsilon_\Lambda V\) implies
\(dE=\epsilon_\Lambda dV\). Comparing with \(dE=-p\,dV\) gives
\(p=-\epsilon_\Lambda\). The energy in the volume grows with the volume;
the local work relation is satisfied by negative pressure.

\hypertarget{why-adding-a-constant-to-a-lagrangian-suddenly-matters}{%
\section{15.7 Why adding a constant to a Lagrangian suddenly
matters}\label{why-adding-a-constant-to-a-lagrangian-suddenly-matters}}

In nongravitational mechanics, adding a constant to \(L\) changes the
action by a fixed amount for a fixed time interval and leaves the
equations of motion unchanged.

In gravity, replace the matter Lagrangian energy density by

\begin{center}
\begingroup
\sbox{\grmathbox}{$\displaystyle 
\mathcal L_m\longrightarrow\mathcal L_m-C,
$}
\typeout{GRDISPLAY|380|\the\wd\grmathbox|\the\linewidth}
\ifdim\wd\grmathbox>\linewidth
\resizebox{\linewidth}{!}{\usebox{\grmathbox}}
\else\usebox{\grmathbox}\fi
\endgroup
\end{center}

where \(C\) is a constant energy density. The matter field equations
remain unchanged, but the metric varies the volume multiplying \(C\).
The stress tensor shifts by

\begin{center}
\begingroup
\sbox{\grmathbox}{$\displaystyle 
T_{\mu\nu}\longrightarrow T_{\mu\nu}-C g_{\mu\nu}.
$}
\typeout{GRDISPLAY|381|\the\wd\grmathbox|\the\linewidth}
\ifdim\wd\grmathbox>\linewidth
\resizebox{\linewidth}{!}{\usebox{\grmathbox}}
\else\usebox{\grmathbox}\fi
\endgroup
\end{center}

Consequently the effective cosmological constant becomes

\begin{center}
\begingroup
\sbox{\grmathbox}{$\displaystyle 
\boxed{\Lambda_{\mathrm{eff}}=\Lambda+\kappa C.}
$}
\typeout{GRDISPLAY|382|\the\wd\grmathbox|\the\linewidth}
\ifdim\wd\grmathbox>\linewidth
\resizebox{\linewidth}{!}{\usebox{\grmathbox}}
\else\usebox{\grmathbox}\fi
\endgroup
\end{center}

That is the conceptual doorway to the cosmological constant problem
discussed in Chapter 23: gravity responds to the metric dependence of a
vacuum-energy term even when nongravitational equations are insensitive
to a constant offset.

A constant \(\Lambda\) is consistent with the Bianchi identity because
\(\nabla g=0\). If someone simply replaces it by an arbitrary function
\(\Lambda(x)\) while keeping \(\nabla^\mu T_{\mu\nu}=0\) and leaving the
rest of the equation unchanged, taking the divergence forces

\begin{center}
\begingroup
\sbox{\grmathbox}{$\displaystyle 
\partial_\nu\Lambda=0.
$}
\typeout{GRDISPLAY|383|\the\wd\grmathbox|\the\linewidth}
\ifdim\wd\grmathbox>\linewidth
\resizebox{\linewidth}{!}{\usebox{\grmathbox}}
\else\usebox{\grmathbox}\fi
\endgroup
\end{center}

A varying dark-energy model therefore needs additional consistent
dynamics or energy exchange. Changing a parameter into a function is not
a free modification of a constrained field theory.

\hypertarget{minimal-coupling-and-what-it-does-not-settle}{%
\section{15.8 Minimal coupling, and what it does not
settle}\label{minimal-coupling-and-what-it-does-not-settle}}

A common way to place a nongravitational field theory on curved
spacetime is to begin with its special-relativistic action and make the
replacements

\begin{center}
\begingroup
\sbox{\grmathbox}{$\displaystyle 
\eta_{\mu\nu}\longrightarrow g_{\mu\nu},
\qquad
 d^4x\longrightarrow\sqrt{-g}\,d^4x,
\qquad
\partial\longrightarrow\nabla
$}
\typeout{GRDISPLAY|384|\the\wd\grmathbox|\the\linewidth}
\ifdim\wd\grmathbox>\linewidth
\resizebox{\linewidth}{!}{\usebox{\grmathbox}}
\else\usebox{\grmathbox}\fi
\endgroup
\end{center}

where the last replacement is appropriate to the kind of field involved.
Scalars have \(\nabla_\mu\phi=\partial_\mu\phi\). For the
electromagnetic field,

\begin{center}
\begingroup
\sbox{\grmathbox}{$\displaystyle 
F_{\mu\nu}=\nabla_\mu A_\nu-\nabla_\nu A_\mu
=\partial_\mu A_\nu-\partial_\nu A_\mu,
$}
\typeout{GRDISPLAY|385|\the\wd\grmathbox|\the\linewidth}
\ifdim\wd\grmathbox>\linewidth
\resizebox{\linewidth}{!}{\usebox{\grmathbox}}
\else\usebox{\grmathbox}\fi
\endgroup
\end{center}

because the two symmetric Levi-Civita connection terms cancel. Spinors
require a local frame and spin connection, introduced in Chapter 21.

This procedure is called minimal coupling. Applying it at the action
level is especially useful: it produces matter equations and a stress
tensor with mutually consistent variational origins.

It is not a uniqueness theorem. In natural units, a term such as

\begin{center}
\begingroup
\sbox{\grmathbox}{$\displaystyle 
-\frac12\zeta R\phi^2
$}
\typeout{GRDISPLAY|386|\the\wd\grmathbox|\the\linewidth}
\ifdim\wd\grmathbox>\linewidth
\resizebox{\linewidth}{!}{\usebox{\grmathbox}}
\else\usebox{\grmathbox}\fi
\endgroup
\end{center}

is also generally covariant, with dimensionless coupling \(\zeta\) in
four dimensions, and vanishes in exactly flat spacetime. A
flat-spacetime theory alone cannot tell you its coefficient. Moreover,
expressions related by commuting derivatives in flat spacetime need not
remain equivalent after covariantization, because covariant derivatives
on tensors do not generally commute.

The equation\textquotesingle s assumptions should now be visible: a
Lorentzian metric, a specified connection structure, chosen
gravitational and matter actions, appropriate boundary data, and a
consistent coupling between them. With those in place, covariance,
curvature identities, local energy-momentum balance, and dynamics
reinforce one another.

That coherence is the achievement. The theory becomes more impressive
when its assumptions are stated than when they are hidden inside a
slogan.

\hypertarget{chapter-16}{%
\chapter{16. Turning geometry into experiments: clocks, light, and
Mercury}\label{chapter-16}}

The Einstein equation has now acquired a meaning, a derivation, and a
respectable collection of indices. An experimentalist is entitled to
ask: what does it make an actual instrument do?

There is a repeatable answer. First solve, or approximate, the field
equation for a metric. Then specify the worldlines of the source,
detector, and light signals. Finally calculate quantities those
observers can measure: elapsed proper time, frequency, angle, or
separation. A coordinate component is an ingredient in that calculation;
it is not automatically an observable.

\hypertarget{two-gravitational-potentials-because-clocks-and-rulers-both-participate}{%
\section{16.1 Two gravitational potentials, because clocks and rulers
both
participate}\label{two-gravitational-potentials-because-clocks-and-rulers-both-participate}}

For a weak, approximately static field with negligible rotation, choose
Cartesian spatial coordinates and write

\begin{center}
\begingroup
\sbox{\grmathbox}{$\displaystyle 
ds^2=-\left(1+\frac{2\Phi}{c^2}\right)c^2dt^2
+\left(1-\frac{2\Psi}{c^2}\right)\delta_{ij}dx^i dx^j.
$}
\typeout{GRDISPLAY|387|\the\wd\grmathbox|\the\linewidth}
\ifdim\wd\grmathbox>\linewidth
\resizebox{\linewidth}{!}{\usebox{\grmathbox}}
\else\usebox{\grmathbox}\fi
\endgroup
\end{center}

Here \(x^0=ct\), and \(|\Phi|/c^2,|\Psi|/c^2\ll1\). Both potentials have
units of velocity squared. The first alters the relation between
coordinate time and clock time. The second alters the relation between
coordinate distances and ruler lengths. We have omitted vector
perturbations associated with mass currents, gravitational waves, and
higher-order terms.

For an isolated, slowly moving, weakly gravitating source in GR, with
negligible anisotropic stress and appropriate boundary conditions,

\begin{center}
\begingroup
\sbox{\grmathbox}{$\displaystyle 
\Phi=\Psi=-\frac{G_NM}{r}
$}
\typeout{GRDISPLAY|388|\the\wd\grmathbox|\the\linewidth}
\ifdim\wd\grmathbox>\linewidth
\resizebox{\linewidth}{!}{\usebox{\grmathbox}}
\else\usebox{\grmathbox}\fi
\endgroup
\end{center}

outside a spherical body. In more general matter systems the equality
needs justification; it is not part of the definition of a gravitational
potential.

Imagine a game whose level editor has separate sliders for clock rates
and spatial distances. Slow projectiles mainly reveal the clock slider.
Light reveals both. The analogy is a way of organizing the metric, not a
claim that nature contains two independently adjustable substances.
Einstein\textquotesingle s equation couples these coefficients.

For a slowly moving freely falling object, the leading spatial geodesic
equation is

\begin{center}
\begingroup
\sbox{\grmathbox}{$\displaystyle 
\frac{d^2x^i}{dt^2}\simeq-c^2\Gamma^i{}_{00}.
$}
\typeout{GRDISPLAY|389|\the\wd\grmathbox|\the\linewidth}
\ifdim\wd\grmathbox>\linewidth
\resizebox{\linewidth}{!}{\usebox{\grmathbox}}
\else\usebox{\grmathbox}\fi
\endgroup
\end{center}

Staticity removes the time derivatives in the connection, giving

\begin{center}
\begingroup
\sbox{\grmathbox}{$\displaystyle 
\Gamma^i{}_{00}\simeq-\frac12\delta^{ij}\partial_jg_{00}
=\frac{1}{c^2}\partial^i\Phi.
$}
\typeout{GRDISPLAY|390|\the\wd\grmathbox|\the\linewidth}
\ifdim\wd\grmathbox>\linewidth
\resizebox{\linewidth}{!}{\usebox{\grmathbox}}
\else\usebox{\grmathbox}\fi
\endgroup
\end{center}

Consequently \(d^2\mathbf x/dt^2=-\symbfit\nabla\Phi\): Newton emerges
because the metric\textquotesingle s clock coefficient has the
appropriate gradient. Terms involving spatial velocities enter at higher
order for this slow particle. They cannot be discarded for light.

Notice the deliberate phrase \textbf{coordinate acceleration}. The
falling object\textquotesingle s accelerometer reads zero. Its covariant
four-acceleration \(a^\mu=u^\nu\nabla_\nu u^\mu\) vanishes. A person
standing on the floor has approximately zero coordinate acceleration in
this chart but nonzero proper acceleration: the floor prevents a
geodesic. Your bathroom scale is measuring the interruption of free
fall.

\hypertarget{what-a-clock-actually-accumulates}{%
\section{16.2 What a clock actually
accumulates}\label{what-a-clock-actually-accumulates}}

Substitute \(dx^i=v^i dt\) and \(ds^2=-c^2d\tau^2\) into the weak-field
metric:

\begin{center}
\begingroup
\sbox{\grmathbox}{$\displaystyle 
\left(\frac{d\tau}{dt}\right)^2
=1+\frac{2\Phi}{c^2}
-\left(1-\frac{2\Psi}{c^2}\right)\frac{v^2}{c^2}.
$}
\typeout{GRDISPLAY|391|\the\wd\grmathbox|\the\linewidth}
\ifdim\wd\grmathbox>\linewidth
\resizebox{\linewidth}{!}{\usebox{\grmathbox}}
\else\usebox{\grmathbox}\fi
\endgroup
\end{center}

Keep terms first order in \(\Phi/c^2\) and \(v^2/c^2\), dropping their
product. Using \(\sqrt{1+q}\simeq1+q/2\) gives

\begin{center}
\begingroup
\sbox{\grmathbox}{$\displaystyle 
\boxed{\frac{d\tau}{dt}\simeq1+\frac{\Phi}{c^2}-\frac{v^2}{2c^2}.}
$}
\typeout{GRDISPLAY|392|\the\wd\grmathbox|\the\linewidth}
\ifdim\wd\grmathbox>\linewidth
\resizebox{\linewidth}{!}{\usebox{\grmathbox}}
\else\usebox{\grmathbox}\fi
\endgroup
\end{center}

The gravitational and kinematic clock corrections now occupy the same
line. A lower, more negative potential reduces accumulated proper time
relative to this coordinate time. Motion also reduces it at this order.
To compare clocks following different routes, integrate each expression
along its own route and specify how their readings are compared.

Near Earth\textquotesingle s surface, two stationary clocks separated
vertically by a small height \(h\) have \(\Delta\Phi\simeq gh\). For
\(h=1\,\mathrm m\),

\begin{center}
\begingroup
\sbox{\grmathbox}{$\displaystyle 
\frac{\Delta\dot\tau}{\dot\tau}\simeq\frac{9.81}{(2.998\times10^8)^2}
\simeq1.09\times10^{-16}.
$}
\typeout{GRDISPLAY|393|\the\wd\grmathbox|\the\linewidth}
\ifdim\wd\grmathbox>\linewidth
\resizebox{\linewidth}{!}{\usebox{\grmathbox}}
\else\usebox{\grmathbox}\fi
\endgroup
\end{center}

Over one day this is approximately \(9.4\) picoseconds. Small? Yes.
Conceptually optional? No. The universe does not round intermediate
calculations to human convenience.

This effect alone does not establish nonzero curvature. Accelerated
observers in flat spacetime can also have systematically different clock
rates. Curvature concerns the obstruction to removing gravitational
effects throughout an extended region, especially tidal effects. It is
stronger information than one pair of differently ticking clocks.

\hypertarget{gravitational-redshift-derived-with-an-honest-definition-of-energy}{%
\section{16.3 Gravitational redshift, derived with an honest definition
of
energy}\label{gravitational-redshift-derived-with-an-honest-definition-of-energy}}

In a stationary region, let \(K^\mu\) be the timelike Killing field
describing time-translation symmetry. Normalize it so that in an
asymptotically flat static chart \(K=\partial/\partial x^0\). Define its
positive norm factor

\begin{center}
\begingroup
\sbox{\grmathbox}{$\displaystyle 
N=\sqrt{-K_\mu K^\mu}.
$}
\typeout{GRDISPLAY|394|\the\wd\grmathbox|\the\linewidth}
\ifdim\wd\grmathbox>\linewidth
\resizebox{\linewidth}{!}{\usebox{\grmathbox}}
\else\usebox{\grmathbox}\fi
\endgroup
\end{center}

An observer remaining on an orbit of this symmetry has four-velocity

\begin{center}
\begingroup
\sbox{\grmathbox}{$\displaystyle 
u^\mu=\frac{cK^\mu}{N}.
$}
\typeout{GRDISPLAY|395|\the\wd\grmathbox|\the\linewidth}
\ifdim\wd\grmathbox>\linewidth
\resizebox{\linewidth}{!}{\usebox{\grmathbox}}
\else\usebox{\grmathbox}\fi
\endgroup
\end{center}

The normalization is not decoration: it ensures \(u_\mu u^\mu=-c^2\).

Let \(p^\mu\) be a photon\textquotesingle s four-momentum, transported
along its null geodesic. The Killing equation is
\(\nabla_{(\mu}K_{\nu)}=0\). Therefore

\begin{center}
\begingroup
\sbox{\grmathbox}{$\displaystyle 
p^\alpha\nabla_\alpha(p_\mu K^\mu)
=p^\alpha p^\mu\nabla_\alpha K_\mu=0.
$}
\typeout{GRDISPLAY|396|\the\wd\grmathbox|\the\linewidth}
\ifdim\wd\grmathbox>\linewidth
\resizebox{\linewidth}{!}{\usebox{\grmathbox}}
\else\usebox{\grmathbox}\fi
\endgroup
\end{center}

The term differentiating \(p_\mu\) vanished by the geodesic equation.
The remaining contraction vanishes because \(p^\alpha p^\mu\) is
symmetric while the relevant part of \(\nabla_\alpha K_\mu\) is
antisymmetric. Thus

\begin{center}
\begingroup
\sbox{\grmathbox}{$\displaystyle 
\mathcal E_K=-p_\mu K^\mu
$}
\typeout{GRDISPLAY|397|\the\wd\grmathbox|\the\linewidth}
\ifdim\wd\grmathbox>\linewidth
\resizebox{\linewidth}{!}{\usebox{\grmathbox}}
\else\usebox{\grmathbox}\fi
\endgroup
\end{center}

is conserved along the light ray. But the energy measured by a
particular observer is

\begin{center}
\begingroup
\sbox{\grmathbox}{$\displaystyle 
E_{\mathrm{local}}=-p_\mu u^\mu=\frac{c\mathcal E_K}{N}.
$}
\typeout{GRDISPLAY|398|\the\wd\grmathbox|\the\linewidth}
\ifdim\wd\grmathbox>\linewidth
\resizebox{\linewidth}{!}{\usebox{\grmathbox}}
\else\usebox{\grmathbox}\fi
\endgroup
\end{center}

One photon, one conserved symmetry charge, different local energies.
This is not a bookkeeping failure. Energy includes a specification of
the observer.

Since photon energy is proportional to measured frequency,

\begin{center}
\begingroup
\sbox{\grmathbox}{$\displaystyle 
\boxed{\frac{\nu_{\mathrm{rec}}}{\nu_{\mathrm{em}}}
=\frac{N_{\mathrm{em}}}{N_{\mathrm{rec}}}.}
$}
\typeout{GRDISPLAY|399|\the\wd\grmathbox|\the\linewidth}
\ifdim\wd\grmathbox>\linewidth
\resizebox{\linewidth}{!}{\usebox{\grmathbox}}
\else\usebox{\grmathbox}\fi
\endgroup
\end{center}

For the static weak-field metric, \(N\simeq1+\Phi/c^2\). Thus

\begin{center}
\begingroup
\sbox{\grmathbox}{$\displaystyle 
z\equiv\frac{\nu_{\mathrm{em}}}{\nu_{\mathrm{rec}}}-1
\simeq\frac{\Phi_{\mathrm{rec}}-\Phi_{\mathrm{em}}}{c^2}.
$}
\typeout{GRDISPLAY|400|\the\wd\grmathbox|\the\linewidth}
\ifdim\wd\grmathbox>\linewidth
\resizebox{\linewidth}{!}{\usebox{\grmathbox}}
\else\usebox{\grmathbox}\fi
\endgroup
\end{center}

Light received higher in the potential is redshifted. The familiar story
that a photon ``spends energy climbing'' gives the right sign, but the
Killing derivation explains precisely what stays conserved and what
changes. In a general time-dependent geometry, the required timelike
Killing symmetry may not exist at all.

\hypertarget{bending-light-finding-the-famous-factor-of-two}{%
\section{16.4 Bending light: finding the famous factor of
two}\label{bending-light-finding-the-famous-factor-of-two}}

For a null trajectory, \(ds^2=0\). Let \(d\ell^2=\delta_{ij}dx^i dx^j\)
denote Euclidean coordinate path length. Then

\begin{center}
\begingroup
\sbox{\grmathbox}{$\displaystyle 
c\,dt=n(\mathbf x)d\ell,
\qquad
n=\sqrt{\frac{1-2\Psi/c^2}{1+2\Phi/c^2}}
\simeq1-\frac{\Phi+\Psi}{c^2}.
$}
\typeout{GRDISPLAY|401|\the\wd\grmathbox|\the\linewidth}
\ifdim\wd\grmathbox>\linewidth
\resizebox{\linewidth}{!}{\usebox{\grmathbox}}
\else\usebox{\grmathbox}\fi
\endgroup
\end{center}

In this static chart, the light path makes the travel-time functional
\(\int n\,d\ell\) stationary. This is the same mathematics as ray optics
in an inhomogeneous refractive medium. No material ether has appeared:
\(n\) is a coordinate description of null geometry, and every local
freely falling observer still measures light speed \(c\).

Take an unperturbed ray traveling along \(z\), passing a mass at
transverse distance \(b\). To first order, the change in its transverse
direction is

\begin{center}
\begingroup
\sbox{\grmathbox}{$\displaystyle 
\Delta\symbfit\theta
\simeq\int_{-\infty}^{\infty}\symbfit\nabla_\perp n\,dz
=-\frac{1}{c^2}\int_{-\infty}^{\infty}
\symbfit\nabla_\perp(\Phi+\Psi)\,dz.
$}
\typeout{GRDISPLAY|402|\the\wd\grmathbox|\the\linewidth}
\ifdim\wd\grmathbox>\linewidth
\resizebox{\linewidth}{!}{\usebox{\grmathbox}}
\else\usebox{\grmathbox}\fi
\endgroup
\end{center}

Why can we integrate along a straight path if the path bends? Because
the bending is already first order in \(G_NM\). Correcting the path
inside this first-order integrand would produce a second-order
correction.

Write \(\Psi=\gamma\Phi\) for a constant comparison parameter. At
transverse position \(b>0\),

\begin{center}
\begingroup
\sbox{\grmathbox}{$\displaystyle 
\partial_b(\Phi+\Psi)
=(1+\gamma)\frac{G_NMb}{(b^2+z^2)^{3/2}}.
$}
\typeout{GRDISPLAY|403|\the\wd\grmathbox|\the\linewidth}
\ifdim\wd\grmathbox>\linewidth
\resizebox{\linewidth}{!}{\usebox{\grmathbox}}
\else\usebox{\grmathbox}\fi
\endgroup
\end{center}

The transverse direction change is negative, toward the mass. Its
magnitude is

\begin{center}
\begingroup
\sbox{\grmathbox}{$\displaystyle 
|\Delta\theta|=
\frac{(1+\gamma)G_NMb}{c^2}
\int_{-\infty}^{\infty}\frac{dz}{(b^2+z^2)^{3/2}}
=\frac{2(1+\gamma)G_NM}{bc^2}.
$}
\typeout{GRDISPLAY|404|\the\wd\grmathbox|\the\linewidth}
\ifdim\wd\grmathbox>\linewidth
\resizebox{\linewidth}{!}{\usebox{\grmathbox}}
\else\usebox{\grmathbox}\fi
\endgroup
\end{center}

The integral equals \(2/b^2\); differentiating \(z/[b^2\sqrt{b^2+z^2}]\)
verifies it. GR gives \(\gamma=1\), hence

\begin{center}
\begingroup
\sbox{\grmathbox}{$\displaystyle 
\boxed{|\Delta\theta|_{\mathrm{GR}}=\frac{4G_NM}{bc^2}.}
$}
\typeout{GRDISPLAY|405|\the\wd\grmathbox|\the\linewidth}
\ifdim\wd\grmathbox>\linewidth
\resizebox{\linewidth}{!}{\usebox{\grmathbox}}
\else\usebox{\grmathbox}\fi
\endgroup
\end{center}

At the solar limb this is about \(1.75\) arcseconds. Keeping only the
clock potential would produce half this value.

A careful gotcha: the split into ``half from time curvature, half from
space curvature'' belongs to this convenient weak-field coordinate
description. It is not an invariant decomposition of four-dimensional
curvature into two independently observable substances. The total
observable deflection is the robust result.

\hypertarget{shapiro-delay-the-other-consequence-of-the-same-optical-geometry}{%
\section{16.5 Shapiro delay: the other consequence of the same optical
geometry}\label{shapiro-delay-the-other-consequence-of-the-same-optical-geometry}}

The effective index also produces an additional travel time. Along a
nearly straight path,

\begin{center}
\begingroup
\sbox{\grmathbox}{$\displaystyle 
\Delta t=-\frac{1}{c^3}\int(\Phi+\Psi)\,d\ell.
$}
\typeout{GRDISPLAY|406|\the\wd\grmathbox|\the\linewidth}
\ifdim\wd\grmathbox>\linewidth
\resizebox{\linewidth}{!}{\usebox{\grmathbox}}
\else\usebox{\grmathbox}\fi
\endgroup
\end{center}

For a point mass in GR, a ray passing with impact parameter \(b\) and
endpoints at longitudinal coordinates \(-z_1\) and \(z_2\) has

\begin{center}
\begingroup
\sbox{\grmathbox}{$\displaystyle 
\Delta t\simeq\frac{2G_NM}{c^3}
\left[\operatorname{arsinh}\frac{z_1}{b}
+\operatorname{arsinh}\frac{z_2}{b}\right].
$}
\typeout{GRDISPLAY|407|\the\wd\grmathbox|\the\linewidth}
\ifdim\wd\grmathbox>\linewidth
\resizebox{\linewidth}{!}{\usebox{\grmathbox}}
\else\usebox{\grmathbox}\fi
\endgroup
\end{center}

For endpoints far from closest approach, \(z_1,z_2\gg b\), use
\(\operatorname{arsinh}q\simeq\ln(2q)\):

\begin{center}
\begingroup
\sbox{\grmathbox}{$\displaystyle 
\Delta t\simeq\frac{2G_NM}{c^3}
\ln\frac{4r_1r_2}{b^2}.
$}
\typeout{GRDISPLAY|408|\the\wd\grmathbox|\the\linewidth}
\ifdim\wd\grmathbox>\linewidth
\resizebox{\linewidth}{!}{\usebox{\grmathbox}}
\else\usebox{\grmathbox}\fi
\endgroup
\end{center}

This is a leading one-way coordinate delay relative to the corresponding
flat path; \(r_1,r_2\) are approximately the endpoint distances from the
mass. An actual radar experiment models the return trip and converts the
result into the tracking station\textquotesingle s proper time.
Gravitational lensing more generally also involves different geometric
path lengths. The observational model must keep both contributions.

The solar coefficient \(2G_NM_\odot/c^3\) is about \(9.85\)
microseconds. A logarithm of order ten turns a geometrically small
correction into a readily meaningful timing signal.

\hypertarget{mercury-a-resonant-correction-that-slowly-rotates-an-ellipse}{%
\section{16.6 Mercury: a resonant correction that slowly rotates an
ellipse}\label{mercury-a-resonant-correction-that-slowly-rotates-an-ellipse}}

Chapter 17 derives the Schwarzschild metric. For a massive test particle
in that geometry, define the conserved specific angular momentum
\(\ell=r^2d\phi/d\tau\) and let \(u=1/r\). The exact equatorial orbit
equation is

\begin{center}
\begingroup
\sbox{\grmathbox}{$\displaystyle 
\boxed{\frac{d^2u}{d\phi^2}+u
=\frac{G_NM}{\ell^2}+\frac{3G_NM}{c^2}u^2.}
$}
\typeout{GRDISPLAY|409|\the\wd\grmathbox|\the\linewidth}
\ifdim\wd\grmathbox>\linewidth
\resizebox{\linewidth}{!}{\usebox{\grmathbox}}
\else\usebox{\grmathbox}\fi
\endgroup
\end{center}

Here \(u\) is reciprocal radius, not four-velocity. The first term gives
the Newtonian ellipse. The second is the relativistic correction. The
same letter is common in both contexts, so the definitions matter more
than the typography.

Set \(p_{\mathrm{orb}}=\ell^2/(G_NM)\) and \(m=G_NM/c^2\). The
unperturbed orbit is

\begin{center}
\begingroup
\sbox{\grmathbox}{$\displaystyle 
u_0=\frac{1}{p_{\mathrm{orb}}}(1+e\cos\phi).
$}
\typeout{GRDISPLAY|410|\the\wd\grmathbox|\the\linewidth}
\ifdim\wd\grmathbox>\linewidth
\resizebox{\linewidth}{!}{\usebox{\grmathbox}}
\else\usebox{\grmathbox}\fi
\endgroup
\end{center}

Insert \(u_0\) into the small correction \(3mu^2\). Its term
proportional to \(\cos\phi\) is \(6me\cos\phi/p_{\mathrm{orb}}^2\). This
drives the same angular frequency as the homogeneous operator
\(d^2/d\phi^2+1\). The resulting resonant particular solution is

\begin{center}
\begingroup
\sbox{\grmathbox}{$\displaystyle 
\delta u_{\mathrm{res}}=\frac{3me}{p_{\mathrm{orb}}^2}\phi\sin\phi,
$}
\typeout{GRDISPLAY|411|\the\wd\grmathbox|\the\linewidth}
\ifdim\wd\grmathbox>\linewidth
\resizebox{\linewidth}{!}{\usebox{\grmathbox}}
\else\usebox{\grmathbox}\fi
\endgroup
\end{center}

because \((d^2/d\phi^2+1)(\phi\sin\phi)=2\cos\phi\). The other forcing
terms produce bounded shape corrections, not the accumulated rotation we
are seeking.

Now expand a slightly shifted oscillation:

\begin{center}
\begingroup
\sbox{\grmathbox}{$\displaystyle 
\frac{e}{p_{\mathrm{orb}}}\cos[(1-\delta)\phi]
\simeq\frac{e}{p_{\mathrm{orb}}}\cos\phi
+\frac{e\delta}{p_{\mathrm{orb}}}\phi\sin\phi.
$}
\typeout{GRDISPLAY|412|\the\wd\grmathbox|\the\linewidth}
\ifdim\wd\grmathbox>\linewidth
\resizebox{\linewidth}{!}{\usebox{\grmathbox}}
\else\usebox{\grmathbox}\fi
\endgroup
\end{center}

Matching coefficients gives \(\delta=3m/p_{\mathrm{orb}}\). One radial
cycle therefore takes slightly more than \(2\pi\) in azimuth:

\begin{center}
\begingroup
\sbox{\grmathbox}{$\displaystyle 
\boxed{\Delta\phi\simeq\frac{6\pi G_NM}
{a_{\mathrm{orb}}(1-e^2)c^2}}
$}
\typeout{GRDISPLAY|413|\the\wd\grmathbox|\the\linewidth}
\ifdim\wd\grmathbox>\linewidth
\resizebox{\linewidth}{!}{\usebox{\grmathbox}}
\else\usebox{\grmathbox}\fi
\endgroup
\end{center}

per orbit, where \(p_{\mathrm{orb}}=a_{\mathrm{orb}}(1-e^2)\) at the
needed Newtonian order. For Mercury, using
\(a_{\mathrm{orb}}\simeq5.79\times10^{10}\,\mathrm m\) and
\(e\simeq0.206\), this gives about \(0.104\) arcseconds per orbit, or
\(43\) arcseconds per century.

This is the relativistic contribution under the approximation of an
isolated spherical Sun. Planetary perturbations, solar structure, and
reference-frame modeling also affect the observed perihelion. The
success lies in calculating the appropriate additional contribution, not
declaring that every observed orbital change is relativistic.

\textbf{The lesson of these experiments:} Newtonian motion, clock
comparison, light bending, and radar delay interrogate related but
different pieces of the geometry. Agreement across them is more
informative than agreement with one attractive number.

\hypertarget{chapter-17}{%
\chapter{17. Black holes: when the causal structure becomes the main
character}\label{chapter-17}}

A black hole is not defined by especially strong acceleration, an
especially dark surface, or even a large local curvature. It is defined
by which events can send signals to the exterior future. To appreciate
how different that is, we will actually solve a piece of the Einstein
equation.

\hypertarget{solving-spherical-vacuum-where-schwarzschild-comes-from}{%
\section{17.1 Solving spherical vacuum: where Schwarzschild comes
from}\label{solving-spherical-vacuum-where-schwarzschild-comes-from}}

Set \(\Lambda=0\). Outside a static spherical source, use the
\textbf{areal radius} \(r\): a symmetry sphere has area \(4\pi r^2\).
This is a geometrically meaningful definition, not a promise that radial
proper distance equals \(r\).

For this subsection the chart is \((t,r,\theta,\phi)\), with \(t\) in
seconds. Write

\begin{center}
\begingroup
\sbox{\grmathbox}{$\displaystyle 
ds^2=-e^{2\alpha(r)}c^2dt^2+e^{2\beta(r)}dr^2+r^2d\Omega^2,
\qquad d\Omega^2=d\theta^2+\sin^2\theta\,d\phi^2.
$}
\typeout{GRDISPLAY|414|\the\wd\grmathbox|\the\linewidth}
\ifdim\wd\grmathbox>\linewidth
\resizebox{\linewidth}{!}{\usebox{\grmathbox}}
\else\usebox{\grmathbox}\fi
\endgroup
\end{center}

Spherical symmetry forbids a preferred angular direction. Staticity
permits a diagonal time-radial form in the exterior region. We have
reduced ten metric components to two unknown functions, without assuming
their values.

A few connection coefficients show the mechanism:

\begin{center}
\begingroup
\sbox{\grmathbox}{$\displaystyle 
\Gamma^t{}_{tr}=\alpha',\qquad
\Gamma^r{}_{tt}=c^2e^{2(\alpha-\beta)}\alpha',\qquad
\Gamma^r{}_{rr}=\beta',\qquad
\Gamma^r{}_{\theta\theta}=-re^{-2\beta}.
$}
\typeout{GRDISPLAY|415|\the\wd\grmathbox|\the\linewidth}
\ifdim\wd\grmathbox>\linewidth
\resizebox{\linewidth}{!}{\usebox{\grmathbox}}
\else\usebox{\grmathbox}\fi
\endgroup
\end{center}

A prime means \(d/dr\). The last expression remembers that spheres
change size as \(r\) changes. Such angular terms are essential even
though the unknown functions depend only on radius.

For example, the Ricci calculation gives

\begin{center}
\begingroup
\sbox{\grmathbox}{$\displaystyle 
R_{tt}=c^2e^{2(\alpha-\beta)}
\left[\alpha''+(\alpha')^2-\alpha'\beta'+\frac{2\alpha'}r\right].
$}
\typeout{GRDISPLAY|416|\the\wd\grmathbox|\the\linewidth}
\ifdim\wd\grmathbox>\linewidth
\resizebox{\linewidth}{!}{\usebox{\grmathbox}}
\else\usebox{\grmathbox}\fi
\endgroup
\end{center}

The second derivative comes from differentiating the connection; the
products come from its \(\Gamma\Gamma\) terms; the \(2/r\) accounts for
the two angular dimensions. Curvature is assembled from precisely the
operations learned earlier.

Define \(f(r)=e^{-2\beta(r)}\). Two convenient mixed Einstein components
are

\begin{center}
\begingroup
\sbox{\grmathbox}{$\displaystyle 
G^t{}_t=\frac{f-1}{r^2}+\frac{f'}r,
\qquad
G^r{}_r=\frac{f-1}{r^2}+\frac{2f\alpha'}r.
$}
\typeout{GRDISPLAY|417|\the\wd\grmathbox|\the\linewidth}
\ifdim\wd\grmathbox>\linewidth
\resizebox{\linewidth}{!}{\usebox{\grmathbox}}
\else\usebox{\grmathbox}\fi
\endgroup
\end{center}

Vacuum requires both to vanish. The first equation can be reorganized as

\begin{center}
\begingroup
\sbox{\grmathbox}{$\displaystyle 
\frac{d}{dr}\big[r(1-f)\big]=0.
$}
\typeout{GRDISPLAY|418|\the\wd\grmathbox|\the\linewidth}
\ifdim\wd\grmathbox>\linewidth
\resizebox{\linewidth}{!}{\usebox{\grmathbox}}
\else\usebox{\grmathbox}\fi
\endgroup
\end{center}

A derivative is zero, so the bracket is a constant: \(r(1-f)=C\). Thus
\(f=1-C/r\). This is the central integration, and it explains the
inverse-radius form instead of merely announcing it.

Subtracting the two vacuum equations gives

\begin{center}
\begingroup
\sbox{\grmathbox}{$\displaystyle 
0=\frac{2f}{r}(\alpha'+\beta').
$}
\typeout{GRDISPLAY|419|\the\wd\grmathbox|\the\linewidth}
\ifdim\wd\grmathbox>\linewidth
\resizebox{\linewidth}{!}{\usebox{\grmathbox}}
\else\usebox{\grmathbox}\fi
\endgroup
\end{center}

In the static exterior \(f\ne0\), hence \(\alpha+\beta\) is constant. A
constant rescaling of \(t\) removes that constant. Therefore
\(e^{2\alpha}=e^{-2\beta}=f\). The remaining angular vacuum equation is
satisfied by these functions; the contracted Bianchi identity explains
why all the apparent equations are not independent.

At large \(r\), match \(g_{tt}=-c^2(1-C/r)\) to the Newtonian clock
coefficient \(-c^2(1-2G_NM/(rc^2))\). This fixes \(C=2G_NM/c^2\). Define
the length \(m=G_NM/c^2\). The solution is

\begin{center}
\begingroup
\sbox{\grmathbox}{$\displaystyle 
\boxed{ds^2=-\left(1-\frac{2m}{r}\right)c^2dt^2
+\frac{dr^2}{1-2m/r}+r^2d\Omega^2.}
$}
\typeout{GRDISPLAY|420|\the\wd\grmathbox|\the\linewidth}
\ifdim\wd\grmathbox>\linewidth
\resizebox{\linewidth}{!}{\usebox{\grmathbox}}
\else\usebox{\grmathbox}\fi
\endgroup
\end{center}

We assumed staticity to make the derivation accessible.
Birkhoff\textquotesingle s theorem says something stronger: a
spherically symmetric vacuum region with \(\Lambda=0\) is locally
Schwarzschild even if the spherical matter boundary moves. A perfectly
spherical pulsating star does not broadcast tensor gravitational waves
into its vacuum exterior. The theorem does not describe a region filled
with an outgoing matter or radiation flux, which is not vacuum. A proof
outline is available in
\href{https://davidtong.org/pdfs/teaching/general-relativity/gr6.pdf}{David
Tong's black-hole lecture notes}.

\hypertarget{two-alarming-radii-two-very-different-problems}{%
\section{17.2 Two alarming radii, two very different
problems}\label{two-alarming-radii-two-very-different-problems}}

At \(r=2m\), the displayed \(g_{rr}\) diverges and \(g_{tt}\) vanishes.
At \(r=0\), multiple expressions fail. Are these the same sort of
disaster?

Compute a coordinate-invariant curvature quantity, the Kretschmann
scalar:

\begin{center}
\begingroup
\sbox{\grmathbox}{$\displaystyle 
\mathcal K=R_{\alpha\beta\gamma\delta}R^{\alpha\beta\gamma\delta}
=\frac{48m^2}{r^6}.
$}
\typeout{GRDISPLAY|421|\the\wd\grmathbox|\the\linewidth}
\ifdim\wd\grmathbox>\linewidth
\resizebox{\linewidth}{!}{\usebox{\grmathbox}}
\else\usebox{\grmathbox}\fi
\endgroup
\end{center}

It is finite at \(r=2m\) and diverges at \(r=0\). Finite scalar
invariants alone do not prove every conceivable spacetime point is
regular, but here an explicit nonsingular chart will establish
regularity at the horizon. At \(r=0\), the divergent invariant proves
the problem cannot be repaired by relabeling coordinates.

Meanwhile \(R_{\mu\nu}=0\) and \(R=0\) everywhere in the vacuum
exterior. There is abundant curvature despite the vanishing Ricci
tensor: this is Weyl curvature. Calling vacuum ``empty'' does not make
it geometrically featureless.

The horizon radius is

\begin{center}
\begingroup
\sbox{\grmathbox}{$\displaystyle 
r_s=2m=\frac{2G_NM}{c^2}\simeq2.95\,\mathrm{km}\left(\frac{M}{M_\odot}\right).
$}
\typeout{GRDISPLAY|422|\the\wd\grmathbox|\the\linewidth}
\ifdim\wd\grmathbox>\linewidth
\resizebox{\linewidth}{!}{\usebox{\grmathbox}}
\else\usebox{\grmathbox}\fi
\endgroup
\end{center}

At the horizon, the tidal scale \(G_NM/r_s^3\) is proportional to
\(M^{-2}\). A larger black hole can have gentler horizon tides.
Event-horizon status and local violence are different questions.

\hypertarget{repairing-the-horizon-with-eddington-finkelstein-coordinates}{%
\section{17.3 Repairing the horizon with Eddington-Finkelstein
coordinates}\label{repairing-the-horizon-with-eddington-finkelstein-coordinates}}

Define the tortoise coordinate

\begin{center}
\begingroup
\sbox{\grmathbox}{$\displaystyle 
r_*=r+2m\ln\left|\frac{r}{2m}-1\right|,
\qquad \frac{dr_*}{dr}=\frac1f,
\qquad f=1-\frac{2m}{r}.
$}
\typeout{GRDISPLAY|423|\the\wd\grmathbox|\the\linewidth}
\ifdim\wd\grmathbox>\linewidth
\resizebox{\linewidth}{!}{\usebox{\grmathbox}}
\else\usebox{\grmathbox}\fi
\endgroup
\end{center}

Then introduce an advanced time \(v=t+r_*/c\), still measured in
seconds. Since \(dt=dv-dr/(cf)\), substitution gives

\begin{center}
\begingroup
\sbox{\grmathbox}{$\displaystyle 
-fc^2\left(dv-\frac{dr}{cf}\right)^2+\frac{dr^2}{f}
=-fc^2dv^2+2c\,dv\,dr.
$}
\typeout{GRDISPLAY|424|\the\wd\grmathbox|\the\linewidth}
\ifdim\wd\grmathbox>\linewidth
\resizebox{\linewidth}{!}{\usebox{\grmathbox}}
\else\usebox{\grmathbox}\fi
\endgroup
\end{center}

The troublesome \(dr^2/f\) terms cancel exactly. The metric becomes

\begin{center}
\begingroup
\sbox{\grmathbox}{$\displaystyle 
\boxed{ds^2=-fc^2dv^2+2c\,dv\,dr+r^2d\Omega^2.}
$}
\typeout{GRDISPLAY|425|\the\wd\grmathbox|\the\linewidth}
\ifdim\wd\grmathbox>\linewidth
\resizebox{\linewidth}{!}{\usebox{\grmathbox}}
\else\usebox{\grmathbox}\fi
\endgroup
\end{center}

Its time-radial block has determinant \(-c^2\), including at \(r=2m\).
The future horizon is a perfectly regular place for these coordinates.

For radial light, set \(d\Omega=0\) and \(ds^2=0\). One family has
\(dv=0\): ingoing rays. The other satisfies

\begin{center}
\begingroup
\sbox{\grmathbox}{$\displaystyle 
\frac{dr}{dv}=\frac c2\left(1-\frac{2m}{r}\right).
$}
\typeout{GRDISPLAY|426|\the\wd\grmathbox|\the\linewidth}
\ifdim\wd\grmathbox>\linewidth
\resizebox{\linewidth}{!}{\usebox{\grmathbox}}
\else\usebox{\grmathbox}\fi
\endgroup
\end{center}

Outside, these outgoing rays increase their radius. At the horizon they
remain on it. Inside, even this outgoing family decreases its areal
radius. Future-directed timelike trajectories lie between the two null
directions and also move toward smaller \(r\).

This is the precise content behind the metaphor ``all roads point
inward.'' It does not mean a rocket engine becomes weak or a photon
slows locally. It means that the future light cones admit no escaping
causal direction in this black-hole interior.

The maximal mathematical extension of eternal Schwarzschild has
additional regions. A black hole produced by stellar collapse need not
contain its white-hole region or second exterior. An exact
metric\textquotesingle s maximal extension and the spacetime of a
particular formation process are distinct objects.

\hypertarget{falling-is-easy-hovering-is-the-expensive-activity}{%
\section{17.4 Falling is easy; hovering is the expensive
activity}\label{falling-is-easy-hovering-is-the-expensive-activity}}

For radial timelike motion, time-translation symmetry gives a
dimensionless conserved energy per unit rest energy,

\begin{center}
\begingroup
\sbox{\grmathbox}{$\displaystyle 
\mathcal E=f\frac{dt}{d\tau}.
$}
\typeout{GRDISPLAY|427|\the\wd\grmathbox|\the\linewidth}
\ifdim\wd\grmathbox>\linewidth
\resizebox{\linewidth}{!}{\usebox{\grmathbox}}
\else\usebox{\grmathbox}\fi
\endgroup
\end{center}

Substitute this into \(g_{\mu\nu}u^\mu u^\nu=-c^2\):

\begin{center}
\begingroup
\sbox{\grmathbox}{$\displaystyle 
\frac{1}{c^2}\left(\frac{dr}{d\tau}\right)^2=\mathcal E^2-f.
$}
\typeout{GRDISPLAY|428|\the\wd\grmathbox|\the\linewidth}
\ifdim\wd\grmathbox>\linewidth
\resizebox{\linewidth}{!}{\usebox{\grmathbox}}
\else\usebox{\grmathbox}\fi
\endgroup
\end{center}

For an object falling from rest at infinity, \(\mathcal E=1\), so

\begin{center}
\begingroup
\sbox{\grmathbox}{$\displaystyle 
\frac{dr}{d\tau}=-c\sqrt{\frac{2m}{r}}.
$}
\typeout{GRDISPLAY|429|\the\wd\grmathbox|\the\linewidth}
\ifdim\wd\grmathbox>\linewidth
\resizebox{\linewidth}{!}{\usebox{\grmathbox}}
\else\usebox{\grmathbox}\fi
\endgroup
\end{center}

Nothing diverges at the horizon. In this idealized classical trajectory,
proper time from horizon to \(r=0\) is

\begin{center}
\begingroup
\sbox{\grmathbox}{$\displaystyle 
\Delta\tau=\frac1c\int_0^{2m}\sqrt{\frac r{2m}}\,dr
=\frac{4m}{3c}.
$}
\typeout{GRDISPLAY|430|\the\wd\grmathbox|\the\linewidth}
\ifdim\wd\grmathbox>\linewidth
\resizebox{\linewidth}{!}{\usebox{\grmathbox}}
\else\usebox{\grmathbox}\fi
\endgroup
\end{center}

This is specific to that radial energy and classical solution, not a
universal countdown for every infaller. Schwarzschild coordinate time
diverges at horizon crossing because that chart fails there. Signals
reaching a distant observer become increasingly delayed and redshifted;
the object does not remain as a permanently bright frozen photograph.

Now hold a rocket at constant \(r>2m\). Its four-velocity has
\(u^t=1/\sqrt f\). Its radial four-acceleration is

\begin{center}
\begingroup
\sbox{\grmathbox}{$\displaystyle 
a^r=\Gamma^r{}_{tt}(u^t)^2=\frac{G_NM}{r^2}.
$}
\typeout{GRDISPLAY|431|\the\wd\grmathbox|\the\linewidth}
\ifdim\wd\grmathbox>\linewidth
\resizebox{\linewidth}{!}{\usebox{\grmathbox}}
\else\usebox{\grmathbox}\fi
\endgroup
\end{center}

This component happens to resemble Newton\textquotesingle s
acceleration, but the accelerometer measures the invariant magnitude

\begin{center}
\begingroup
\sbox{\grmathbox}{$\displaystyle 
\boxed{\sqrt{a_\mu a^\mu}=\frac{G_NM}{r^2\sqrt{1-2m/r}}.}
$}
\typeout{GRDISPLAY|432|\the\wd\grmathbox|\the\linewidth}
\ifdim\wd\grmathbox>\linewidth
\resizebox{\linewidth}{!}{\usebox{\grmathbox}}
\else\usebox{\grmathbox}\fi
\endgroup
\end{center}

It diverges on approach to the horizon. The diverging quantity belongs
to the family of observers trying to remain static. It does not imply a
freely falling observer measures infinite curvature there. At the
horizon, being static would require following a null worldline; no
massive rocket can do that.

\hypertarget{the-photon-sphere-and-isco-three-radii-you-should-never-merge}{%
\section{17.5 The photon sphere and ISCO: three radii you should never
merge}\label{the-photon-sphere-and-isco-three-radii-you-should-never-merge}}

Spherical symmetry lets a geodesic lie in an equatorial plane. Define
\(\ell=r^2d\phi/d\tau\). The timelike normalization becomes

\begin{center}
\begingroup
\sbox{\grmathbox}{$\displaystyle 
\frac{\dot r^2}{c^2}+V_{\mathrm{eff}}(r)=\mathcal E^2,
\qquad
V_{\mathrm{eff}}=\left(1-\frac{2m}{r}\right)
\left(1+\frac{\ell^2}{c^2r^2}\right).
$}
\typeout{GRDISPLAY|433|\the\wd\grmathbox|\the\linewidth}
\ifdim\wd\grmathbox>\linewidth
\resizebox{\linewidth}{!}{\usebox{\grmathbox}}
\else\usebox{\grmathbox}\fi
\endgroup
\end{center}

Dots here mean proper-time derivatives. Expanding the potential reveals
a term proportional to \(-1/r^3\), absent from the Newtonian effective
potential. It changes the centrifugal barrier near the hole.

A circular orbit requires \(V'_{\mathrm{eff}}=0\). Solving this
algebraic condition gives

\begin{center}
\begingroup
\sbox{\grmathbox}{$\displaystyle 
\frac{\ell^2}{c^2}=\frac{mr^2}{r-3m}.
$}
\typeout{GRDISPLAY|434|\the\wd\grmathbox|\the\linewidth}
\ifdim\wd\grmathbox>\linewidth
\resizebox{\linewidth}{!}{\usebox{\grmathbox}}
\else\usebox{\grmathbox}\fi
\endgroup
\end{center}

A small radial displacement is stable only when
\(V''_{\mathrm{eff}}>0\). At a circular orbit,

\begin{center}
\begingroup
\sbox{\grmathbox}{$\displaystyle 
V''_{\mathrm{eff}}=\frac{2m(r-6m)}{r^3(r-3m)}.
$}
\typeout{GRDISPLAY|435|\the\wd\grmathbox|\the\linewidth}
\ifdim\wd\grmathbox>\linewidth
\resizebox{\linewidth}{!}{\usebox{\grmathbox}}
\else\usebox{\grmathbox}\fi
\endgroup
\end{center}

Thus circular timelike geodesics exist for \(r>3m\) and are stable for
\(r>6m\). The marginal boundary \(r=6m\) is the \textbf{innermost stable
circular orbit}, or ISCO. Unstable circular orbits can exist between
\(3m\) and \(6m\); ``unstable'' does not mean ``algebraically
nonexistent.''

For null geodesics, the effective potential is proportional to
\(f/r^2\). Its derivative vanishes at \(r=3m\), a maximum. That is the
\textbf{photon sphere}. Its circular light orbits are unstable.

\begingroup\small\setlength{\tabcolsep}{5pt}\renewcommand{\arraystretch}{1.13}

\begin{longtable}[]{@{}
  >{\raggedright\arraybackslash}p{(\columnwidth - 2\tabcolsep) * \real{0.3000}}
  >{\raggedright\arraybackslash}p{(\columnwidth - 2\tabcolsep) * \real{0.7000}}@{}}
\toprule\noalign{}
\begin{minipage}[b]{\linewidth}\raggedright
Radius in Schwarzschild
\end{minipage} & \begin{minipage}[b]{\linewidth}\raggedright
Meaning
\end{minipage} \\
\midrule\noalign{}
\endhead
\bottomrule\noalign{}
\endlastfoot
\(2m\) & Event horizon of the black-hole solution \\
\(3m\) & Unstable circular null orbits; photon sphere \\
\(6m\) & Marginally stable circular timelike orbit; ISCO \\
\end{longtable}

\endgroup

The photon sphere is not a material surface. Nor is a black-hole image a
direct photograph of the horizon\textquotesingle s coordinate radius:
lensing, emission, absorption, and observer geometry intervene.

For completeness, the effective potential also produces the orbit
equation used in Chapter 16. Put \(u=1/r\) and use
\(\dot r=-\ell\,du/d\phi\). Differentiate the radial energy equation
with respect to \(\phi\), cancel the common first-derivative factor
where it is nonzero, and extend smoothly through turning points. The
result is \(u''+u=G_NM/\ell^2+3mu^2\). The apparently magical precession
term is just the same Schwarzschild geometry expressed as an orbit
shape.

\hypertarget{an-event-horizon-knows-about-the-future}{%
\section{17.6 An event horizon knows about the
future}\label{an-event-horizon-knows-about-the-future}}

In an asymptotically flat spacetime, the black-hole region is the set of
events that cannot send a future-directed causal signal to future null
infinity. Its boundary is the event horizon. The qualifier ``future''
means the entire future development matters.

A sufficiently small freely falling laboratory generally cannot
determine by purely local experiments whether it has crossed an event
horizon. It can measure curvature and tidal forces, but horizon
membership is a global causal statement. Locally defined trapped
surfaces and foliation-dependent apparent horizons provide useful
related diagnostics; they are not interchangeable definitions.

An analogy is a shipping port that eventually closes permanently.
Whether a departure can still reach the open sea depends on the future
geometry of the route, not just the waves measured beside the boat. The
limitation is that a black hole\textquotesingle s obstruction is causal
geometry itself, not an authority changing a schedule.

\hypertarget{rotation-kerr-frame-dragging-and-carefully-qualified-no-hair-claims}{%
\section{17.7 Rotation: Kerr, frame dragging, and carefully qualified
no-hair
claims}\label{rotation-kerr-frame-dragging-and-carefully-qualified-no-hair-claims}}

Realistic rotating black-hole models use the Kerr solution. Deriving it
fully is substantially harder than the spherical calculation; the
following is a guided reading of its geometry, not a concealed claim of
a complete derivation.

Define

\begin{center}
\begingroup
\sbox{\grmathbox}{$\displaystyle 
m=\frac{G_NM}{c^2},\qquad
a_K=\frac{J}{Mc},\qquad
\chi_{\mathrm{spin}}=\frac{a_K}{m}=\frac{cJ}{G_NM^2},
$}
\typeout{GRDISPLAY|436|\the\wd\grmathbox|\the\linewidth}
\ifdim\wd\grmathbox>\linewidth
\resizebox{\linewidth}{!}{\usebox{\grmathbox}}
\else\usebox{\grmathbox}\fi
\endgroup
\end{center}

and

\begin{center}
\begingroup
\sbox{\grmathbox}{$\displaystyle 
\Sigma=r^2+a_K^2\cos^2\theta,
\qquad \Delta=r^2-2mr+a_K^2.
$}
\typeout{GRDISPLAY|437|\the\wd\grmathbox|\the\linewidth}
\ifdim\wd\grmathbox>\linewidth
\resizebox{\linewidth}{!}{\usebox{\grmathbox}}
\else\usebox{\grmathbox}\fi
\endgroup
\end{center}

In Boyer-Lindquist coordinates, with \(t\) in seconds, the metric is

\begin{center}
\begingroup
\sbox{\grmathbox}{$\displaystyle 
\begin{aligned}
ds^2={}&-\left(1-\frac{2mr}{\Sigma}\right)c^2dt^2
-\frac{4ma_Kr\sin^2\theta}{\Sigma}\,c\,dt\,d\phi
+\frac{\Sigma}{\Delta}dr^2+\Sigma d\theta^2\\
&+\left(r^2+a_K^2+
\frac{2ma_K^2r\sin^2\theta}{\Sigma}\right)\sin^2\theta\,d\phi^2.
\end{aligned}
$}
\typeout{GRDISPLAY|438|\the\wd\grmathbox|\the\linewidth}
\ifdim\wd\grmathbox>\linewidth
\resizebox{\linewidth}{!}{\usebox{\grmathbox}}
\else\usebox{\grmathbox}\fi
\endgroup
\end{center}

Setting \(a_K=0\) recovers Schwarzschild. The new \(dt\,d\phi\) term
mixes time evolution with angular motion. Because a cross term in
\(ds^2\) is \(2g_{t\phi}dt\,d\phi\), its displayed coefficient is twice
the metric component. Missing that factor produces an impressively wrong
frame-dragging calculation.

An observer with zero conserved axial angular momentum satisfies

\begin{center}
\begingroup
\sbox{\grmathbox}{$\displaystyle 
p_\phi=0
\quad\Longrightarrow\quad
\frac{d\phi}{dt}=-\frac{g_{t\phi}}{g_{\phi\phi}}.
$}
\typeout{GRDISPLAY|439|\the\wd\grmathbox|\the\linewidth}
\ifdim\wd\grmathbox>\linewidth
\resizebox{\linewidth}{!}{\usebox{\grmathbox}}
\else\usebox{\grmathbox}\fi
\endgroup
\end{center}

Zero angular momentum therefore does not mean zero coordinate angular
velocity. This is one operational expression of \textbf{frame dragging}.
It is an off-diagonal geometric effect, not viscous friction against an
invisible fluid.

The roots of \(\Delta=0\) are

\begin{center}
\begingroup
\sbox{\grmathbox}{$\displaystyle 
r_\pm=m\pm\sqrt{m^2-a_K^2}.
$}
\typeout{GRDISPLAY|440|\the\wd\grmathbox|\the\linewidth}
\ifdim\wd\grmathbox>\linewidth
\resizebox{\linewidth}{!}{\usebox{\grmathbox}}
\else\usebox{\grmathbox}\fi
\endgroup
\end{center}

For the Kerr black-hole family, \(|\chi_{\mathrm{spin}}|\le1\). The
outer root is the event-horizon radius. The surface where the stationary
Killing field becomes null is instead \(g_{tt}=0\):

\begin{center}
\begingroup
\sbox{\grmathbox}{$\displaystyle 
r_{\mathrm{ergo}}(\theta)=m+\sqrt{m^2-a_K^2\cos^2\theta}.
$}
\typeout{GRDISPLAY|441|\the\wd\grmathbox|\the\linewidth}
\ifdim\wd\grmathbox>\linewidth
\resizebox{\linewidth}{!}{\usebox{\grmathbox}}
\else\usebox{\grmathbox}\fi
\endgroup
\end{center}

Between this surface and the outer horizon lies the \textbf{ergoregion}.
There, remaining at fixed spatial coordinates is impossible for a
timelike observer, although escape can still be possible. At the poles
the two surfaces meet. Elsewhere they differ: inability to remain
stationary is weaker than inability to escape.

``No hair'' is not a theorem that every possible gravitating theory has
only two black-hole parameters. Kerr uniqueness results apply under
substantial assumptions about vacuum Einstein gravity, stationarity,
asymptotics, horizon structure, and regularity. For example, a rigorous
result establishes Kerr uniqueness within a class of connected,
nondegenerate, analytic regular vacuum black holes. Additional fields,
different asymptotics, or dynamical settings change the question. See
the primary mathematical result,
\href{https://arxiv.org/abs/0806.0016}{Chruściel and Costa, \emph{On
uniqueness of stationary vacuum black holes}}.

The useful physical idea is that an isolated black hole settling into
the appropriate stationary vacuum state is described by very few
exterior parameters. The qualification is what turns a slogan into a
scientific statement.

\hypertarget{chapter-18}{%
\chapter{18. Gravitational waves: curvature can carry a
message}\label{chapter-18}}

The metric responds to matter, but it is not required to follow matter
instantaneously. Einstein\textquotesingle s equation is a dynamical
field equation. Once disturbed, geometry has propagating degrees of
freedom of its own.

The key conceptual distinction is between a field\textquotesingle s
\textbf{source} and the \textbf{field already present}.
Maxwell\textquotesingle s equations permit light in a charge-free
region. Einstein\textquotesingle s equation permits gravitational waves
in a matter-free region. ``The source is zero here'' does not imply
``the solution is zero here.''

\hypertarget{linearization-is-a-controlled-approximation-not-a-different-theory}{%
\section{18.1 Linearization is a controlled approximation, not a
different
theory}\label{linearization-is-a-controlled-approximation-not-a-different-theory}}

Use Cartesian coordinates \(x^0=ct\) on a Minkowski background, take
\(\Lambda=0\), and write

\begin{center}
\begingroup
\sbox{\grmathbox}{$\displaystyle 
g_{\mu\nu}=\eta_{\mu\nu}+h_{\mu\nu},
\qquad |h_{\mu\nu}|\ll1.
$}
\typeout{GRDISPLAY|442|\the\wd\grmathbox|\the\linewidth}
\ifdim\wd\grmathbox>\linewidth
\resizebox{\linewidth}{!}{\usebox{\grmathbox}}
\else\usebox{\grmathbox}\fi
\endgroup
\end{center}

Smallness is asserted in a suitable background-adapted coordinate
system; a wild coordinate transformation can make components large
without creating strong physical gravity. We keep first-order terms in
\(h\). The inverse metric is

\begin{center}
\begingroup
\sbox{\grmathbox}{$\displaystyle 
g^{\mu\nu}=\eta^{\mu\nu}-h^{\mu\nu}+O(h^2),
$}
\typeout{GRDISPLAY|443|\the\wd\grmathbox|\the\linewidth}
\ifdim\wd\grmathbox>\linewidth
\resizebox{\linewidth}{!}{\usebox{\grmathbox}}
\else\usebox{\grmathbox}\fi
\endgroup
\end{center}

where perturbation indices are raised with \(\eta\). Multiplying the two
metrics verifies the minus sign: the first-order cross terms cancel.

The linearized connection is

\begin{center}
\begingroup
\sbox{\grmathbox}{$\displaystyle 
\Gamma^{(1)\rho}{}_{\mu\nu}
=\frac12\eta^{\rho\sigma}
(\partial_\mu h_{\sigma\nu}+\partial_\nu h_{\sigma\mu}
-\partial_\sigma h_{\mu\nu}).
$}
\typeout{GRDISPLAY|444|\the\wd\grmathbox|\the\linewidth}
\ifdim\wd\grmathbox>\linewidth
\resizebox{\linewidth}{!}{\usebox{\grmathbox}}
\else\usebox{\grmathbox}\fi
\endgroup
\end{center}

Why not use the full inverse metric in this expression? Its correction
is already first order, so multiplying it by \(\partial h\) would give a
second-order term. Why omit \(\Gamma\Gamma\) from linearized curvature?
For the same reason: each connection is first order around this constant
background.

Contracting the derivative terms in Riemann gives

\begin{center}
\begingroup
\sbox{\grmathbox}{$\displaystyle 
R^{(1)}_{\mu\nu}=\frac12\left(
\partial_\rho\partial_\mu h^\rho{}_\nu
+\partial_\rho\partial_\nu h^\rho{}_\mu
-\Box h_{\mu\nu}-\partial_\mu\partial_\nu h\right),
$}
\typeout{GRDISPLAY|445|\the\wd\grmathbox|\the\linewidth}
\ifdim\wd\grmathbox>\linewidth
\resizebox{\linewidth}{!}{\usebox{\grmathbox}}
\else\usebox{\grmathbox}\fi
\endgroup
\end{center}

with

\begin{center}
\begingroup
\sbox{\grmathbox}{$\displaystyle 
h=\eta^{\mu\nu}h_{\mu\nu},
\qquad
\Box=\eta^{\mu\nu}\partial_\mu\partial_\nu
=-\frac1{c^2}\partial_t^2+\nabla^2.
$}
\typeout{GRDISPLAY|446|\the\wd\grmathbox|\the\linewidth}
\ifdim\wd\grmathbox>\linewidth
\resizebox{\linewidth}{!}{\usebox{\grmathbox}}
\else\usebox{\grmathbox}\fi
\endgroup
\end{center}

The trace is

\begin{center}
\begingroup
\sbox{\grmathbox}{$\displaystyle 
R^{(1)}=\partial_\mu\partial_\nu h^{\mu\nu}-\Box h.
$}
\typeout{GRDISPLAY|447|\the\wd\grmathbox|\the\linewidth}
\ifdim\wd\grmathbox>\linewidth
\resizebox{\linewidth}{!}{\usebox{\grmathbox}}
\else\usebox{\grmathbox}\fi
\endgroup
\end{center}

Read the Ricci formula slowly. Two terms involve divergences of \(h\),
one is a wave operator acting on each component, and one differentiates
the trace. That structure suggests choosing variables and coordinates
that isolate the wave operator.

\hypertarget{trace-reversal-and-lorenz-gauge-cleaning-the-algebra-without-deleting-physics}{%
\section{18.2 Trace reversal and Lorenz gauge: cleaning the algebra
without deleting
physics}\label{trace-reversal-and-lorenz-gauge-cleaning-the-algebra-without-deleting-physics}}

Define the trace-reversed perturbation

\begin{center}
\begingroup
\sbox{\grmathbox}{$\displaystyle 
\bar h_{\mu\nu}=h_{\mu\nu}-\frac12\eta_{\mu\nu}h.
$}
\typeout{GRDISPLAY|448|\the\wd\grmathbox|\the\linewidth}
\ifdim\wd\grmathbox>\linewidth
\resizebox{\linewidth}{!}{\usebox{\grmathbox}}
\else\usebox{\grmathbox}\fi
\endgroup
\end{center}

In four dimensions its trace is \(\bar h=-h\), since contracting
\(\eta_{\mu\nu}\) produces four. Consequently the inverse operation has
the same form:

\begin{center}
\begingroup
\sbox{\grmathbox}{$\displaystyle 
h_{\mu\nu}=\bar h_{\mu\nu}-\frac12\eta_{\mu\nu}\bar h.
$}
\typeout{GRDISPLAY|449|\the\wd\grmathbox|\the\linewidth}
\ifdim\wd\grmathbox>\linewidth
\resizebox{\linewidth}{!}{\usebox{\grmathbox}}
\else\usebox{\grmathbox}\fi
\endgroup
\end{center}

The name refers to the sign of the trace. It does not mean reversing
every tensor component.

The linearized Einstein tensor becomes

\begin{center}
\begingroup
\sbox{\grmathbox}{$\displaystyle 
G^{(1)}_{\mu\nu}
=-\frac12\Box\bar h_{\mu\nu}
+\partial_{(\mu}\partial^\alpha\bar h_{\nu)\alpha}
-\frac12\eta_{\mu\nu}\partial^\alpha\partial^\beta\bar h_{\alpha\beta}.
$}
\typeout{GRDISPLAY|450|\the\wd\grmathbox|\the\linewidth}
\ifdim\wd\grmathbox>\linewidth
\resizebox{\linewidth}{!}{\usebox{\grmathbox}}
\else\usebox{\grmathbox}\fi
\endgroup
\end{center}

Symmetrization includes its factor of \(1/2\). The unhelpful terms are
now explicitly divergences of \(\bar h\).

Under an infinitesimal coordinate change \(x'^\mu=x^\mu+\xi^\mu\), the
first-order perturbation changes by

\begin{center}
\begingroup
\sbox{\grmathbox}{$\displaystyle 
h'_{\mu\nu}=h_{\mu\nu}-\partial_\mu\xi_\nu-\partial_\nu\xi_\mu.
$}
\typeout{GRDISPLAY|451|\the\wd\grmathbox|\the\linewidth}
\ifdim\wd\grmathbox>\linewidth
\resizebox{\linewidth}{!}{\usebox{\grmathbox}}
\else\usebox{\grmathbox}\fi
\endgroup
\end{center}

This is a different coordinate description of the same metric, to the
stated order. Its trace-reversed divergence transforms as

\begin{center}
\begingroup
\sbox{\grmathbox}{$\displaystyle 
\partial^\mu\bar h'_{\mu\nu}
=\partial^\mu\bar h_{\mu\nu}-\Box\xi_\nu.
$}
\typeout{GRDISPLAY|452|\the\wd\grmathbox|\the\linewidth}
\ifdim\wd\grmathbox>\linewidth
\resizebox{\linewidth}{!}{\usebox{\grmathbox}}
\else\usebox{\grmathbox}\fi
\endgroup
\end{center}

Locally, with suitable initial and boundary conditions, solve
\(\Box\xi_\nu=\partial^\mu\bar h_{\mu\nu}\). In the resulting
coordinates,

\begin{center}
\begingroup
\sbox{\grmathbox}{$\displaystyle 
\partial^\mu\bar h_{\mu\nu}=0.
$}
\typeout{GRDISPLAY|453|\the\wd\grmathbox|\the\linewidth}
\ifdim\wd\grmathbox>\linewidth
\resizebox{\linewidth}{!}{\usebox{\grmathbox}}
\else\usebox{\grmathbox}\fi
\endgroup
\end{center}

This is usually called Lorenz, harmonic, or de Donder gauge in this
context. It removes four coordinate-dependent combinations, not four
physical forces.

The field equation collapses to

\begin{center}
\begingroup
\sbox{\grmathbox}{$\displaystyle 
\boxed{\Box\bar h_{\mu\nu}=-\frac{16\pi G_N}{c^4}T_{\mu\nu}.}
$}
\typeout{GRDISPLAY|454|\the\wd\grmathbox|\the\linewidth}
\ifdim\wd\grmathbox>\linewidth
\resizebox{\linewidth}{!}{\usebox{\grmathbox}}
\else\usebox{\grmathbox}\fi
\endgroup
\end{center}

The coefficient follows directly from
\(G^{(1)}_{\mu\nu}=-\Box\bar h_{\mu\nu}/2\) and
Einstein\textquotesingle s \(8\pi G_N/c^4\). If you lose a factor of two
here, every predicted wave amplitude will faithfully preserve your
mistake.

Taking a divergence requires \(\partial^\mu T_{\mu\nu}=0\) at this
order. The source cannot be chosen arbitrarily; its leading dynamics
must be consistent with energy-momentum conservation. For a
self-gravitating compact system, systematically including the
gravitational contribution requires the appropriate perturbative
expansion rather than inserting an inconsistent prescribed matter
motion.

\hypertarget{why-there-are-two-polarizations-not-ten}{%
\section{18.3 Why there are two polarizations, not
ten}\label{why-there-are-two-polarizations-not-ten}}

In vacuum, \(\Box\bar h_{\mu\nu}=0\). A plane wave has the form

\begin{center}
\begingroup
\sbox{\grmathbox}{$\displaystyle 
\bar h_{\mu\nu}=\operatorname{Re}
\left[A_{\mu\nu}e^{ik_\alpha x^\alpha}\right].
$}
\typeout{GRDISPLAY|455|\the\wd\grmathbox|\the\linewidth}
\ifdim\wd\grmathbox>\linewidth
\resizebox{\linewidth}{!}{\usebox{\grmathbox}}
\else\usebox{\grmathbox}\fi
\endgroup
\end{center}

Applying \(\Box\) multiplies this by \(-k_\alpha k^\alpha\), so a
nontrivial wave requires

\begin{center}
\begingroup
\sbox{\grmathbox}{$\displaystyle 
k_\alpha k^\alpha=0.
$}
\typeout{GRDISPLAY|456|\the\wd\grmathbox|\the\linewidth}
\ifdim\wd\grmathbox>\linewidth
\resizebox{\linewidth}{!}{\usebox{\grmathbox}}
\else\usebox{\grmathbox}\fi
\endgroup
\end{center}

Its wave covector is null; for a wave propagating along \(z\), the phase
depends on \(t-z/c\). Gravitational disturbances propagate on the
background light cone at this order.

The gauge condition gives \(k^\mu A_{\mu\nu}=0\), four restrictions on a
symmetric tensor\textquotesingle s ten components. But this gauge is not
fully fixed: transformations satisfying \(\Box\xi^\mu=0\) preserve it.
For a nonzero vacuum plane wave, four residual gauge choices remove four
further amplitude combinations. What remains are two independent
radiative polarizations.

This \(10-4-4=2\) count is an on-shell plane-wave count. It should not
be applied blindly to arbitrary metric perturbations with matter,
boundaries, or constrained nonradiative components.

A convenient representative is \textbf{transverse-traceless}, or TT,
gauge. For propagation along \(z\),

\begin{center}
\begingroup
\sbox{\grmathbox}{$\displaystyle 
h^{\mathrm{TT}}_{0\mu}=0,\qquad
h^{\mathrm{TT}}_{ij}=
\begin{pmatrix}
h_+&h_\times&0\\
h_\times&-h_+&0\\
0&0&0
\end{pmatrix},
\qquad h_+=h_+(t-z/c),\quad h_\times=h_\times(t-z/c).
$}
\typeout{GRDISPLAY|457|\the\wd\grmathbox|\the\linewidth}
\ifdim\wd\grmathbox>\linewidth
\resizebox{\linewidth}{!}{\usebox{\grmathbox}}
\else\usebox{\grmathbox}\fi
\endgroup
\end{center}

``Transverse'' means the perturbation has no component along the
propagation direction. ``Traceless'' means its first-order expansion
along one transverse axis is accompanied by contraction along the other.

The plus polarization stretches an initially circular ring of free test
particles along one axis and compresses it along the perpendicular axis.
Half a cycle later the roles reverse. The cross polarization does the
same with axes rotated by \(45^\circ\). Superpositions produce
elliptical or circular polarization.

An ordinary vector\textquotesingle s components mix by an angle
\(\theta\) when transverse axes rotate by \(\theta\). These two
polarization amplitudes mix by \(2\theta\). That doubled angular
response is a classical signature of the spin-2 character of the field.
It does not require a quantum detector to see the relevant
transformation rule.

\hypertarget{but-the-particles-do-not-move-in-tt-coordinates-exactly}{%
\section{18.4 ``But the particles do not move in TT coordinates!''
Exactly}\label{but-the-particles-do-not-move-in-tt-coordinates-exactly}}

For initially stationary free test particles in TT coordinates,
\(\Gamma^i{}_{00}=0\) at first order. Their spatial coordinates can
remain constant. Does that mean the wave is an illusion?

No. Coordinates are labels; the metric tells us the distance between
labels. Along a short arm directed along \(x\),

\begin{center}
\begingroup
\sbox{\grmathbox}{$\displaystyle 
L_x(t)=\int_0^{L_0}\sqrt{1+h_+}\,dx
\simeq L_0\left(1+\frac12h_+\right).
$}
\typeout{GRDISPLAY|458|\the\wd\grmathbox|\the\linewidth}
\ifdim\wd\grmathbox>\linewidth
\resizebox{\linewidth}{!}{\usebox{\grmathbox}}
\else\usebox{\grmathbox}\fi
\endgroup
\end{center}

Along \(y\), the corresponding change is \(-h_+/2\). This expression
assumes an arm short compared with the wavelength, so the field is
approximately uniform across it on the chosen time slice.

The same observable effect can be obtained from curvature, avoiding the
impression that moving grid lines create physics. In a local inertial
frame associated with the detector,

\begin{center}
\begingroup
\sbox{\grmathbox}{$\displaystyle 
R_{i0j0}=-\frac{1}{2c^2}\frac{\partial^2h^{\mathrm{TT}}_{ij}}{\partial t^2}.
$}
\typeout{GRDISPLAY|459|\the\wd\grmathbox|\the\linewidth}
\ifdim\wd\grmathbox>\linewidth
\resizebox{\linewidth}{!}{\usebox{\grmathbox}}
\else\usebox{\grmathbox}\fi
\endgroup
\end{center}

The geodesic-deviation equation therefore gives

\begin{center}
\begingroup
\sbox{\grmathbox}{$\displaystyle 
\frac{d^2\xi^i}{dt^2}
=-c^2R^i{}_{0j0}\xi^j
=\frac12\ddot h^{\mathrm{TT\,i}}{}_j\xi^j.
$}
\typeout{GRDISPLAY|460|\the\wd\grmathbox|\the\linewidth}
\ifdim\wd\grmathbox>\linewidth
\resizebox{\linewidth}{!}{\usebox{\grmathbox}}
\else\usebox{\grmathbox}\fi
\endgroup
\end{center}

To first order, with initial conditions corresponding to undisturbed
separations before a passing wave,

\begin{center}
\begingroup
\sbox{\grmathbox}{$\displaystyle 
\xi^i(t)\simeq\left(\delta^i{}_j+\frac12h^{\mathrm{TT\,i}}{}_j(t)\right)\xi_0^j.
$}
\typeout{GRDISPLAY|461|\the\wd\grmathbox|\the\linewidth}
\ifdim\wd\grmathbox>\linewidth
\resizebox{\linewidth}{!}{\usebox{\grmathbox}}
\else\usebox{\grmathbox}\fi
\endgroup
\end{center}

In one coordinate description the masses stay at fixed coordinates and
the metric changes their separation. In another their coordinates
respond to tidal acceleration. They predict the same instrument
response. Linearized Riemann is unchanged by a pure linearized gauge
transformation about flat spacetime because the extra third derivatives
cancel.

A laser interferometer compares the phases accumulated by light
traversing differently oriented arms and returning to a common observer.
In the ideal long-wavelength, favorably oriented case,

\begin{center}
\begingroup
\sbox{\grmathbox}{$\displaystyle 
\frac{\delta L_x-\delta L_y}{L_0}=h_+.
$}
\typeout{GRDISPLAY|462|\the\wd\grmathbox|\the\linewidth}
\ifdim\wd\grmathbox>\linewidth
\resizebox{\linewidth}{!}{\usebox{\grmathbox}}
\else\usebox{\grmathbox}\fi
\endgroup
\end{center}

For \(L_0=4\,\mathrm{km}\) and \(h_+=10^{-21}\), the differential
equivalent length is \(4\times10^{-18}\,\mathrm m\). Each individual
arm\textquotesingle s change in this idealized example is half that
magnitude with opposite sign. Actual responses include source direction,
polarization, optical configuration, and frequency-dependent light
travel effects.

Why does the light not ``stretch exactly with the apparatus'' and erase
the measurement? Because an interferometer measures the relation between
null propagation, mirror worldlines, and a clock at the beamsplitter.
The wave creates a time-dependent anisotropic tidal geometry. There is
no universal rescaling of every relevant relation that turns this into
an unchanged experiment. The full light-travel calculation agrees with
the simple strain description in its regime of validity.

\hypertarget{retarded-solutions-the-field-receives-yesterdays-news}{%
\section{18.5 Retarded solutions: the field receives
yesterday\textquotesingle s
news}\label{retarded-solutions-the-field-receives-yesterdays-news}}

With no incoming radiation and an appropriate localized weak source, the
retarded solution is

\begin{center}
\begingroup
\sbox{\grmathbox}{$\displaystyle 
\bar h_{\mu\nu}(t,\mathbf x)
=\frac{4G_N}{c^4}\int
\frac{T_{\mu\nu}\left(t-|\mathbf x-\mathbf x'|/c,\mathbf x'\right)}
{|\mathbf x-\mathbf x'|}\,d^3x'.
$}
\typeout{GRDISPLAY|463|\the\wd\grmathbox|\the\linewidth}
\ifdim\wd\grmathbox>\linewidth
\resizebox{\linewidth}{!}{\usebox{\grmathbox}}
\else\usebox{\grmathbox}\fi
\endgroup
\end{center}

Every source element contributes at its own retarded time. The
denominator spreads the disturbance over distance; the time argument
enforces finite propagation. This is the gravitational cousin of a
retarded electromagnetic potential, with a tensor source and tensor
response.

For a source much smaller than its characteristic gravitational
wavelength, observed far away at distance \(D\), approximate the
denominator by \(D\) and the leading source time by \(t-D/c\). Then

\begin{center}
\begingroup
\sbox{\grmathbox}{$\displaystyle 
\bar h_{ij}\simeq\frac{4G_N}{c^4D}\int T_{ij}(t-D/c,\mathbf x')\,d^3x'.
$}
\typeout{GRDISPLAY|464|\the\wd\grmathbox|\the\linewidth}
\ifdim\wd\grmathbox>\linewidth
\resizebox{\linewidth}{!}{\usebox{\grmathbox}}
\else\usebox{\grmathbox}\fi
\endgroup
\end{center}

The spatial stresses look inconvenient. Conservation converts them into
a more intuitive changing mass distribution. Define the leading mass
quadrupole moment before trace removal,

\begin{center}
\begingroup
\sbox{\grmathbox}{$\displaystyle 
I_{ij}=\frac1{c^2}\int T^{00}x_i x_j\,d^3x
\simeq\int\rho\,x_i x_j\,d^3x.
$}
\typeout{GRDISPLAY|465|\the\wd\grmathbox|\the\linewidth}
\ifdim\wd\grmathbox>\linewidth
\resizebox{\linewidth}{!}{\usebox{\grmathbox}}
\else\usebox{\grmathbox}\fi
\endgroup
\end{center}

Here \(\rho=\epsilon/c^2\) at leading nonrelativistic order. Spatial
index positions in these Euclidean Cartesian formulas are
interchangeable.

Use \(\partial_tT^{00}=-c\partial_kT^{k0}\) and integrate by parts.
Surface terms vanish for the localized source:

\begin{center}
\begingroup
\sbox{\grmathbox}{$\displaystyle 
\dot I_{ij}=\frac1c\int(T_i{}^0x_j+T_j{}^0x_i)\,d^3x.
$}
\typeout{GRDISPLAY|466|\the\wd\grmathbox|\the\linewidth}
\ifdim\wd\grmathbox>\linewidth
\resizebox{\linewidth}{!}{\usebox{\grmathbox}}
\else\usebox{\grmathbox}\fi
\endgroup
\end{center}

Differentiate again, use \(\partial_tT_i{}^0=-c\partial_kT_i{}^k\), and
integrate by parts again:

\begin{center}
\begingroup
\sbox{\grmathbox}{$\displaystyle 
\boxed{\ddot I_{ij}=2\int T_{ij}\,d^3x.}
$}
\typeout{GRDISPLAY|467|\the\wd\grmathbox|\the\linewidth}
\ifdim\wd\grmathbox>\linewidth
\resizebox{\linewidth}{!}{\usebox{\grmathbox}}
\else\usebox{\grmathbox}\fi
\endgroup
\end{center}

The factor of two comes from the two positions in the symmetric product
\(x_i x_j\). This manipulation is doing physical work: conservation
relates momentum transport to changes in the shape of the mass
distribution.

Define its trace-free part

\begin{center}
\begingroup
\sbox{\grmathbox}{$\displaystyle 
Q_{ij}=I_{ij}-\frac13\delta_{ij}\delta^{kl}I_{kl}.
$}
\typeout{GRDISPLAY|468|\the\wd\grmathbox|\the\linewidth}
\ifdim\wd\grmathbox>\linewidth
\resizebox{\linewidth}{!}{\usebox{\grmathbox}}
\else\usebox{\grmathbox}\fi
\endgroup
\end{center}

The leading radiative field is

\begin{center}
\begingroup
\sbox{\grmathbox}{$\displaystyle 
\boxed{h^{\mathrm{TT}}_{ij}(t,\mathbf x)
=\frac{2G_N}{c^4D}\,\mathcal P_{ij}{}^{kl}(\mathbf n)
\ddot Q_{kl}(t-D/c).}
$}
\typeout{GRDISPLAY|469|\the\wd\grmathbox|\the\linewidth}
\ifdim\wd\grmathbox>\linewidth
\resizebox{\linewidth}{!}{\usebox{\grmathbox}}
\else\usebox{\grmathbox}\fi
\endgroup
\end{center}

The unit vector \(\mathbf n\) points toward the observer. The TT
projector is

\begin{center}
\begingroup
\sbox{\grmathbox}{$\displaystyle 
P_{ij}=\delta_{ij}-n_i n_j,
\qquad
\mathcal P_{ij}{}^{kl}=P_i{}^kP_j{}^l-\frac12P_{ij}P^{kl},
$}
\typeout{GRDISPLAY|470|\the\wd\grmathbox|\the\linewidth}
\ifdim\wd\grmathbox>\linewidth
\resizebox{\linewidth}{!}{\usebox{\grmathbox}}
\else\usebox{\grmathbox}\fi
\endgroup
\end{center}

acting here on symmetric tensors. The first operation projects both
indices into the observer\textquotesingle s transverse plane. The second
removes the trace within that two-dimensional plane, hence \(1/2\)
rather than the \(1/3\) used to remove a three-dimensional trace. That
difference is a useful check that the geometry remains attached to the
algebra.

\hypertarget{why-a-violently-breathing-sphere-still-does-not-radiate-tensor-waves}{%
\section{18.6 Why a violently breathing sphere still does not radiate
tensor
waves}\label{why-a-violently-breathing-sphere-still-does-not-radiate-tensor-waves}}

An isolated source\textquotesingle s leading mass monopole is its
conserved total mass-energy. Its leading mass dipole describes its
center of mass, whose velocity is fixed by conserved total momentum. A
changing monopole or accelerating isolated center of mass is unavailable
as a radiative degree of freedom. The leading current dipole is tied to
conserved angular momentum. The first available leading tensor radiation
comes from a changing quadrupole.

Two masses orbiting each other continually change their quadrupole even
if their total mass and center of mass remain constant. A perfectly
spherical body expanding and contracting has no trace-free mass
quadrupole. Exact spherical symmetry also forbids tensor gravitational
radiation beyond this approximation, consistent with
Birkhoff\textquotesingle s theorem in the vacuum exterior.

The tempting rule ``anything accelerating emits gravitational waves'' is
therefore too crude. Radiation depends on the collective multipole
structure and conservation laws, not a tally of individual
accelerations. Other theories with additional radiative fields can have
different multipole channels.

\hypertarget{wave-energy-lives-one-order-beyond-the-linear-equation}{%
\section{18.7 Wave energy lives one order beyond the linear
equation}\label{wave-energy-lives-one-order-beyond-the-linear-equation}}

The field equation is nonlinear. Expand it schematically as

\begin{center}
\begingroup
\sbox{\grmathbox}{$\displaystyle 
G[g_{\mathrm{background}}+h]
=G[g_{\mathrm{background}}]+G^{(1)}[h]+G^{(2)}[h,h]+\cdots.
$}
\typeout{GRDISPLAY|471|\the\wd\grmathbox|\the\linewidth}
\ifdim\wd\grmathbox>\linewidth
\resizebox{\linewidth}{!}{\usebox{\grmathbox}}
\else\usebox{\grmathbox}\fi
\endgroup
\end{center}

The linear term describes propagation. Quadratic terms describe how the
waves themselves influence the slowly varying background. When the
wavelength is short compared with the background curvature scale, one
can average over several wavelengths while remaining local relative to
that background. This produces an effective wave stress-energy tensor.

The scale separation is essential. This is not an exact, generally
covariant pointwise gravitational-energy tensor for arbitrary
geometries. Isaacson\textquotesingle s high-frequency treatment
explicitly constructs the appropriate averaged effective description.
See \href{https://doi.org/10.1103/PhysRev.166.1272}{Isaacson,
\emph{Gravitational Radiation in the Limit of High Frequency. II}}.

For a nearly planar wave in a locally flat wave zone, the average flux
is

\begin{center}
\begingroup
\sbox{\grmathbox}{$\displaystyle 
F=\frac{c^3}{32\pi G_N}
\left\langle\dot h^{\mathrm{TT}}_{ij}\dot h_{\mathrm{TT}}^{ij}\right\rangle
=\frac{c^3}{16\pi G_N}
\left\langle\dot h_+^2+\dot h_\times^2\right\rangle.
$}
\typeout{GRDISPLAY|472|\the\wd\grmathbox|\the\linewidth}
\ifdim\wd\grmathbox>\linewidth
\resizebox{\linewidth}{!}{\usebox{\grmathbox}}
\else\usebox{\grmathbox}\fi
\endgroup
\end{center}

Angle brackets denote the averaging. The second equality uses the fact
that the polarization matrix contributes twice each squared amplitude.
Flux has units of power per area; \(c^3/G_N\) times an inverse time
squared has exactly those units.

The leading total radiated power from a slow isolated source is

\begin{center}
\begingroup
\sbox{\grmathbox}{$\displaystyle 
P_{\mathrm{GW}}=\frac{G_N}{5c^5}
\left\langle\dddot Q_{ij}\dddot Q^{ij}\right\rangle.
$}
\typeout{GRDISPLAY|473|\the\wd\grmathbox|\the\linewidth}
\ifdim\wd\grmathbox>\linewidth
\resizebox{\linewidth}{!}{\usebox{\grmathbox}}
\else\usebox{\grmathbox}\fi
\endgroup
\end{center}

The extra time derivative relative to the strain reflects that energy
depends on the rate of change of the wave. We can even uncover the
mysterious \(1/5\) rather than asking you to trust a coefficient that
apparently wandered in from another textbook.

Write \(A_{ij}=\dddot Q_{ij}\) at a fixed retarded time. Substituting
the quadrupole strain into the flux and integrating over a large sphere
gives

\begin{center}
\begingroup
\sbox{\grmathbox}{$\displaystyle 
P_{\mathrm{GW}}=\frac{G_N}{8\pi c^5}\int d\Omega\,
\left\langle A^{\mathrm{TT}}_{ij}A_{\mathrm{TT}}^{ij}\right\rangle.
$}
\typeout{GRDISPLAY|474|\the\wd\grmathbox|\the\linewidth}
\ifdim\wd\grmathbox>\linewidth
\resizebox{\linewidth}{!}{\usebox{\grmathbox}}
\else\usebox{\grmathbox}\fi
\endgroup
\end{center}

The distance cancels: wave amplitude falls as \(1/D\), flux as
\(1/D^2\), and sphere area grows as \(D^2\). The remaining calculation
asks what fraction of a trace-free tensor survives transverse
projection, averaged over viewing directions.

Define \(S=A_{ij}A^{ij}\), \(B_i=A_{ij}n^j\), and \(s=A_{ij}n^i n^j\).
Expanding the projector gives

\begin{center}
\begingroup
\sbox{\grmathbox}{$\displaystyle 
A^{\mathrm{TT}}_{ij}A_{\mathrm{TT}}^{ij}=S-2B_iB^i+\frac12s^2.
$}
\typeout{GRDISPLAY|475|\the\wd\grmathbox|\the\linewidth}
\ifdim\wd\grmathbox>\linewidth
\resizebox{\linewidth}{!}{\usebox{\grmathbox}}
\else\usebox{\grmathbox}\fi
\endgroup
\end{center}

For a uniform average over directions, isotropy requires

\begin{center}
\begingroup
\sbox{\grmathbox}{$\displaystyle 
\langle n_i n_j\rangle_\Omega=\frac{\delta_{ij}}3,
\qquad
\langle n_i n_j n_k n_l\rangle_\Omega
=\frac{\delta_{ij}\delta_{kl}+\delta_{ik}\delta_{jl}
+\delta_{il}\delta_{jk}}{15}.
$}
\typeout{GRDISPLAY|476|\the\wd\grmathbox|\the\linewidth}
\ifdim\wd\grmathbox>\linewidth
\resizebox{\linewidth}{!}{\usebox{\grmathbox}}
\else\usebox{\grmathbox}\fi
\endgroup
\end{center}

Why these forms? There is no preferred direction, so only Kronecker
deltas can appear. Symmetry fixes the combinations. Contracting indices
and using \(n_i n^i=1\) fixes the denominators. Consequently
\(\langle B_iB^i\rangle_\Omega=S/3\) and, since \(A\) is trace-free,
\(\langle s^2\rangle_\Omega=2S/15\). Therefore

\begin{center}
\begingroup
\sbox{\grmathbox}{$\displaystyle 
\langle A^{\mathrm{TT}}_{ij}A_{\mathrm{TT}}^{ij}\rangle_\Omega
=S-\frac{2S}3+\frac S{15}=\frac{2S}5.
$}
\typeout{GRDISPLAY|477|\the\wd\grmathbox|\the\linewidth}
\ifdim\wd\grmathbox>\linewidth
\resizebox{\linewidth}{!}{\usebox{\grmathbox}}
\else\usebox{\grmathbox}\fi
\endgroup
\end{center}

Multiplying by the sphere\textquotesingle s solid angle \(4\pi\)
produces \(8\pi S/5\). That cancels the flux prefactor\textquotesingle s
\(8\pi\) and leaves precisely \(G_N/(5c^5)\). The time average and the
angular average serve different purposes; their labels keep those
operations distinct.

The \(c^{-5}\) makes ordinary laboratory gravitational radiation
staggeringly weak. Large masses, rapid asymmetric motion, and compact
configurations help overcome that suppression.

For compact bodies, the leading quadrupole law can describe their slow
orbital dynamics even when gravity inside each body is strong; its
derivation must then be embedded in a consistent approximation for the
effective orbital source. It is not a demand that each black hole itself
be a weak-field object.

For a circular binary with total mass \(M\), reduced mass
\(\mu=m_1m_2/M\), and separation \(r\), the leading result is

\begin{center}
\begingroup
\sbox{\grmathbox}{$\displaystyle 
P_{\mathrm{GW}}=\frac{32}{5}\frac{G_N^4\mu^2M^3}{c^5r^5}.
$}
\typeout{GRDISPLAY|478|\the\wd\grmathbox|\the\linewidth}
\ifdim\wd\grmathbox>\linewidth
\resizebox{\linewidth}{!}{\usebox{\grmathbox}}
\else\usebox{\grmathbox}\fi
\endgroup
\end{center}

To see the binary coefficient, put the relative position at
\(\mathbf r=r(\cos\Omega t,\sin\Omega t,0)\) in its center-of-mass
frame. Then \(I_{ij}=\mu r_i r_j\). The changing quadrupole components
are

\begin{center}
\begingroup
\sbox{\grmathbox}{$\displaystyle 
\begin{aligned}
Q_{xx}&=\mu r^2\left(\frac16+\frac12\cos2\Omega t\right),\\
Q_{yy}&=\mu r^2\left(\frac16-\frac12\cos2\Omega t\right),\\
Q_{xy}&=Q_{yx}=\frac12\mu r^2\sin2\Omega t,\qquad
Q_{zz}=-\frac13\mu r^2.
\end{aligned}
$}
\typeout{GRDISPLAY|479|\the\wd\grmathbox|\the\linewidth}
\ifdim\wd\grmathbox>\linewidth
\resizebox{\linewidth}{!}{\usebox{\grmathbox}}
\else\usebox{\grmathbox}\fi
\endgroup
\end{center}

Three time derivatives remove the constant terms and give amplitudes
\(4\mu r^2\Omega^3\). Contracting the tensor counts both \(xy\) and
\(yx\) and yields

\begin{center}
\begingroup
\sbox{\grmathbox}{$\displaystyle 
\dddot Q_{ij}\dddot Q^{ij}=32\mu^2r^4\Omega^6.
$}
\typeout{GRDISPLAY|480|\the\wd\grmathbox|\the\linewidth}
\ifdim\wd\grmathbox>\linewidth
\resizebox{\linewidth}{!}{\usebox{\grmathbox}}
\else\usebox{\grmathbox}\fi
\endgroup
\end{center}

Use \(\Omega^2=G_NM/r^3\) in the quadrupole power law to obtain the
displayed binary luminosity. This leading calculation treats the orbit
as approximately circular and nearly unchanged over one cycle; the slow
inspiral is then included through energy balance.

Its Newtonian binding energy is \(E=-G_N\mu M/(2r)\). Because
\(dE/dt=-P_{\mathrm{GW}}<0\), the energy becomes more negative and \(r\)
decreases. Kepler\textquotesingle s relation \(\Omega^2=G_NM/r^3\) then
makes the orbital frequency increase. The system loses energy and speeds
up: a bound gravitational system has just ambushed everyday friction
intuition.

The dominant wave frequency is twice the orbital frequency,
\(f=\Omega/\pi\). Rewriting \(E\) and \(P\) in terms of \(f\),
differentiating \(E(f)\), and applying energy balance yields

\begin{center}
\begingroup
\sbox{\grmathbox}{$\displaystyle 
\dot f=\frac{96}{5}\pi^{8/3}
\left(\frac{G_N\mathcal M}{c^3}\right)^{5/3}f^{11/3},
\qquad
\mathcal M=\mu^{3/5}M^{2/5}.
$}
\typeout{GRDISPLAY|481|\the\wd\grmathbox|\the\linewidth}
\ifdim\wd\grmathbox>\linewidth
\resizebox{\linewidth}{!}{\usebox{\grmathbox}}
\else\usebox{\grmathbox}\fi
\endgroup
\end{center}

The combination \(\mathcal M\) is the \textbf{chirp mass}. Its name is
operational: the measured rate at which the signal\textquotesingle s
pitch rises strongly constrains it. This is a leading inspiral formula,
requiring slow enough orbital motion for the approximation. Near merger,
higher-order analytic methods and numerical solutions of
Einstein\textquotesingle s equation become necessary.

\hypertarget{a-real-event-that-turned-these-symbols-into-data}{%
\section{18.8 A real event that turned these symbols into
data}\label{a-real-event-that-turned-these-symbols-into-data}}

On September 14, 2015, LIGO detected GW150914. The discovery report
described a signal rising from approximately \(35\) to \(250\) Hz with
peak strain about \(10^{-21}\). Its inferred source was a merging binary
black hole; the initial analysis estimated that roughly three solar
masses of energy were radiated. These are findings of the original
analysis, with model-dependent parameter estimates and uncertainties,
rather than exact source properties.
\href{https://arxiv.org/abs/1602.03837}{LIGO Scientific Collaboration
and Virgo Collaboration, \emph{Observation of Gravitational Waves from a
Binary Black Hole Merger}}.

The conceptual achievement is broader than hearing a cosmic chirp. A
theory derived from a variational principle and local geometric
identities predicted a timed sequence of tidal distortions caused by
distant dynamical curvature. Instruments measured that sequence.
Agreement constrains alternatives; it does not logically prove that no
other theory can agree in the tested regime.

\hypertarget{chapter-19}{%
\chapter{19. Cosmology: Einstein\textquotesingle s equation for the
large-scale universe}\label{chapter-19}}

Black holes exploit isolation. Cosmology exploits symmetry of a
different kind: on sufficiently large scales, approximate homogeneity
and isotropy. This does not mean the universe has no galaxies. It means
we construct a smooth background description whose departures can then
be studied as perturbations.

Homogeneity says no spatial location is special in the background model.
Isotropy says no spatial direction is special for its fundamental
observers. Neither assumption requires time independence. The universe
is allowed to evolve while treating every background location
equivalently.

\hypertarget{flrw-geometry-with-the-units-declared-before-the-equations-multiply}{%
\section{19.1 FLRW geometry, with the units declared before the
equations
multiply}\label{flrw-geometry-with-the-units-declared-before-the-equations-multiply}}

Use the chart \((t,\chi,\theta,\phi)\), with \(t\) in seconds and the
comoving spatial coordinates dimensionless. Choose

\begin{center}
\begingroup
\sbox{\grmathbox}{$\displaystyle 
\boxed{ds^2=-c^2dt^2+a^2(t)
\left[\frac{d\chi^2}{1-k\chi^2}+\chi^2d\Omega^2\right],}
$}
\typeout{GRDISPLAY|482|\the\wd\grmathbox|\the\linewidth}
\ifdim\wd\grmathbox>\linewidth
\resizebox{\linewidth}{!}{\usebox{\grmathbox}}
\else\usebox{\grmathbox}\fi
\endgroup
\end{center}

where \(a(t)\) has units of length and \(k\in\{-1,0,+1\}\) is
dimensionless. This is the Friedmann-Lemaître-Robertson-Walker metric.
The spatial slices have scalar curvature

\begin{center}
\begingroup
\sbox{\grmathbox}{$\displaystyle 
{}^{(3)}R=\frac{6k}{a^2}.
$}
\typeout{GRDISPLAY|483|\the\wd\grmathbox|\the\linewidth}
\ifdim\wd\grmathbox>\linewidth
\resizebox{\linewidth}{!}{\usebox{\grmathbox}}
\else\usebox{\grmathbox}\fi
\endgroup
\end{center}

Positive \(k\) describes positive constant spatial curvature, negative
\(k\) negative curvature, and zero \(k\) flat spatial slices. Local
curvature does not by itself settle every global topology question.

Some books instead make \(a\) dimensionless and place length units in
the spatial coordinates or curvature parameter. Either convention works.
Combining the formulas without converting conventions does not. Here,
when we want a dimensionless normalized scale factor, we will explicitly
write \(A(t)=a(t)/a(t_0)\).

Observers at fixed comoving coordinates have \(d\tau=dt\). The cosmic
time is their proper time; it is not a new absolute Newtonian time
available in every spacetime. It is singled out by this geometry and
matter congruence.

The Hubble parameter is

\begin{center}
\begingroup
\sbox{\grmathbox}{$\displaystyle 
H=\frac{\dot a}{a},
$}
\typeout{GRDISPLAY|484|\the\wd\grmathbox|\the\linewidth}
\ifdim\wd\grmathbox>\linewidth
\resizebox{\linewidth}{!}{\usebox{\grmathbox}}
\else\usebox{\grmathbox}\fi
\endgroup
\end{center}

with units of inverse time. It measures the fractional expansion rate.
An expansion factor is a ratio of scale factors, while \(H\) is a rate
of change. Confusing those is like confusing the current size of your
bank balance with its interest rate.

\hypertarget{deriving-the-curvature-without-hiding-the-time-dependence}{%
\section{19.2 Deriving the curvature without hiding the time
dependence}\label{deriving-the-curvature-without-hiding-the-time-dependence}}

Write the spatial bracket as \(\gamma_{ij}dx^i dx^j\), so
\(g_{ij}=a^2\gamma_{ij}\) and \(\gamma\) has unit constant curvature
\(k\). The useful connection coefficients are

\begin{center}
\begingroup
\sbox{\grmathbox}{$\displaystyle 
\Gamma^t{}_{ij}=\frac{a\dot a}{c^2}\gamma_{ij},
\qquad
\Gamma^i{}_{tj}=H\delta^i{}_j,
\qquad
\Gamma^i{}_{jk}={}^{(3)}\Gamma^i{}_{jk}[\gamma].
$}
\typeout{GRDISPLAY|485|\the\wd\grmathbox|\the\linewidth}
\ifdim\wd\grmathbox>\linewidth
\resizebox{\linewidth}{!}{\usebox{\grmathbox}}
\else\usebox{\grmathbox}\fi
\endgroup
\end{center}

The first says spatial motion participates in time evolution because
spatial distances depend on time. The second says a spatial basis
carried through cosmic time changes its scale. The third contains the
ordinary intrinsic connection of the constant-curvature spatial
geometry.

For the time-time Ricci component, \(\Gamma^i{}_{ti}=3H\). The relevant
contraction gives

\begin{center}
\begingroup
\sbox{\grmathbox}{$\displaystyle 
R_{tt}=-\partial_t(3H)-3H^2
=-3\frac{\ddot a}{a}.
$}
\typeout{GRDISPLAY|486|\the\wd\grmathbox|\the\linewidth}
\ifdim\wd\grmathbox>\linewidth
\resizebox{\linewidth}{!}{\usebox{\grmathbox}}
\else\usebox{\grmathbox}\fi
\endgroup
\end{center}

The cancellation uses \(\dot H=\ddot a/a-H^2\). There are three equal
spatial contributions because there are three equivalent spatial
directions.

For the spatial Ricci components, the intrinsic curvature contributes
\(2k\gamma_{ij}\). The time-dependent terms contribute
\((a\ddot a+2\dot a^2)\gamma_{ij}/c^2\). Thus

\begin{center}
\begingroup
\sbox{\grmathbox}{$\displaystyle 
R_{ij}=\left(\frac{a\ddot a+2\dot a^2}{c^2}+2k\right)\gamma_{ij}.
$}
\typeout{GRDISPLAY|487|\the\wd\grmathbox|\the\linewidth}
\ifdim\wd\grmathbox>\linewidth
\resizebox{\linewidth}{!}{\usebox{\grmathbox}}
\else\usebox{\grmathbox}\fi
\endgroup
\end{center}

Contract with \(g^{tt}=-1/c^2\) and \(g^{ij}=\gamma^{ij}/a^2\):

\begin{center}
\begingroup
\sbox{\grmathbox}{$\displaystyle 
R=\frac6{c^2}\left(\frac{\ddot a}{a}+H^2+\frac{kc^2}{a^2}\right).
$}
\typeout{GRDISPLAY|488|\the\wd\grmathbox|\the\linewidth}
\ifdim\wd\grmathbox>\linewidth
\resizebox{\linewidth}{!}{\usebox{\grmathbox}}
\else\usebox{\grmathbox}\fi
\endgroup
\end{center}

Notice that \(k=0\) does not generally make this vanish. Spatial
flatness is not spacetime flatness. Evolving distances generate
spacetime curvature even when every spatial slice is intrinsically
Euclidean.

Combining Ricci and its trace gives

\begin{center}
\begingroup
\sbox{\grmathbox}{$\displaystyle 
G_{tt}=3\left(H^2+\frac{kc^2}{a^2}\right),
$}
\typeout{GRDISPLAY|489|\the\wd\grmathbox|\the\linewidth}
\ifdim\wd\grmathbox>\linewidth
\resizebox{\linewidth}{!}{\usebox{\grmathbox}}
\else\usebox{\grmathbox}\fi
\endgroup
\end{center}

and

\begin{center}
\begingroup
\sbox{\grmathbox}{$\displaystyle 
G^i{}_j=-\frac1{c^2}
\left(2\frac{\ddot a}{a}+H^2+\frac{kc^2}{a^2}\right)\delta^i{}_j.
$}
\typeout{GRDISPLAY|490|\the\wd\grmathbox|\the\linewidth}
\ifdim\wd\grmathbox>\linewidth
\resizebox{\linewidth}{!}{\usebox{\grmathbox}}
\else\usebox{\grmathbox}\fi
\endgroup
\end{center}

The \(tt\) component has inverse-time-squared units because this chart
uses \(t\), not \(ct\). The mixed spatial components have
inverse-length-squared units. This is consistent tensor dimensional
bookkeeping, not an inconsistency to be repaired by arbitrarily
inserting \(c\).

\hypertarget{the-friedmann-equations-what-controls-expansion-and-acceleration}{%
\section{19.3 The Friedmann equations: what controls expansion and
acceleration}\label{the-friedmann-equations-what-controls-expansion-and-acceleration}}

Homogeneity and isotropy select a perfect-fluid background stress-energy
tensor. Let \(\epsilon\) be physical rest-frame energy density and \(p\)
pressure. In the comoving chart,

\begin{center}
\begingroup
\sbox{\grmathbox}{$\displaystyle 
T^t{}_t=-\epsilon,\qquad T^i{}_j=p\delta^i{}_j,
\qquad T_{tt}=\epsilon c^2.
$}
\typeout{GRDISPLAY|491|\the\wd\grmathbox|\the\linewidth}
\ifdim\wd\grmathbox>\linewidth
\resizebox{\linewidth}{!}{\usebox{\grmathbox}}
\else\usebox{\grmathbox}\fi
\endgroup
\end{center}

Substitute the time-time components into
\(G_{\mu\nu}+\Lambda g_{\mu\nu}=8\pi G_NT_{\mu\nu}/c^4\):

\begin{center}
\begingroup
\sbox{\grmathbox}{$\displaystyle 
3\left(H^2+\frac{kc^2}{a^2}\right)-\Lambda c^2
=\frac{8\pi G_N}{c^2}\epsilon.
$}
\typeout{GRDISPLAY|492|\the\wd\grmathbox|\the\linewidth}
\ifdim\wd\grmathbox>\linewidth
\resizebox{\linewidth}{!}{\usebox{\grmathbox}}
\else\usebox{\grmathbox}\fi
\endgroup
\end{center}

Divide by three and rearrange:

\begin{center}
\begingroup
\sbox{\grmathbox}{$\displaystyle 
\boxed{H^2=\frac{8\pi G_N}{3c^2}\epsilon
-\frac{kc^2}{a^2}+\frac{\Lambda c^2}{3}.}
$}
\typeout{GRDISPLAY|493|\the\wd\grmathbox|\the\linewidth}
\ifdim\wd\grmathbox>\linewidth
\resizebox{\linewidth}{!}{\usebox{\grmathbox}}
\else\usebox{\grmathbox}\fi
\endgroup
\end{center}

This is the first Friedmann equation. In mass-equivalent density
\(\rho=\epsilon/c^2\), the matter term is \(8\pi G_N\rho/3\). That
conversion explains many apparently different textbook versions.

The spatial equation gives

\begin{center}
\begingroup
\sbox{\grmathbox}{$\displaystyle 
2\frac{\ddot a}{a}+H^2+\frac{kc^2}{a^2}
=-\frac{8\pi G_N}{c^2}p+\Lambda c^2.
$}
\typeout{GRDISPLAY|494|\the\wd\grmathbox|\the\linewidth}
\ifdim\wd\grmathbox>\linewidth
\resizebox{\linewidth}{!}{\usebox{\grmathbox}}
\else\usebox{\grmathbox}\fi
\endgroup
\end{center}

Eliminate \(H^2+kc^2/a^2\) using the first equation:

\begin{center}
\begingroup
\sbox{\grmathbox}{$\displaystyle 
\boxed{\frac{\ddot a}{a}
=-\frac{4\pi G_N}{3c^2}(\epsilon+3p)
+\frac{\Lambda c^2}{3}.}
$}
\typeout{GRDISPLAY|495|\the\wd\grmathbox|\the\linewidth}
\ifdim\wd\grmathbox>\linewidth
\resizebox{\linewidth}{!}{\usebox{\grmathbox}}
\else\usebox{\grmathbox}\fi
\endgroup
\end{center}

This is the acceleration equation. Energy density and isotropic pressure
both gravitate. The three pressures arise from the three spatial
directions, not from an arbitrary correction inserted to surprise
Newton.

A positive expansion rate \(H>0\) does not imply accelerating expansion
\(\ddot a>0\). A ball thrown upward moves upward while slowing.
Likewise, a matter-filled model can grow in size while its growth rate
decreases.

Conversely, a positive cosmological constant contributes positively to
\(\ddot a/a\). If it dominates, expansion accelerates. The equations
make the condition quantitative instead of relying on the ambiguous
phrase ``repulsive gravity.''

\hypertarget{conservation-becomes-the-universes-first-law-of-thermodynamics}{%
\section{19.4 Conservation becomes the universe\textquotesingle s first
law of
thermodynamics}\label{conservation-becomes-the-universes-first-law-of-thermodynamics}}

The time component of \(\nabla_\mu T^{\mu\nu}=0\) gives

\begin{center}
\begingroup
\sbox{\grmathbox}{$\displaystyle 
\boxed{\dot\epsilon+3H(\epsilon+p)=0.}
$}
\typeout{GRDISPLAY|496|\the\wd\grmathbox|\the\linewidth}
\ifdim\wd\grmathbox>\linewidth
\resizebox{\linewidth}{!}{\usebox{\grmathbox}}
\else\usebox{\grmathbox}\fi
\endgroup
\end{center}

For a fixed comoving volume, its physical volume is proportional to
\(a^3\). Multiply the conservation equation by \(a^3\):

\begin{center}
\begingroup
\sbox{\grmathbox}{$\displaystyle 
\frac{d}{dt}(\epsilon a^3)=-p\frac{d}{dt}(a^3).
$}
\typeout{GRDISPLAY|497|\the\wd\grmathbox|\the\linewidth}
\ifdim\wd\grmathbox>\linewidth
\resizebox{\linewidth}{!}{\usebox{\grmathbox}}
\else\usebox{\grmathbox}\fi
\endgroup
\end{center}

This has the familiar form \(dE=-p\,dV\). As the volume expands,
positive pressure reduces the energy within that comoving volume. The
fluid does expansion work in this local continuum sense.

These equations are not three independent pieces of information.
Differentiate the first Friedmann equation and use the continuity
equation; away from a turning point, dividing by \(H\) recovers the
acceleration equation. At \(H=0\), use the original Einstein and
conservation equations rather than dividing by zero. Bianchi consistency
is doing its job even when an algebraic shortcut is unavailable.

For a separately conserved component with constant equation-of-state
parameter

\begin{center}
\begingroup
\sbox{\grmathbox}{$\displaystyle 
w=\frac p\epsilon,
$}
\typeout{GRDISPLAY|498|\the\wd\grmathbox|\the\linewidth}
\ifdim\wd\grmathbox>\linewidth
\resizebox{\linewidth}{!}{\usebox{\grmathbox}}
\else\usebox{\grmathbox}\fi
\endgroup
\end{center}

the continuity equation becomes

\begin{center}
\begingroup
\sbox{\grmathbox}{$\displaystyle 
\frac{d\epsilon}{\epsilon}=-3(1+w)\frac{da}{a}.
$}
\typeout{GRDISPLAY|499|\the\wd\grmathbox|\the\linewidth}
\ifdim\wd\grmathbox>\linewidth
\resizebox{\linewidth}{!}{\usebox{\grmathbox}}
\else\usebox{\grmathbox}\fi
\endgroup
\end{center}

Integrate:

\begin{center}
\begingroup
\sbox{\grmathbox}{$\displaystyle 
\boxed{\epsilon=\epsilon_0A^{-3(1+w)},\qquad A=\frac{a}{a_0}.}
$}
\typeout{GRDISPLAY|500|\the\wd\grmathbox|\the\linewidth}
\ifdim\wd\grmathbox>\linewidth
\resizebox{\linewidth}{!}{\usebox{\grmathbox}}
\else\usebox{\grmathbox}\fi
\endgroup
\end{center}

The integration is a statement about dilution and work, not an
additional gravitational law. If components exchange energy, each gets
an exchange term and need not obey this separate scaling; the total
still obeys conservation.

\begingroup\small\setlength{\tabcolsep}{5pt}\renewcommand{\arraystretch}{1.13}

\begin{longtable}[]{@{}
  >{\raggedright\arraybackslash}p{(\columnwidth - 6\tabcolsep) * \real{0.2308}}
  >{\raggedleft\arraybackslash}p{(\columnwidth - 6\tabcolsep) * \real{0.3077}}
  >{\raggedright\arraybackslash}p{(\columnwidth - 6\tabcolsep) * \real{0.2308}}
  >{\raggedright\arraybackslash}p{(\columnwidth - 6\tabcolsep) * \real{0.2308}}@{}}
\toprule\noalign{}
\begin{minipage}[b]{\linewidth}\raggedright
Component
\end{minipage} & \begin{minipage}[b]{\linewidth}\raggedleft
Approximate \(w\)
\end{minipage} & \begin{minipage}[b]{\linewidth}\raggedright
Energy-density scaling
\end{minipage} & \begin{minipage}[b]{\linewidth}\raggedright
Physical reason
\end{minipage} \\
\midrule\noalign{}
\endhead
\bottomrule\noalign{}
\endlastfoot
Nonrelativistic matter, or dust & \(0\) & \(A^{-3}\) & Approximately
fixed rest energy per particle, diluted by volume \\
Radiation & \(1/3\) & \(A^{-4}\) & Volume dilution plus redshift of each
quantum\textquotesingle s energy \\
Cosmological constant as a fluid & \(-1\) & Constant & Negative pressure
exactly offsets dilution in the continuity equation \\
\end{longtable}

\endgroup

The dust model does not mean microscopic dust grains specifically. It
means negligible pressure relative to energy density at the scale and
accuracy being modeled.

If \(\Lambda\) is moved to the matter side, its effective density and
pressure are

\begin{center}
\begingroup
\sbox{\grmathbox}{$\displaystyle 
\epsilon_\Lambda=\frac{\Lambda c^4}{8\pi G_N},
\qquad p_\Lambda=-\epsilon_\Lambda.
$}
\typeout{GRDISPLAY|501|\the\wd\grmathbox|\the\linewidth}
\ifdim\wd\grmathbox>\linewidth
\resizebox{\linewidth}{!}{\usebox{\grmathbox}}
\else\usebox{\grmathbox}\fi
\endgroup
\end{center}

Either keep \(\Lambda\) explicit in the Friedmann equations or include
this component in total \(\epsilon,p\) and remove the explicit term.
Doing both counts the same effect twice.

A comoving volume filled with this effective component gains total
energy as it grows, since its energy density remains constant. This does
not violate the continuity equation; its negative pressure makes the
right-hand side \(-p\,dV\) positive. In a general expanding spacetime
there is no global timelike translation symmetry supplying a universally
conserved total energy of the elementary mechanics kind.

\hypertarget{solving-for-the-scale-factor-three-recognizable-cosmic-personalities}{%
\section{19.5 Solving for the scale factor: three recognizable cosmic
personalities}\label{solving-for-the-scale-factor-three-recognizable-cosmic-personalities}}

Consider a spatially flat universe dominated by a single separately
conserved constant-\(w\) component, with no additional explicit
\(\Lambda\). For \(w>-1\), choose the expanding branch. Combining the
first Friedmann equation with the density scaling gives

\begin{center}
\begingroup
\sbox{\grmathbox}{$\displaystyle 
\frac{\dot A}{A}=C A^{-3(1+w)/2},
$}
\typeout{GRDISPLAY|502|\the\wd\grmathbox|\the\linewidth}
\ifdim\wd\grmathbox>\linewidth
\resizebox{\linewidth}{!}{\usebox{\grmathbox}}
\else\usebox{\grmathbox}\fi
\endgroup
\end{center}

where \(C>0\) is a constant with inverse-time units. Move the power of
\(A\) to the left and integrate:

\begin{center}
\begingroup
\sbox{\grmathbox}{$\displaystyle 
A^{3(1+w)/2}\propto t-t_B,
\qquad
\boxed{a(t)\propto(t-t_B)^{2/[3(1+w)]}.}
$}
\typeout{GRDISPLAY|503|\the\wd\grmathbox|\the\linewidth}
\ifdim\wd\grmathbox>\linewidth
\resizebox{\linewidth}{!}{\usebox{\grmathbox}}
\else\usebox{\grmathbox}\fi
\endgroup
\end{center}

The time \(t_B\) is the integration constant locating \(a=0\) in this
classical idealized solution. The approximation is not a license to
extrapolate an arbitrarily chosen matter model into the quantum-gravity
regime.

For dust, \(a\propto t^{2/3}\); for radiation, \(a\propto t^{1/2}\)
after shifting \(t_B\) to zero. Both expand while decelerating.
Radiation decelerates more strongly because its pressure also
contributes to the acceleration equation.

For a positive cosmological constant alone in a spatially flat expanding
slicing,

\begin{center}
\begingroup
\sbox{\grmathbox}{$\displaystyle 
H=\sqrt{\frac{\Lambda c^2}{3}}=\text{constant},
\qquad a(t)\propto e^{Ht}.
$}
\typeout{GRDISPLAY|504|\the\wd\grmathbox|\the\linewidth}
\ifdim\wd\grmathbox>\linewidth
\resizebox{\linewidth}{!}{\usebox{\grmathbox}}
\else\usebox{\grmathbox}\fi
\endgroup
\end{center}

The power-law derivation excluded \(w=-1\), so the exponential must be
obtained separately. Plugging \(w=-1\) into the power-law exponent and
celebrating infinity would be mathematics performing a distress signal.

More generally, a positive-density single component with \(w<-1/3\)
produces acceleration in the flat model. That criterion depends on the
total effective \(\epsilon+3p\) when multiple components are present.

A realistic background calculation combines components with different
scalings. Radiation fades fastest, matter more slowly, and a
cosmological-constant density stays fixed. Different terms can therefore
dominate at different epochs without any of them abruptly changing its
fundamental identity.

\hypertarget{cosmological-redshift-derived-from-neighboring-wave-crests}{%
\section{19.6 Cosmological redshift, derived from neighboring wave
crests}\label{cosmological-redshift-derived-from-neighboring-wave-crests}}

Define a radial comoving distance coordinate along a ray by

\begin{center}
\begingroup
\sbox{\grmathbox}{$\displaystyle 
d\psi=\frac{d\chi}{\sqrt{1-k\chi^2}}.
$}
\typeout{GRDISPLAY|505|\the\wd\grmathbox|\the\linewidth}
\ifdim\wd\grmathbox>\linewidth
\resizebox{\linewidth}{!}{\usebox{\grmathbox}}
\else\usebox{\grmathbox}\fi
\endgroup
\end{center}

Radial null propagation gives \(c\,dt=\pm a(t)d\psi\). For a fixed
comoving source and receiver, the comoving distance traversed by a light
signal is

\begin{center}
\begingroup
\sbox{\grmathbox}{$\displaystyle 
\Delta\psi=\int_{t_{\mathrm{em}}}^{t_{\mathrm{rec}}}\frac{c\,dt}{a(t)}.
$}
\typeout{GRDISPLAY|506|\the\wd\grmathbox|\the\linewidth}
\ifdim\wd\grmathbox>\linewidth
\resizebox{\linewidth}{!}{\usebox{\grmathbox}}
\else\usebox{\grmathbox}\fi
\endgroup
\end{center}

A neighboring wave crest leaves at
\(t_{\mathrm{em}}+\delta t_{\mathrm{em}}\) and arrives at
\(t_{\mathrm{rec}}+\delta t_{\mathrm{rec}}\). It crosses the same
comoving separation. Subtract the two integrals, treating the periods as
short compared with the expansion time:

\begin{center}
\begingroup
\sbox{\grmathbox}{$\displaystyle 
\frac{\delta t_{\mathrm{rec}}}{a(t_{\mathrm{rec}})}
=\frac{\delta t_{\mathrm{em}}}{a(t_{\mathrm{em}})}.
$}
\typeout{GRDISPLAY|507|\the\wd\grmathbox|\the\linewidth}
\ifdim\wd\grmathbox>\linewidth
\resizebox{\linewidth}{!}{\usebox{\grmathbox}}
\else\usebox{\grmathbox}\fi
\endgroup
\end{center}

The comoving clocks measure these coordinate intervals as proper
periods. Frequency is inverse period, so

\begin{center}
\begingroup
\sbox{\grmathbox}{$\displaystyle 
\boxed{1+z=\frac{\nu_{\mathrm{em}}}{\nu_{\mathrm{rec}}}
=\frac{a(t_{\mathrm{rec}})}{a(t_{\mathrm{em}})}.}
$}
\typeout{GRDISPLAY|508|\the\wd\grmathbox|\the\linewidth}
\ifdim\wd\grmathbox>\linewidth
\resizebox{\linewidth}{!}{\usebox{\grmathbox}}
\else\usebox{\grmathbox}\fi
\endgroup
\end{center}

A photon observed at redshift \(z=2\) was emitted when the scale factor
was one third its value at observation. Its observed wavelength is three
times its emitted wavelength, assuming no additional peculiar-motion or
local gravitational shifts.

Photon energy is proportional to frequency, so it scales as \(a^{-1}\).
Combined with number-density dilution \(a^{-3}\), this independently
explains radiation\textquotesingle s \(a^{-4}\) energy-density law. The
geodesic calculation and the fluid conservation calculation agree: two
conceptual roads reach the same equation.

This redshift differs from comparing stationary observers in a static
potential. Generic FLRW spacetime has no corresponding global timelike
Killing symmetry. A useful alternative interpretation builds the
redshift from many small local Doppler shifts between neighboring
comoving observers. What one should not do is pretend all widely
separated cosmological observers share one global special-relativistic
inertial frame.

\hypertarget{conformal-time-changing-the-graph-paper-to-straighten-light-rays}{%
\section{19.7 Conformal time: changing the graph paper to straighten
light
rays}\label{conformal-time-changing-the-graph-paper-to-straighten-light-rays}}

Define dimensionless conformal time by

\begin{center}
\begingroup
\sbox{\grmathbox}{$\displaystyle 
d\eta=\frac{c\,dt}{a(t)}.
$}
\typeout{GRDISPLAY|509|\the\wd\grmathbox|\the\linewidth}
\ifdim\wd\grmathbox>\linewidth
\resizebox{\linewidth}{!}{\usebox{\grmathbox}}
\else\usebox{\grmathbox}\fi
\endgroup
\end{center}

Then

\begin{center}
\begingroup
\sbox{\grmathbox}{$\displaystyle 
ds^2=a^2(\eta)\left[-d\eta^2+d\psi^2+
S_k^2(\psi)d\Omega^2\right],
$}
\typeout{GRDISPLAY|510|\the\wd\grmathbox|\the\linewidth}
\ifdim\wd\grmathbox>\linewidth
\resizebox{\linewidth}{!}{\usebox{\grmathbox}}
\else\usebox{\grmathbox}\fi
\endgroup
\end{center}

where \(S_{+1}(\psi)=\sin\psi\), \(S_0(\psi)=\psi\), and
\(S_{-1}(\psi)=\sinh\psi\) in their appropriate coordinate domains.
Radial light rays now satisfy \(d\psi=\pm d\eta\).

The transformation has made causal bookkeeping much easier. Multiplying
a metric by a positive conformal factor preserves its null cones. It
does not preserve proper times, physical lengths, or affine parameters
of null geodesics. The graph paper can straighten the light rays without
stopping the universe\textquotesingle s clocks and rulers from evolving.

This distinction becomes especially useful for horizons: what matters is
how much conformal time has elapsed or remains, not merely whether
today\textquotesingle s expansion rate sounds large.

\hypertarget{three-distances-called-a-horizon-far-too-casually}{%
\section{19.8 Three distances called a horizon far too
casually}\label{three-distances-called-a-horizon-far-too-casually}}

At fixed cosmic time, radial proper distance from the origin along a
spatial slice is \(D=a(t)\psi\). A comoving object has fixed \(\psi\),
so

\begin{center}
\begingroup
\sbox{\grmathbox}{$\displaystyle 
\dot D=HD.
$}
\typeout{GRDISPLAY|511|\the\wd\grmathbox|\the\linewidth}
\ifdim\wd\grmathbox>\linewidth
\resizebox{\linewidth}{!}{\usebox{\grmathbox}}
\else\usebox{\grmathbox}\fi
\endgroup
\end{center}

This rate can exceed \(c\) at sufficiently large \(D\). It is a rate of
change of a nonlocal separation defined using cosmic simultaneity, not
the velocity measured when one object passes another in the same local
inertial frame. Special relativity\textquotesingle s local causal limit
remains intact.

For a radial light ray,

\begin{center}
\begingroup
\sbox{\grmathbox}{$\displaystyle 
\dot D=HD\pm c.
$}
\typeout{GRDISPLAY|512|\the\wd\grmathbox|\the\linewidth}
\ifdim\wd\grmathbox>\linewidth
\resizebox{\linewidth}{!}{\usebox{\grmathbox}}
\else\usebox{\grmathbox}\fi
\endgroup
\end{center}

An inward-directed ray can initially have increasing proper distance
when \(HD>c\). Whether it later approaches us depends on the subsequent
expansion history. This makes the Hubble radius a useful instantaneous
scale but not generally an event horizon.

\begingroup\small\setlength{\tabcolsep}{5pt}\renewcommand{\arraystretch}{1.13}

\begin{longtable}[]{@{}
  >{\raggedright\arraybackslash}p{(\columnwidth - 4\tabcolsep) * \real{0.2500}}
  >{\raggedright\arraybackslash}p{(\columnwidth - 4\tabcolsep) * \real{0.3600}}
  >{\raggedright\arraybackslash}p{(\columnwidth - 4\tabcolsep) * \real{0.3900}}@{}}
\toprule\noalign{}
\begin{minipage}[b]{\linewidth}\raggedright
Distance at cosmic time \(t\)
\end{minipage} & \begin{minipage}[b]{\linewidth}\raggedright
Formula
\end{minipage} & \begin{minipage}[b]{\linewidth}\raggedright
Question it answers
\end{minipage} \\
\midrule\noalign{}
\endhead
\bottomrule\noalign{}
\endlastfoot
Hubble radius & \(D_H=c/H\) & Where does recession rate \(HD\) equal
\(c\) on this cosmic-time slice? \\
Particle-horizon distance & \(D_p=a(t)\int_{t_B}^{t}c\,dt'/a(t')\) & How
far could light have traveled to us since the model\textquotesingle s
initial boundary? \\
Event-horizon distance &
\(D_e=a(t)\int_t^{t_{\mathrm{max}}}c\,dt'/a(t')\) & Which comoving
sources can ever communicate with us in the modeled future? \\
\end{longtable}

\endgroup

The relevant horizon exists with finite distance only when the
corresponding integral converges, with global topology and the
spacetime\textquotesingle s actual domain also taken into account. The
event-horizon upper limit is the future endpoint, often infinity; its
existence therefore depends on future evolution.

For an ideal flat dust universe \(a\propto t^{2/3}\) beginning at
\(t=0\),

\begin{center}
\begingroup
\sbox{\grmathbox}{$\displaystyle 
D_p=3ct,\qquad D_H=\frac32ct.
$}
\typeout{GRDISPLAY|513|\the\wd\grmathbox|\the\linewidth}
\ifdim\wd\grmathbox>\linewidth
\resizebox{\linewidth}{!}{\usebox{\grmathbox}}
\else\usebox{\grmathbox}\fi
\endgroup
\end{center}

The future event-horizon integral diverges if that evolution continues
forever: there is no cosmological event horizon in this model. Two
distinct present-day distances and one absent horizon have emerged from
a single simple scale factor.

For exponentially expanding flat de Sitter slicing, the future event
horizon is \(D_e=c/H\). Its equality with the Hubble radius is a
property of that special evolution, not a universal identity.

\hypertarget{does-expansion-stretch-your-atoms-and-what-exactly-is-dark-energy}{%
\section{19.9 Does expansion stretch your atoms? And what exactly is
dark
energy?}\label{does-expansion-stretch-your-atoms-and-what-exactly-is-dark-energy}}

FLRW describes a smoothed cosmological background. A bound atom,
planetary system, or galaxy is a local solution with its own stresses
and gravitational field. You cannot obtain its size evolution by
multiplying every internal distance by the background scale factor while
ignoring the forces that bind it.

Cosmological effects can appear as tiny tidal terms in suitable local
approximations. Whether they matter is determined by comparing them with
binding dynamics. ``Everything stretches'' is not the field equation;
neither is ``cosmology can never affect a bound system.'' Specify the
system and compare scales.

Dark matter and dark energy also play different dynamical roles. In the
usual cosmological modeling, cold dark matter behaves approximately as
pressureless matter, contributes to gravitational clustering, and has
background density scaling approximately as \(a^{-3}\). Dark energy
labels the component or effective physics invoked for accelerated
expansion; a cosmological constant is the simplest constant-\(w=-1\)
realization. The equations alone do not identify dark
matter\textquotesingle s particle properties or prove that dark energy
is exactly a cosmological constant.

Historically, distant Type Ia supernova measurements supplied evidence
for accelerated expansion through the relation between observed redshift
and inferred luminosity distance. A primary account is
\href{https://arxiv.org/abs/astro-ph/9805201}{Riess and collaborators,
\emph{Observational Evidence from Supernovae for an Accelerating
Universe and a Cosmological Constant}}. Such an inference connects
calibrated observations to a model for light propagation and cosmic
evolution; it is not a direct photograph of negative pressure.

The modern task is to confront expansion, lensing, clustering, and other
observables together while checking systematics and assumptions. A
successful fit within GR supports that description. It does not
establish that every alternative gravitational theory is mathematically
incapable of producing the same particular observations.

\hypertarget{a-curvature-comparison-that-should-permanently-cure-one-misconception}{%
\section{19.10 A curvature comparison that should permanently cure one
misconception}\label{a-curvature-comparison-that-should-permanently-cure-one-misconception}}

FLRW is conformally flat: its Weyl tensor vanishes. Its curvature is
entirely in the Ricci part. The Schwarzschild vacuum exterior has the
opposite pattern: its Ricci tensor vanishes while its Weyl tensor
carries the tidal field.

There is an additional surprise. A radiation-filled FLRW solution with
\(\Lambda=0\) has \(T=-\epsilon+3p=0\), so the traced Einstein equation
gives \(R=0\). Yet \(R_{\mu\nu}\) is nonzero. Even vanishing scalar
curvature is a remarkably weak statement about the full geometry.

\begingroup\small\setlength{\tabcolsep}{5pt}\renewcommand{\arraystretch}{1.13}

\begin{longtable}[]{@{}
  >{\raggedright\arraybackslash}p{(\columnwidth - 6\tabcolsep) * \real{0.2500}}
  >{\raggedright\arraybackslash}p{(\columnwidth - 6\tabcolsep) * \real{0.2500}}
  >{\raggedright\arraybackslash}p{(\columnwidth - 6\tabcolsep) * \real{0.2500}}
  >{\raggedright\arraybackslash}p{(\columnwidth - 6\tabcolsep) * \real{0.2500}}@{}}
\toprule\noalign{}
\begin{minipage}[b]{\linewidth}\raggedright
Geometry or region
\end{minipage} & \begin{minipage}[b]{\linewidth}\raggedright
Ricci tensor
\end{minipage} & \begin{minipage}[b]{\linewidth}\raggedright
Weyl tensor
\end{minipage} & \begin{minipage}[b]{\linewidth}\raggedright
What this teaches
\end{minipage} \\
\midrule\noalign{}
\endhead
\bottomrule\noalign{}
\endlastfoot
Minkowski spacetime & Zero & Zero & Full spacetime curvature vanishes \\
Schwarzschild vacuum exterior & Zero & Nonzero & Vacuum can contain
tidal curvature \\
Nonempty radiation FLRW, \(\Lambda=0\) & Nonzero, with scalar trace
\(R=0\) & Zero & Even zero scalar curvature need not mean zero Ricci
curvature \\
de Sitter spacetime & \(R_{\mu\nu}=\Lambda g_{\mu\nu}\) & Zero & A
cosmological constant curves spacetime without Weyl tides \\
\end{longtable}

\endgroup

The Einstein equation controls a particular contraction of curvature.
Matter-filled cosmology, vacuum black holes, and vacuum gravitational
waves demonstrate why that distinction matters. Geometry has both
locally sourced structure and dynamical information carried through the
spacetime solution.

Vanishing Weyl curvature does not mean vanishing geodesic deviation. De
Sitter spacetime, for example, has isotropic relative acceleration of
neighboring comoving geodesics. The Ricci-Weyl split distinguishes parts
of the tidal geometry; it does not assign all measurable gravitational
effects to Weyl alone.

\hypertarget{chapter-20}{%
\chapter{20. Making spacetime run: initial data, constraints, and
numerical relativity}\label{chapter-20}}

\hypertarget{an-equation-is-not-yet-a-prediction}{%
\section{20.1 An equation is not yet a
prediction}\label{an-equation-is-not-yet-a-prediction}}

An equation can be magnificent and still fail to tell you what happens
next. Newton\textquotesingle s \(m\ddot x=F\) needs an initial position
and velocity. Maxwell\textquotesingle s equations need an
electromagnetic field whose initial divergence satisfies the charge
constraints. Einstein\textquotesingle s equation also needs initial data
- but now the object being initialized includes the geometry used to
define ``initial.''

Imagine a video editor handed a single frame of a film and asked to
generate the next frame. A photograph alone cannot determine motion. In
GR the spatial geometry is the photograph; a second geometric object,
\textbf{extrinsic curvature}, supplies the relevant velocity
information. Matter fields bring their own initial data. The catch is
that these ingredients cannot be chosen independently: they must already
satisfy four Einstein equations.

For this chapter set \(c=1\), so time and length have the same units.
Keep \(G_N\) explicit.

\hypertarget{slicing-spacetime-without-claiming-that-nature-has-a-preferred-slicer}{%
\section{20.2 Slicing spacetime without claiming that nature has a
preferred
slicer}\label{slicing-spacetime-without-claiming-that-nature-has-a-preferred-slicer}}

Choose spacelike hypersurfaces \(\Sigma_t\) labeled by a coordinate
\(t\). Locally, wherever this foliation is valid, write the metric in
\textbf{3+1 form}:

\begin{center}
\begingroup
\sbox{\grmathbox}{$\displaystyle 
ds^2=-N^2dt^2+\gamma_{ij}(dx^i+\beta^i dt)(dx^j+\beta^jdt).
$}
\typeout{GRDISPLAY|514|\the\wd\grmathbox|\the\linewidth}
\ifdim\wd\grmathbox>\linewidth
\resizebox{\linewidth}{!}{\usebox{\grmathbox}}
\else\usebox{\grmathbox}\fi
\endgroup
\end{center}

Here are the three ingredients.

\begingroup\small\setlength{\tabcolsep}{5pt}\renewcommand{\arraystretch}{1.13}

\begin{longtable}[]{@{}
  >{\raggedright\arraybackslash}p{(\columnwidth - 4\tabcolsep) * \real{0.2500}}
  >{\raggedright\arraybackslash}p{(\columnwidth - 4\tabcolsep) * \real{0.3600}}
  >{\raggedright\arraybackslash}p{(\columnwidth - 4\tabcolsep) * \real{0.3900}}@{}}
\toprule\noalign{}
\begin{minipage}[b]{\linewidth}\raggedright
Object
\end{minipage} & \begin{minipage}[b]{\linewidth}\raggedright
What it describes
\end{minipage} & \begin{minipage}[b]{\linewidth}\raggedright
What it does not mean
\end{minipage} \\
\midrule\noalign{}
\endhead
\bottomrule\noalign{}
\endlastfoot
\(\gamma_{ij}\), the spatial metric & Distances measured within a slice
& The whole spacetime metric \\
\(N>0\), the lapse & Proper time separation between nearby slices along
their unit normals: \(d\tau=Ndt\) & A universal cosmic clock \\
\(\beta^i\), the shift & How the coordinate grid slides sideways between
slices & Matter necessarily moving through space \\
\end{longtable}

\endgroup

Let \(n^\mu\) be the future-directed unit normal, \(n^\mu n_\mu=-1\).
The coordinate time vector decomposes as

\begin{center}
\begingroup
\sbox{\grmathbox}{$\displaystyle 
\partial_t=Nn+\beta^i\partial_i,
\qquad
n^\mu=\left(\frac1N,-\frac{\beta^i}{N}\right).
$}
\typeout{GRDISPLAY|515|\the\wd\grmathbox|\the\linewidth}
\ifdim\wd\grmathbox>\linewidth
\resizebox{\linewidth}{!}{\usebox{\grmathbox}}
\else\usebox{\grmathbox}\fi
\endgroup
\end{center}

The sign of the shift has operational meaning: someone following a
normal worldline has \(dx^i/dt=-\beta^i\). The coordinates drift
relative to that person.

A useful analogy is slicing a loaf while moving the cutting board
sideways. Slice thickness corresponds to lapse; sideways movement
corresponds to shift. Its limitation is decisive: spacetime is
Lorentzian, and these slices are conventions for organizing causal
evolution, not pieces cut from an object sitting in a larger room.

There need not be a convenient global slicing of an arbitrary spacetime.
We will meet the additional causal conditions that support one in
Chapter 22.

\hypertarget{the-velocity-of-space-is-extrinsic-curvature}{%
\section{20.3 The ``velocity of space'' is extrinsic
curvature}\label{the-velocity-of-space-is-extrinsic-curvature}}

Extend the spatial projector to spacetime:

\begin{center}
\begingroup
\sbox{\grmathbox}{$\displaystyle 
\gamma_{\mu\nu}=g_{\mu\nu}+n_\mu n_\nu.
$}
\typeout{GRDISPLAY|516|\the\wd\grmathbox|\the\linewidth}
\ifdim\wd\grmathbox>\linewidth
\resizebox{\linewidth}{!}{\usebox{\grmathbox}}
\else\usebox{\grmathbox}\fi
\endgroup
\end{center}

It removes a vector\textquotesingle s component normal to the slice.
Define extrinsic curvature with the following sign convention:

\begin{center}
\begingroup
\sbox{\grmathbox}{$\displaystyle 
K_{\mu\nu}
=-\gamma_\mu{}^\alpha\gamma_\nu{}^\beta\nabla_\alpha n_\beta.
$}
\typeout{GRDISPLAY|517|\the\wd\grmathbox|\the\linewidth}
\ifdim\wd\grmathbox>\linewidth
\resizebox{\linewidth}{!}{\usebox{\grmathbox}}
\else\usebox{\grmathbox}\fi
\endgroup
\end{center}

The derivative asks how the normals change as we move along the slice.
The projectors retain the part visible within the slice. A plane has
parallel normals; a curved surface generally does not.

For tangent indices this becomes

\begin{center}
\begingroup
\sbox{\grmathbox}{$\displaystyle 
K_{ij}=-\frac12\mathcal L_n\gamma_{ij}
=-\frac1{2N}\left(\partial_t\gamma_{ij}
-D_i\beta_j-D_j\beta_i\right),
$}
\typeout{GRDISPLAY|518|\the\wd\grmathbox|\the\linewidth}
\ifdim\wd\grmathbox>\linewidth
\resizebox{\linewidth}{!}{\usebox{\grmathbox}}
\else\usebox{\grmathbox}\fi
\endgroup
\end{center}

where \(D_i\) is the Levi-Civita derivative of \(\gamma_{ij}\). The Lie
derivative \(\mathcal L\) measures change under the flow of a vector
field. In this equation it subtracts change caused merely by sliding the
coordinates.

Equivalently,

\begin{center}
\begingroup
\sbox{\grmathbox}{$\displaystyle 
\partial_t\gamma_{ij}=-2NK_{ij}+\mathcal L_\beta\gamma_{ij}.
$}
\typeout{GRDISPLAY|519|\the\wd\grmathbox|\the\linewidth}
\ifdim\wd\grmathbox>\linewidth
\resizebox{\linewidth}{!}{\usebox{\grmathbox}}
\else\usebox{\grmathbox}\fi
\endgroup
\end{center}

This is a precise version of ``extrinsic curvature contains the velocity
of the spatial metric.'' The qualification matters: it is velocity
relative to the chosen slicing, after correcting for coordinate drift.

\textbf{Sign trap:} Chapter 14 used the boundary convention
\(K_{\mathrm{boundary}}=h^{\mu\nu}\nabla_\mu n_\nu\). For the same
spacelike hypersurface and the same normal, the present ADM convention
gives \(K=-K_{\mathrm{boundary}}\). Other textbooks also differ in this
choice. Every equation containing an odd number of \(K\) factors must be
translated consistently. A sign difference here is not a disagreement
about expanding universes.

\textbf{Geometry trap:} extrinsic curvature need not indicate spacetime
curvature. Curved slices can be drawn inside flat Minkowski spacetime.
Intrinsic spatial curvature, extrinsic curvature, and four-dimensional
spacetime curvature are related objects, not synonyms.

\hypertarget{four-equations-your-initial-photograph-must-already-obey}{%
\section{20.4 Four equations your initial photograph must already
obey}\label{four-equations-your-initial-photograph-must-already-obey}}

Define the matter energy and momentum seen by the normal observers:

\begin{center}
\begingroup
\sbox{\grmathbox}{$\displaystyle 
E=T_{\mu\nu}n^\mu n^\nu,
\qquad
j_i=-\gamma_i{}^\mu n^\nu T_{\mu\nu}.
$}
\typeout{GRDISPLAY|520|\the\wd\grmathbox|\the\linewidth}
\ifdim\wd\grmathbox>\linewidth
\resizebox{\linewidth}{!}{\usebox{\grmathbox}}
\else\usebox{\grmathbox}\fi
\endgroup
\end{center}

\(E\) is not automatically the fluid\textquotesingle s rest-frame
density \(\epsilon\): a moving fluid has extra energy in the normal
frame. The minus sign in \(j_i\) compensates for the timelike metric
sign so that its ordinary local interpretation is momentum density.

The normal-normal projection of Einstein\textquotesingle s equation is
the \textbf{Hamiltonian constraint}:

\begin{center}
\begingroup
\sbox{\grmathbox}{$\displaystyle 
\boxed{{}^{(3)}R+K^2-K_{ij}K^{ij}=16\pi G_N E+2\Lambda.}
$}
\typeout{GRDISPLAY|521|\the\wd\grmathbox|\the\linewidth}
\ifdim\wd\grmathbox>\linewidth
\resizebox{\linewidth}{!}{\usebox{\grmathbox}}
\else\usebox{\grmathbox}\fi
\endgroup
\end{center}

Here \(K=\gamma^{ij}K_{ij}\) and \({}^{(3)}R\) is the scalar curvature
of the spatial metric. The three mixed normal-spatial projections give
the \textbf{momentum constraints}:

\begin{center}
\begingroup
\sbox{\grmathbox}{$\displaystyle 
\boxed{D_j\left(K^{ij}-\gamma^{ij}K\right)=8\pi G_Nj^i.}
$}
\typeout{GRDISPLAY|522|\the\wd\grmathbox|\the\linewidth}
\ifdim\wd\grmathbox>\linewidth
\resizebox{\linewidth}{!}{\usebox{\grmathbox}}
\else\usebox{\grmathbox}\fi
\endgroup
\end{center}

Why these combinations? The Gauss relation connects curvature measured
entirely within a slice to spacetime curvature plus products of
\(K_{ij}\). After contraction it gives

\begin{center}
\begingroup
\sbox{\grmathbox}{$\displaystyle 
2G_{\mu\nu}n^\mu n^\nu
={}^{(3)}R+K^2-K_{ij}K^{ij}.
$}
\typeout{GRDISPLAY|523|\the\wd\grmathbox|\the\linewidth}
\ifdim\wd\grmathbox>\linewidth
\resizebox{\linewidth}{!}{\usebox{\grmathbox}}
\else\usebox{\grmathbox}\fi
\endgroup
\end{center}

Projecting \(G_{\mu\nu}+\Lambda g_{\mu\nu}=8\pi G_NT_{\mu\nu}\) twice
along \(n\) contributes \(-\Lambda\), because \(g(n,n)=-1\). Move it
across and multiply by two: the \(+2\Lambda\) in the constraint follows.
The Codazzi relation similarly turns the mixed projection into a spatial
divergence of \(K\); our definition of \(j_i\) fixes the momentum
equation\textquotesingle s sign.

These equations contain no second time derivative of the geometry. They
constrain what can consistently exist on one slice. The 3+1 projection
and this extrinsic-curvature convention are developed systematically in
\href{https://arxiv.org/abs/gr-qc/0703035}{Éric
Gourgoulhon\textquotesingle s author-written notes on the 3+1
formalism}.

Suppose you invent a lumpy matter distribution, decree that space is
perfectly Euclidean, and also decree that it is instantaneously
unchanging, \(K_{ij}=0\). With \(\Lambda=0\), the Hamiltonian constraint
says \(E=0\). Your proposed universe has failed its entrance exam.

This is analogous to specifying an electric field with zero divergence
everywhere while also inserting a charge. Evolution cannot repair an
inconsistent starting point without changing the data.

\hypertarget{worked-check-the-friedmann-equation-is-an-initial-data-constraint}{%
\section{20.5 Worked check: the Friedmann equation is an initial-data
constraint}\label{worked-check-the-friedmann-equation-is-an-initial-data-constraint}}

For homogeneous, isotropic slices,

\begin{center}
\begingroup
\sbox{\grmathbox}{$\displaystyle 
\gamma_{ij}=a^2(t)\bar\gamma_{ij},
\qquad N=1,\qquad \beta^i=0,
\qquad {}^{(3)}R=\frac{6k}{a^2}.
$}
\typeout{GRDISPLAY|524|\the\wd\grmathbox|\the\linewidth}
\ifdim\wd\grmathbox>\linewidth
\resizebox{\linewidth}{!}{\usebox{\grmathbox}}
\else\usebox{\grmathbox}\fi
\endgroup
\end{center}

Take normal observers comoving with the fluid, so \(E=\epsilon\). From
the definition,

\begin{center}
\begingroup
\sbox{\grmathbox}{$\displaystyle 
K_{ij}=-\frac12\partial_t(a^2\bar\gamma_{ij})
=-H\gamma_{ij},
\qquad H=\frac{\dot a}{a}.
$}
\typeout{GRDISPLAY|525|\the\wd\grmathbox|\the\linewidth}
\ifdim\wd\grmathbox>\linewidth
\resizebox{\linewidth}{!}{\usebox{\grmathbox}}
\else\usebox{\grmathbox}\fi
\endgroup
\end{center}

Therefore \(K=-3H\), while \(K_{ij}K^{ij}=3H^2\). The difference is
\(9H^2-3H^2=6H^2\). The Hamiltonian constraint becomes

\begin{center}
\begingroup
\sbox{\grmathbox}{$\displaystyle 
\frac{6k}{a^2}+6H^2=16\pi G_N\epsilon+2\Lambda,
$}
\typeout{GRDISPLAY|526|\the\wd\grmathbox|\the\linewidth}
\ifdim\wd\grmathbox>\linewidth
\resizebox{\linewidth}{!}{\usebox{\grmathbox}}
\else\usebox{\grmathbox}\fi
\endgroup
\end{center}

or

\begin{center}
\begingroup
\sbox{\grmathbox}{$\displaystyle 
\boxed{H^2+\frac{k}{a^2}
=\frac{8\pi G_N}{3}\epsilon+\frac\Lambda3.}
$}
\typeout{GRDISPLAY|527|\the\wd\grmathbox|\the\linewidth}
\ifdim\wd\grmathbox>\linewidth
\resizebox{\linewidth}{!}{\usebox{\grmathbox}}
\else\usebox{\grmathbox}\fi
\endgroup
\end{center}

The familiar cosmological equation is the statement that the initial
geometry, expansion rate, and matter density fit together. Homogeneity
made the constraint algebraic. In a binary-black-hole calculation,
solving its spatially varying version is substantial work.

\hypertarget{evolution-gauge-and-why-ten-metric-components-do-not-mean-ten-gravitons}{%
\section{20.6 Evolution, gauge, and why ten metric components do not
mean ten
gravitons}\label{evolution-gauge-and-why-ten-metric-components-do-not-mean-ten-gravitons}}

The spatial projections provide evolution equations. In vacuum with
\(\Lambda=0\), one useful form is

\begin{center}
\begingroup
\sbox{\grmathbox}{$\displaystyle 
(\partial_t-\mathcal L_\beta)K_{ij}
=-D_iD_jN
+N\left({}^{(3)}R_{ij}+KK_{ij}-2K_{ik}K^k{}_j\right).
$}
\typeout{GRDISPLAY|528|\the\wd\grmathbox|\the\linewidth}
\ifdim\wd\grmathbox>\linewidth
\resizebox{\linewidth}{!}{\usebox{\grmathbox}}
\else\usebox{\grmathbox}\fi
\endgroup
\end{center}

Read it together with the equation for \(\partial_t\gamma_{ij}\).
Spatial curvature and extrinsic-curvature products govern the next
change of \(K\); lapse gradients describe the acceleration associated
with the chosen normal observers. Matter adds appropriate spatial stress
and energy terms.

The Einstein-Hilbert action, after separating its boundary contribution,
contains the bulk combination

\begin{center}
\begingroup
\sbox{\grmathbox}{$\displaystyle 
S_{\mathrm{bulk}}=\frac1{16\pi G_N}\int dt\,d^3x\,
N\sqrt\gamma\left({}^{(3)}R+K_{ij}K^{ij}-K^2-2\Lambda\right).
$}
\typeout{GRDISPLAY|529|\the\wd\grmathbox|\the\linewidth}
\ifdim\wd\grmathbox>\linewidth
\resizebox{\linewidth}{!}{\usebox{\grmathbox}}
\else\usebox{\grmathbox}\fi
\endgroup
\end{center}

There are no independent time derivatives of lapse and shift. In the
Hamiltonian formulation they act as multipliers enforcing the
constraints, rather than adding propagating gravitational polarizations.

Now count carefully. The symmetric \(\gamma_{ij}\) has six components.
Its six conjugate momenta make twelve \textbf{phase-space} variables per
spatial point. Four first-class constraints each remove one phase-space
direction by imposing an equation and one by identifying
gauge-equivalent descriptions:

\begin{center}
\begingroup
\sbox{\grmathbox}{$\displaystyle 
12-2\times4=4\quad\text{physical phase-space dimensions}.
$}
\typeout{GRDISPLAY|530|\the\wd\grmathbox|\the\linewidth}
\ifdim\wd\grmathbox>\linewidth
\resizebox{\linewidth}{!}{\usebox{\grmathbox}}
\else\usebox{\grmathbox}\fi
\endgroup
\end{center}

That means \textbf{two configuration degrees of freedom}, each with its
conjugate momentum. In weak gravitational waves, these become the two
familiar polarizations. This is a local count in ordinary
four-dimensional GR; boundaries, topology, special backgrounds, and
matter require additional care.

``Ten minus four equals six'' does not perform this count. It subtracts
coordinate functions while overlooking the constrained dynamical
structure.

\hypertarget{a-stable-spacetime-simulator-needs-a-good-coordinate-policy}{%
\section{20.7 A stable spacetime simulator needs a good coordinate
policy}\label{a-stable-spacetime-simulator-needs-a-good-coordinate-policy}}

Einstein\textquotesingle s equations contain gauge freedom, so their
unreduced component form is not simply ten independent wave equations. A
coordinate condition can expose the wave structure. In harmonic
coordinates, for example,

\begin{center}
\begingroup
\sbox{\grmathbox}{$\displaystyle 
\Box_g x^\mu=0,
$}
\typeout{GRDISPLAY|531|\the\wd\grmathbox|\the\linewidth}
\ifdim\wd\grmathbox>\linewidth
\resizebox{\linewidth}{!}{\usebox{\grmathbox}}
\else\usebox{\grmathbox}\fi
\endgroup
\end{center}

and the principal, highest-derivative part of the reduced metric
equations is schematically

\begin{center}
\begingroup
\sbox{\grmathbox}{$\displaystyle 
g^{\alpha\beta}\partial_\alpha\partial_\beta g_{\mu\nu}
=\text{lower-derivative geometric terms and matter sources}.
$}
\typeout{GRDISPLAY|532|\the\wd\grmathbox|\the\linewidth}
\ifdim\wd\grmathbox>\linewidth
\resizebox{\linewidth}{!}{\usebox{\grmathbox}}
\else\usebox{\grmathbox}\fi
\endgroup
\end{center}

The metric itself supplies the coefficients that determine wave
propagation: the stage also controls its own signal speed. Such a system
is called \textbf{quasilinear}.

A well-posed formulation needs existence, uniqueness in the appropriate
sense, and continuous dependence on initial data. Tiny numerical errors
must not instantly become arbitrary disasters. A formulation can be
mathematically equivalent on exact constraint-satisfying solutions yet
behave very differently when roundoff and discretization introduce small
constraint violations. Generalized harmonic, BSSN, and related
formulations make different choices for variables, gauge evolution, and
controlling those violations.

The contracted Bianchi identity supplies constraint-propagation
relations. With consistent matter evolution, exact constraints that hold
initially continue to hold in a suitable exact evolution. A computer
approximates that theorem; it does not receive a magical exemption from
error analysis.

Initial conditions and boundary conditions play different roles. Initial
data describe a spatial slice. Boundaries of a finite simulation also
require treatment of incoming characteristic fields, gravitational
radiation, and gauge or constraint modes. For an isolated system one
often approximates an asymptotically flat exterior; arbitrary reflective
boundary conditions would instead build a gravitational echo chamber.

Solving an elliptic constraint across a slice does not transmit a
physical signal instantly. It constructs a mutually compatible initial
state. Subsequent physical disturbances propagate according to the
causal equations.

\hypertarget{the-hole-argument-predicting-geometry-without-predicting-coordinate-labels}{%
\section{20.8 The hole argument: predicting geometry without predicting
coordinate
labels}\label{the-hole-argument-predicting-geometry-without-predicting-coordinate-labels}}

Imagine a smooth relabeling of spacetime that is exactly the identity
near the initial slice but changes labels inside a later empty region -
the ``hole.'' Apply it to the metric and all physical fields. General
covariance produces a new coordinate description satisfying the same
initial data. Does that destroy determinism?

Only if you assume that bare manifold points already possess observable
identities independently of every field. The two descriptions preserve
coincidences: where a detector meets a pulse, how much proper time its
clock records, what curvature its instruments measure. They represent
the same physical solution when related by the appropriate gauge
diffeomorphism.

A prediction should concern ``the curvature measured when this clock
reads this value,'' not ``the curvature at a label whose attachment to
any physical event I am free to change.'' Boundary symmetries require
care: transformations acting nontrivially on prescribed asymptotic data
can carry physical charges and are not all disposable gauge.

For suitable constraint-satisfying vacuum data, and for appropriate
well-posed matter systems, the relevant uniqueness statement is
uniqueness of the maximal globally hyperbolic development \textbf{up to
diffeomorphism}. ``Maximal'' does not promise geodesic completeness or a
nonsingular future. This is the landmark result of
\href{https://projecteuclid.org/journals/communications-in-mathematical-physics/volume-14/issue-4/Global-aspects-of-the-Cauchy-problem-in-general-relativity/cmp/1103841822.pdf}{Choquet-Bruhat
and Geroch\textquotesingle s original Cauchy-problem paper}.

GR does predict a spacetime. It declines to endow your
spreadsheet\textquotesingle s row labels with additional physical
reality.

\hypertarget{chapter-21}{%
\chapter{21. A local laboratory at every point: tetrads, forms, and the
gauge viewpoint}\label{chapter-21}}

\hypertarget{taking-a-square-root-of-the-metric}{%
\section{21.1 Taking a square root of the
metric}\label{taking-a-square-root-of-the-metric}}

Coordinate bases are versatile, but they need not look like a
laboratory\textquotesingle s orthogonal ruler-and-clock axes. In
spherical coordinates, a change of one radian is not a change of one
meter. For a local experiment we often want an orthonormal basis
instead.

Introduce four one-forms

\begin{center}
\begingroup
\sbox{\grmathbox}{$\displaystyle 
e^a=e^a{}_\mu dx^\mu,
$}
\typeout{GRDISPLAY|533|\the\wd\grmathbox|\the\linewidth}
\ifdim\wd\grmathbox>\linewidth
\resizebox{\linewidth}{!}{\usebox{\grmathbox}}
\else\usebox{\grmathbox}\fi
\endgroup
\end{center}

such that

\begin{center}
\begingroup
\sbox{\grmathbox}{$\displaystyle 
\boxed{g_{\mu\nu}=\eta_{ab}e^a{}_\mu e^b{}_\nu,\qquad
\eta_{ab}=\operatorname{diag}(-1,1,1,1).}
$}
\typeout{GRDISPLAY|534|\the\wd\grmathbox|\the\linewidth}
\ifdim\wd\grmathbox>\linewidth
\resizebox{\linewidth}{!}{\usebox{\grmathbox}}
\else\usebox{\grmathbox}\fi
\endgroup
\end{center}

In this chapter only, \(a,b,c,d=0,1,2,3\) label \textbf{internal
orthonormal-frame directions}, not the spatial indices \(i,j,k\) used
elsewhere. Greek indices still label spacetime coordinate components.

The \(e^a{}_\mu\) are called a \textbf{tetrad}, \textbf{vierbein}, or
\textbf{coframe}. Its inverse \(e_a{}^\mu\) defines basis vectors
\(e_a=e_a{}^\mu\partial_\mu\). Their duality says
\(e^a(e_b)=\delta^a_b\). A vector\textquotesingle s laboratory
components are

\begin{center}
\begingroup
\sbox{\grmathbox}{$\displaystyle 
V^a=e^a{}_\mu V^\mu.
$}
\typeout{GRDISPLAY|535|\the\wd\grmathbox|\the\linewidth}
\ifdim\wd\grmathbox>\linewidth
\resizebox{\linewidth}{!}{\usebox{\grmathbox}}
\else\usebox{\grmathbox}\fi
\endgroup
\end{center}

The metric is a kind of matrix square built from \(e\), but not an
ordinary unique positive square root. Six choices remain free: at every
point we may rotate the three laboratory axes and boost the
laboratory\textquotesingle s time axis.

More precisely, if \(L^a{}_b(x)\) obeys \(L^T\eta L=\eta\), then

\begin{center}
\begingroup
\sbox{\grmathbox}{$\displaystyle 
e'^a=L^a{}_b e^b
$}
\typeout{GRDISPLAY|536|\the\wd\grmathbox|\the\linewidth}
\ifdim\wd\grmathbox>\linewidth
\resizebox{\linewidth}{!}{\usebox{\grmathbox}}
\else\usebox{\grmathbox}\fi
\endgroup
\end{center}

produces the same metric. Sixteen tetrad components minus six local
Lorentz freedoms leave the metric\textquotesingle s ten components.
Coordinate freedom remains as well; this count describes the additional
frame representation, not the physical degrees of freedom counted in
Chapter 20.

\textbf{Deep gotcha:} writing the metric as \(\eta_{ab}\) in an
orthonormal frame does not flatten spacetime. The frame changes from
point to point, and its comparison law contains the geometry. Writing
every bank balance in dollars does not make all accounts contain the
same amount of money.

\hypertarget{differential-forms-antisymmetry-earns-its-keep}{%
\section{21.2 Differential forms: antisymmetry earns its
keep}\label{differential-forms-antisymmetry-earns-its-keep}}

A one-form takes one vector and returns a number. A two-form takes two
vectors and returns an antisymmetric number, naturally measuring
oriented area. A \(p\)-form generalizes this to \(p\) directions.

The \textbf{wedge product} antisymmetrizes:

\begin{center}
\begingroup
\sbox{\grmathbox}{$\displaystyle 
dx\wedge dy=-dy\wedge dx,
\qquad dx\wedge dx=0.
$}
\typeout{GRDISPLAY|537|\the\wd\grmathbox|\the\linewidth}
\ifdim\wd\grmathbox>\linewidth
\resizebox{\linewidth}{!}{\usebox{\grmathbox}}
\else\usebox{\grmathbox}\fi
\endgroup
\end{center}

For one-forms \(\alpha\) and \(\beta\),

\begin{center}
\begingroup
\sbox{\grmathbox}{$\displaystyle 
(\alpha\wedge\beta)(V,W)
=\alpha(V)\beta(W)-\alpha(W)\beta(V).
$}
\typeout{GRDISPLAY|538|\the\wd\grmathbox|\the\linewidth}
\ifdim\wd\grmathbox>\linewidth
\resizebox{\linewidth}{!}{\usebox{\grmathbox}}
\else\usebox{\grmathbox}\fi
\endgroup
\end{center}

This is the determinant of a two-by-two array. It measures the signed
parallelogram area seen by the two measuring devices. Exchange the sides
and orientation reverses; use the same side twice and the area vanishes.

For a \(p\)-form and a \(q\)-form,

\begin{center}
\begingroup
\sbox{\grmathbox}{$\displaystyle 
\alpha\wedge\beta=(-1)^{pq}\beta\wedge\alpha.
$}
\typeout{GRDISPLAY|539|\the\wd\grmathbox|\the\linewidth}
\ifdim\wd\grmathbox>\linewidth
\resizebox{\linewidth}{!}{\usebox{\grmathbox}}
\else\usebox{\grmathbox}\fi
\endgroup
\end{center}

Thus two two-forms commute under the wedge product, while two one-forms
anticommute. The sign follows from moving \(p\) directions past \(q\)
directions: \(pq\) exchanges.

The \textbf{exterior derivative} raises the form degree by one. For a
scalar,

\begin{center}
\begingroup
\sbox{\grmathbox}{$\displaystyle 
df=\partial_\mu f\,dx^\mu.
$}
\typeout{GRDISPLAY|540|\the\wd\grmathbox|\the\linewidth}
\ifdim\wd\grmathbox>\linewidth
\resizebox{\linewidth}{!}{\usebox{\grmathbox}}
\else\usebox{\grmathbox}\fi
\endgroup
\end{center}

For a one-form \(A=A_\mu dx^\mu\),

\begin{center}
\begingroup
\sbox{\grmathbox}{$\displaystyle 
dA=\frac12(\partial_\mu A_\nu-\partial_\nu A_\mu)
\,dx^\mu\wedge dx^\nu.
$}
\typeout{GRDISPLAY|541|\the\wd\grmathbox|\the\linewidth}
\ifdim\wd\grmathbox>\linewidth
\resizebox{\linewidth}{!}{\usebox{\grmathbox}}
\else\usebox{\grmathbox}\fi
\endgroup
\end{center}

The factor \(1/2\) prevents counting each antisymmetric pair twice. For
example, if \(A=x^2dy\), then \(dA=2x\,dx\wedge dy\).

Two rules do most of the work:

\begin{center}
\begingroup
\sbox{\grmathbox}{$\displaystyle 
d^2=0,
\qquad
d(\alpha\wedge\beta)=d\alpha\wedge\beta+(-1)^p\alpha\wedge d\beta.
$}
\typeout{GRDISPLAY|542|\the\wd\grmathbox|\the\linewidth}
\ifdim\wd\grmathbox>\linewidth
\resizebox{\linewidth}{!}{\usebox{\grmathbox}}
\else\usebox{\grmathbox}\fi
\endgroup
\end{center}

Why does \(d^2f=0\)? The second partial derivatives are symmetric under
index exchange, while the wedge basis is antisymmetric. Their
contraction cancels. This is the common engine behind ``the curl of a
gradient vanishes'' and ``the divergence of a curl vanishes.''

The exterior derivative needs no metric or connection. It does not tell
you the full directional variation of an arbitrary tensor. It extracts a
particular antisymmetric derivative of differential forms. It therefore
complements \(\nabla\); it does not replace it.

Forms also make Stokes\textquotesingle{} theorem compact:

\begin{center}
\begingroup
\sbox{\grmathbox}{$\displaystyle 
\int_\Omega d\alpha=\int_{\partial\Omega}\alpha.
$}
\typeout{GRDISPLAY|543|\the\wd\grmathbox|\the\linewidth}
\ifdim\wd\grmathbox>\linewidth
\resizebox{\linewidth}{!}{\usebox{\grmathbox}}
\else\usebox{\grmathbox}\fi
\endgroup
\end{center}

The boundary integral and interior derivative are two descriptions of
the same accumulated oriented change. Appropriate orientations and
smoothness are part of the statement.

\hypertarget{why-an-orthonormal-frame-needs-a-spin-connection}{%
\section{21.3 Why an orthonormal frame needs a spin
connection}\label{why-an-orthonormal-frame-needs-a-spin-connection}}

Two nearby laboratories may choose differently rotated or boosted axes.
Differentiating their component lists naively confuses physical change
with a change of reference frame. That is precisely the problem that
connections solve.

Introduce the matrix of one-forms

\begin{center}
\begingroup
\sbox{\grmathbox}{$\displaystyle 
\omega^a{}_b=\omega^a{}_{b\mu}dx^\mu,
$}
\typeout{GRDISPLAY|544|\the\wd\grmathbox|\the\linewidth}
\ifdim\wd\grmathbox>\linewidth
\resizebox{\linewidth}{!}{\usebox{\grmathbox}}
\else\usebox{\grmathbox}\fi
\endgroup
\end{center}

and define, for internal vector components,

\begin{center}
\begingroup
\sbox{\grmathbox}{$\displaystyle 
D_\mu V^a=\partial_\mu V^a+\omega^a{}_{b\mu}V^b.
$}
\typeout{GRDISPLAY|545|\the\wd\grmathbox|\the\linewidth}
\ifdim\wd\grmathbox>\linewidth
\resizebox{\linewidth}{!}{\usebox{\grmathbox}}
\else\usebox{\grmathbox}\fi
\endgroup
\end{center}

This is the \textbf{spin connection}, despite its usefulness even before
spinors appear. Metric compatibility in the orthonormal frame implies

\begin{center}
\begingroup
\sbox{\grmathbox}{$\displaystyle 
\omega_{ab}=-\omega_{ba},
\qquad \omega_{ab}=\eta_{ac}\omega^c{}_b.
$}
\typeout{GRDISPLAY|546|\the\wd\grmathbox|\the\linewidth}
\ifdim\wd\grmathbox>\linewidth
\resizebox{\linewidth}{!}{\usebox{\grmathbox}}
\else\usebox{\grmathbox}\fi
\endgroup
\end{center}

Antisymmetry applies after lowering the first frame index. A boost
component can therefore satisfy \(\omega^0{}_1=\omega^1{}_0\) with both
indices in the displayed mixed positions. The minus sign hiding in
\(\eta_{00}\) is doing real work.

The relationship to Christoffel symbols follows by demanding that
conversion between coordinate and frame components commute with
differentiation:

\begin{center}
\begingroup
\sbox{\grmathbox}{$\displaystyle 
\boxed{
\partial_\mu e^a{}_\nu
-\Gamma^\rho{}_{\mu\nu}e^a{}_\rho
+\omega^a{}_{b\mu}e^b{}_\nu=0.
}
$}
\typeout{GRDISPLAY|547|\the\wd\grmathbox|\the\linewidth}
\ifdim\wd\grmathbox>\linewidth
\resizebox{\linewidth}{!}{\usebox{\grmathbox}}
\else\usebox{\grmathbox}\fi
\endgroup
\end{center}

This is the \textbf{tetrad postulate}. It is a compatibility relation
among two representations of the same connection and the map between
their bases. Once the tetrad and Levi-Civita connection are specified,
it determines the associated spin connection.

Under \(V'=LV\), require \(DV\) to transform as \(D'V'=L(DV)\).
Expanding \(d(LV)\) produces an unwanted \((dL)V\) term. The connection
cancels it precisely if

\begin{center}
\begingroup
\sbox{\grmathbox}{$\displaystyle 
\boxed{\omega'=L\omega L^{-1}-(dL)L^{-1}.}
$}
\typeout{GRDISPLAY|548|\the\wd\grmathbox|\the\linewidth}
\ifdim\wd\grmathbox>\linewidth
\resizebox{\linewidth}{!}{\usebox{\grmathbox}}
\else\usebox{\grmathbox}\fi
\endgroup
\end{center}

This is why a connection transforms inhomogeneously. The extra term is
the bookkeeping fee for changing frames differently at different places.

\hypertarget{cartans-equations-package-geometry-into-two-lines}{%
\section{21.4 Cartan\textquotesingle s equations package geometry into
two lines}\label{cartans-equations-package-geometry-into-two-lines}}

The torsion two-form is

\begin{center}
\begingroup
\sbox{\grmathbox}{$\displaystyle 
\mathcal T^a=de^a+\omega^a{}_b\wedge e^b.
$}
\typeout{GRDISPLAY|549|\the\wd\grmathbox|\the\linewidth}
\ifdim\wd\grmathbox>\linewidth
\resizebox{\linewidth}{!}{\usebox{\grmathbox}}
\else\usebox{\grmathbox}\fi
\endgroup
\end{center}

For the Levi-Civita connection of ordinary GR, \(\mathcal T^a=0\), so

\begin{center}
\begingroup
\sbox{\grmathbox}{$\displaystyle 
\boxed{de^a+\omega^a{}_b\wedge e^b=0.}
$}
\typeout{GRDISPLAY|550|\the\wd\grmathbox|\the\linewidth}
\ifdim\wd\grmathbox>\linewidth
\resizebox{\linewidth}{!}{\usebox{\grmathbox}}
\else\usebox{\grmathbox}\fi
\endgroup
\end{center}

This is the torsion-free first Cartan structure equation. Together with
metric compatibility it determines \(\omega\) from the tetrad. It is
often a much quicker calculation than a page of Christoffel symbols.

The curvature two-form is

\begin{center}
\begingroup
\sbox{\grmathbox}{$\displaystyle 
\boxed{\mathcal R^a{}_b=d\omega^a{}_b
+\omega^a{}_c\wedge\omega^c{}_b.}
$}
\typeout{GRDISPLAY|551|\the\wd\grmathbox|\the\linewidth}
\ifdim\wd\grmathbox>\linewidth
\resizebox{\linewidth}{!}{\usebox{\grmathbox}}
\else\usebox{\grmathbox}\fi
\endgroup
\end{center}

This is the second structure equation. The matrix product includes a sum
over \(c\) and a wedge product of its one-form entries. Curvature has
components

\begin{center}
\begingroup
\sbox{\grmathbox}{$\displaystyle 
\mathcal R^a{}_b=\frac12R^a{}_{b\mu\nu}
\,dx^\mu\wedge dx^\nu.
$}
\typeout{GRDISPLAY|552|\the\wd\grmathbox|\the\linewidth}
\ifdim\wd\grmathbox>\linewidth
\resizebox{\linewidth}{!}{\usebox{\grmathbox}}
\else\usebox{\grmathbox}\fi
\endgroup
\end{center}

With our Riemann convention these agree with the curvature defined by
\([\nabla_\mu,\nabla_\nu]\). Applying \(D=d+\omega\) twice to a
vector-valued zero-form gives \(D^2V=\mathcal R V\): the leftover is
exactly the failure of parallel transport around an infinitesimal loop
to return a vector unchanged.

The connection\textquotesingle s inhomogeneous transformation disappears
from curvature:

\begin{center}
\begingroup
\sbox{\grmathbox}{$\displaystyle 
\mathcal R'=L\mathcal R L^{-1}.
$}
\typeout{GRDISPLAY|553|\the\wd\grmathbox|\the\linewidth}
\ifdim\wd\grmathbox>\linewidth
\resizebox{\linewidth}{!}{\usebox{\grmathbox}}
\else\usebox{\grmathbox}\fi
\endgroup
\end{center}

This does not mean its component matrix is invariant. It means curvature
transforms tensorially, so observers can compare the same geometric
object in their different frames. For a detailed frame-based
development, see
\href{https://davidtong.org/teaching/general-relativity/grhtml/S3}{David
Tong\textquotesingle s author-written chapter on connections and Cartan
geometry}.

\hypertarget{worked-example-a-rotating-frame-can-have-connection-without-curvature}{%
\section{21.5 Worked example: a rotating frame can have connection
without
curvature}\label{worked-example-a-rotating-frame-can-have-connection-without-curvature}}

Consider the Euclidean plane, away from the origin, in polar
coordinates:

\begin{center}
\begingroup
\sbox{\grmathbox}{$\displaystyle 
dl^2=dr^2+r^2d\phi^2,
\qquad e^1=dr,\quad e^2=r\,d\phi.
$}
\typeout{GRDISPLAY|554|\the\wd\grmathbox|\the\linewidth}
\ifdim\wd\grmathbox>\linewidth
\resizebox{\linewidth}{!}{\usebox{\grmathbox}}
\else\usebox{\grmathbox}\fi
\endgroup
\end{center}

This two-dimensional example uses only spatial frame indices. Compute

\begin{center}
\begingroup
\sbox{\grmathbox}{$\displaystyle 
de^1=0,
\qquad de^2=dr\wedge d\phi.
$}
\typeout{GRDISPLAY|555|\the\wd\grmathbox|\the\linewidth}
\ifdim\wd\grmathbox>\linewidth
\resizebox{\linewidth}{!}{\usebox{\grmathbox}}
\else\usebox{\grmathbox}\fi
\endgroup
\end{center}

Choose

\begin{center}
\begingroup
\sbox{\grmathbox}{$\displaystyle 
\omega^2{}_1=d\phi,
\qquad\omega^1{}_2=-d\phi.
$}
\typeout{GRDISPLAY|556|\the\wd\grmathbox|\the\linewidth}
\ifdim\wd\grmathbox>\linewidth
\resizebox{\linewidth}{!}{\usebox{\grmathbox}}
\else\usebox{\grmathbox}\fi
\endgroup
\end{center}

The second first-structure equation reads

\begin{center}
\begingroup
\sbox{\grmathbox}{$\displaystyle 
de^2+\omega^2{}_1\wedge e^1
=dr\wedge d\phi+d\phi\wedge dr=0.
$}
\typeout{GRDISPLAY|557|\the\wd\grmathbox|\the\linewidth}
\ifdim\wd\grmathbox>\linewidth
\resizebox{\linewidth}{!}{\usebox{\grmathbox}}
\else\usebox{\grmathbox}\fi
\endgroup
\end{center}

The first equation vanishes too, because \(d\phi\wedge d\phi=0\). The
nonzero connection records how the radial and angular unit vectors
rotate as \(\phi\) changes.

But

\begin{center}
\begingroup
\sbox{\grmathbox}{$\displaystyle 
\mathcal R^1{}_2=d(-d\phi)=0
$}
\typeout{GRDISPLAY|558|\the\wd\grmathbox|\the\linewidth}
\ifdim\wd\grmathbox>\linewidth
\resizebox{\linewidth}{!}{\usebox{\grmathbox}}
\else\usebox{\grmathbox}\fi
\endgroup
\end{center}

on a regular angular coordinate patch; the matrix wedge contribution
also vanishes here. The plane is flat. The polar frame is undefined at
the origin, but a bad frame does not create a physical curvature
singularity there.

Now replace the plane by a sphere of radius \(a\):

\begin{center}
\begingroup
\sbox{\grmathbox}{$\displaystyle 
e^1=a\,d\theta,
\qquad e^2=a\sin\theta\,d\phi.
$}
\typeout{GRDISPLAY|559|\the\wd\grmathbox|\the\linewidth}
\ifdim\wd\grmathbox>\linewidth
\resizebox{\linewidth}{!}{\usebox{\grmathbox}}
\else\usebox{\grmathbox}\fi
\endgroup
\end{center}

Then \(de^2=a\cos\theta\,d\theta\wedge d\phi\). The same cancellation
method gives

\begin{center}
\begingroup
\sbox{\grmathbox}{$\displaystyle 
\omega^2{}_1=\cos\theta\,d\phi,
\qquad\omega^1{}_2=-\cos\theta\,d\phi.
$}
\typeout{GRDISPLAY|560|\the\wd\grmathbox|\the\linewidth}
\ifdim\wd\grmathbox>\linewidth
\resizebox{\linewidth}{!}{\usebox{\grmathbox}}
\else\usebox{\grmathbox}\fi
\endgroup
\end{center}

Now the connection\textquotesingle s coefficient varies with latitude:

\begin{center}
\begingroup
\sbox{\grmathbox}{$\displaystyle 
\mathcal R^1{}_2
=d(-\cos\theta\,d\phi)
=\sin\theta\,d\theta\wedge d\phi
=\frac1{a^2}e^1\wedge e^2.
$}
\typeout{GRDISPLAY|561|\the\wd\grmathbox|\the\linewidth}
\ifdim\wd\grmathbox>\linewidth
\resizebox{\linewidth}{!}{\usebox{\grmathbox}}
\else\usebox{\grmathbox}\fi
\endgroup
\end{center}

Its Gaussian curvature is \(1/a^2\), and its scalar curvature is
\(2/a^2\). The same two-line machinery has distinguished a rotating
choice of axes from actual curved geometry.

That distinction is one of the central survival skills of GR: a
connection can look busy while spacetime is doing nothing
gravitationally interesting.

\hypertarget{what-gravity-shares-with-gauge-theory---and-what-it-adds}{%
\section{21.6 What gravity shares with gauge theory - and what it
adds}\label{what-gravity-shares-with-gauge-theory---and-what-it-adds}}

Electromagnetism uses a potential one-form \(A\) and field strength
\(F=dA\). A non-Abelian gauge theory uses a matrix-valued connection
with schematic curvature \(F=dA+A\wedge A\). The spin connection obeys
the same geometric pattern.

The shared idea is a freedom to choose a local reference convention,
accompanied by a connection that compares neighboring conventions. It is
not an assertion that gravity is ordinary electromagnetism with a larger
alphabet.

In GR, the tetrad ties the internal Lorentz frame to actual tangent
directions: it connects the gauge description to rods, clocks, causal
cones, and volume. The Einstein-Hilbert action is linear in curvature,
while the usual Yang-Mills action is quadratic in its field strength.
Their symmetry, variables, and dynamics therefore differ in crucial
ways.

Tetrads also let us couple spin-\(1/2\) matter to gravity. Spinors
transform under the spin representation associated with the local
Lorentz group; they are not ordinary spacetime vectors with an unusual
number of entries. Locally, choose gamma matrices satisfying

\begin{center}
\begingroup
\sbox{\grmathbox}{$\displaystyle 
\{\gamma^a,\gamma^b\}=2\eta^{ab}I.
$}
\typeout{GRDISPLAY|562|\the\wd\grmathbox|\the\linewidth}
\ifdim\wd\grmathbox>\linewidth
\resizebox{\linewidth}{!}{\usebox{\grmathbox}}
\else\usebox{\grmathbox}\fi
\endgroup
\end{center}

A compatible spinor derivative is

\begin{center}
\begingroup
\sbox{\grmathbox}{$\displaystyle 
D_\mu\psi=\partial_\mu\psi
+\frac14\omega_{ab\mu}\gamma^a\gamma^b\psi.
$}
\typeout{GRDISPLAY|563|\the\wd\grmathbox|\the\linewidth}
\ifdim\wd\grmathbox>\linewidth
\resizebox{\linewidth}{!}{\usebox{\grmathbox}}
\else\usebox{\grmathbox}\fi
\endgroup
\end{center}

The gamma matrices supply the Lorentz generators in the spinor
representation. The connection is still solving the same problem: how to
compare field components defined using different local frames. A global
spinor theory also requires suitable global topology - a spin structure
- not merely a locally clever notation.

\hypertarget{torsion-nonmetricity-and-the-limits-of-the-palatini-shortcut}{%
\section{21.7 Torsion, nonmetricity, and the limits of the Palatini
shortcut}\label{torsion-nonmetricity-and-the-limits-of-the-palatini-shortcut}}

Three different geometric properties deserve three different names:

\begingroup\small\setlength{\tabcolsep}{5pt}\renewcommand{\arraystretch}{1.13}

\begin{longtable}[]{@{}
  >{\raggedright\arraybackslash}p{(\columnwidth - 4\tabcolsep) * \real{0.2500}}
  >{\raggedright\arraybackslash}p{(\columnwidth - 4\tabcolsep) * \real{0.3600}}
  >{\raggedright\arraybackslash}p{(\columnwidth - 4\tabcolsep) * \real{0.3900}}@{}}
\toprule\noalign{}
\begin{minipage}[b]{\linewidth}\raggedright
Property
\end{minipage} & \begin{minipage}[b]{\linewidth}\raggedright
Representative definition
\end{minipage} & \begin{minipage}[b]{\linewidth}\raggedright
What it measures
\end{minipage} \\
\midrule\noalign{}
\endhead
\bottomrule\noalign{}
\endlastfoot
Curvature & \(R^\rho{}_{\sigma\mu\nu}\) & Failure of infinitesimal
parallel transport around loops to agree \\
Torsion &
\(T^\rho{}_{\mu\nu}=\Gamma^\rho{}_{\mu\nu}-\Gamma^\rho{}_{\nu\mu}\) in a
coordinate basis & Antisymmetric part of the connection; covariantly,
\(T(X,Y)=\nabla_XY-\nabla_YX-[X,Y]\) \\
Nonmetricity & \(Q_{\rho\mu\nu}=-\nabla_\rho g_{\mu\nu}\), with this
chosen sign & Failure of the connection to preserve the metric under
parallel transport \\
\end{longtable}

\endgroup

Standard GR uses a torsion-free, metric-compatible connection, while
allowing curvature. More general theories can change these assumptions.
Calling a curved spacetime ``twisted'' in ordinary speech does not imply
nonzero mathematical torsion.

In the \textbf{Palatini approach}, vary the metric and connection
independently in an Einstein-Hilbert-type action. Under the usual
assumptions - four dimensions, a nondegenerate metric, a torsion-free
independent connection, and matter action independent of that connection
- the connection equation enforces the Levi-Civita connection.
Substituting it back recovers metric GR.

Why does this work? The action\textquotesingle s curvature is linear in
derivatives of the connection. Integrating by parts transfers those
derivatives onto \(\sqrt{-g}g^{\mu\nu}\). The resulting equation
requires compatibility of the connection with that metric density; in
dimensions above two, under these assumptions, it reduces to metric
compatibility. Torsion freedom then selects the unique Levi-Civita
connection.

There are important exceptions to the slogan ``independent connection
variation always gives GR.'' Allowing completely general connections
introduces projective gauge subtleties. Connection-dependent matter
changes the connection equation; spinor matter can source torsion in
Einstein-Cartan formulations. Replacing \(R\) by a nonlinear function
\(f(R)\) generally makes metric and Palatini variation different
theories. The careful equivalence statement is analyzed in
\href{https://arxiv.org/abs/1010.0869}{Dadhich and
Pons\textquotesingle s paper on Einstein-Hilbert and Einstein-Palatini
formulations}.

Using forms, a corresponding first-order gravitational action can be
written, with \(c=1\) and a consistently chosen orientation,

\begin{center}
\begingroup
\sbox{\grmathbox}{$\displaystyle 
S=\frac1{32\pi G_N}\int
\varepsilon_{abcd}\,e^a\wedge e^b\wedge
\left(\mathcal R^{cd}-\frac\Lambda6e^c\wedge e^d\right).
$}
\typeout{GRDISPLAY|564|\the\wd\grmathbox|\the\linewidth}
\ifdim\wd\grmathbox>\linewidth
\resizebox{\linewidth}{!}{\usebox{\grmathbox}}
\else\usebox{\grmathbox}\fi
\endgroup
\end{center}

Here \(\varepsilon_{0123}=+1\) is the internal alternating symbol and
\(\mathcal R^{cd}=\eta^{de}\mathcal R^c{}_e\). For an invertible tetrad,
a Lorentz-compatible independent connection, and no torsion-sourcing
matter, varying the connection imposes zero torsion; varying the tetrad
gives Einstein\textquotesingle s equation. The apparent change of
language has exposed a new organization of the same dynamics.

The payoff is larger than elegant notation. We now understand a metric
formulation, a moving-laboratory formulation, and a gauge-connection
formulation as different ways of asking the same questions about
comparison, motion, and curvature.

\hypertarget{chapter-22}{%
\chapter{22. When geodesics crowd together: focusing, singularities, and
black-hole thermodynamics}\label{chapter-22}}

\hypertarget{from-one-falling-observer-to-a-cloud-of-them}{%
\section{22.1 From one falling observer to a cloud of
them}\label{from-one-falling-observer-to-a-cloud-of-them}}

A geodesic describes one freely falling observer. Geodesic deviation
describes the changing separation of nearby observers. A
\textbf{congruence} is a smooth family of worldlines filling a region
without crossing there. Think of an enormous cloud of tiny spacecraft,
each with its engines off, carrying rulers to monitor its neighbors.

The cloud can expand, distort, and rotate. Those are different effects.
A sphere becoming a larger sphere expands. A sphere becoming a
same-volume ellipsoid shears. An ellipsoid turning without changing
shape rotates.

Use \(c=1\) through the focusing discussion and let the congruence have
unit tangent \(u^\mu\), with \(u^\mu u_\mu=-1\). Its local rest-space
metric is

\begin{center}
\begingroup
\sbox{\grmathbox}{$\displaystyle 
h_{\mu\nu}=g_{\mu\nu}+u_\mu u_\nu.
$}
\typeout{GRDISPLAY|565|\the\wd\grmathbox|\the\linewidth}
\ifdim\wd\grmathbox>\linewidth
\resizebox{\linewidth}{!}{\usebox{\grmathbox}}
\else\usebox{\grmathbox}\fi
\endgroup
\end{center}

Define the projected velocity-gradient tensor

\begin{center}
\begingroup
\sbox{\grmathbox}{$\displaystyle 
B_{\mu\nu}=h_\mu{}^\alpha h_\nu{}^\beta\nabla_\beta u_\alpha.
$}
\typeout{GRDISPLAY|566|\the\wd\grmathbox|\the\linewidth}
\ifdim\wd\grmathbox>\linewidth
\resizebox{\linewidth}{!}{\usebox{\grmathbox}}
\else\usebox{\grmathbox}\fi
\endgroup
\end{center}

It tells neighboring observers how their relative velocity depends, to
first order, on their spatial separation. Split this three-dimensional
linear map into its trace, symmetric traceless part, and antisymmetric
part:

\begin{center}
\begingroup
\sbox{\grmathbox}{$\displaystyle 
\boxed{B_{\mu\nu}=\frac13\theta h_{\mu\nu}
+\sigma_{\mu\nu}+\omega_{\mu\nu}.}
$}
\typeout{GRDISPLAY|567|\the\wd\grmathbox|\the\linewidth}
\ifdim\wd\grmathbox>\linewidth
\resizebox{\linewidth}{!}{\usebox{\grmathbox}}
\else\usebox{\grmathbox}\fi
\endgroup
\end{center}

Explicitly,

\begin{center}
\begingroup
\sbox{\grmathbox}{$\displaystyle 
\theta=\nabla_\mu u^\mu,
\qquad
\sigma_{\mu\nu}=B_{(\mu\nu)}-\frac13\theta h_{\mu\nu},
\qquad
\omega_{\mu\nu}=B_{[\mu\nu]}.
$}
\typeout{GRDISPLAY|568|\the\wd\grmathbox|\the\linewidth}
\ifdim\wd\grmathbox>\linewidth
\resizebox{\linewidth}{!}{\usebox{\grmathbox}}
\else\usebox{\grmathbox}\fi
\endgroup
\end{center}

The equality between the projected trace and \(\nabla_\mu u^\mu\)
follows from differentiating the fixed normalization of \(u\).

\(\theta\) is \textbf{expansion}, \(\sigma\) is \textbf{shear}, and
\(\omega\) is \textbf{vorticity} or twist. For an infinitesimal comoving
volume \(\mathcal V\),

\begin{center}
\begingroup
\sbox{\grmathbox}{$\displaystyle 
\theta=\frac1{\mathcal V}\frac{d\mathcal V}{d\tau}
=\frac{d}{d\tau}\ln\mathcal V.
$}
\typeout{GRDISPLAY|569|\the\wd\grmathbox|\the\linewidth}
\ifdim\wd\grmathbox>\linewidth
\resizebox{\linewidth}{!}{\usebox{\grmathbox}}
\else\usebox{\grmathbox}\fi
\endgroup
\end{center}

The trace is the fractional volume-growth rate because the first-order
fractional change of a determinant is the trace of the underlying linear
deformation. That is the same determinant identity that appeared when
varying \(\sqrt{-g}\).

There is also a direct bridge to geodesic deviation. For a commuting
family of geodesics with separation vector \(\xi\), \([u,\xi]=0\)
implies

\begin{center}
\begingroup
\sbox{\grmathbox}{$\displaystyle 
\nabla_u\xi=\nabla_\xi u.
$}
\typeout{GRDISPLAY|570|\the\wd\grmathbox|\the\linewidth}
\ifdim\wd\grmathbox>\linewidth
\resizebox{\linewidth}{!}{\usebox{\grmathbox}}
\else\usebox{\grmathbox}\fi
\endgroup
\end{center}

In a transported rest frame, this is ``separation velocity equals \(B\)
times separation.'' Differentiating again and using the
geodesic-deviation equation gives a matrix evolution law with a term
\(-B^2\) and a curvature term. Raychaudhuri\textquotesingle s equation
is its trace.

\hypertarget{deriving-the-raychaudhuri-equation-one-logical-move-at-a-time}{%
\section{22.2 Deriving the Raychaudhuri equation, one logical move at a
time}\label{deriving-the-raychaudhuri-equation-one-logical-move-at-a-time}}

Start with the expansion and differentiate along the flow:

\begin{center}
\begingroup
\sbox{\grmathbox}{$\displaystyle 
\frac{d\theta}{d\tau}
=u^\rho\nabla_\rho(\nabla_\mu u^\mu).
$}
\typeout{GRDISPLAY|571|\the\wd\grmathbox|\the\linewidth}
\ifdim\wd\grmathbox>\linewidth
\resizebox{\linewidth}{!}{\usebox{\grmathbox}}
\else\usebox{\grmathbox}\fi
\endgroup
\end{center}

Commute the derivatives, using the Riemann convention of this book. The
contraction contributes a minus Ricci term. Then apply the product rule:

\begin{center}
\begingroup
\sbox{\grmathbox}{$\displaystyle 
\frac{d\theta}{d\tau}
=\nabla_\mu(u^\rho\nabla_\rho u^\mu)
-(\nabla_\mu u^\rho)(\nabla_\rho u^\mu)
-R_{\mu\nu}u^\mu u^\nu.
$}
\typeout{GRDISPLAY|572|\the\wd\grmathbox|\the\linewidth}
\ifdim\wd\grmathbox>\linewidth
\resizebox{\linewidth}{!}{\usebox{\grmathbox}}
\else\usebox{\grmathbox}\fi
\endgroup
\end{center}

The first term is \(\nabla_\mu a^\mu\), where \(a^\mu=\nabla_u u^\mu\).
It vanishes for a geodesic congruence. The second is the trace of the
square of the velocity-gradient map. For geodesic flow its rest-space
decomposition gives

\begin{center}
\begingroup
\sbox{\grmathbox}{$\displaystyle 
\operatorname{tr}(B^2)
=\frac13\theta^2+\sigma_{\mu\nu}\sigma^{\mu\nu}
-\omega_{\mu\nu}\omega^{\mu\nu}.
$}
\typeout{GRDISPLAY|573|\the\wd\grmathbox|\the\linewidth}
\ifdim\wd\grmathbox>\linewidth
\resizebox{\linewidth}{!}{\usebox{\grmathbox}}
\else\usebox{\grmathbox}\fi
\endgroup
\end{center}

Why the vorticity minus sign? A real antisymmetric matrix squares to a
matrix with nonpositive trace. In two dimensions, the rotation generator

\begin{center}
\begingroup
\sbox{\grmathbox}{$\displaystyle 
\begin{pmatrix}0&-w\\w&0\end{pmatrix}
$}
\typeout{GRDISPLAY|574|\the\wd\grmathbox|\the\linewidth}
\ifdim\wd\grmathbox>\linewidth
\resizebox{\linewidth}{!}{\usebox{\grmathbox}}
\else\usebox{\grmathbox}\fi
\endgroup
\end{center}

squares to \(-w^2I\). The trace-free symmetric matrix has a positive sum
of squared eigenvalues. Cross terms vanish because a symmetric tensor
contracted with an antisymmetric tensor is zero, and shear has zero
trace.

Therefore, for a timelike geodesic congruence in four spacetime
dimensions,

\begin{center}
\begingroup
\sbox{\grmathbox}{$\displaystyle 
\boxed{
\frac{d\theta}{d\tau}
=-\frac13\theta^2
-\sigma_{\mu\nu}\sigma^{\mu\nu}
+\omega_{\mu\nu}\omega^{\mu\nu}
-R_{\mu\nu}u^\mu u^\nu.
}
$}
\typeout{GRDISPLAY|575|\the\wd\grmathbox|\the\linewidth}
\ifdim\wd\grmathbox>\linewidth
\resizebox{\linewidth}{!}{\usebox{\grmathbox}}
\else\usebox{\grmathbox}\fi
\endgroup
\end{center}

For a normalized accelerated congruence, add \(+\nabla_\mu a^\mu\) to
this expression, with the same projected definitions. Rockets can alter
the cloud\textquotesingle s behavior; geodesic focusing theorems do not
silently include thrust.

Every term has a story:

\begin{itemize}
\tightlist
\item
  \(-\theta^2/3\): convergence can reinforce itself even without
  curvature.
\item
  \(-\sigma^2\): stretching along some directions can accelerate volume
  focusing.
\item
  \(+\omega^2\): rotation opposes the simple focusing tendency.
\item
  \(-R_{\mu\nu}u^\mu u^\nu\): curvature directly changes the trace of
  relative acceleration.
\end{itemize}

The Weyl tensor is absent from the explicit last term, but it can
generate shear, which then affects expansion through \(-\sigma^2\).
Vacuum curvature can matter enormously even when \(R_{\mu\nu}=0\).

\hypertarget{how-an-inequality-becomes-a-finite-time-prediction}{%
\section{22.3 How an inequality becomes a finite-time
prediction}\label{how-an-inequality-becomes-a-finite-time-prediction}}

Suppose the congruence is geodesic and hypersurface-orthogonal, so
\(\omega_{\mu\nu}=0\). The equivalence between vanishing twist and local
orthogonality to hypersurfaces is an application of the Frobenius
integrability theorem. It is a condition on the flow, not a claim that
every family of geodesics has zero vorticity.

Also suppose the \textbf{timelike convergence condition} holds:

\begin{center}
\begingroup
\sbox{\grmathbox}{$\displaystyle 
R_{\mu\nu}u^\mu u^\nu\ge0.
$}
\typeout{GRDISPLAY|576|\the\wd\grmathbox|\the\linewidth}
\ifdim\wd\grmathbox>\linewidth
\resizebox{\linewidth}{!}{\usebox{\grmathbox}}
\else\usebox{\grmathbox}\fi
\endgroup
\end{center}

Raychaudhuri then implies

\begin{center}
\begingroup
\sbox{\grmathbox}{$\displaystyle 
\frac{d\theta}{d\tau}\le-\frac13\theta^2.
$}
\typeout{GRDISPLAY|577|\the\wd\grmathbox|\the\linewidth}
\ifdim\wd\grmathbox>\linewidth
\resizebox{\linewidth}{!}{\usebox{\grmathbox}}
\else\usebox{\grmathbox}\fi
\endgroup
\end{center}

If the cloud starts converging, \(\theta_0<0\), divide by \(\theta^2>0\)
and differentiate its reciprocal:

\begin{center}
\begingroup
\sbox{\grmathbox}{$\displaystyle 
\frac{d}{d\tau}\left(\frac1\theta\right)\ge\frac13.
$}
\typeout{GRDISPLAY|578|\the\wd\grmathbox|\the\linewidth}
\ifdim\wd\grmathbox>\linewidth
\resizebox{\linewidth}{!}{\usebox{\grmathbox}}
\else\usebox{\grmathbox}\fi
\endgroup
\end{center}

If the geodesics extend through the required interval, this inequality
forces the initially negative reciprocal toward zero within proper time
no larger than \(3/|\theta_0|\): \(\theta\) becomes unboundedly
negative, and the smooth congruence develops a focal point or
\textbf{caustic} by that bound. An earlier incomplete geodesic endpoint
is another possibility. The focusing argument does not by itself
guarantee that the geodesics exist for the entire interval.

\textbf{This is not yet a spacetime singularity.} Aim a family of
straight worldlines at the same event in Minkowski spacetime. Their
congruence focuses; each worldline continues perfectly well. What breaks
is the single-valued smooth velocity field used to describe that
overlapping family.

The jump from a local focusing theorem to a global singularity theorem
requires additional causal and topological arguments. Skipping that jump
is like proving that two highways intersect and declaring that the Earth
ends there.

\hypertarget{energy-conditions-are-assumptions-with-jobs-to-do}{%
\section{22.4 Energy conditions are assumptions with jobs to
do}\label{energy-conditions-are-assumptions-with-jobs-to-do}}

Einstein\textquotesingle s equation converts curvature conditions into
matter conditions, but it does not itself require ordinary matter to
satisfy those conditions. They are additional hypotheses.

For a perfect fluid, useful pointwise conditions are:

\begingroup\small\setlength{\tabcolsep}{5pt}\renewcommand{\arraystretch}{1.13}

\begin{longtable}[]{@{}
  >{\raggedright\arraybackslash}p{(\columnwidth - 4\tabcolsep) * \real{0.2500}}
  >{\raggedright\arraybackslash}p{(\columnwidth - 4\tabcolsep) * \real{0.3600}}
  >{\raggedright\arraybackslash}p{(\columnwidth - 4\tabcolsep) * \real{0.3900}}@{}}
\toprule\noalign{}
\begin{minipage}[b]{\linewidth}\raggedright
Condition
\end{minipage} & \begin{minipage}[b]{\linewidth}\raggedright
General idea
\end{minipage} & \begin{minipage}[b]{\linewidth}\raggedright
Perfect-fluid inequalities
\end{minipage} \\
\midrule\noalign{}
\endhead
\bottomrule\noalign{}
\endlastfoot
Null energy condition, NEC & \(T_{\mu\nu}k^\mu k^\nu\ge0\) for every
null \(k\) & \(\epsilon+p\ge0\) \\
Weak energy condition, WEC & Every timelike observer measures
nonnegative local energy density & \(\epsilon\ge0\),
\(\epsilon+p\ge0\) \\
Dominant energy condition, DEC & Energy density is nonnegative and its
flux is causal & \(\epsilon\ge\lvert p\rvert\) \\
Strong energy condition, SEC &
\((T_{\mu\nu}-\tfrac12Tg_{\mu\nu})v^\mu v^\nu\ge0\) for all timelike
\(v\) & \(\epsilon+p\ge0\), \(\epsilon+3p\ge0\) \\
\end{longtable}

\endgroup

For \(\Lambda=0\), the SEC implies timelike convergence. If \(\Lambda\)
remains on the geometric side, however,

\begin{center}
\begingroup
\sbox{\grmathbox}{$\displaystyle 
R_{\mu\nu}u^\mu u^\nu
=8\pi G_N\left(T_{\mu\nu}u^\mu u^\nu+\frac12T\right)-\Lambda.
$}
\typeout{GRDISPLAY|579|\the\wd\grmathbox|\the\linewidth}
\ifdim\wd\grmathbox>\linewidth
\resizebox{\linewidth}{!}{\usebox{\grmathbox}}
\else\usebox{\grmathbox}\fi
\endgroup
\end{center}

For a comoving perfect-fluid observer this is
\(4\pi G_N(\epsilon+3p)-\Lambda\). Positive \(\Lambda\) can defeat
timelike focusing. Equivalently, move it into an effective vacuum stress
tensor with \(p_\Lambda=-\epsilon_\Lambda\): it violates the SEC when
its density is positive, while saturating the NEC.

For null vectors, the trace and cosmological terms vanish because
\(g_{\mu\nu}k^\mu k^\nu=0\). The NEC therefore implies null convergence
in GR even with \(\Lambda\).

Classical scalar potentials can violate the SEC, and quantum fields can
violate classical pointwise energy conditions more broadly. These are
reasons to inspect a theorem\textquotesingle s hypotheses carefully, not
to call the theorem mistaken.

\hypertarget{trapped-surfaces-and-what-penrose-actually-proved}{%
\section{22.5 Trapped surfaces and what Penrose actually
proved}\label{trapped-surfaces-and-what-penrose-actually-proved}}

For an affinely parametrized null geodesic congruence, the transverse
screen has two dimensions. Its Raychaudhuri equation has the
corresponding coefficient:

\begin{center}
\begingroup
\sbox{\grmathbox}{$\displaystyle 
\frac{d\theta}{d\lambda}
=-\frac12\theta^2-\sigma_{\mu\nu}\sigma^{\mu\nu}
+\omega_{\mu\nu}\omega^{\mu\nu}
-R_{\mu\nu}k^\mu k^\nu.
$}
\typeout{GRDISPLAY|580|\the\wd\grmathbox|\the\linewidth}
\ifdim\wd\grmathbox>\linewidth
\resizebox{\linewidth}{!}{\usebox{\grmathbox}}
\else\usebox{\grmathbox}\fi
\endgroup
\end{center}

The shear and twist here live on the positive-definite two-dimensional
screen transverse to the rays. An affine parameter is essential; a
nonaffine parameter introduces an additional term proportional to
\(\theta\).

Take a closed spacelike two-surface and send future light rays
orthogonally away from it in both null-normal directions. For an
ordinary sphere in flat space, the outward bundle grows in area and the
inward bundle shrinks. A \textbf{future trapped surface} has negative
expansion in both directions. Even the outward-directed light bundle
initially loses cross-sectional area.

This is a local geometric condition on the surface and its null normals.
It does not say that a photon locally travels more slowly than light,
nor does it require a coordinate speed to become negative.

One standard form of Penrose\textquotesingle s theorem says that a
sufficiently regular spacetime with a noncompact Cauchy hypersurface,
null convergence, and a closed future trapped surface must be future
null geodesically incomplete. The theorem combines focusing with global
causal geometry; it does not assume spherical symmetry. The original
result is
\href{https://link.aps.org/doi/10.1103/PhysRevLett.14.57}{Penrose\textquotesingle s
1965 paper, ``Gravitational Collapse and Space-Time Singularities''}.

\textbf{Geodesic incompleteness} means at least one inextendible
geodesic has finite affine length in the relevant direction; for
timelike geodesics, proper time is the physical parameter. It does not
universally mean that a curvature scalar tends to infinity. The theorem
does not supply a location, a topology, or a detailed microscopic
description of ``the singularity.''

Even incompleteness must be interpreted carefully: deleting one point
from otherwise regular Minkowski spacetime creates incomplete geodesics
artificially. One must consider extendibility and which spacetime has
actually been specified. Conversely, a coordinate singularity that
disappears in a larger smooth chart is not evidence that physics has
ended.

The genuinely unsettling message is that broad geometric conditions can
force classical evolution to confront an endpoint without our being able
to remove it by assuming that realistic matter is a little less
symmetric.

\hypertarget{horizons-are-about-causal-access-predictability-needs-another-definition}{%
\section{22.6 Horizons are about causal access; predictability needs
another
definition}\label{horizons-are-about-causal-access-predictability-needs-another-definition}}

In an asymptotically flat spacetime, the black-hole region consists of
events unable to send a causal signal to future null infinity
\(\mathscr I^+\), the ideal destination of escaping light. Its boundary
is the \textbf{event horizon}:

\begin{center}
\begingroup
\sbox{\grmathbox}{$\displaystyle 
\mathcal H^+=\partial J^-(\mathscr I^+).
$}
\typeout{GRDISPLAY|581|\the\wd\grmathbox|\the\linewidth}
\ifdim\wd\grmathbox>\linewidth
\resizebox{\linewidth}{!}{\usebox{\grmathbox}}
\else\usebox{\grmathbox}\fi
\endgroup
\end{center}

\(J^-(\mathscr I^+)\) denotes the causal past of that destination. This
definition depends on the entire future spacetime. An event horizon is
not generally locatable using measurements made at one point and one
instant.

A marginally outer trapped surface instead has vanishing outward null
expansion, usually with negative inward expansion in the black-hole
setting. Under suitable conditions an apparent horizon is the outer
boundary of the trapped region on a chosen slice. These are valuable
tools for dynamical calculations, but they depend on slicing and need
not coincide with the event horizon.

For predictability, define the future \textbf{domain of dependence}
\(D^+(\Sigma)\): an event belongs to it if every past-inextendible
causal curve through that event intersects \(\Sigma\). No causal
influence can arrive there without passing through the supplied initial
data.

A \textbf{Cauchy surface} meets every inextendible causal curve exactly
once. A spacetime admitting such a surface is globally hyperbolic, in
the usual boundary-free setting. This is the natural arena for the
initial-value description of Chapter 20.

A \textbf{Cauchy horizon} bounds a domain of dependence. Beyond it,
initial data on the chosen surface no longer control every possible
incoming influence. It is different from an event horizon: crossing a
black-hole event horizon need not destroy local predictability. Inner
horizons in idealized charged or rotating solutions motivate the
question of whether such extendible boundaries survive generic
perturbations. Strong cosmic censorship studies this issue, and its
precise claims depend on the matter model and the allowed regularity of
extensions.

A final geometric tool helps organize these questions. Under a smooth
positive conformal rescaling \(g_{\mu\nu}\mapsto\Omega^2g_{\mu\nu}\),
null cones are unchanged. Proper times and distances change. Conformal
diagrams exploit this separation to compress enormous regions while
preserving causal relations; they are maps of who can signal whom, not
faithful scale drawings.

\hypertarget{a-horizon-gets-a-temperature}{%
\section{22.7 A horizon gets a
temperature}\label{a-horizon-gets-a-temperature}}

Restore \(c\) and \(\hbar\). For a Schwarzschild black hole,

\begin{center}
\begingroup
\sbox{\grmathbox}{$\displaystyle 
r_s=\frac{2G_NM}{c^2},
\qquad A=4\pi r_s^2.
$}
\typeout{GRDISPLAY|582|\the\wd\grmathbox|\the\linewidth}
\ifdim\wd\grmathbox>\linewidth
\resizebox{\linewidth}{!}{\usebox{\grmathbox}}
\else\usebox{\grmathbox}\fi
\endgroup
\end{center}

A short derivation reveals why a temperature appears. Near the horizon
write \(r=r_s+x\), with \(x\ll r_s\), so \(1-r_s/r\approx x/r_s\).
Continue the stationary time coordinate to imaginary time
\(t=-i\tau_E\). The near-horizon radial-time metric becomes

\begin{center}
\begingroup
\sbox{\grmathbox}{$\displaystyle 
ds_E^2\approx\frac{x}{r_s}c^2d\tau_E^2+\frac{r_s}{x}dx^2.
$}
\typeout{GRDISPLAY|583|\the\wd\grmathbox|\the\linewidth}
\ifdim\wd\grmathbox>\linewidth
\resizebox{\linewidth}{!}{\usebox{\grmathbox}}
\else\usebox{\grmathbox}\fi
\endgroup
\end{center}

Introduce a radial proper-distance coordinate \(\varrho=2\sqrt{r_sx}\).
Direct substitution gives

\begin{center}
\begingroup
\sbox{\grmathbox}{$\displaystyle 
ds_E^2\approx d\varrho^2
+\varrho^2\left(\frac{c\,d\tau_E}{2r_s}\right)^2.
$}
\typeout{GRDISPLAY|584|\the\wd\grmathbox|\the\linewidth}
\ifdim\wd\grmathbox>\linewidth
\resizebox{\linewidth}{!}{\usebox{\grmathbox}}
\else\usebox{\grmathbox}\fi
\endgroup
\end{center}

This is a flat plane in polar coordinates. Its angular coordinate is
\(c\tau_E/(2r_s)\). Smoothness at the origin requires that angle to have
period \(2\pi\), so imaginary time has period

\begin{center}
\begingroup
\sbox{\grmathbox}{$\displaystyle 
\Delta\tau_E=\frac{4\pi r_s}{c}
=\frac{8\pi G_NM}{c^3}.
$}
\typeout{GRDISPLAY|585|\the\wd\grmathbox|\the\linewidth}
\ifdim\wd\grmathbox>\linewidth
\resizebox{\linewidth}{!}{\usebox{\grmathbox}}
\else\usebox{\grmathbox}\fi
\endgroup
\end{center}

Quantum statistical mechanics identifies thermal equilibrium with
imaginary-time period \(\hbar/(k_BT)\). Equating these periods yields

\begin{center}
\begingroup
\sbox{\grmathbox}{$\displaystyle 
\boxed{T_H=\frac{\hbar c^3}{8\pi G_NMk_B}.}
$}
\typeout{GRDISPLAY|586|\the\wd\grmathbox|\the\linewidth}
\ifdim\wd\grmathbox>\linewidth
\resizebox{\linewidth}{!}{\usebox{\grmathbox}}
\else\usebox{\grmathbox}\fi
\endgroup
\end{center}

This Euclidean argument characterizes the regular stationary thermal
construction. Deriving the outgoing flux in a collapse spacetime
involves a quantum-field calculation with the appropriate state and
boundary conditions; it yields the same Hawking temperature at infinity,
with frequency-dependent transmission factors modifying an ideal
blackbody spectrum. The calculation is semiclassical, not a complete
quantization of spacetime. See
\href{https://arxiv.org/abs/gr-qc/9912119}{Wald\textquotesingle s
research review of black-hole thermodynamics}.

The commonly illustrated story of a particle pair appearing exactly on
the horizon is a heuristic, not this derivation. Hawking radiation
depends on quantum field modes, the state, and the global relation
between early and late notions of positive frequency. It is not
adequately explained by assigning ordinary local particle trajectories
to vacuum fluctuations.

\hypertarget{entropy-is-written-in-area}{%
\section{22.8 Entropy is written in
area}\label{entropy-is-written-in-area}}

For the nonrotating, uncharged family, the first law is

\begin{center}
\begingroup
\sbox{\grmathbox}{$\displaystyle 
d(Mc^2)=T_H\,dS_{\mathrm{BH}}.
$}
\typeout{GRDISPLAY|587|\the\wd\grmathbox|\the\linewidth}
\ifdim\wd\grmathbox>\linewidth
\resizebox{\linewidth}{!}{\usebox{\grmathbox}}
\else\usebox{\grmathbox}\fi
\endgroup
\end{center}

Substitute \(T_H\) and solve for the entropy change:

\begin{center}
\begingroup
\sbox{\grmathbox}{$\displaystyle 
dS_{\mathrm{BH}}=\frac{8\pi k_BG_NM}{\hbar c}\,dM.
$}
\typeout{GRDISPLAY|588|\the\wd\grmathbox|\the\linewidth}
\ifdim\wd\grmathbox>\linewidth
\resizebox{\linewidth}{!}{\usebox{\grmathbox}}
\else\usebox{\grmathbox}\fi
\endgroup
\end{center}

Integrating gives, with the conventional additive normalization,

\begin{center}
\begingroup
\sbox{\grmathbox}{$\displaystyle 
S_{\mathrm{BH}}=\frac{4\pi k_BG_NM^2}{\hbar c}
=\boxed{\frac{k_BAc^3}{4\hbar G_N}
=\frac{k_BA}{4\ell_P^2}},
\qquad
\ell_P=\sqrt{\frac{\hbar G_N}{c^3}}.
$}
\typeout{GRDISPLAY|589|\the\wd\grmathbox|\the\linewidth}
\ifdim\wd\grmathbox>\linewidth
\resizebox{\linewidth}{!}{\usebox{\grmathbox}}
\else\usebox{\grmathbox}\fi
\endgroup
\end{center}

Entropy proportional to area is startling because ordinary extensive
systems usually organize entropy by volume. This result says the
gravitational problem has a different information accounting. It does
not, by itself, prove that spacetime is made from literal square pixels.

Since \(T_H\propto M^{-1}\), a Schwarzschild black hole becomes hotter
as it loses mass. Its heat capacity is negative. Ordinary
canonical-ensemble intuition - put it in a bath and expect a stable
equilibrium - therefore needs caution.

The classical horizon-area theorem requires the relevant convergence and
global regularity assumptions. Hawking evaporation does not contradict
it: the quantum stress tensor need not satisfy the classical energy
hypothesis, and the horizon area can decrease. The thermodynamic
quantity then involves generalized entropy,

\begin{center}
\begingroup
\sbox{\grmathbox}{$\displaystyle 
S_{\mathrm{gen}}=\frac{k_BA}{4\ell_P^2}+S_{\mathrm{outside}},
$}
\typeout{GRDISPLAY|590|\the\wd\grmathbox|\the\linewidth}
\ifdim\wd\grmathbox>\linewidth
\resizebox{\linewidth}{!}{\usebox{\grmathbox}}
\else\usebox{\grmathbox}\fi
\endgroup
\end{center}

with the quantum-field entropy and gravitational parameters treated
consistently under renormalization. The generalized second law has
substantial support and proofs in specified settings; it should not be
promoted without qualifications to a theorem covering every unknown
quantum-gravitational process.

Curvature, causality, quantum theory, and entropy now meet in one
calculation. That intersection is a clue about the depth of gravity,
even though it is not yet a finished microscopic explanation.

\hypertarget{chapter-23}{%
\chapter{23. Einstein\textquotesingle s equation as a low-energy
masterpiece: effective theory and the frontier}\label{chapter-23}}

\hypertarget{a-theory-can-be-incomplete-without-being-unreliable}{%
\section{23.1 A theory can be incomplete without being
unreliable}\label{a-theory-can-be-incomplete-without-being-unreliable}}

The phrase ``we need quantum gravity'' can create the impression that
quantum mechanics and GR cannot be used together at all. That impression
is false. There is a systematic, predictive framework for quantum
gravitational effects at sufficiently low energies: \textbf{effective
field theory}, or EFT.

An effective theory describes the degrees of freedom accessible at a
chosen resolution. It represents unresolved shorter-distance physics
through coefficients multiplying allowed local interactions. It does not
require us to know every microscopic detail before predicting a
long-wavelength experiment.

Think about sound in a solid. At long wavelengths, elasticity uses
displacement, density, and a few elastic constants. Its predictions can
be excellent without tracking every electron. If you demand wavelengths
comparable to atomic separations, the continuum expansion loses its
organizing advantage. Gravity\textquotesingle s microscopic completion
need not resemble a crystal; the analogy concerns separation of scales,
not a claim that spacetime is an atomic material.

The quantum EFT treatment separates calculable long-distance effects
from unknown short-distance coefficients. This is the central result of
\href{https://arxiv.org/abs/gr-qc/9405057}{Donoghue\textquotesingle s
original work on general relativity as an effective field theory}.

\hypertarget{why-unresolved-heavy-physics-becomes-derivatives}{%
\section{23.2 Why unresolved heavy physics becomes
derivatives}\label{why-unresolved-heavy-physics-becomes-derivatives}}

Use natural units \(c=\hbar=1\) in this subsection. A simple algebraic
model explains the expansion. Suppose a heavy field \(X\) responds to a
source \(J\) through

\begin{center}
\begingroup
\sbox{\grmathbox}{$\displaystyle 
(M_*^2-\Box)X=J.
$}
\typeout{GRDISPLAY|591|\the\wd\grmathbox|\the\linewidth}
\ifdim\wd\grmathbox>\linewidth
\resizebox{\linewidth}{!}{\usebox{\grmathbox}}
\else\usebox{\grmathbox}\fi
\endgroup
\end{center}

Formally,

\begin{center}
\begingroup
\sbox{\grmathbox}{$\displaystyle 
X=\frac1{M_*^2-\Box}J
=\frac1{M_*^2}\left(1+\frac\Box{M_*^2}
+\frac{\Box^2}{M_*^4}+\cdots\right)J.
$}
\typeout{GRDISPLAY|592|\the\wd\grmathbox|\the\linewidth}
\ifdim\wd\grmathbox>\linewidth
\resizebox{\linewidth}{!}{\usebox{\grmathbox}}
\else\usebox{\grmathbox}\fi
\endgroup
\end{center}

This is the geometric series \(1/(1-z)=1+z+z^2+\cdots\), now applied to
a differential operator. It is useful when the source varies on scales
for which the relevant derivatives are small compared with \(M_*^2\). At
higher frequencies the expansion fails, and the heavy
field\textquotesingle s independent dynamics must be restored.

This toy response suppresses boundary-condition and propagator details
to isolate the derivative expansion. Actual quantum matching also
involves loops, symmetries, and the available light fields.

For a generally covariant metric theory, local gravitational
interactions must be scalar combinations of curvature and covariant
derivatives, integrated with the invariant volume. In these units,
define the \textbf{reduced Planck mass} by

\begin{center}
\begingroup
\sbox{\grmathbox}{$\displaystyle 
M_{\mathrm{Pl}}^2=\frac1{8\pi G_N}.
$}
\typeout{GRDISPLAY|593|\the\wd\grmathbox|\the\linewidth}
\ifdim\wd\grmathbox>\linewidth
\resizebox{\linewidth}{!}{\usebox{\grmathbox}}
\else\usebox{\grmathbox}\fi
\endgroup
\end{center}

A schematic local EFT action is

\begin{center}
\begingroup
\sbox{\grmathbox}{$\displaystyle 
S_{\mathrm{local}}=\int d^4x\sqrt{-g}\left[
\frac{M_{\mathrm{Pl}}^2}{2}(R-2\Lambda)
+a_1R^2+a_2R_{\mu\nu}R^{\mu\nu}
+a_3R_{\mu\nu\rho\sigma}R^{\mu\nu\rho\sigma}
+\sum_i\frac{b_i}{M_*^2}\mathcal O_i^{(6)}+\cdots
\right]+S_{\mathrm{light}}.
$}
\typeout{GRDISPLAY|594|\the\wd\grmathbox|\the\linewidth}
\ifdim\wd\grmathbox>\linewidth
\resizebox{\linewidth}{!}{\usebox{\grmathbox}}
\else\usebox{\grmathbox}\fi
\endgroup
\end{center}

The symbols \(\mathcal O_i^{(6)}\) denote local scalar operators of mass
dimension six, such as suitable cubic-curvature contractions.
\(S_{\mathrm{light}}\) contains the light matter fields retained
explicitly.

Check every dimension. For this dimensional count, choose length-valued
local coordinates and a dimensionless metric. A coordinate then has mass
dimension \(-1\), a derivative has dimension \(+1\), and curvature has
dimension \(+2\). The measure \(d^4x\) has dimension \(-4\). Thus
\(M_{\mathrm{Pl}}^2R\), \(R^2\), and \(\mathcal O^{(6)}/M_*^2\) all have
dimension \(+4\), as required for a dimensionless action. The \(a_i\)
and \(b_i\) are dimensionless in this notation. \(M_*\) is a
heavy-physics or cutoff scale; it need not equal \(M_{\mathrm{Pl}}\).

Generic derivative power counting compares curvature-squared terms with
the Einstein term at relative order
\(a_i\mathcal R_*/M_{\mathrm{Pl}}^2\), where \(\mathcal R_*\) denotes a
characteristic magnitude of curvature components in a physically
specified orthonormal frame, with mass dimension two. This estimate
organizes possible corrections; specific operators can vanish on
particular backgrounds. In four dimensions, constant-coefficient local
curvature-squared terms produce no bulk correction when evaluated on a
Ricci-flat vacuum solution: variations of \(R^2\) and
\(R_{\mu\nu}R^{\mu\nu}\) vanish there, and the remaining quadratic
contraction is related to them by the Gauss-Bonnet combination below.
Nonzero Weyl curvature still matters for higher operators and for EFT
validity. The actual coefficients determine the suppression scale. ``Low
energy'' is a quantitative hierarchy, not a promise that every
coefficient is conveniently small.

The displayed curvature-squared basis is intentionally redundant. In
four dimensions the constant-coefficient Gauss-Bonnet combination

\begin{center}
\begingroup
\sbox{\grmathbox}{$\displaystyle 
R_{\mu\nu\rho\sigma}R^{\mu\nu\rho\sigma}
-4R_{\mu\nu}R^{\mu\nu}+R^2
$}
\typeout{GRDISPLAY|595|\the\wd\grmathbox|\the\linewidth}
\ifdim\wd\grmathbox>\linewidth
\resizebox{\linewidth}{!}{\usebox{\grmathbox}}
\else\usebox{\grmathbox}\fi
\endgroup
\end{center}

does not change local bulk equations under the appropriate variational
boundary conditions. Field redefinitions can remove additional redundant
operators in specified calculations, sometimes moving their effects into
matter interactions. Counting written terms is not the same as counting
measurable new parameters.

Massless quantum fields also produce nonlocal contributions,
schematically involving expressions such as \(R\log(-\Box/\mu^2)R\).
They cannot all be hidden in a finite list of local constants: massless
particles propagate over long distances. Here \(\mu\) is a
renormalization scale, with corresponding coefficient dependence
arranged so physical predictions do not depend on this arbitrary
bookkeeping choice.

\hypertarget{nonrenormalizable-does-not-mean-nonpredictive}{%
\section{23.3 Nonrenormalizable does not mean
nonpredictive}\label{nonrenormalizable-does-not-mean-nonpredictive}}

A perturbatively renormalizable theory can absorb ultraviolet
divergences into a fixed finite set of couplings at all orders. Einstein
gravity, treated as a quantum theory about a suitable background,
requires successively higher-order operators. It is not perturbatively
renormalizable in that narrow sense.

EFT asks a different question: \textbf{how many parameters contribute at
the accuracy of this experiment?} At a fixed order in the low-energy
expansion, only finitely many operators contribute. Determine their
coefficients by measurement or matching to a more microscopic theory,
and the remaining predictions at that order follow.

It resembles approximating a smooth function by a Taylor series. An
arbitrary function contains infinitely many coefficients, but a
controlled second-order approximation does not require knowing the
coefficient of \(x^{47}\). The crucial requirement is a valid small
expansion parameter and an honest estimate of neglected terms.

Quantum gravitational loop corrections often carry powers of energy
divided by the Planck scale, along with loop factors. In a long-distance
problem, a characteristic quantum ratio is

\begin{center}
\begingroup
\sbox{\grmathbox}{$\displaystyle 
\frac{\ell_P^2}{L^2}
=\frac{\hbar G_N}{c^3L^2}.
$}
\typeout{GRDISPLAY|596|\the\wd\grmathbox|\the\linewidth}
\ifdim\wd\grmathbox>\linewidth
\resizebox{\linewidth}{!}{\usebox{\grmathbox}}
\else\usebox{\grmathbox}\fi
\endgroup
\end{center}

A classical strong-gravity ratio is instead

\begin{center}
\begingroup
\sbox{\grmathbox}{$\displaystyle 
\frac{G_NM}{c^2L}.
$}
\typeout{GRDISPLAY|597|\the\wd\grmathbox|\the\linewidth}
\ifdim\wd\grmathbox>\linewidth
\resizebox{\linewidth}{!}{\usebox{\grmathbox}}
\else\usebox{\grmathbox}\fi
\endgroup
\end{center}

These are different. Near the horizon of a large black hole, the second
can be order unity while the first is tiny. Strong classical gravity is
not automatically Planckian quantum gravity. Nor must one expand about
flat space to use low-energy reasoning; curved backgrounds can be
treated when their physical scales and the quantum state allow a
controlled approximation. See
\href{https://arxiv.org/abs/2211.09902}{Donoghue\textquotesingle s
review of quantum GR and its effective-theory limits}.

\textbf{Higher-derivative trap:} if we truncate an EFT and then solve
its higher-derivative equations exactly at arbitrarily high frequency,
we may find extra runaway or ghostlike solutions. That extrapolates the
truncated expression beyond the expansion that justified it. Consistent
EFT calculations treat higher-order corrections perturbatively, using
methods such as order reduction or field redefinition where appropriate.
An extra physical pole genuinely below the proposed cutoff would require
reexamining the field content, not dismissing it by slogan. These
distinctions are developed in
\href{https://arxiv.org/abs/1709.09695}{Solomon and
Trodden\textquotesingle s research on higher derivatives in EFT}.

\hypertarget{the-spin-2-route-gravity-is-forced-to-notice-its-own-bookkeeping}{%
\section{23.4 The spin-2 route: gravity is forced to notice its own
bookkeeping}\label{the-spin-2-route-gravity-is-forced-to-notice-its-own-bookkeeping}}

There is another way to approach Einstein\textquotesingle s equation.
Begin with a free massless spin-2 field \(h_{\mu\nu}\) in Minkowski
spacetime. Its linear gauge freedom has the form

\begin{center}
\begingroup
\sbox{\grmathbox}{$\displaystyle 
h_{\mu\nu}\mapsto h_{\mu\nu}
+\partial_\mu\xi_\nu+\partial_\nu\xi_\mu.
$}
\typeout{GRDISPLAY|598|\the\wd\grmathbox|\the\linewidth}
\ifdim\wd\grmathbox>\linewidth
\resizebox{\linewidth}{!}{\usebox{\grmathbox}}
\else\usebox{\grmathbox}\fi
\endgroup
\end{center}

This removes unphysical components and is related to the two propagating
polarizations. Couple \(h\) to matter schematically through

\begin{center}
\begingroup
\sbox{\grmathbox}{$\displaystyle 
S_{\mathrm{int}}\propto\int d^4x\,h_{\mu\nu}T^{\mu\nu}.
$}
\typeout{GRDISPLAY|599|\the\wd\grmathbox|\the\linewidth}
\ifdim\wd\grmathbox>\linewidth
\resizebox{\linewidth}{!}{\usebox{\grmathbox}}
\else\usebox{\grmathbox}\fi
\endgroup
\end{center}

Under the gauge transformation, integration by parts makes its variation
proportional to \(\xi_\nu\partial_\mu T^{\mu\nu}\). Linear gauge
consistency therefore wants the source to be conserved.

But once matter interacts with gravity, it exchanges energy and momentum
with the gravitational field. Matter stress alone cannot continue to
serve as an independently conserved source in the naive flat-background
equation. The field must also respond to its own contribution. That
changes its dynamics, which changes its contribution again. The
nonlinear completion reorganizes this self-coupling into GR under the
relevant assumptions.

The logic is powerful: a universally interacting massless spin-2 field
cannot consistently behave as if its own interactions are invisible to
its source equation. It also explains why simple superposition fails. A
primary presentation of the consistency construction is
\href{https://arxiv.org/abs/gr-qc/0411023}{Deser\textquotesingle s
``Self-Interaction and Gauge Invariance''}.

This is not a theorem that any imaginable spin-2 system must equal pure
GR at all energies. The argument relies on assumptions including
locality, Lorentz-compatible dynamics, appropriate gauge consistency,
field content, and the leading derivative structure. Field redefinitions
and stress-tensor improvements affect intermediate expressions.
Additional fields, higher-derivative terms, nonlocality, or different
backgrounds require separate analysis. A Minkowski-background
construction also presupposes an appropriate flat-background limit; it
does not determine an arbitrary cosmological constant from nothing.

\hypertarget{what-the-uniqueness-of-einsteins-equation-actually-says}{%
\section{23.5 What the uniqueness of Einstein\textquotesingle s equation
actually
says}\label{what-the-uniqueness-of-einsteins-equation-actually-says}}

In four dimensions, the Lovelock classification implies that a natural,
symmetric, divergence-free rank-two tensor built locally from the metric
and at most its second derivatives has the Einstein tensor and metric as
the available gravitational ingredients, under the
theorem\textquotesingle s hypotheses. Consequently, a metric-only
second-order field equation of this type takes the
Einstein-plus-cosmological form, up to constants. A precise mathematical
statement appears in
\href{https://arxiv.org/html/1005.2386v4}{``Lovelock\textquotesingle s
theorem revisited''}.

The assumptions are the engine of the conclusion. Add another field,
permit higher derivatives, change dimension, or change locality, and the
conclusion changes. Thus GR\textquotesingle s distinguished simplicity
and EFT\textquotesingle s higher-order corrections are compatible
claims. One concerns a restricted class of exact equations; the other
organizes small departures when that class is not assumed exact at every
scale.

Einstein\textquotesingle s equation is remarkably constrained at its
leading level. That makes its success intelligible without making the
unfinished parts of physics disappear.

\hypertarget{vacuum-energy-the-term-that-refuses-to-be-a-small-correction}{%
\section{23.6 Vacuum energy: the term that refuses to be a small
correction}\label{vacuum-energy-the-term-that-refuses-to-be-a-small-correction}}

The cosmological constant is a special challenge because it multiplies
the zero-derivative volume term. It is not automatically suppressed by
the long-wavelength expansion that weakens higher-curvature terms.

A Lorentz-invariant vacuum has stress tensor

\begin{center}
\begingroup
\sbox{\grmathbox}{$\displaystyle 
T^{\mathrm{vac}}_{\mu\nu}=-\epsilon_{\mathrm{vac}}g_{\mu\nu}.
$}
\typeout{GRDISPLAY|600|\the\wd\grmathbox|\the\linewidth}
\ifdim\wd\grmathbox>\linewidth
\resizebox{\linewidth}{!}{\usebox{\grmathbox}}
\else\usebox{\grmathbox}\fi
\endgroup
\end{center}

It has positive energy density if \(\epsilon_{\mathrm{vac}}>0\), and
pressure \(p=-\epsilon_{\mathrm{vac}}\). In Einstein\textquotesingle s
equation its effect has exactly the form of a cosmological constant.
Schematically, restoring \(c\),

\begin{center}
\begingroup
\sbox{\grmathbox}{$\displaystyle 
\Lambda_{\mathrm{effective}}
=\Lambda_{\mathrm{gravitational}}
+\frac{8\pi G_N}{c^4}\epsilon_{\mathrm{vac}},
$}
\typeout{GRDISPLAY|601|\the\wd\grmathbox|\the\linewidth}
\ifdim\wd\grmathbox>\linewidth
\resizebox{\linewidth}{!}{\usebox{\grmathbox}}
\else\usebox{\grmathbox}\fi
\endgroup
\end{center}

with the separate terms understood within a consistent renormalization
prescription. Only their physical combination is measurable.

Why does quantum theory make the small observed combination surprising?
In natural units, a free bosonic field suggests a zero-point
contribution

\begin{center}
\begingroup
\sbox{\grmathbox}{$\displaystyle 
\epsilon_{\mathrm{zero\ point}}
=\frac12\int\frac{d^3k}{(2\pi)^3}\sqrt{k^2+m^2}.
$}
\typeout{GRDISPLAY|602|\the\wd\grmathbox|\the\linewidth}
\ifdim\wd\grmathbox>\linewidth
\resizebox{\linewidth}{!}{\usebox{\grmathbox}}
\else\usebox{\grmathbox}\fi
\endgroup
\end{center}

A large-momentum cutoff \(M_*\) makes this grow roughly as \(M_*^4\).
You can see the fourth power without doing the integral: the
three-dimensional momentum measure contributes three powers, and the
high-momentum oscillator energy contributes one more.

That cutoff estimate is not a unique, covariant prediction of the
measured cosmological constant. Renormalization, the regulator, masses,
interactions, phase transitions, and the gravitational vacuum parameter
all matter. The notorious ``roughly 120 orders of magnitude'' comparison
uses a Planck-scale heuristic; it should not be presented as an exact
regulator-independent prediction that an experiment simply refuted.
\href{https://arxiv.org/abs/1205.3365}{Jérôme Martin\textquotesingle s
review of the cosmological constant problem}.

Nevertheless, the problem survives the correction to the slogan. In the
standard cosmological interpretation, the effective dark-energy density
corresponds to an energy scale of only a few millielectronvolts raised
to the fourth power. Contributions associated with much higher known
particle-physics scales naturally dwarf that. Why does the renormalized
combination stay so small when such contributions change? That is a
radiative-stability and naturalness question, not merely an unremoved
divergent integral.

GR permits a small cosmological constant. It does not explain its
observed value, and ordinary EFT bookkeeping does not supply the missing
explanation by itself.

\hypertarget{the-semiclassical-equation-and-its-limits}{%
\section{23.7 The semiclassical equation and its
limits}\label{the-semiclassical-equation-and-its-limits}}

An intermediate framework keeps the metric classical while treating
matter quantum mechanically:

\begin{center}
\begingroup
\sbox{\grmathbox}{$\displaystyle 
G_{\mu\nu}+\Lambda g_{\mu\nu}
+H^{\mathrm{higher\ curvature}}_{\mu\nu}
=\frac{8\pi G_N}{c^4}
\langle T_{\mu\nu}\rangle_{\mathrm{ren}}.
$}
\typeout{GRDISPLAY|603|\the\wd\grmathbox|\the\linewidth}
\ifdim\wd\grmathbox>\linewidth
\resizebox{\linewidth}{!}{\usebox{\grmathbox}}
\else\usebox{\grmathbox}\fi
\endgroup
\end{center}

The higher-curvature terms are included because renormalizing quantum
matter on a curved background requires corresponding gravitational
couplings. The state-dependent expectation value is a mean stress
tensor, not a literal classical list of particles.

This approximation is most credible when the relevant curvature and
momenta are below its cutoff, the quantum state is suitable, and
neglected metric fluctuations or stress fluctuations do not undermine
the mean-field description. Small curvature alone is not a universal
certificate of validity. Long evolution, delicate quantum correlations,
unusual states, or large fluctuations can raise additional issues.

Hawking radiation inhabits this framework. The endpoint of evaporation
and the complete accounting of information generally do not follow just
by extending the leading approximation until a black hole becomes
arbitrarily small. A calculation can advertise the boundary of its own
reliability; that is useful scientific information.

\hypertarget{what-experiments-test---and-what-dark-matter-and-dark-energy-mean}{%
\section{23.8 What experiments test - and what dark matter and dark
energy
mean}\label{what-experiments-test---and-what-dark-matter-and-dark-energy-mean}}

An observation does not compare ``all of GR'' with ``all alternatives''
in a single stroke. It constrains particular effects over particular
scales and source conditions.

\begingroup\small\setlength{\tabcolsep}{5pt}\renewcommand{\arraystretch}{1.13}

\begin{longtable}[]{@{}
  >{\raggedright\arraybackslash}p{(\columnwidth - 2\tabcolsep) * \real{0.3000}}
  >{\raggedright\arraybackslash}p{(\columnwidth - 2\tabcolsep) * \real{0.7000}}@{}}
\toprule\noalign{}
\begin{minipage}[b]{\linewidth}\raggedright
Measurement family
\end{minipage} & \begin{minipage}[b]{\linewidth}\raggedright
Examples of what it can constrain
\end{minipage} \\
\midrule\noalign{}
\endhead
\bottomrule\noalign{}
\endlastfoot
Freely falling bodies, clocks, and local laboratory tests & Composition
dependence, local Lorentz behavior, gravitational redshift \\
Planetary motion, timing, and lensing & Weak-field metric structure and
specified deviations from it \\
Binary pulsars & Strongly self-gravitating bodies, orbital dynamics,
radiative energy loss \\
Gravitational-wave signals & Wave generation, propagation, polarization,
and remnant dynamics within a chosen analysis \\
Cosmological expansion and structure & The joint behavior of gravity,
matter content, initial conditions, and large-scale evolution \\
\end{longtable}

\endgroup

For example, the LIGO-Virgo-KAGRA GWTC-5.0 analysis compares waveform
residuals, polarizations, generation, and remnant properties and reports
no overall evidence for physics beyond GR in those tests. That is a
strong set of constrained comparisons, with stated statistical and
modeling limits. It is not a proof that every possible modification at
every scale has vanished.
\href{https://arxiv.org/abs/2607.19293}{LVK\textquotesingle s primary
GWTC-5.0 tests paper}.

\textbf{Dark matter} and \textbf{dark energy} also name different
explanatory roles. In the usual cosmological model, dark matter behaves
approximately as clustering, nearly pressureless matter on large scales.
Dark energy denotes a component producing the observed accelerated
expansion; a cosmological constant is its simplest standard
representation. They are not two names for vacuum energy, nor does
either term alone establish that Einstein\textquotesingle s geometric
equation is wrong.

The observational inference always depends on a combined model:
gravitational laws, visible and invisible sources, their interactions,
and initial conditions. A successful alternative must fit the web of
measurements together. Matching one galaxy curve or one expansion
history is a starting point, not the entire examination.

\hypertarget{the-open-questions-are-sharper-than-space-is-mysterious}{%
\section{23.9 The open questions are sharper than ``space is
mysterious''}\label{the-open-questions-are-sharper-than-space-is-mysterious}}

The unresolved frontier contains concrete questions:

\begin{itemize}
\tightlist
\item
  What microscopic or nonperturbative description remains predictive
  where the gravitational low-energy expansion fails?
\item
  How do smooth causal geometry and approximately local fields emerge,
  if they are not fundamental at every scale?
\item
  What counts the black-hole entropy in sufficiently general situations,
  and how is information represented through formation and evaporation?
\item
  What mechanism, if any, explains the small effective cosmological
  constant and its stability under quantum corrections?
\item
  Which genuinely quantum properties of gravity can be isolated
  experimentally, rather than inferred solely from a classical
  gravitational fit?
\end{itemize}

Different research programs offer different partial answers and
controlled special cases. No derivation in this book establishes that
spacetime is a lattice, that consciousness creates geometry, or that a
particular microscopic proposal is experimentally selected.

The modern achievement is already substantial: the same geometry can be
understood as a dynamical constrained system, a local-frame gauge
structure, a theory of causal focusing, a thermodynamic participant, and
a predictive low-energy quantum field theory. Those are independent
pressures on the same equation. A future theory must explain why this
structure works so well, as well as where its limits lie.

\hypertarget{chapter-24}{%
\chapter{24. Bringing the whole machine together}\label{chapter-24}}

\hypertarget{a-calculation-is-a-chain-of-questions}{%
\section{24.1 A calculation is a chain of
questions}\label{a-calculation-is-a-chain-of-questions}}

After a long journey through geometry, it is possible to forget which
object answers which question. Here is the operating manual.

\begingroup\small\setlength{\tabcolsep}{5pt}\renewcommand{\arraystretch}{1.13}

\begin{longtable}[]{@{}
  >{\raggedright\arraybackslash}p{(\columnwidth - 4\tabcolsep) * \real{0.2500}}
  >{\raggedright\arraybackslash}p{(\columnwidth - 4\tabcolsep) * \real{0.3600}}
  >{\raggedright\arraybackslash}p{(\columnwidth - 4\tabcolsep) * \real{0.3900}}@{}}
\toprule\noalign{}
\begin{minipage}[b]{\linewidth}\raggedright
Stage
\end{minipage} & \begin{minipage}[b]{\linewidth}\raggedright
Question
\end{minipage} & \begin{minipage}[b]{\linewidth}\raggedright
Mathematical object
\end{minipage} \\
\midrule\noalign{}
\endhead
\bottomrule\noalign{}
\endlastfoot
Specify the model & Which fields exist, and what dynamics are assumed? &
Action, matter model, coupling constants \\
Choose a description & How are events labeled? What symmetry can be
used? & Coordinates, chart domain, ansatz, gauge \\
Measure locally & What are intervals, clocks, cones, and observer rest
spaces? & Metric and orthonormal frames \\
Compare nearby & What does it mean to differentiate or transport a
vector? & Connection and covariant derivative \\
Detect irreducible gravity & Do neighboring free-fall trajectories
develop tidal acceleration? & Riemann tensor \\
Impose gravitational dynamics & Is this geometry compatible with this
matter? & Einstein equation and matter equations \\
Specify a particular solution & Which physical history is being
described? & Constraint-satisfying initial data and appropriate
boundary/asymptotic information \\
Predict an experiment & What does a specified observer or detector
record? & Proper times, frequency ratios, tidal response, scattering
data \\
\end{longtable}

\endgroup

A component formula is usually a middle step, not the final observable.
An impressively complicated \(g_{00}\) is no substitute for specifying
who carries the clock.

\hypertarget{one-last-complete-example-clocks-can-disagree-in-flat-spacetime}{%
\section{24.2 One last complete example: clocks can disagree in flat
spacetime}\label{one-last-complete-example-clocks-can-disagree-in-flat-spacetime}}

This example deliberately combines several ideas that are easy to
confuse.

Start in Minkowski spacetime, with inertial coordinates \((T,X,Y,Z)\):

\begin{center}
\begingroup
\sbox{\grmathbox}{$\displaystyle 
ds^2=-c^2dT^2+dX^2+dY^2+dZ^2.
$}
\typeout{GRDISPLAY|604|\the\wd\grmathbox|\the\linewidth}
\ifdim\wd\grmathbox>\linewidth
\resizebox{\linewidth}{!}{\usebox{\grmathbox}}
\else\usebox{\grmathbox}\fi
\endgroup
\end{center}

Introduce accelerated coordinates \((t,x,y,z)\) by

\begin{center}
\begingroup
\sbox{\grmathbox}{$\displaystyle 
cT=\left(\frac{c^2}{a_0}+z\right)
\sinh\left(\frac{a_0t}{c}\right),
$}
\typeout{GRDISPLAY|605|\the\wd\grmathbox|\the\linewidth}
\ifdim\wd\grmathbox>\linewidth
\resizebox{\linewidth}{!}{\usebox{\grmathbox}}
\else\usebox{\grmathbox}\fi
\endgroup
\end{center}

\begin{center}
\begingroup
\sbox{\grmathbox}{$\displaystyle 
Z+\frac{c^2}{a_0}
=\left(\frac{c^2}{a_0}+z\right)
\cosh\left(\frac{a_0t}{c}\right),
\qquad X=x,\quad Y=y,
$}
\typeout{GRDISPLAY|606|\the\wd\grmathbox|\the\linewidth}
\ifdim\wd\grmathbox>\linewidth
\resizebox{\linewidth}{!}{\usebox{\grmathbox}}
\else\usebox{\grmathbox}\fi
\endgroup
\end{center}

where \(a_0>0\) is a constant acceleration scale. Work in the wedge with

\begin{center}
\begingroup
\sbox{\grmathbox}{$\displaystyle 
q(z)=1+\frac{a_0z}{c^2}>0.
$}
\typeout{GRDISPLAY|607|\the\wd\grmathbox|\the\linewidth}
\ifdim\wd\grmathbox>\linewidth
\resizebox{\linewidth}{!}{\usebox{\grmathbox}}
\else\usebox{\grmathbox}\fi
\endgroup
\end{center}

Differentiate the transformation. The mixed \(dt\,dz\) terms from
\(-c^2dT^2\) and \(dZ^2\) cancel, because the hyperbolic sine and cosine
enter symmetrically. The remaining terms simplify using
\(\cosh^2\eta-\sinh^2\eta=1\):

\begin{center}
\begingroup
\sbox{\grmathbox}{$\displaystyle 
\boxed{ds^2=-q(z)^2c^2dt^2+dx^2+dy^2+dz^2.}
$}
\typeout{GRDISPLAY|608|\the\wd\grmathbox|\the\linewidth}
\ifdim\wd\grmathbox>\linewidth
\resizebox{\linewidth}{!}{\usebox{\grmathbox}}
\else\usebox{\grmathbox}\fi
\endgroup
\end{center}

This is the Rindler metric. We have changed coordinates, so spacetime
must still be flat. But the metric now varies with position, and
stationary clocks in these coordinates run at different rates relative
to \(t\):

\begin{center}
\begingroup
\sbox{\grmathbox}{$\displaystyle 
d\tau=q(z)\,dt.
$}
\typeout{GRDISPLAY|609|\the\wd\grmathbox|\the\linewidth}
\ifdim\wd\grmathbox>\linewidth
\resizebox{\linewidth}{!}{\usebox{\grmathbox}}
\else\usebox{\grmathbox}\fi
\endgroup
\end{center}

Does that contradict anything we learned? No. ``Stationary in Rindler
coordinates'' describes an accelerated family of observers. The question
includes the observers, not only the spacetime.

\hypertarget{calculate-the-connection}{%
\subsection{Calculate the connection}\label{calculate-the-connection}}

Here the time coordinate is \(t\) in seconds, so \(g_{tt}=-c^2q^2\). The
nonzero connection coefficients involving the \((t,z)\) sector are

\begin{center}
\begingroup
\sbox{\grmathbox}{$\displaystyle 
\Gamma^t{}_{tz}=\Gamma^t{}_{zt}=\frac{q'}{q},
\qquad
\Gamma^z{}_{tt}=c^2qq'=a_0q,
$}
\typeout{GRDISPLAY|610|\the\wd\grmathbox|\the\linewidth}
\ifdim\wd\grmathbox>\linewidth
\resizebox{\linewidth}{!}{\usebox{\grmathbox}}
\else\usebox{\grmathbox}\fi
\endgroup
\end{center}

where \(q'=dq/dz=a_0/c^2\).

The first coefficient describes how the time basis changes as we move in
\(z\). The second is the coordinate acceleration term entering the \(z\)
geodesic equation.

An observer held at fixed \(z\) has \(u^t=dt/d\tau=1/q\) and \(u^z=0\).
Their four-acceleration is

\begin{center}
\begingroup
\sbox{\grmathbox}{$\displaystyle 
a^z=\Gamma^z{}_{tt}(u^t)^2=\frac{a_0}{q}.
$}
\typeout{GRDISPLAY|611|\the\wd\grmathbox|\the\linewidth}
\ifdim\wd\grmathbox>\linewidth
\resizebox{\linewidth}{!}{\usebox{\grmathbox}}
\else\usebox{\grmathbox}\fi
\endgroup
\end{center}

Because \(g_{zz}=1\) and the acceleration is spatial in this
observer\textquotesingle s rest frame, the accelerometer magnitude is

\begin{center}
\begingroup
\sbox{\grmathbox}{$\displaystyle 
\boxed{a_{\mathrm{proper}}(z)=\frac{a_0}{q(z)}.}
$}
\typeout{GRDISPLAY|612|\the\wd\grmathbox|\the\linewidth}
\ifdim\wd\grmathbox>\linewidth
\resizebox{\linewidth}{!}{\usebox{\grmathbox}}
\else\usebox{\grmathbox}\fi
\endgroup
\end{center}

The observers at different heights require different proper
accelerations to maintain this stationary arrangement. They are not a
collection of freely falling observers.

\hypertarget{calculate-the-curvature}{%
\subsection{Calculate the curvature}\label{calculate-the-curvature}}

Compute a representative component directly:

\begin{center}
\begingroup
\sbox{\grmathbox}{$\displaystyle 
R^z{}_{tzt}
=\partial_z\Gamma^z{}_{tt}
-\Gamma^z{}_{tt}\Gamma^t{}_{zt}.
$}
\typeout{GRDISPLAY|613|\the\wd\grmathbox|\the\linewidth}
\ifdim\wd\grmathbox>\linewidth
\resizebox{\linewidth}{!}{\usebox{\grmathbox}}
\else\usebox{\grmathbox}\fi
\endgroup
\end{center}

The other terms vanish for this static diagonal metric. Substitute:

\begin{center}
\begingroup
\sbox{\grmathbox}{$\displaystyle 
R^z{}_{tzt}
=c^2\bigl[(q')^2+qq''\bigr]
-c^2(q')^2
=c^2qq''=0.
$}
\typeout{GRDISPLAY|614|\the\wd\grmathbox|\the\linewidth}
\ifdim\wd\grmathbox>\linewidth
\resizebox{\linewidth}{!}{\usebox{\grmathbox}}
\else\usebox{\grmathbox}\fi
\endgroup
\end{center}

The cancellation matters. A derivative of a connection coefficient alone
is not generally curvature. The quadratic connection terms are part of
the definition. Here \(q\) is linear in \(z\), so \(q''=0\), and every
Riemann component vanishes.

The Einstein tensor is therefore zero. In the idealized test-apparatus
limit, this metric is compatible with vacuum and \(\Lambda=0\).

\hypertarget{calculate-the-frequency-shift}{%
\subsection{Calculate the frequency
shift}\label{calculate-the-frequency-shift}}

The metric is stationary, so the covariant photon momentum component
\(p_t\) is conserved along a null geodesic. A stationary observer
measures

\begin{center}
\begingroup
\sbox{\grmathbox}{$\displaystyle 
E(z)=-p_\mu u^\mu=-\frac{p_t}{q(z)}.
$}
\typeout{GRDISPLAY|615|\the\wd\grmathbox|\the\linewidth}
\ifdim\wd\grmathbox>\linewidth
\resizebox{\linewidth}{!}{\usebox{\grmathbox}}
\else\usebox{\grmathbox}\fi
\endgroup
\end{center}

Since photon energy is proportional to frequency, a photon emitted at
\(z_e\) and received at \(z_r\) satisfies

\begin{center}
\begingroup
\sbox{\grmathbox}{$\displaystyle 
\boxed{\frac{\nu_r}{\nu_e}=\frac{q(z_e)}{q(z_r)}.}
$}
\typeout{GRDISPLAY|616|\the\wd\grmathbox|\the\linewidth}
\ifdim\wd\grmathbox>\linewidth
\resizebox{\linewidth}{!}{\usebox{\grmathbox}}
\else\usebox{\grmathbox}\fi
\endgroup
\end{center}

If the receiver is at larger \(z\), the received frequency is lower.
There is a frequency shift between accelerated observers even though
spacetime curvature is identically zero.

The inertial description interprets the same experiment in terms of the
observers\textquotesingle{} motion during the light exchange. The
accelerated description interprets it using a position-dependent lapse.
Both predict the same detector readings.

\begin{quote}
\textbf{Capstone trap door:} Nonconstant metric components, nonzero
Christoffel symbols, acceleration readings, and frequency shifts between
a specified family of observers do not individually prove nonzero
spacetime curvature. Tidal curvature requires the appropriate invariant
geometric test.
\end{quote}

The chart covers the wedge \(Z+c^2/a_0>|cT|\). Its boundaries
\(Z+c^2/a_0=\pm cT\) are null acceleration horizons for the stationary
Rindler observers. They are neither curvature singularities nor
black-hole horizons in Minkowski spacetime. A coordinate system can have
a limited view without the universe having a wound.

\hypertarget{how-to-calculate-a-spacetime-without-getting-lost}{%
\section{24.3 How to calculate a spacetime without getting
lost}\label{how-to-calculate-a-spacetime-without-getting-lost}}

Suppose someone hands you a metric and asks you to interpret it. Use
this sequence.

\begin{enumerate}
\def\labelenumi{\arabic{enumi}.}
\tightlist
\item
  \textbf{Check the chart and domain.} Which coordinate is time? Are
  angles dimensionless? Where is the matrix nondegenerate? Do any
  apparent singularities occur only at a chart boundary?
\item
  \textbf{Invert the matrix.} Verify
  \(g^{\mu\alpha}g_{\alpha\nu}=\delta^\mu{}_{\nu}\). For a nondiagonal
  metric, componentwise reciprocals are incorrect.
\item
  \textbf{Compute the determinant.} This controls the volume element and
  often reveals where a coordinate chart fails.
\item
  \textbf{Identify symmetries before differentiating.} Independence of a
  coordinate can provide a Killing vector and conserved quantities.
  Symmetry can save pages of algebra.
\item
  \textbf{Compute \(\Gamma\) from \(g\) and \(\partial g\).} Exploit its
  lower-index symmetry only for the Levi-Civita connection in a
  coordinate basis.
\item
  \textbf{Compute Riemann before trusting a curvature slogan.} Form the
  derivative terms and both quadratic terms with the chosen convention.
\item
  \textbf{Contract carefully.} Obtain Ricci, scalar curvature, and
  Einstein tensor; free indices must remain in the right places.
\item
  \textbf{Compare with a physically admissible stress tensor.}
  Conservation, matter equations, and an equation of state may rule out
  an apparently convenient interpretation.
\item
  \textbf{Choose observers.} Build four-velocities or a local
  orthonormal frame and project coordinate tensors into quantities those
  observers measure.
\item
  \textbf{Check a known limit.} Flat space, weak fields, small
  velocities, spherical symmetry, or an independently known invariant
  can reveal an error that elegant notation concealed.
\end{enumerate}

A computer can perform steps 2-7 beautifully while understanding none of
steps 1, 8, or 9. That is why a hundred lines of symbolic output can be
less trustworthy than one well-chosen physical limit.

\hypertarget{four-calibration-geometries}{%
\section{24.4 Four calibration
geometries}\label{four-calibration-geometries}}

These examples are excellent tests for a hand calculation or a computer
implementation. Here \(A\) is the radius of a sphere; in the
Schwarzschild row, \(m=G_NM/c^2\); in the FLRW row, use \(c=1\) and a
spatially flat cosmology with \(H=\dot a/a\).

\begingroup\small\setlength{\tabcolsep}{5pt}\renewcommand{\arraystretch}{1.13}

\begin{longtable}[]{@{}
  >{\raggedright\arraybackslash}p{(\columnwidth - 4\tabcolsep) * \real{0.2500}}
  >{\raggedright\arraybackslash}p{(\columnwidth - 4\tabcolsep) * \real{0.3600}}
  >{\raggedright\arraybackslash}p{(\columnwidth - 4\tabcolsep) * \real{0.3900}}@{}}
\toprule\noalign{}
\begin{minipage}[b]{\linewidth}\raggedright
Geometry
\end{minipage} & \begin{minipage}[b]{\linewidth}\raggedright
A result that must emerge
\end{minipage} & \begin{minipage}[b]{\linewidth}\raggedright
Error it often catches
\end{minipage} \\
\midrule\noalign{}
\endhead
\bottomrule\noalign{}
\endlastfoot
Euclidean plane, \(ds^2=dr^2+r^2d\theta^2\) & Nonzero \(\Gamma\), but
all \(R^a{}_{bcd}=0\) & Mistaking curvilinear coordinates for
curvature \\
Round sphere, \(ds^2=A^2(d\theta^2+\sin^2\theta\,d\phi^2)\) &
\(R=2/A^2\) & Riemann sign and Ricci contraction errors \\
Schwarzschild exterior & \(R_{\mu\nu}=0\), but
\(R_{\alpha\beta\gamma\delta}R^{\alpha\beta\gamma\delta}=48m^2/r^6\) &
Mistaking Ricci-flatness for flatness \\
Flat FLRW & \(G_{tt}=3H^2\) and \(R=6(\ddot a/a+H^2)\) & Time-sign
mistakes and missing quadratic connection terms \\
\end{longtable}

\endgroup

The FLRW expression also shows that scalar curvature need not be
positive even when the spatial slices are flat. ``Spatially flat'' is
not ``spacetime flat.''

\hypertarget{why-substituting-a-solution-into-the-action-too-early-can-destroy-the-question}{%
\section{24.5 Why substituting a solution into the action too early can
destroy the
question}\label{why-substituting-a-solution-into-the-action-too-early-can-destroy-the-question}}

For a vacuum Ricci-flat solution with \(\Lambda=0\), the
Einstein-Hilbert bulk integrand \(\sqrt{-g}R\) vanishes. Does this mean
the action could not possibly determine that solution?

No. A function can vanish at a point without having zero derivative in
every direction there. More specifically, a field configuration can
satisfy \(R=0\) while nearby off-shell configurations do not.

\textbf{On shell} means the fields satisfy their equations of motion.
\textbf{Off shell} means we allow variations that need not satisfy those
equations. A variational principle compares nearby off-shell
configurations and asks whether the first change in the total action
vanishes.

If you impose the equations before varying, you discard the directions
that the variation is supposed to test. It is like replacing a function
by its value at the minimum and then trying to recover the slope from
that single number.

Boundary contributions add another reason not to identify a vanishing
bulk integrand with a physically empty action. In gravitational
problems, boundary conditions and boundary terms can carry decisive
information.

\hypertarget{a-compact-mental-reconstruction-of-the-einstein-equation}{%
\section{24.6 A compact mental reconstruction of the Einstein
equation}\label{a-compact-mental-reconstruction-of-the-einstein-equation}}

Close the book for a moment and try to rebuild the logic.

A metric defines intervals and local causal structure. Its Levi-Civita
connection supplies a derivative that respects the metric and has no
torsion. The connection\textquotesingle s failure to return transported
vectors unchanged around infinitesimal loops is curvature.
Curvature\textquotesingle s contractions give the Ricci tensor and
scalar.

A local gravitational action with the usual metric-only, two-derivative
assumptions contains the scalar curvature and a constant term,
integrated with the invariant volume measure. Matter has its own action
and defines stress-energy through its response to a metric variation.

Vary the inverse metric. The curvature variation contributes
\(R_{\mu\nu}\) and a boundary divergence. The measure variation
contributes \(-\tfrac12Rg_{\mu\nu}\). The constant term contributes
\(\Lambda g_{\mu\nu}\). The matter variation contributes \(-T_{\mu\nu}\)
with the convention-dependent normalization already fixed.

Stationarity for arbitrary permitted interior variations gives

\begin{center}
\begingroup
\sbox{\grmathbox}{$\displaystyle 
G_{\mu\nu}+\Lambda g_{\mu\nu}
=\frac{8\pi G_N}{c^4}T_{\mu\nu}.
$}
\typeout{GRDISPLAY|617|\the\wd\grmathbox|\the\linewidth}
\ifdim\wd\grmathbox>\linewidth
\resizebox{\linewidth}{!}{\usebox{\grmathbox}}
\else\usebox{\grmathbox}\fi
\endgroup
\end{center}

The Bianchi identity makes its left-hand side covariantly
divergence-free. Diffeomorphism symmetry organizes the corresponding
identity and matter balance law. The Newtonian limit fixes the coupling.
Initial data and constraints select a physical history, up to gauge.
Matter equations and the metric evolve together.

Then you ask an actual observer to make a measurement.

That final step is what keeps the entire geometric construction a theory
of physics.

\hypertarget{the-highest-value-change-in-intuition}{%
\section{24.7 The highest-value change in
intuition}\label{the-highest-value-change-in-intuition}}

Before studying general relativity, you may ask, ``What force pulls the
body along that curve?''

After studying it, you have more precise questions available:

\begin{itemize}
\tightlist
\item
  Is that curve a geodesic of the physical metric?
\item
  Is an accelerometer measuring nonzero proper acceleration?
\item
  Does a neighboring geodesic reveal tidal curvature?
\item
  Is the apparent effect a coordinate feature, an observer choice, or an
  invariant property?
\item
  Which degrees of freedom are fixed by matter, and which arrive as
  gravitational initial or radiative data?
\item
  Which approximation connects this mathematical model to the
  experiment?
\end{itemize}

You have not merely acquired a new answer about gravity. You have
acquired a better set of questions.

\hypertarget{appendix-a}{%
\chapter{Appendix A. Thirty exercises that turn recognition into
understanding}\label{appendix-a}}

These are not speed tests. Several are designed so that the tempting
answer is wrong. Try the problem before reading the solution. If you
obtain a different sign, first compare conventions; if you obtain a
different physical prediction after doing so, investigate.

\hypertarget{a.1-a-vector-and-a-covector-change-their-clothes}{%
\section{A.1 A vector and a covector change their
clothes}\label{a.1-a-vector-and-a-covector-change-their-clothes}}

\textbf{Problem.} In two dimensions, let \(x'=2x+y\) and \(y'=y\). A
vector has components \((V^x,V^y)=(3,4)\) and a covector has components
\((\omega_x,\omega_y)=(5,-2)\). Transform both and check their
contraction.

\textbf{Solution.} The Jacobian is

\begin{center}
\begingroup
\sbox{\grmathbox}{$\displaystyle 
J=\begin{pmatrix}2&1\\0&1\end{pmatrix},
\qquad J^{-1}=\begin{pmatrix}1/2&-1/2\\0&1\end{pmatrix}.
$}
\typeout{GRDISPLAY|618|\the\wd\grmathbox|\the\linewidth}
\ifdim\wd\grmathbox>\linewidth
\resizebox{\linewidth}{!}{\usebox{\grmathbox}}
\else\usebox{\grmathbox}\fi
\endgroup
\end{center}

A vector transforms with \(J\), so \(V'=(10,4)\). A covector,
represented as a row, transforms with \(J^{-1}\), giving
\(\omega'=(5/2,-9/2)\). Thus

\begin{center}
\begingroup
\sbox{\grmathbox}{$\displaystyle 
\omega_\mu V^\mu=15-8=7,
\qquad
\omega'_{\mu'}V'^{\mu'}=25-18=7.
$}
\typeout{GRDISPLAY|619|\the\wd\grmathbox|\the\linewidth}
\ifdim\wd\grmathbox>\linewidth
\resizebox{\linewidth}{!}{\usebox{\grmathbox}}
\else\usebox{\grmathbox}\fi
\endgroup
\end{center}

The components changed while the pairing did not. Using the same
transformation rule for both objects would lose this invariance. The
upper and lower indices encode how the two kinds of components
compensate each other.

\hypertarget{a.2-an-inverse-metric-is-a-matrix-inverse}{%
\section{A.2 An inverse metric is a matrix
inverse}\label{a.2-an-inverse-metric-is-a-matrix-inverse}}

\textbf{Problem.} For the positive-definite metric

\begin{center}
\begingroup
\sbox{\grmathbox}{$\displaystyle 
g_{ij}=\begin{pmatrix}2&1\\1&2\end{pmatrix}
$}
\typeout{GRDISPLAY|620|\the\wd\grmathbox|\the\linewidth}
\ifdim\wd\grmathbox>\linewidth
\resizebox{\linewidth}{!}{\usebox{\grmathbox}}
\else\usebox{\grmathbox}\fi
\endgroup
\end{center}

and \(V^i=(1,2)\), find \(g^{ij}\), \(V_i\), and \(g(V,V)\).

\textbf{Solution.} Since the determinant is \(3\),

\begin{center}
\begingroup
\sbox{\grmathbox}{$\displaystyle 
g^{ij}=\frac13\begin{pmatrix}2&-1\\-1&2\end{pmatrix}.
$}
\typeout{GRDISPLAY|621|\the\wd\grmathbox|\the\linewidth}
\ifdim\wd\grmathbox>\linewidth
\resizebox{\linewidth}{!}{\usebox{\grmathbox}}
\else\usebox{\grmathbox}\fi
\endgroup
\end{center}

Lowering gives \(V_i=(4,5)\), and the norm squared is
\(V^iV_i=1\cdot4+2\cdot5=14\). Raising again returns \((1,2)\).
Replacing each matrix entry by its reciprocal would fail even the
identity check \(g^{ik}g_{kj}=\delta^i{}_j\).

\hypertarget{a.3-energy-depends-on-the-observer-even-when-it-is-a-scalar}{%
\section{A.3 Energy depends on the observer, even when it is a
scalar}\label{a.3-energy-depends-on-the-observer-even-when-it-is-a-scalar}}

\textbf{Problem.} A massive particle moves at \(v=0.6c\) relative to an
inertial observer. Find its energy and momentum. What energy does an
observer moving with the particle assign to it?

\textbf{Solution.} \(\gamma=1/\sqrt{1-0.36}=1.25\). The first observer
measures \(E=1.25mc^2\) and \(p=\gamma mv=0.75mc\). The comoving
observer measures \(mc^2\) and zero spatial momentum.

There is no contradiction with \(E_{(U)}=-p_\mu U^\mu\) being a scalar.
A coordinate change transforms both \(p\) and the same observer \(U\),
leaving that number unchanged. Choosing a different observer replaces
\(U\) by a different vector and defines a different measurement.

\hypertarget{a.4-the-hessian-knows-about-polar-coordinates}{%
\section{A.4 The Hessian knows about polar
coordinates}\label{a.4-the-hessian-knows-about-polar-coordinates}}

\textbf{Problem.} On the Euclidean plane with
\(ds^2=dr^2+r^2d\theta^2\), let \(f=r\). Although
\(\partial_\theta\partial_\theta f=0\), show that
\(\nabla_\theta\nabla_\theta f=r\) and \(\nabla^2f=1/r\) for \(r>0\).

\textbf{Solution.} The first derivative of a scalar is the covector
\(df\). Differentiating it requires a connection:

\begin{center}
\begingroup
\sbox{\grmathbox}{$\displaystyle 
\nabla_\theta\nabla_\theta f
=\partial_\theta\partial_\theta f
-\Gamma^r{}_{\theta\theta}\partial_rf
=0-(-r)(1)=r.
$}
\typeout{GRDISPLAY|622|\the\wd\grmathbox|\the\linewidth}
\ifdim\wd\grmathbox>\linewidth
\resizebox{\linewidth}{!}{\usebox{\grmathbox}}
\else\usebox{\grmathbox}\fi
\endgroup
\end{center}

The radial Hessian component vanishes. Contracting with
\(g^{\theta\theta}=1/r^2\) gives \(\nabla^2r=r/r^2=1/r\). The second
covariant derivative of a scalar is not generally just a matrix of
ordinary second derivatives.

\hypertarget{a.5-a-divergence-free-field-can-hide-a-puncture}{%
\section{A.5 A divergence-free field can hide a
puncture}\label{a.5-a-divergence-free-field-can-hide-a-puncture}}

\textbf{Problem.} On the same plane, take \(V^r=1/r\) and
\(V^\theta=0\). Compute its divergence for \(r>0\). Why is the flux
through a circle nonzero?

\textbf{Solution.} Since \(\sqrt{g}=r\),

\begin{center}
\begingroup
\sbox{\grmathbox}{$\displaystyle 
\nabla_iV^i=\frac1r\partial_r(rV^r)=0
$}
\typeout{GRDISPLAY|623|\the\wd\grmathbox|\the\linewidth}
\ifdim\wd\grmathbox>\linewidth
\resizebox{\linewidth}{!}{\usebox{\grmathbox}}
\else\usebox{\grmathbox}\fi
\endgroup
\end{center}

away from the origin. But a circle of radius \(r\) has outward flux
\((1/r)(2\pi r)=2\pi\).

The field is undefined at the origin. Applying the divergence theorem to
a disk without accounting for that singularity violates its smoothness
assumptions. On an annulus, the outer and inner boundary fluxes cancel
with their appropriate orientations. In a distributional extension to
the whole plane, the divergence is \(2\pi\delta^{(2)}(\mathbf{x})\).

The lesson applies far beyond this example: local equations plus domain
assumptions determine whether an integral inference is legitimate.

\hypertarget{a.6-constant-cartesian-arrows-have-changing-polar-components}{%
\section{A.6 Constant Cartesian arrows have changing polar
components}\label{a.6-constant-cartesian-arrows-have-changing-polar-components}}

\textbf{Problem.} A constant Cartesian vector pointing in the \(x\)
direction has polar components \(V^r=\cos\theta\) and
\(V^\theta=-\sin\theta/r\). Show that \(\nabla_\theta V^r=0\).

\textbf{Solution.} The partial derivative is
\(\partial_\theta V^r=-\sin\theta\), which by itself falsely suggests
the vector changes. The connection contribution is

\begin{center}
\begingroup
\sbox{\grmathbox}{$\displaystyle 
\Gamma^r{}_{\theta\theta}V^\theta
=(-r)\left(-\frac{\sin\theta}{r}\right)=\sin\theta.
$}
\typeout{GRDISPLAY|624|\the\wd\grmathbox|\the\linewidth}
\ifdim\wd\grmathbox>\linewidth
\resizebox{\linewidth}{!}{\usebox{\grmathbox}}
\else\usebox{\grmathbox}\fi
\endgroup
\end{center}

They cancel. This is a concrete example of a covariant derivative
correcting the change in a coordinate basis. No physical bending of
space was needed.

\hypertarget{a.7-the-connection-can-vanish-while-curvature-survives}{%
\section{A.7 The connection can vanish while curvature
survives}\label{a.7-the-connection-can-vanish-while-curvature-survives}}

\textbf{Problem.} At an event \(P\), use normal coordinates so that
\(\Gamma^\rho{}_{\mu\nu}(P)=0\). Why does this not imply
\(R^\rho{}_{\sigma\mu\nu}(P)=0\)?

\textbf{Solution.} At \(P\), the terms quadratic in \(\Gamma\) vanish,
but

\begin{center}
\begingroup
\sbox{\grmathbox}{$\displaystyle 
R^\rho{}_{\sigma\mu\nu}(P)
=\partial_\mu\Gamma^\rho{}_{\nu\sigma}(P)
-\partial_\nu\Gamma^\rho{}_{\mu\sigma}(P)
$}
\typeout{GRDISPLAY|625|\the\wd\grmathbox|\the\linewidth}
\ifdim\wd\grmathbox>\linewidth
\resizebox{\linewidth}{!}{\usebox{\grmathbox}}
\else\usebox{\grmathbox}\fi
\endgroup
\end{center}

can remain nonzero. A function vanishing at a point need not have
vanishing derivatives there. In normal coordinates the
metric\textquotesingle s first derivatives vanish at \(P\), but its
second derivatives contain curvature information.

\hypertarget{a.8-a-spheres-circumference-confesses-its-curvature}{%
\section{A.8 A sphere\textquotesingle s circumference confesses its
curvature}\label{a.8-a-spheres-circumference-confesses-its-curvature}}

\textbf{Problem.} On a sphere of radius \(A\), a geodesic circle at
geodesic distance \(s\) from the north pole has circumference
\(C(s)=2\pi A\sin(s/A)\). Expand it for small \(s\) and compare it with
a flat circle.

\textbf{Solution.} Taylor expansion gives

\begin{center}
\begingroup
\sbox{\grmathbox}{$\displaystyle 
C(s)=2\pi s\left(1-\frac{s^2}{6A^2}+O(s^4/A^4)\right).
$}
\typeout{GRDISPLAY|626|\the\wd\grmathbox|\the\linewidth}
\ifdim\wd\grmathbox>\linewidth
\resizebox{\linewidth}{!}{\usebox{\grmathbox}}
\else\usebox{\grmathbox}\fi
\endgroup
\end{center}

It is smaller than \(2\pi s\). Intrinsic measurements of radius and
circumference reveal positive curvature without any view from an
embedding space. The Gaussian curvature is \(1/A^2\); in two dimensions
the scalar curvature is twice that, \(R=2/A^2\).

\hypertarget{a.9-newtonian-tides-know-the-sign-convention}{%
\section{A.9 Newtonian tides know the sign
convention}\label{a.9-newtonian-tides-know-the-sign-convention}}

\textbf{Problem.} For \(\Phi=-G_NM/r\), derive the radial and transverse
relative accelerations of nearby freely falling particles in the
Newtonian limit.

\textbf{Solution.} In local Cartesian directions aligned with the radial
axis, the Hessian of the potential has eigenvalues

\begin{center}
\begingroup
\sbox{\grmathbox}{$\displaystyle 
\left(-\frac{2G_NM}{r^3},\frac{G_NM}{r^3},\frac{G_NM}{r^3}\right).
$}
\typeout{GRDISPLAY|627|\the\wd\grmathbox|\the\linewidth}
\ifdim\wd\grmathbox>\linewidth
\resizebox{\linewidth}{!}{\usebox{\grmathbox}}
\else\usebox{\grmathbox}\fi
\endgroup
\end{center}

Relative acceleration is
\(\delta\ddot x^i=-\delta^{ik}\partial_k\partial_j\Phi\,\delta x^j\), so
the corresponding tidal eigenvalues are

\begin{center}
\begingroup
\sbox{\grmathbox}{$\displaystyle 
\left(\frac{2G_NM}{r^3},-\frac{G_NM}{r^3},-\frac{G_NM}{r^3}\right).
$}
\typeout{GRDISPLAY|628|\the\wd\grmathbox|\the\linewidth}
\ifdim\wd\grmathbox>\linewidth
\resizebox{\linewidth}{!}{\usebox{\grmathbox}}
\else\usebox{\grmathbox}\fi
\endgroup
\end{center}

Radial separations stretch; transverse separations compress. The trace
vanishes outside the source, consistent with the vacuum Poisson
equation. In the relativistic convention used here,
\(R^i{}_{0j0}\simeq\delta^{ik}\partial_k\partial_j\Phi/c^2\), and
geodesic deviation supplies the minus sign.

\hypertarget{a.10-trace-reversal-produces-the-missing-factor-of-two}{%
\section{A.10 Trace reversal produces the missing factor of
two}\label{a.10-trace-reversal-produces-the-missing-factor-of-two}}

\textbf{Problem.} With \(\Lambda=0\), contract the four-dimensional
Einstein equation and rewrite it as an equation for \(R_{\mu\nu}\).
Evaluate its \(00\) source for slowly moving dust.

\textbf{Solution.} The trace is \(R-2R=\kappa T\), so \(R=-\kappa T\),
where \(\kappa=8\pi G_N/c^4\). Therefore

\begin{center}
\begingroup
\sbox{\grmathbox}{$\displaystyle 
R_{\mu\nu}=\kappa\left(T_{\mu\nu}-\frac12g_{\mu\nu}T\right).
$}
\typeout{GRDISPLAY|629|\the\wd\grmathbox|\the\linewidth}
\ifdim\wd\grmathbox>\linewidth
\resizebox{\linewidth}{!}{\usebox{\grmathbox}}
\else\usebox{\grmathbox}\fi
\endgroup
\end{center}

For slow dust, \(T_{00}\simeq\rho c^2\), \(T\simeq-\rho c^2\), and
\(g_{00}\simeq-1\). The parenthesis in the \(00\) component becomes
\(\rho c^2/2\). Thus \(R_{00}\simeq4\pi G_N\rho/c^2\), consistent with
\(R_{00}\simeq\nabla^2\Phi/c^2\) and Newton\textquotesingle s Poisson
equation.

\hypertarget{a.11-a-zero-scalar-curvature-proves-much-less-than-you-think}{%
\section{A.11 A zero scalar curvature proves much less than you
think}\label{a.11-a-zero-scalar-curvature-proves-much-less-than-you-think}}

\textbf{Problem.} Does \(R=0\) imply vacuum? Does vacuum with
\(\Lambda=0\) imply no gravitational waves?

\textbf{Solution.} Both answers are no. A classical electromagnetic
field in four dimensions has a trace-free stress tensor.
Einstein\textquotesingle s trace equation gives \(R=0\) when
\(\Lambda=0\), even though \(T_{\mu\nu}\) and \(R_{\mu\nu}\) may be
nonzero. Vacuum implies \(R_{\mu\nu}=0\), but the Weyl tensor can still
describe gravitational waves or an exterior tidal field.

Three distinct statements are being separated: scalar-flat, Ricci-flat,
and Riemann-flat. Do not let a shorter word count persuade you that they
are interchangeable.

\hypertarget{a.12-pressure-enters-the-local-curvature-source}{%
\section{A.12 Pressure enters the local curvature
source}\label{a.12-pressure-enters-the-local-curvature-source}}

\textbf{Problem.} In a perfect fluid\textquotesingle s local orthonormal
rest frame, find \(T\) and the trace-reversed \(00\) source.

\textbf{Solution.} With
\(T_{\hat\mu\hat\nu}=\operatorname{diag}(\epsilon,p,p,p)\),

\begin{center}
\begingroup
\sbox{\grmathbox}{$\displaystyle 
T=-\epsilon+3p,
$}
\typeout{GRDISPLAY|630|\the\wd\grmathbox|\the\linewidth}
\ifdim\wd\grmathbox>\linewidth
\resizebox{\linewidth}{!}{\usebox{\grmathbox}}
\else\usebox{\grmathbox}\fi
\endgroup
\end{center}

and

\begin{center}
\begingroup
\sbox{\grmathbox}{$\displaystyle 
T_{\hat0\hat0}-\frac12\eta_{\hat0\hat0}T
=\epsilon+\frac12(-\epsilon+3p)
=\frac12(\epsilon+3p).
$}
\typeout{GRDISPLAY|631|\the\wd\grmathbox|\the\linewidth}
\ifdim\wd\grmathbox>\linewidth
\resizebox{\linewidth}{!}{\usebox{\grmathbox}}
\else\usebox{\grmathbox}\fi
\endgroup
\end{center}

That is one precise sense in which pressure gravitates. It does not
license adding \(3p/c^2\) to the mass of every isolated box while
ignoring stresses in its walls. A complete system includes its stresses,
binding, and appropriate global mass definition.

\hypertarget{a.13-the-determinant-supplies-the-trace-term}{%
\section{A.13 The determinant supplies the trace
term}\label{a.13-the-determinant-supplies-the-trace-term}}

\textbf{Problem.} Starting from
\(\delta\ln|\det M|=\operatorname{tr}(M^{-1}\delta M)\), show why
varying \(g^{\mu\nu}\) produces a minus sign in \(\delta\sqrt{-g}\).

\textbf{Solution.} Varying
\(g^{\mu\alpha}g_{\alpha\nu}=\delta^\mu{}_\nu\) gives

\begin{center}
\begingroup
\sbox{\grmathbox}{$\displaystyle 
\delta g_{\mu\nu}=-g_{\mu\alpha}g_{\nu\beta}\delta g^{\alpha\beta}.
$}
\typeout{GRDISPLAY|632|\the\wd\grmathbox|\the\linewidth}
\ifdim\wd\grmathbox>\linewidth
\resizebox{\linewidth}{!}{\usebox{\grmathbox}}
\else\usebox{\grmathbox}\fi
\endgroup
\end{center}

Then

\begin{center}
\begingroup
\sbox{\grmathbox}{$\displaystyle 
\delta\sqrt{-g}
=\frac12\sqrt{-g}\,g^{\mu\nu}\delta g_{\mu\nu}
=-\frac12\sqrt{-g}\,g_{\alpha\beta}\delta g^{\alpha\beta}.
$}
\typeout{GRDISPLAY|633|\the\wd\grmathbox|\the\linewidth}
\ifdim\wd\grmathbox>\linewidth
\resizebox{\linewidth}{!}{\usebox{\grmathbox}}
\else\usebox{\grmathbox}\fi
\endgroup
\end{center}

The sign is a consequence of inverse-matrix variation. It is not an
arbitrary special rule for gravity.

\hypertarget{a.14-a-constant-lagrangian-can-affect-gravity}{%
\section{A.14 A constant Lagrangian can affect
gravity}\label{a.14-a-constant-lagrangian-can-affect-gravity}}

\textbf{Problem.} In natural units \(c=1\), shift the matter Lagrangian
by a constant: \(\mathcal L_m\mapsto\mathcal L_m-V_0\). Find the stress
tensor change.

\textbf{Solution.} The action changes by \(-V_0\int\sqrt{-g}\,d^4x\).
Its variation is

\begin{center}
\begingroup
\sbox{\grmathbox}{$\displaystyle 
\delta S_{\mathrm{shift}}
=\frac12\int\sqrt{-g}\,V_0g_{\mu\nu}\delta g^{\mu\nu}\,d^4x.
$}
\typeout{GRDISPLAY|634|\the\wd\grmathbox|\the\linewidth}
\ifdim\wd\grmathbox>\linewidth
\resizebox{\linewidth}{!}{\usebox{\grmathbox}}
\else\usebox{\grmathbox}\fi
\endgroup
\end{center}

Comparing with
\(\delta S_m=-\tfrac12\int\sqrt{-g}\,T_{\mu\nu}\delta g^{\mu\nu}\,d^4x\)
yields

\begin{center}
\begingroup
\sbox{\grmathbox}{$\displaystyle 
\Delta T_{\mu\nu}=-V_0g_{\mu\nu}.
$}
\typeout{GRDISPLAY|635|\the\wd\grmathbox|\the\linewidth}
\ifdim\wd\grmathbox>\linewidth
\resizebox{\linewidth}{!}{\usebox{\grmathbox}}
\else\usebox{\grmathbox}\fi
\endgroup
\end{center}

Every timelike observer measures energy density \(V_0\) and isotropic
pressure \(-V_0\); this stress tensor selects no preferred rest frame. A
constant can be dynamically irrelevant to nongravitational equations on
a fixed background while affecting the metric equation when geometry is
dynamical.

\hypertarget{a.15-why-a-boundary-condition-on-the-metric-is-not-automatically-enough}{%
\section{A.15 Why a boundary condition on the metric is not
automatically
enough}\label{a.15-why-a-boundary-condition-on-the-metric-is-not-automatically-enough}}

\textbf{Problem.} Suppose \(\delta g_{\mu\nu}=0\) on a boundary. Must
the normal derivative of \(\delta g_{\mu\nu}\) vanish there too?

\textbf{Solution.} No. The elementary function \(f(x)=x\) satisfies
\(f(0)=0\) but \(f'(0)=1\). The same distinction applies componentwise
to a field variation. The Einstein-Hilbert bulk action produces boundary
terms containing derivatives of the metric variation. For a suitable
fixed-induced-metric problem on a non-null smooth boundary, the
Gibbons-Hawking-York term supplies the required cancellation. Setting
more boundary data to zero than the problem calls for can conceal rather
than solve this issue.

\hypertarget{a.16-a-killing-vector-manufactures-a-conserved-current}{%
\section{A.16 A Killing vector manufactures a conserved
current}\label{a.16-a-killing-vector-manufactures-a-conserved-current}}

\textbf{Problem.} Let \(T^{\mu\nu}\) be symmetric and covariantly
conserved. If \(\xi^\mu\) obeys \(\nabla_{(\mu}\xi_{\nu)}=0\), prove
that \(J^\mu=T^{\mu\nu}\xi_\nu\) is conserved.

\textbf{Solution.} The product rule gives

\begin{center}
\begingroup
\sbox{\grmathbox}{$\displaystyle 
\nabla_\mu J^\mu
=(\nabla_\mu T^{\mu\nu})\xi_\nu
+T^{\mu\nu}\nabla_\mu\xi_\nu.
$}
\typeout{GRDISPLAY|636|\the\wd\grmathbox|\the\linewidth}
\ifdim\wd\grmathbox>\linewidth
\resizebox{\linewidth}{!}{\usebox{\grmathbox}}
\else\usebox{\grmathbox}\fi
\endgroup
\end{center}

The first term vanishes. A symmetric tensor contracts only the symmetric
part of the second factor, so the second becomes
\(T^{\mu\nu}\nabla_{(\mu}\xi_{\nu)}=0\). Integrating this local
conservation law into a conserved charge also requires an appropriate
region, hypersurfaces, and control of boundary fluxes. The Killing field
supplies an actual symmetry, not merely a preferred-looking coordinate
label.

\hypertarget{a.17-a-height-difference-changes-a-clocks-rate}{%
\section{A.17 A height difference changes a clock\textquotesingle s
rate}\label{a.17-a-height-difference-changes-a-clocks-rate}}

\textbf{Problem.} Estimate the fractional clock-rate difference between
stationary clocks separated vertically by \(100\,\mathrm m\) in a weak
approximately uniform field with
\(g_{\mathrm{local}}=9.8\,\mathrm{m/s^2}\). Estimate the accumulated
difference over one day.

\textbf{Solution.} The potential difference is
\(\Delta\Phi\simeq g_{\mathrm{local}}h=980\,\mathrm{m^2/s^2}\). Thus

\begin{center}
\begingroup
\sbox{\grmathbox}{$\displaystyle 
\frac{\Delta(d\tau/dt)}{d\tau/dt}
\simeq\frac{\Delta\Phi}{c^2}
\simeq1.09\times10^{-14}.
$}
\typeout{GRDISPLAY|637|\the\wd\grmathbox|\the\linewidth}
\ifdim\wd\grmathbox>\linewidth
\resizebox{\linewidth}{!}{\usebox{\grmathbox}}
\else\usebox{\grmathbox}\fi
\endgroup
\end{center}

Over \(86{,}400\) seconds, the higher clock gains approximately
\(9.4\times10^{-10}\) seconds, or \(0.94\) nanoseconds. This estimate
isolates the gravitational height effect; Earth\textquotesingle s
rotation, the actual potential, motion, and the comparison protocol
matter in a precision experiment.

\hypertarget{a.18-a-horizon-is-not-a-place-where-every-clock-stops}{%
\section{A.18 A horizon is not a place where every clock
stops}\label{a.18-a-horizon-is-not-a-place-where-every-clock-stops}}

\textbf{Problem.} In Schwarzschild spacetime, a static clock obeys
\(d\tau=\sqrt{1-2m/r}\,dt\). Why can\textquotesingle t we infer that a
freely falling clock physically stops at \(r=2m\)?

\textbf{Solution.} The formula applies to worldlines with fixed
Schwarzschild \(r,\theta,\phi\). Such worldlines require increasing
proper acceleration as \(r\) approaches \(2m\) from outside. The
limiting static worldline is not a timelike observer sitting on the
horizon. A falling observer follows a different worldline, has a finite
proper-time crossing for ordinary infall, and can use a regular
horizon-crossing chart. A relation between one coordinate time and one
family of clocks is not a universal claim about all clocks.

\hypertarget{a.19-doubling-a-black-hole-changes-horizon-tides}{%
\section{A.19 Doubling a black hole changes horizon
tides}\label{a.19-doubling-a-black-hole-changes-horizon-tides}}

\textbf{Problem.} The Schwarzschild Kretschmann scalar is
\(K=48m^2/r^6\), with \(m=G_NM/c^2\). Evaluate it at \(r=2m\) and
determine its mass scaling.

\textbf{Solution.} At the horizon,

\begin{center}
\begingroup
\sbox{\grmathbox}{$\displaystyle 
K_H=\frac{48m^2}{64m^6}=\frac{3}{4m^4}\propto M^{-4}.
$}
\typeout{GRDISPLAY|638|\the\wd\grmathbox|\the\linewidth}
\ifdim\wd\grmathbox>\linewidth
\resizebox{\linewidth}{!}{\usebox{\grmathbox}}
\else\usebox{\grmathbox}\fi
\endgroup
\end{center}

Representative orthonormal curvature components scale as
\(m/r^3\sim M^{-2}\) at the horizon. Larger nonrotating black holes can
therefore have weaker horizon-scale tidal curvature. ``Bigger black
hole'' does not mean ``more violent local horizon crossing.'' This says
nothing by itself about the singularity deeper inside.

\hypertarget{a.20-gravitational-wave-strain-is-a-relative-measurement}{%
\section{A.20 Gravitational-wave strain is a relative
measurement}\label{a.20-gravitational-wave-strain-is-a-relative-measurement}}

\textbf{Problem.} A plus-polarized wave has \(h_+=10^{-21}\). For an
ideal freely falling separation \(L=4\,\mathrm{km}\) along one principal
axis, estimate the leading displacement amplitude. What is the ideal
differential change between orthogonal axes?

\textbf{Solution.} To leading order in the long-wavelength
approximation,

\begin{center}
\begingroup
\sbox{\grmathbox}{$\displaystyle 
\delta L_x\simeq\frac12h_+L=2\times10^{-18}\,\mathrm m,
\qquad
\delta L_y\simeq-\frac12h_+L.
$}
\typeout{GRDISPLAY|639|\the\wd\grmathbox|\the\linewidth}
\ifdim\wd\grmathbox>\linewidth
\resizebox{\linewidth}{!}{\usebox{\grmathbox}}
\else\usebox{\grmathbox}\fi
\endgroup
\end{center}

The differential change is
\(\delta L_x-\delta L_y\simeq h_+L=4\times10^{-18}\,\mathrm m\). A real
interferometer measures optical phase with a frequency-dependent
response; this elementary estimate captures the geometric strain scale,
not the full instrument transfer function.

\hypertarget{a.21-why-an-isolated-source-has-no-leading-mass-dipole-gravitational-radiation}{%
\section{A.21 Why an isolated source has no leading mass-dipole
gravitational
radiation}\label{a.21-why-an-isolated-source-has-no-leading-mass-dipole-gravitational-radiation}}

\textbf{Problem.} In the slow-motion weak-field approximation, define
\(D^i=\int\rho x^i\,d^3x\). Explain why its second time derivative
vanishes for an isolated system and what that implies.

\textbf{Solution.} The first derivative is total momentum,
\(\dot D^i=P^i\), assuming suitable mass conservation and vanishing
boundary flux. For an isolated system \(\dot P^i=0\), hence
\(\ddot D^i=0\). A putative leading radiative term based on that second
derivative cannot carry changing dipolar structure. Together with
conservation of the monopole at the relevant leading order, this helps
explain why the mass quadrupole is the first generic radiative
contribution in GR\textquotesingle s slow-source expansion. The claim
has an approximation and isolation regime; it is not a statement about
arbitrary additional fields in modified gravity.

\hypertarget{a.22-an-expanding-box-gives-the-fluid-equation}{%
\section{A.22 An expanding box gives the fluid
equation}\label{a.22-an-expanding-box-gives-the-fluid-equation}}

\textbf{Problem.} For a homogeneous perfect fluid in a comoving cell
with physical volume \(V\propto a^3\), use \(d(\epsilon V)=-p\,dV\) to
derive the continuity equation and the scaling for \(p=w\epsilon\), with
constant \(w\).

\textbf{Solution.} Expand the differential:

\begin{center}
\begingroup
\sbox{\grmathbox}{$\displaystyle 
V\,d\epsilon+\epsilon\,dV=-p\,dV.
$}
\typeout{GRDISPLAY|640|\the\wd\grmathbox|\the\linewidth}
\ifdim\wd\grmathbox>\linewidth
\resizebox{\linewidth}{!}{\usebox{\grmathbox}}
\else\usebox{\grmathbox}\fi
\endgroup
\end{center}

Divide by \(Vdt\) and use \(\dot V/V=3H\):

\begin{center}
\begingroup
\sbox{\grmathbox}{$\displaystyle 
\dot\epsilon+3H(\epsilon+p)=0.
$}
\typeout{GRDISPLAY|641|\the\wd\grmathbox|\the\linewidth}
\ifdim\wd\grmathbox>\linewidth
\resizebox{\linewidth}{!}{\usebox{\grmathbox}}
\else\usebox{\grmathbox}\fi
\endgroup
\end{center}

For constant \(w\), \(d\ln\epsilon=-3(1+w)d\ln a\), giving
\(\epsilon\propto a^{-3(1+w)}\). Dust has \(a^{-3}\), radiation
\(a^{-4}\), and a cosmological-constant fluid \(a^0\). This equation is
a local covariant balance law specialized to the symmetry, not evidence
for a universally defined conserved total cosmic energy.

\hypertarget{a.23-derive-a-universes-power-law-growth}{%
\section{A.23 Derive a universe\textquotesingle s power-law
growth}\label{a.23-derive-a-universes-power-law-growth}}

\textbf{Problem.} In a spatially flat, single-fluid universe with
\(\Lambda=0\) and constant \(w>-1\), combine \(H^2\propto\epsilon\) with
the previous result to derive \(a(t)\).

\textbf{Solution.} We have

\begin{center}
\begingroup
\sbox{\grmathbox}{$\displaystyle 
\left(\frac{\dot a}{a}\right)^2\propto a^{-3(1+w)},
\qquad
\dot a\propto a^{-(1+3w)/2}.
$}
\typeout{GRDISPLAY|642|\the\wd\grmathbox|\the\linewidth}
\ifdim\wd\grmathbox>\linewidth
\resizebox{\linewidth}{!}{\usebox{\grmathbox}}
\else\usebox{\grmathbox}\fi
\endgroup
\end{center}

Integrating the expanding branch gives \(a^{3(1+w)/2}\propto t-t_0\),
hence

\begin{center}
\begingroup
\sbox{\grmathbox}{$\displaystyle 
a(t)\propto(t-t_0)^{2/[3(1+w)]}.
$}
\typeout{GRDISPLAY|643|\the\wd\grmathbox|\the\linewidth}
\ifdim\wd\grmathbox>\linewidth
\resizebox{\linewidth}{!}{\usebox{\grmathbox}}
\else\usebox{\grmathbox}\fi
\endgroup
\end{center}

Dust gives \(t^{2/3}\), radiation \(t^{1/2}\). The \(w=-1\) case must be
solved separately: constant positive energy density gives constant \(H\)
and exponential expansion. Substituting \(w=-1\) into the power-law
exponent is not a valid limiting derivation.

\hypertarget{a.24-when-does-pressure-accelerate-expansion}{%
\section{A.24 When does pressure accelerate
expansion?}\label{a.24-when-does-pressure-accelerate-expansion}}

\textbf{Problem.} With no separately written cosmological constant and
\(p=w\epsilon\), determine when a positive-density FLRW fluid drives
\(\ddot a>0\).

\textbf{Solution.} The acceleration equation is

\begin{center}
\begingroup
\sbox{\grmathbox}{$\displaystyle 
\frac{\ddot a}{a}
=-\frac{4\pi G_N}{3c^2}(\epsilon+3p)
=-\frac{4\pi G_N\epsilon}{3c^2}(1+3w).
$}
\typeout{GRDISPLAY|644|\the\wd\grmathbox|\the\linewidth}
\ifdim\wd\grmathbox>\linewidth
\resizebox{\linewidth}{!}{\usebox{\grmathbox}}
\else\usebox{\grmathbox}\fi
\endgroup
\end{center}

For \(\epsilon>0\), acceleration requires \(w<-1/3\). Ordinary positive
pressure contributes toward deceleration in this equation. Sufficiently
negative pressure reverses the sign. This cosmological statement must
not be replaced by the indiscriminate slogan ``pressure is repulsive''
or ``pressure is always attractive.''

\hypertarget{a.25-four-constraints-do-not-mean-four-lost-metric-components}{%
\section{A.25 Four constraints do not mean four lost metric
components}\label{a.25-four-constraints-do-not-mean-four-lost-metric-components}}

\textbf{Problem.} Explain why GR has two local propagating metric
degrees of freedom in four spacetime dimensions using ADM phase space.

\textbf{Solution.} The spatial metric has six independent components,
and its conjugate momentum adds six, for twelve phase-space variables
per spatial point. Four first-class constraints each remove one
constrained phase-space direction and one associated gauge direction.
Thus \(12-2\times4=4\) physical phase-space dimensions remain,
corresponding to two configuration degrees of freedom and their
conjugate momenta. Lapse and shift act as gauge multipliers rather than
additional propagating fields in this count. This is a local count for
ordinary GR; boundaries and global sectors require further care.

\hypertarget{a.26-the-adm-constraint-recognizes-an-expanding-universe}{%
\section{A.26 The ADM constraint recognizes an expanding
universe}\label{a.26-the-adm-constraint-recognizes-an-expanding-universe}}

\textbf{Problem.} Use \(c=1\), a flat FLRW slice, and the convention
\(K_{ij}=-\tfrac12\mathcal L_n\gamma_{ij}\). Show that the Hamiltonian
constraint reproduces the first Friedmann equation.

\textbf{Solution.} For comoving proper-time slicing,
\(K_{ij}=-H\gamma_{ij}\). Hence \(K=-3H\) and \(K_{ij}K^{ij}=3H^2\).
Since the slice has \({}^{(3)}R=0\),

\begin{center}
\begingroup
\sbox{\grmathbox}{$\displaystyle 
{}^{(3)}R+K^2-K_{ij}K^{ij}=6H^2.
$}
\typeout{GRDISPLAY|645|\the\wd\grmathbox|\the\linewidth}
\ifdim\wd\grmathbox>\linewidth
\resizebox{\linewidth}{!}{\usebox{\grmathbox}}
\else\usebox{\grmathbox}\fi
\endgroup
\end{center}

Set this equal to \(16\pi G_N\epsilon+2\Lambda\) to obtain

\begin{center}
\begingroup
\sbox{\grmathbox}{$\displaystyle 
H^2=\frac{8\pi G_N}{3}\epsilon+\frac{\Lambda}{3}.
$}
\typeout{GRDISPLAY|646|\the\wd\grmathbox|\the\linewidth}
\ifdim\wd\grmathbox>\linewidth
\resizebox{\linewidth}{!}{\usebox{\grmathbox}}
\else\usebox{\grmathbox}\fi
\endgroup
\end{center}

The expansion lives in the extrinsic curvature of the slices even when
their intrinsic spatial curvature vanishes.

\hypertarget{a.27-raychaudhuri-gives-a-deadline-for-focusing}{%
\section{A.27 Raychaudhuri gives a deadline for
focusing}\label{a.27-raychaudhuri-gives-a-deadline-for-focusing}}

\textbf{Problem.} For a geodesic, hypersurface-orthogonal timelike
congruence in four dimensions with \(R_{\mu\nu}u^\mu u^\nu\ge0\), use
\(c=1\) and show that negative initial expansion \(\theta_0<0\) must
diverge to negative infinity within proper time at most
\(3/|\theta_0|\), provided the congruence remains defined up to that
point.

\textbf{Solution.} Raychaudhuri gives

\begin{center}
\begingroup
\sbox{\grmathbox}{$\displaystyle 
\frac{d\theta}{d\tau}\le-\frac13\theta^2
$}
\typeout{GRDISPLAY|647|\the\wd\grmathbox|\the\linewidth}
\ifdim\wd\grmathbox>\linewidth
\resizebox{\linewidth}{!}{\usebox{\grmathbox}}
\else\usebox{\grmathbox}\fi
\endgroup
\end{center}

because shear contributes nonpositively and vorticity vanishes. For
\(\theta<0\),

\begin{center}
\begingroup
\sbox{\grmathbox}{$\displaystyle 
\frac{d}{d\tau}\left(\frac1\theta\right)
=-\frac{\dot\theta}{\theta^2}\ge\frac13.
$}
\typeout{GRDISPLAY|648|\the\wd\grmathbox|\the\linewidth}
\ifdim\wd\grmathbox>\linewidth
\resizebox{\linewidth}{!}{\usebox{\grmathbox}}
\else\usebox{\grmathbox}\fi
\endgroup
\end{center}

Starting from \(1/\theta_0<0\), the reciprocal must reach zero no later
than \(3/|\theta_0|\), corresponding to focusing. A caustic is not by
itself a spacetime singularity: geodesics can cross in perfectly regular
spacetime. Singularity theorems require additional global assumptions to
infer incompleteness.

\hypertarget{a.28-a-pure-frame-rotation-is-not-curvature}{%
\section{A.28 A pure frame rotation is not
curvature}\label{a.28-a-pure-frame-rotation-is-not-curvature}}

\textbf{Problem.} On the flat plane, use the orthonormal coframe
\(e^1=dr\), \(e^2=r\,d\theta\). Why can the spin connection be nonzero
while the curvature two-form vanishes?

\textbf{Solution.} Since \(de^2=dr\wedge d\theta\), the torsion-free
Cartan equation requires \(\omega^2{}_1=d\theta\) and
\(\omega^1{}_2=-d\theta\). These one-forms record the rotation of the
polar frame. But \(d(d\theta)=0\) on a regular angular chart, and the
relevant connection wedge products vanish in this two-dimensional
example. Thus
\(\Omega^a{}_b=d\omega^a{}_b+\omega^a{}_c\wedge\omega^c{}_b=0\). The
polar chart and frame fail at the origin; their behavior there must not
be confused with a curved plane.

\hypertarget{a.29-why-black-hole-entropy-has-the-right-dimensions}{%
\section{A.29 Why black-hole entropy has the right
dimensions}\label{a.29-why-black-hole-entropy-has-the-right-dimensions}}

\textbf{Problem.} Verify that \(S_{\mathrm{BH}}=k_BA/(4\ell_P^2)\) has
entropy units, where \(\ell_P^2=\hbar G_N/c^3\). For a Schwarzschild
black hole, derive the \(M^2\) dependence of entropy and check that
\(T_HdS=d(Mc^2)\).

\textbf{Solution.} \(A/\ell_P^2\) is dimensionless, so \(S\) has units
of \(k_B\). With \(A=16\pi G_N^2M^2/c^4\),

\begin{center}
\begingroup
\sbox{\grmathbox}{$\displaystyle 
S_{\mathrm{BH}}=\frac{4\pi k_BG_NM^2}{\hbar c},
\qquad
\frac{dS_{\mathrm{BH}}}{dM}
=\frac{8\pi k_BG_NM}{\hbar c}.
$}
\typeout{GRDISPLAY|649|\the\wd\grmathbox|\the\linewidth}
\ifdim\wd\grmathbox>\linewidth
\resizebox{\linewidth}{!}{\usebox{\grmathbox}}
\else\usebox{\grmathbox}\fi
\endgroup
\end{center}

Multiply by \(T_H=\hbar c^3/(8\pi G_NMk_B)\) to obtain \(T_HdS=c^2dM\).
This is a consistency check within the semiclassical result for an
uncharged, nonrotating hole, not a derivation of its microscopic degrees
of freedom.

\hypertarget{a.30-the-effective-theory-expansion-has-a-speed-limit}{%
\section{A.30 The effective-theory expansion has a speed
limit}\label{a.30-the-effective-theory-expansion-has-a-speed-limit}}

\textbf{Problem.} In units \(c=\hbar=1\), consider a schematic
gravitational Lagrangian

\begin{center}
\begingroup
\sbox{\grmathbox}{$\displaystyle 
\mathcal L\sim M_*^2\left[R+\frac{a}{M_*^2}R^2+\cdots\right],
$}
\typeout{GRDISPLAY|650|\the\wd\grmathbox|\the\linewidth}
\ifdim\wd\grmathbox>\linewidth
\resizebox{\linewidth}{!}{\usebox{\grmathbox}}
\else\usebox{\grmathbox}\fi
\endgroup
\end{center}

with dimensionless coefficient \(a\) of order unity. On a slowly varying
geometry with typical curvature scale \(R\sim L^{-2}\), estimate the
relative size of the correction. State why \(R=0\) alone is not a
sufficient validity check.

\textbf{Solution.} Relative to the \(R\) term, the displayed correction
scales as \(a/(M_*L)^2\). It is small when \(M_*L\gg1\). But an actual
effective action contains other independent curvature contractions and
derivative operators. The scalar \(R\) may vanish while Riemann or Weyl
curvature is nonzero, as in a Schwarzschild exterior. Validity requires
control of the physically relevant curvature components, invariant
scales, frequencies, and state-dependent effects, not a single
convenient scalar. The effective theory can be predictive below its
cutoff without claiming validity at arbitrarily short distances.

\hypertarget{appendix-b}{%
\chapter{Appendix B. A working reference sheet}\label{appendix-b}}

A reference sheet is useful once an equation has a meaning. Before that,
it is a decorative wall of Greek letters. Use this appendix as a
retrieval aid, and return to the associated chapter whenever a formula
feels suspiciously effortless.

\hypertarget{b.1-the-geometry-ladder}{%
\section{B.1 The geometry ladder}\label{b.1-the-geometry-ladder}}

\textbf{Interval and clock time}

\begin{center}
\begingroup
\sbox{\grmathbox}{$\displaystyle 
ds^2=g_{\mu\nu}dx^\mu dx^\nu,
\qquad d\tau^2=-ds^2/c^2\quad\text{on a timelike worldline}.
$}
\typeout{GRDISPLAY|651|\the\wd\grmathbox|\the\linewidth}
\ifdim\wd\grmathbox>\linewidth
\resizebox{\linewidth}{!}{\usebox{\grmathbox}}
\else\usebox{\grmathbox}\fi
\endgroup
\end{center}

\textbf{Inverse and index conversion}

\begin{center}
\begingroup
\sbox{\grmathbox}{$\displaystyle 
g^{\mu\alpha}g_{\alpha\nu}=\delta^\mu{}_\nu,
\qquad V_\mu=g_{\mu\nu}V^\nu,
\qquad V^\mu=g^{\mu\nu}V_\nu.
$}
\typeout{GRDISPLAY|652|\the\wd\grmathbox|\the\linewidth}
\ifdim\wd\grmathbox>\linewidth
\resizebox{\linewidth}{!}{\usebox{\grmathbox}}
\else\usebox{\grmathbox}\fi
\endgroup
\end{center}

\textbf{Levi-Civita connection in a coordinate basis}

\begin{center}
\begingroup
\sbox{\grmathbox}{$\displaystyle 
\Gamma^\rho{}_{\mu\nu}
=\frac12g^{\rho\sigma}
\left(\partial_\mu g_{\sigma\nu}
+\partial_\nu g_{\sigma\mu}
-\partial_\sigma g_{\mu\nu}\right).
$}
\typeout{GRDISPLAY|653|\the\wd\grmathbox|\the\linewidth}
\ifdim\wd\grmathbox>\linewidth
\resizebox{\linewidth}{!}{\usebox{\grmathbox}}
\else\usebox{\grmathbox}\fi
\endgroup
\end{center}

\textbf{Covariant derivatives}

\begin{center}
\begingroup
\sbox{\grmathbox}{$\displaystyle 
\nabla_\mu f=\partial_\mu f,
\qquad
\nabla_\mu V^\nu=\partial_\mu V^\nu
+\Gamma^\nu{}_{\mu\rho}V^\rho,
$}
\typeout{GRDISPLAY|654|\the\wd\grmathbox|\the\linewidth}
\ifdim\wd\grmathbox>\linewidth
\resizebox{\linewidth}{!}{\usebox{\grmathbox}}
\else\usebox{\grmathbox}\fi
\endgroup
\end{center}

\begin{center}
\begingroup
\sbox{\grmathbox}{$\displaystyle 
\nabla_\mu\omega_\nu=\partial_\mu\omega_\nu
-\Gamma^\rho{}_{\mu\nu}\omega_\rho.
$}
\typeout{GRDISPLAY|655|\the\wd\grmathbox|\the\linewidth}
\ifdim\wd\grmathbox>\linewidth
\resizebox{\linewidth}{!}{\usebox{\grmathbox}}
\else\usebox{\grmathbox}\fi
\endgroup
\end{center}

An upper tensor index gets a plus connection term; a lower index gets a
minus term. This rule applies to ordinary tensors; tensor densities have
an additional weight term.

\textbf{Geodesic with affine parameter \(\lambda\)}

\begin{center}
\begingroup
\sbox{\grmathbox}{$\displaystyle 
\frac{d^2x^\mu}{d\lambda^2}
+\Gamma^\mu{}_{\alpha\beta}
\frac{dx^\alpha}{d\lambda}
\frac{dx^\beta}{d\lambda}=0.
$}
\typeout{GRDISPLAY|656|\the\wd\grmathbox|\the\linewidth}
\ifdim\wd\grmathbox>\linewidth
\resizebox{\linewidth}{!}{\usebox{\grmathbox}}
\else\usebox{\grmathbox}\fi
\endgroup
\end{center}

Proper time is affine on a timelike geodesic. Null geodesics have no
proper-time parameter.

\textbf{Riemann curvature}

\begin{center}
\begingroup
\sbox{\grmathbox}{$\displaystyle 
R^\rho{}_{\sigma\mu\nu}
=\partial_\mu\Gamma^\rho{}_{\nu\sigma}
-\partial_\nu\Gamma^\rho{}_{\mu\sigma}
+\Gamma^\rho{}_{\mu\lambda}\Gamma^\lambda{}_{\nu\sigma}
-\Gamma^\rho{}_{\nu\lambda}\Gamma^\lambda{}_{\mu\sigma}.
$}
\typeout{GRDISPLAY|657|\the\wd\grmathbox|\the\linewidth}
\ifdim\wd\grmathbox>\linewidth
\resizebox{\linewidth}{!}{\usebox{\grmathbox}}
\else\usebox{\grmathbox}\fi
\endgroup
\end{center}

\textbf{Contractions}

\begin{center}
\begingroup
\sbox{\grmathbox}{$\displaystyle 
R_{\mu\nu}=R^\rho{}_{\mu\rho\nu},
\qquad R=g^{\mu\nu}R_{\mu\nu},
\qquad G_{\mu\nu}=R_{\mu\nu}-\frac12Rg_{\mu\nu}.
$}
\typeout{GRDISPLAY|658|\the\wd\grmathbox|\the\linewidth}
\ifdim\wd\grmathbox>\linewidth
\resizebox{\linewidth}{!}{\usebox{\grmathbox}}
\else\usebox{\grmathbox}\fi
\endgroup
\end{center}

\textbf{Tidal acceleration}

\begin{center}
\begingroup
\sbox{\grmathbox}{$\displaystyle 
\frac{D^2\xi^\mu}{d\tau^2}
=-R^\mu{}_{\alpha\nu\beta}u^\alpha\xi^\nu u^\beta.
$}
\typeout{GRDISPLAY|659|\the\wd\grmathbox|\the\linewidth}
\ifdim\wd\grmathbox>\linewidth
\resizebox{\linewidth}{!}{\usebox{\grmathbox}}
\else\usebox{\grmathbox}\fi
\endgroup
\end{center}

This form assumes a neighboring family of affinely parameterized
geodesics with the usual commuting tangent and separation fields.

\textbf{Volume, divergence, scalar wave operator}

\begin{center}
\begingroup
\sbox{\grmathbox}{$\displaystyle 
dV_4=\sqrt{-g}\,d^4x,
\qquad
\nabla_\mu V^\mu=
\frac1{\sqrt{-g}}\partial_\mu(\sqrt{-g}V^\mu),
$}
\typeout{GRDISPLAY|660|\the\wd\grmathbox|\the\linewidth}
\ifdim\wd\grmathbox>\linewidth
\resizebox{\linewidth}{!}{\usebox{\grmathbox}}
\else\usebox{\grmathbox}\fi
\endgroup
\end{center}

\begin{center}
\begingroup
\sbox{\grmathbox}{$\displaystyle 
\Box f=\nabla_\mu\nabla^\mu f
=\frac1{\sqrt{-g}}\partial_\mu
\left(\sqrt{-g}g^{\mu\nu}\partial_\nu f\right).
$}
\typeout{GRDISPLAY|661|\the\wd\grmathbox|\the\linewidth}
\ifdim\wd\grmathbox>\linewidth
\resizebox{\linewidth}{!}{\usebox{\grmathbox}}
\else\usebox{\grmathbox}\fi
\endgroup
\end{center}

In flat inertial coordinates, \(\Box=-c^{-2}\partial_t^2+\nabla^2\).

\hypertarget{b.2-matter-action-and-dynamics}{%
\section{B.2 Matter, action, and
dynamics}\label{b.2-matter-action-and-dynamics}}

\textbf{Observer energy}

\begin{center}
\begingroup
\sbox{\grmathbox}{$\displaystyle 
E_{(U)}=-p_\mu U^\mu,
\qquad U^\mu U_\mu=-c^2.
$}
\typeout{GRDISPLAY|662|\the\wd\grmathbox|\the\linewidth}
\ifdim\wd\grmathbox>\linewidth
\resizebox{\linewidth}{!}{\usebox{\grmathbox}}
\else\usebox{\grmathbox}\fi
\endgroup
\end{center}

\textbf{Perfect fluid}

\begin{center}
\begingroup
\sbox{\grmathbox}{$\displaystyle 
T^{\mu\nu}=
\frac{\epsilon+p}{c^2}u^\mu u^\nu+pg^{\mu\nu},
\qquad T=-\epsilon+3p.
$}
\typeout{GRDISPLAY|663|\the\wd\grmathbox|\the\linewidth}
\ifdim\wd\grmathbox>\linewidth
\resizebox{\linewidth}{!}{\usebox{\grmathbox}}
\else\usebox{\grmathbox}\fi
\endgroup
\end{center}

\textbf{Einstein equation and trace reversal in four dimensions}

\begin{center}
\begingroup
\sbox{\grmathbox}{$\displaystyle 
G_{\mu\nu}+\Lambda g_{\mu\nu}=\kappa T_{\mu\nu},
\qquad \kappa=\frac{8\pi G_N}{c^4},
$}
\typeout{GRDISPLAY|664|\the\wd\grmathbox|\the\linewidth}
\ifdim\wd\grmathbox>\linewidth
\resizebox{\linewidth}{!}{\usebox{\grmathbox}}
\else\usebox{\grmathbox}\fi
\endgroup
\end{center}

\begin{center}
\begingroup
\sbox{\grmathbox}{$\displaystyle 
R=4\Lambda-\kappa T,
\qquad
R_{\mu\nu}
=\kappa\left(T_{\mu\nu}-\frac12Tg_{\mu\nu}\right)
+\Lambda g_{\mu\nu}.
$}
\typeout{GRDISPLAY|665|\the\wd\grmathbox|\the\linewidth}
\ifdim\wd\grmathbox>\linewidth
\resizebox{\linewidth}{!}{\usebox{\grmathbox}}
\else\usebox{\grmathbox}\fi
\endgroup
\end{center}

\textbf{Action and metric variation, with \(x^0=ct\)}

\begin{center}
\begingroup
\sbox{\grmathbox}{$\displaystyle 
S_{\mathrm{EH}}=\frac{c^3}{16\pi G_N}
\int d^4x\sqrt{-g}(R-2\Lambda),
$}
\typeout{GRDISPLAY|666|\the\wd\grmathbox|\the\linewidth}
\ifdim\wd\grmathbox>\linewidth
\resizebox{\linewidth}{!}{\usebox{\grmathbox}}
\else\usebox{\grmathbox}\fi
\endgroup
\end{center}

\begin{center}
\begingroup
\sbox{\grmathbox}{$\displaystyle 
\delta\sqrt{-g}=-\frac12\sqrt{-g}g_{\mu\nu}\delta g^{\mu\nu},
\qquad
\delta S_m=-\frac1{2c}\int d^4x\sqrt{-g}
T_{\mu\nu}\delta g^{\mu\nu}.
$}
\typeout{GRDISPLAY|667|\the\wd\grmathbox|\the\linewidth}
\ifdim\wd\grmathbox>\linewidth
\resizebox{\linewidth}{!}{\usebox{\grmathbox}}
\else\usebox{\grmathbox}\fi
\endgroup
\end{center}

The total variational problem also includes its appropriate boundary
terms and boundary conditions. The bulk expression alone is not a
universal boundary prescription.

\textbf{Identities and balance laws}

\begin{center}
\begingroup
\sbox{\grmathbox}{$\displaystyle 
\nabla_\mu G^{\mu\nu}=0,
\qquad
\nabla_\mu T^{\mu\nu}=0
$}
\typeout{GRDISPLAY|668|\the\wd\grmathbox|\the\linewidth}
\ifdim\wd\grmathbox>\linewidth
\resizebox{\linewidth}{!}{\usebox{\grmathbox}}
\else\usebox{\grmathbox}\fi
\endgroup
\end{center}

for matter compatible with the Einstein equation and constant
\(\Lambda\). The first is a geometric identity; the second is an
on-shell matter balance law in the standard coupled theory. Neither is
the assertion that all components are constant.

\textbf{Weak Newtonian limit}

\begin{center}
\begingroup
\sbox{\grmathbox}{$\displaystyle 
g_{00}\simeq-\left(1+\frac{2\Phi}{c^2}\right),
\qquad
\frac{d^2\mathbf{x}}{dt^2}\simeq-\symbfit\nabla\Phi,
\qquad
\nabla^2\Phi=4\pi G_N\rho
$}
\typeout{GRDISPLAY|669|\the\wd\grmathbox|\the\linewidth}
\ifdim\wd\grmathbox>\linewidth
\resizebox{\linewidth}{!}{\usebox{\grmathbox}}
\else\usebox{\grmathbox}\fi
\endgroup
\end{center}

for slow test motion, weak approximately static fields, negligible
pressure, and negligible \(\Lambda\) on the scales considered.

\hypertarget{b.3-vacuum-is-a-family-of-possibilities}{%
\section{B.3 Vacuum is a family of
possibilities}\label{b.3-vacuum-is-a-family-of-possibilities}}

\begingroup\small\setlength{\tabcolsep}{5pt}\renewcommand{\arraystretch}{1.13}

\begin{longtable}[]{@{}
  >{\raggedright\arraybackslash}p{(\columnwidth - 4\tabcolsep) * \real{0.2500}}
  >{\raggedright\arraybackslash}p{(\columnwidth - 4\tabcolsep) * \real{0.3600}}
  >{\raggedright\arraybackslash}p{(\columnwidth - 4\tabcolsep) * \real{0.3900}}@{}}
\toprule\noalign{}
\begin{minipage}[b]{\linewidth}\raggedright
Condition
\end{minipage} & \begin{minipage}[b]{\linewidth}\raggedright
Consequence in four-dimensional GR
\end{minipage} & \begin{minipage}[b]{\linewidth}\raggedright
What is still possible
\end{minipage} \\
\midrule\noalign{}
\endhead
\bottomrule\noalign{}
\endlastfoot
\(T_{\mu\nu}=0\), \(\Lambda=0\) & \(R_{\mu\nu}=0\), \(R=0\) & Weyl
curvature, black-hole exterior tides, gravitational waves \\
\(T_{\mu\nu}=0\), \(\Lambda\ne0\) & \(R_{\mu\nu}=\Lambda g_{\mu\nu}\),
\(R=4\Lambda\) & Additional Weyl curvature; maximal symmetry is an extra
condition \\
\(T=0\), \(\Lambda=0\) & \(R=0\) & Nonzero Ricci from trace-free matter,
such as a classical Maxwell field \\
\(R^\rho{}_{\sigma\mu\nu}=0\) on an open region & Locally flat geometry
& Curvilinear/accelerated coordinates; possible global topology beyond a
local chart \\
\end{longtable}

\endgroup

\hypertarget{b.4-six-distinctions-to-keep-beside-your-notebook}{%
\section{B.4 Six distinctions to keep beside your
notebook}\label{b.4-six-distinctions-to-keep-beside-your-notebook}}

\begingroup\small\setlength{\tabcolsep}{5pt}\renewcommand{\arraystretch}{1.13}

\begin{longtable}[]{@{}
  >{\raggedright\arraybackslash}p{(\columnwidth - 2\tabcolsep) * \real{0.3000}}
  >{\raggedright\arraybackslash}p{(\columnwidth - 2\tabcolsep) * \real{0.7000}}@{}}
\toprule\noalign{}
\begin{minipage}[b]{\linewidth}\raggedright
Do not merge
\end{minipage} & \begin{minipage}[b]{\linewidth}\raggedright
The distinction
\end{minipage} \\
\midrule\noalign{}
\endhead
\bottomrule\noalign{}
\endlastfoot
Coordinates and observers & A chart labels events; an observer is a
physical timelike worldline, with additional frame choices for
measurements \\
Connection and curvature & The connection compares nearby tangent
spaces; curvature measures a particular failure of path-independent
transport \\
Spatial curvature and spacetime curvature & Spatial slices depend on a
slicing; the four-dimensional curvature describes spacetime \\
Local conservation and global conserved energy & A covariant balance law
does not automatically provide a global time-translation symmetry or a
total energy \\
Coordinate singularity and geometric singularity & A failing chart can
be replaced; spacetime incompleteness requires a different analysis \\
Mathematical identity and equation of motion & An identity holds for
every field of the relevant class; an equation of motion restricts the
physically allowed fields \\
\end{longtable}

\endgroup

\hypertarget{appendix-c}{%
\chapter{Appendix C. A plain-language glossary}\label{appendix-c}}

\begingroup\small\setlength{\tabcolsep}{5pt}\renewcommand{\arraystretch}{1.13}

\begin{longtable}[]{@{}
  >{\raggedright\arraybackslash}p{(\columnwidth - 2\tabcolsep) * \real{0.3000}}
  >{\raggedright\arraybackslash}p{(\columnwidth - 2\tabcolsep) * \real{0.7000}}@{}}
\toprule\noalign{}
\begin{minipage}[b]{\linewidth}\raggedright
Term
\end{minipage} & \begin{minipage}[b]{\linewidth}\raggedright
Meaning
\end{minipage} \\
\midrule\noalign{}
\endhead
\bottomrule\noalign{}
\endlastfoot
Affine parameter & A geodesic parameter for which the tangent transports
parallel to itself; linear rescalings preserve this property \\
Atlas & A collection of overlapping coordinate charts covering a
manifold \\
Boundary condition & A restriction on fields or their behavior at a
boundary or asymptotic region \\
Causal curve & A timelike or null curve; with a chosen time orientation,
it can represent future-directed causal influence \\
Chart & A smooth local system of coordinates; it need not cover the
whole spacetime \\
Christoffel symbols & Coordinate-basis coefficients of the Levi-Civita
connection; the coefficients alone are not tensor components \\
Connection & A rule for differentiation and infinitesimal parallel
transport; in standard metric GR it is the torsion-free,
metric-compatible one \\
Constraint & An equation restricting permissible initial data rather
than independently specifying all their time derivatives \\
Covector & A linear map from a tangent vector to a scalar; the
differential \(df\), with components \(\partial_\mu f\), is an example.
Obtaining the gradient vector requires raising its index with the
metric \\
Covariant derivative & A derivative that accounts for the connection so
its output transforms as the intended tensor \\
Diffeomorphism & A smooth invertible map with smooth inverse; it
underlies coordinate changes and GR\textquotesingle s field-relabeling
gauge symmetry \\
Effective field theory & An expansion valid below a specified scale,
organized so that finitely many parameters control predictions to a
chosen accuracy \\
Einstein tensor & The divergence-free contraction of curvature that
appears in Einstein\textquotesingle s field equation \\
Equivalence principle & A family of experimentally meaningful statements
about universal free fall and the local behavior of nongravitational
physics; its precise version matters \\
Event horizon & A global causal boundary separating events that can send
signals to the relevant future infinity from those that cannot \\
Extrinsic curvature & How a hypersurface sits within a
higher-dimensional geometry; in a spacetime slicing it encodes important
information about the slices\textquotesingle{} evolution \\
Gauge freedom & Redundancy in a mathematical description; gauge-related
descriptions represent the same physical situation when the
transformations are treated as redundancies with the appropriate
boundary conditions \\
Geodesic & A curve whose tangent is parallel-transported along itself
with an affine parameter \\
Geodesic deviation & The evolution of separation between neighboring
geodesics, governed by curvature \\
Global hyperbolicity & A strong causality condition associated with a
well-defined Cauchy initial-value formulation \\
Holonomy & The net transformation produced by parallel transport around
a closed loop \\
Killing vector & A vector field generating a continuous metric symmetry,
with vanishing Lie derivative of the metric \\
Lapse & In a spacetime slicing, the factor converting coordinate-time
separation into proper separation along the slice normal \\
Lie derivative & Change of a field under the flow generated by a vector
field; it is distinct from covariant differentiation along that
vector \\
Manifold & A space locally describable by ordinary coordinate
neighborhoods with compatible smooth overlaps \\
Metric & A smoothly varying nondegenerate bilinear form defining
intervals, causal character, and index conversion \\
Nonmetricity & Failure of a connection to preserve the metric under
covariant differentiation \\
Null & Having zero spacetime norm; a nonzero null vector is not the zero
vector \\
On shell & Satisfying the equations of motion; off shell means not
imposing them \\
Orthonormal frame & A local basis in which the metric has the standard
diagonal Minkowski form \\
Parallel transport & Moving a vector along a path so its covariant
derivative along that path vanishes \\
Proper acceleration & The invariant magnitude of four-acceleration,
measured by an ideal accelerometer \\
Proper time & The time accumulated by an ideal clock along its timelike
worldline \\
Ricci curvature & A contraction of Riemann curvature, directly tied to
stress-energy through Einstein\textquotesingle s equation \\
Riemann curvature & The tensor encoding local curvature, including tidal
and parallel-transport information \\
Scalar curvature & A further contraction of Ricci; useful but
insufficient to describe all curvature \\
Shift & The tangential part of the coordinate-time flow relative to a
chosen slicing \\
Stress-energy tensor & The local tensor encoding energy, momentum, their
fluxes, and mechanical stresses \\
Tangent space & The vector space of possible tangent directions at one
event \\
Tensor & A multilinear geometric object with components transforming by
the appropriate Jacobian factors \\
Tetrad & A local orthonormal frame or its dual coframe in
four-dimensional spacetime \\
Torsion & The antisymmetric failure of a connection\textquotesingle s
directional derivatives to match the Lie bracket; it vanishes for the
Levi-Civita connection \\
Weyl curvature & The trace-free part of Riemann curvature; in four
dimensions it can remain nonzero in Ricci-flat vacuum \\
Worldline & A path through spacetime representing the history of a
localized object or observer \\
\end{longtable}

\endgroup

\hypertarget{appendix-d}{%
\chapter{Appendix D. Where to go next}\label{appendix-d}}

\hypertarget{d.1-foundational-lecture-notes}{%
\section{D.1 Foundational lecture
notes}\label{d.1-foundational-lecture-notes}}

Two substantial, freely accessible courses make good companions after
this tutorial:

\begin{itemize}
\tightlist
\item
  \href{https://davidtong.org/teaching/general-relativity/}{David Tong,
  General Relativity}: a connected course spanning geometry,
  Einstein\textquotesingle s equation, weak fields, waves, and black
  holes, with problem sheets.
\item
  \href{https://preposterousuniverse.com/grnotes/}{Sean Carroll, Lecture
  Notes on General Relativity}: a full set of introductory graduate
  notes, including special relativity, manifolds, curvature, dynamics,
  and applications. Carroll\textquotesingle s page also flags that the
  older notes are not maintained as comprehensively as the later
  textbook.
\end{itemize}

Do not compare a single curvature equation across books before comparing
their signature, Riemann definition, Ricci contraction, and units. The
physics can agree while several intermediate signs differ.

\hypertarget{d.2-choose-the-next-subject-by-the-question-that-bothers-you}{%
\section{D.2 Choose the next subject by the question that bothers
you}\label{d.2-choose-the-next-subject-by-the-question-that-bothers-you}}

\begingroup\small\setlength{\tabcolsep}{5pt}\renewcommand{\arraystretch}{1.13}

\begin{longtable}[]{@{}
  >{\raggedright\arraybackslash}p{(\columnwidth - 4\tabcolsep) * \real{0.2500}}
  >{\raggedright\arraybackslash}p{(\columnwidth - 4\tabcolsep) * \real{0.3600}}
  >{\raggedright\arraybackslash}p{(\columnwidth - 4\tabcolsep) * \real{0.3900}}@{}}
\toprule\noalign{}
\begin{minipage}[b]{\linewidth}\raggedright
If your question is\ldots{}
\end{minipage} & \begin{minipage}[b]{\linewidth}\raggedright
Study next
\end{minipage} & \begin{minipage}[b]{\linewidth}\raggedright
What to learn to do
\end{minipage} \\
\midrule\noalign{}
\endhead
\bottomrule\noalign{}
\endlastfoot
``What does curvature do beyond symmetric examples?'' & Differential
geometry and Jacobi fields & Move between abstract, coordinate, and
frame formulations \\
``How do we evolve a binary black hole?'' & Initial-value GR and
numerical relativity & Solve constraints, choose gauge, understand
stability and waveform extraction \\
``How do I calculate what an observer sees?'' & Relativistic
astrophysics, ray tracing, and geometric optics & Connect tetrads, null
geodesics, emission models, and detector quantities \\
``How do waves carry energy?'' & Perturbation theory, asymptotic
structure, and radiation theory & Distinguish local approximations from
asymptotic flux definitions \\
``How can a horizon have entropy?'' & Quantum fields in curved spacetime
and black-hole thermodynamics & Separate classical horizon geometry from
quantum state and detector effects \\
``Why this action rather than another?'' & Effective field theory and
gravitational amplitudes & State symmetry and scale assumptions and
calculate controlled corrections \\
``What is conserved in a universe without a preferred time?'' &
Symmetries, Hamiltonian GR, and covariant phase space & Construct
charges with explicit boundary and symmetry assumptions \\
``Why did this take Einstein so long?'' & History based on notebooks and
correspondence & Separate a modern textbook derivation from the actual
process of discovery \\
\end{longtable}

\endgroup

\hypertarget{d.3-research-and-historical-sources-cited-in-the-chapters}{%
\section{D.3 Research and historical sources cited in the
chapters}\label{d.3-research-and-historical-sources-cited-in-the-chapters}}

This list gathers the sources linked at the point of use. A link to a
paper does not mean every interpretation of its subject is settled. The
chapter text specifies whether a statement is classical, perturbative,
semiclassical, observational, or an open problem.

\begin{enumerate}
\def\labelenumi{\arabic{enumi}.}
\tightlist
\item
  \href{https://www.mpiwg-berlin.mpg.de/Preprints/P264.PDF}{Janssen and
  Renn, \emph{Untying the Knot}}
\item
  \href{https://online.ucpress.edu/hsns/article/14/2/253/47626/How-Einstein-Found-His-Field-Equations-1912-1915}{Norton,
  \emph{How Einstein Found His Field Equations: 1912-1915}}
\item
  \href{https://davidtong.org/pdfs/teaching/general-relativity/gr2.pdf}{David
  Tong\textquotesingle s differential-geometry chapter}
\item
  \href{https://sites.pitt.edu/~jdnorton/teaching/Einstein_graduate/pdfs/Einstein_STR_1905_English.pdf}{Einstein\textquotesingle s
  1905 paper, in English translation}
\item
  \href{https://davidtong.org/pdfs/teaching/general-relativity/gr3.pdf}{Tong\textquotesingle s
  discussion of the metric volume form}
\item
  \href{https://arxiv.org/abs/1403.7377}{Clifford Will\textquotesingle s
  review of tests of gravitation}
\item
  \href{https://davidtong.org/pdfs/teaching/general-relativity/gr1.pdf}{Tong\textquotesingle s
  treatment of the equivalence principle and Rindler motion}
\item
  \href{https://arxiv.org/pdf/gr-qc/9712019}{Sean
  Carroll\textquotesingle s notes, in the geodesics section}
\item
  \href{https://ned.ipac.caltech.edu/level5/March01/Carroll3/Carroll3.html}{Sean
  Carroll\textquotesingle s university lecture notes, ``Curvature''}
\item
  \href{https://davidtong.org/pdfs/teaching/general-relativity/gr.pdf}{David
  Tong\textquotesingle s general relativity notes}
\item
  \href{https://doi.org/10.1063/1.1724316}{Manasse and
  Misner\textquotesingle s original Fermi-coordinate paper}
\item
  \href{https://arxiv.org/abs/gr-qc/0209024}{Pravda, Pravdova, Coley,
  and Milson}
\item
  \href{https://link.aps.org/doi/10.1103/PhysRev.116.1045}{Misner and
  Putnam in ``Active Gravitational Mass''}
\item
  \href{https://arxiv.org/abs/1609.00207}{Lehner, Myers, Poisson, and
  Sorkin, ``Gravitational action with null boundaries''}
\item
  \href{https://arxiv.org/abs/gr-qc/9209012}{``Quasilocal Energy and
  Conserved Charges Derived from the Gravitational Action''}
\item
  \href{https://davidtong.org/pdfs/teaching/general-relativity/gr6.pdf}{David
  Tong's black-hole lecture notes}
\item
  \href{https://arxiv.org/abs/0806.0016}{Chruściel and Costa, \emph{On
  uniqueness of stationary vacuum black holes}}
\item
  \href{https://doi.org/10.1103/PhysRev.166.1272}{Isaacson,
  \emph{Gravitational Radiation in the Limit of High Frequency. II}}
\item
  \href{https://arxiv.org/abs/1602.03837}{LIGO Scientific Collaboration
  and Virgo Collaboration, \emph{Observation of Gravitational Waves from
  a Binary Black Hole Merger}}
\item
  \href{https://arxiv.org/abs/astro-ph/9805201}{Riess and collaborators,
  \emph{Observational Evidence from Supernovae for an Accelerating
  Universe and a Cosmological Constant}}
\item
  \href{https://arxiv.org/abs/gr-qc/0703035}{Éric
  Gourgoulhon\textquotesingle s author-written notes on the 3+1
  formalism}
\item
  \href{https://projecteuclid.org/journals/communications-in-mathematical-physics/volume-14/issue-4/Global-aspects-of-the-Cauchy-problem-in-general-relativity/cmp/1103841822.pdf}{Choquet-Bruhat
  and Geroch\textquotesingle s original Cauchy-problem paper}
\item
  \href{https://davidtong.org/teaching/general-relativity/grhtml/S3}{David
  Tong\textquotesingle s author-written chapter on connections and
  Cartan geometry}
\item
  \href{https://arxiv.org/abs/1010.0869}{Dadhich and
  Pons\textquotesingle s paper on Einstein-Hilbert and Einstein-Palatini
  formulations}
\item
  \href{https://link.aps.org/doi/10.1103/PhysRevLett.14.57}{Penrose\textquotesingle s
  1965 paper, ``Gravitational Collapse and Space-Time Singularities''}
\item
  \href{https://arxiv.org/abs/gr-qc/9912119}{Wald\textquotesingle s
  research review of black-hole thermodynamics}
\item
  \href{https://arxiv.org/abs/gr-qc/9405057}{Donoghue\textquotesingle s
  original work on general relativity as an effective field theory}
\item
  \href{https://arxiv.org/abs/2211.09902}{Donoghue\textquotesingle s
  review of quantum GR and its effective-theory limits}
\item
  \href{https://arxiv.org/abs/1709.09695}{Solomon and
  Trodden\textquotesingle s research on higher derivatives in EFT}
\item
  \href{https://arxiv.org/abs/gr-qc/0411023}{Deser\textquotesingle s
  ``Self-Interaction and Gauge Invariance''}
\item
  \href{https://arxiv.org/html/1005.2386v4}{``Lovelock\textquotesingle s
  theorem revisited''}
\item
  \href{https://arxiv.org/abs/1205.3365}{Jérôme Martin\textquotesingle s
  review of the cosmological constant problem}
\item
  \href{https://arxiv.org/abs/2607.19293}{LVK\textquotesingle s primary
  GWTC-5.0 tests paper}
\item
  \href{https://davidtong.org/teaching/general-relativity/}{David Tong,
  General Relativity}
\item
  \href{https://preposterousuniverse.com/grnotes/}{Sean Carroll, Lecture
  Notes on General Relativity}
\end{enumerate}

\textbf{One final challenge:} explain the Einstein equation to a friend
without saying ``mass bends a rubber sheet.'' Use a clock, two
neighboring freely falling laboratories, a rule for comparing their
directions, and an action whose stationary points determine the
geometry. If you can do that - and explain why empty spacetime can still
carry waves - you have moved well beyond recognizing the symbols.

\backmatter
\end{document}
